% Generated mechanically from the frozen cohomology.tex target.
% Do not edit this generated file; edit the bound locale source instead.

% OK, start here.
%

\setcounter{chapter}{19}
\chapter{层上同调}



\phantomsection
\label{cohomology-section-phantom}


\section{引言}
\label{cohomology-section-introduction}

\noindent
本章讨论拓扑空间上层的上同调若干专题。我们主要在环层上的模这一一般框架下
工作，并研究环空间的态射。关于站点上的层会出现什么情形，见《站点上的
上同调》第 \ref{sites-cohomology-section-introduction} 节。基本参考文献为
\cite{Godement} 与 \cite{Iversen}。





\section{层上同调}
\label{cohomology-section-cohomology-sheaves}

\noindent
设 $X$ 为拓扑空间，$\mathcal{F}$ 为阿贝尔层。我们已知 $X$ 上阿贝尔层的
范畴有足够多的内射对象；见《内射对象》引理
\ref{injectives-lemma-abelian-sheaves-space}。因此可以选取内射消解
$\mathcal{F}[0] \to \mathcal{I}^\bullet$。按照通常约定，定义
\begin{equation}
\label{cohomology-equation-cohomology}
H^i(X, \mathcal{F}) = H^i(\Gamma(X, \mathcal{I}^\bullet))
\end{equation}
并称之为阿贝尔层 $\mathcal{F}$ 的{\it 第 $i$ 个上同调群}。
函子族 $H^i(X, -)$ 构成从 $\textit{Ab}(X)$ 到 $\textit{Ab}$ 的
泛 $\delta$-函子。

\medskip\noindent
设 $f : X \to Y$ 为拓扑空间的连续映射。沿用上面的
$\mathcal{F}[0] \to \mathcal{I}^\bullet$，定义
\begin{equation}
\label{cohomology-equation-higher-direct-image}
R^if_*\mathcal{F} = H^i(f_*\mathcal{I}^\bullet)
\end{equation}
并称之为 $\mathcal{F}$ 的{\it 第 $i$ 个高阶直像}。
函子族 $R^if_*$ 构成从 $\textit{Ab}(X)$ 到 $\textit{Ab}(Y)$ 的
泛 $\delta$-函子。

\medskip\noindent
设 $(X, \mathcal{O}_X)$ 为环空间，$\mathcal{F}$ 为
$\mathcal{O}_X$-模。我们已知 $X$ 上 $\mathcal{O}_X$-模的范畴有足够多的
内射对象；见《内射对象》引理
\ref{injectives-lemma-sheaves-modules-space}。因此可以选取内射消解
$\mathcal{F}[0] \to \mathcal{I}^\bullet$。按照通常约定，定义
\begin{equation}
\label{cohomology-equation-cohomology-modules}
H^i(X, \mathcal{F}) = H^i(\Gamma(X, \mathcal{I}^\bullet))
\end{equation}
并称之为 $\mathcal{F}$ 的{\it 第 $i$ 个上同调群}。
函子族 $H^i(X, -)$ 构成从 $\textit{Mod}(\mathcal{O}_X)$ 到
$\text{Mod}_{\mathcal{O}_X(X)}$ 的泛 $\delta$-函子。

\medskip\noindent
设 $f : (X, \mathcal{O}_X) \to (Y, \mathcal{O}_Y)$ 为环空间的态射。
沿用上面的 $\mathcal{F}[0] \to \mathcal{I}^\bullet$，定义
\begin{equation}
\label{cohomology-equation-higher-direct-image-modules}
R^if_*\mathcal{F} = H^i(f_*\mathcal{I}^\bullet)
\end{equation}
并称之为 $\mathcal{F}$ 的{\it 第 $i$ 个高阶直像}。
函子族 $R^if_*$ 构成从 $\textit{Mod}(\mathcal{O}_X)$ 到
$\textit{Mod}(\mathcal{O}_Y)$ 的泛 $\delta$-函子。





\section{导出函子}
\label{cohomology-section-derived-functors}

\noindent
下面简要说明如何利用消解函子得到右导出函子。关于无界导出函子，见
第 \ref{cohomology-section-unbounded} 节。

\medskip\noindent
设 $(X, \mathcal{O}_X)$ 为环空间。范畴
$\textit{Mod}(\mathcal{O}_X)$ 是阿贝尔范畴；见《模层》引理
\ref{modules-lemma-abelian}。本章记
$$
K(\mathcal{O}_X) = K(\textit{Mod}(\mathcal{O}_X))
\quad
\text{以及}
\quad
D(\mathcal{O}_X) = D(\textit{Mod}(\mathcal{O}_X)).
$$
对于《导出范畴》定义 \ref{derived-definition-complexes-notation} 与定义
\ref{derived-definition-unbounded-derived-category} 中引入的三角范畴，
其有界版本采用类似记号。由《导出范畴》注
\ref{derived-remark-big-abelian-category}，存在消解函子
$$
j = j_X :
K^{+}(\textit{Mod}(\mathcal{O}_X))
\longrightarrow
K^{+}(\mathcal{I})
$$
其中 $\mathcal{I}$ 是 $\textit{Mod}(\mathcal{O}_X)$ 中由内射层构成的
严格满加性子范畴。对任意取值于阿贝尔范畴 $\mathcal{B}$ 的左正合函子
$F : \textit{Mod}(\mathcal{O}_X) \to \mathcal{B}$
，以 $RF$ 表示《导出范畴》第 \ref{derived-section-right-derived-functor} 节
所述、利用刚才的消解函子 $j_X$ 构造的右导出函子：
\begin{equation}
\label{cohomology-equation-RF}
RF = F \circ j_X' : D^{+}(X) \longrightarrow D^{+}(\mathcal{B})
\end{equation}
记号见《导出范畴》引理 \ref{derived-lemma-right-derived-functor}。
注意，视具体情形，可把 $RF$ 看作定义于
$\textit{Mod}(\mathcal{O}_X)$、
$\text{Comp}^{+}(\textit{Mod}(\mathcal{O}_X))$,
$K^{+}(X)$ 或 $D^{+}(X)$ 上。按照《导出范畴》定义
\ref{derived-definition-higher-derived-functors}，得到第 $i$ 个右导出函子
\begin{equation}
\label{cohomology-equation-RFi}
R^iF = H^i \circ RF : \textit{Mod}(\mathcal{O}_X) \longrightarrow \mathcal{B}
\end{equation}
从而 $R^0F = F$，且 $\{R^iF, \delta\}_{i \geq 0}$ 是泛
$\delta$-函子；见《导出范畴》引理
\ref{derived-lemma-higher-derived-functors}。

\medskip\noindent
下面给出这一构造的两个特殊情形。给定环 $R$，记
$K(R) = K(\text{Mod}_R)$ 与 $D(R) = D(\text{Mod}_R)$，有界版本也采用
类似记号。对任意开集 $U \subset X$，有左正合函子
$
\Gamma(U, -) :
\textit{Mod}(\mathcal{O}_X)
\longrightarrow
\text{Mod}_{\mathcal{O}_X(U)}
$
，由上述讨论得到
\begin{equation}
\label{cohomology-equation-total-derived-cohomology}
R\Gamma(U, -) :
D^{+}(X)
\longrightarrow
D^{+}(\mathcal{O}_X(U))
\end{equation}
。置 $H^i(U, -) = R^i\Gamma(U, -)$。当 $U = X$ 时，便恢复
(\ref{cohomology-equation-cohomology-modules})。若 $f : X \to Y$ 是环空间的态射，
则有左正合函子
$
f_* :
\textit{Mod}(\mathcal{O}_X)
\longrightarrow
\textit{Mod}(\mathcal{O}_Y)
$
，由此得到{\it 导出直像}
\begin{equation}
\label{cohomology-equation-total-derived-direct-image}
Rf_* :
D^{+}(X)
\longrightarrow
D^{+}(Y)
\end{equation}
。$Rf_*\mathcal{F}^\bullet$ 的第 $i$ 个上同调层记为
$R^if_*\mathcal{F}^\bullet$，并依照
(\ref{cohomology-equation-higher-direct-image-modules}) 称为第 $i$ 个{\it 高阶直像}。
上述两个显示式中的函子都是导出范畴之间的正合函子。

\medskip\noindent
{\bf 记号的滥用：}当函子 $Rf_*$ 或任何其他导出函子作用于 $X$ 上的层
$\mathcal{F}$ 或层复形时，约定先把 $\mathcal{F}$ 替换为一个适当的消解。
为方便进行这类操作，对给定对象
$\mathcal{F}^\bullet \in D(\mathcal{O}_X)$，若存在拟同构
$\mathcal{F}^\bullet \to \mathcal{I}^\bullet$，其中
$\mathcal{I}^\bullet$ 是由 $\textit{Mod}(\mathcal{O}_X)$ 中内射对象组成的
下有界复形，则称 $\mathcal{I}^\bullet$ 在导出范畴中{\it 表示
$\mathcal{F}^\bullet$}。同理，“设
$\alpha : \mathcal{F}^\bullet \to \mathcal{G}^\bullet$ 为
$D(\mathcal{O}_X)$ 的态射”并不意味着 $\alpha$ 由复形态射表示。若所指确为
复形态射，我们会明确说明。









\section{第一上同调与挠子}
\label{cohomology-section-h1-torsors}

\begin{definition}
\label{cohomology-definition-torsor}
设 $X$ 为拓扑空间，$\mathcal{G}$ 为 $X$ 上的（未必交换的）群层。
所谓{\it 伪挠子}，更准确地说，所谓{\it 伪 $\mathcal{G}$-挠子}，是指
$X$ 上的集合层 $\mathcal{F}$，配备作用
$\mathcal{G} \times \mathcal{F} \to \mathcal{F}$，使得
\begin{enumerate}
\item 只要 $\mathcal{F}(U)$ 非空，作用
$\mathcal{G}(U) \times \mathcal{F}(U) \to \mathcal{F}(U)$
便是单可迁的；
\end{enumerate}
伪 $\mathcal{G}$-挠子的态射 $\mathcal{F} \to \mathcal{F}'$ 是与
$\mathcal{G}$-作用相容的集合层态射。所谓{\it 挠子}，更准确地说，所谓
{\it $\mathcal{G}$-挠子}，是还满足下列条件的伪 $\mathcal{G}$-挠子：
\begin{enumerate}
\item[(2)] 对每个 $x \in X$，茎 $\mathcal{F}_x$ 非空。
\end{enumerate}
$\mathcal{G}$-挠子的态射是伪 $\mathcal{G}$-挠子的态射。
{\it 平凡 $\mathcal{G}$-挠子}是配备显然的左 $\mathcal{G}$-作用的层
$\mathcal{G}$。
\end{definition}

\noindent
显然，挠子之间的态射自动是同构。

\begin{lemma}
\label{cohomology-lemma-trivial-torsor}
设 $X$ 为拓扑空间，$\mathcal{G}$ 为 $X$ 上的（未必交换的）群层。
$\mathcal{G}$-挠子 $\mathcal{F}$ 平凡，当且仅当
$\mathcal{F}(X) \not = \emptyset$。
\end{lemma}

\begin{proof}
从略。
\end{proof}

\begin{lemma}
\label{cohomology-lemma-torsors-h1}
设 $X$ 为拓扑空间，$\mathcal{H}$ 为 $X$ 上的阿贝尔层。
$\mathcal{H}$-挠子的同构类集合与 $H^1(X, \mathcal{H})$ 之间有典范双射。
\end{lemma}

\begin{proof}
设 $\mathcal{F}$ 为 $\mathcal{H}$-挠子。考虑由 $\mathcal{F}$ 生成的
自由阿贝尔层 $\mathbf{Z}[\mathcal{F}]$。它是下述规则的层化：对开集
$U \subset X$，取所有有限形式和 $\sum n_i[s_i]$，其中
$n_i \in \mathbf{Z}$ 且 $s_i \in \mathcal{F}(U)$。存在自然映射
$$
\sigma : \mathbf{Z}[\mathcal{F}] \longrightarrow \underline{\mathbf{Z}}
$$
把局部截面 $\sum n_i[s_i]$ 映到 $\sum n_i$。$\sigma$ 的核由形如
$[s] - [s']$ 的局部截面生成。存在典范映射
$a : \Ker(\sigma) \to \mathcal{H}$
，把 $[s] - [s']$ 映到满足 $h \cdot s' = s$ 的 $\mathcal{H}$ 的局部截面
$h$。考虑推出图
$$
\xymatrix{
0 \ar[r] &
\Ker(\sigma) \ar[r] \ar[d]^a &
\mathbf{Z}[\mathcal{F}] \ar[r] \ar[d] &
\underline{\mathbf{Z}} \ar[r] \ar[d] &
0 \\
0 \ar[r] &
\mathcal{H} \ar[r] &
\mathcal{E} \ar[r] &
\underline{\mathbf{Z}} \ar[r] &
0
}
$$
其中 $\mathcal{E}$ 是通过推出得到的扩张。对下方短正合列应用相应的上同调
长正合列，并把连接同态作用于
$1 \in H^0(X, \underline{\mathbf{Z}})$，得到元素
$\xi = \xi_\mathcal{F} \in H^1(X, \mathcal{H})$
。

\medskip\noindent
反过来，给定 $\xi \in H^1(X, \mathcal{H})$，可按如下方式为它配上一个
挠子。选取 $\mathcal{H}$ 到内射阿贝尔层 $\mathcal{I}$ 中的嵌入
$\mathcal{H} \to \mathcal{I}$。置
$\mathcal{Q} = \mathcal{I}/\mathcal{H}$，于是有短正合列
$$
\xymatrix{
0 \ar[r] &
\mathcal{H} \ar[r] &
\mathcal{I} \ar[r] &
\mathcal{Q} \ar[r] &
0
}
$$
因为 $H^1(X, \mathcal{I}) = 0$（见《导出范畴》引理
\ref{derived-lemma-higher-derived-functors}），元素 $\xi$ 是某个整体截面
$q \in H^0(X, \mathcal{Q})$ 的像。令 $\mathcal{F} \subset \mathcal{I}$
为由在层 $\mathcal{Q}$ 中映到 $q$ 的截面组成的子层（视为集合层）。容易验证
$\mathcal{F}$ 是挠子。

\medskip\noindent
关于上述两个构造互为逆构造的验证从略。
\end{proof}







\section{第一上同调与扩张}
\label{cohomology-section-h1-extensions}

\begin{lemma}
\label{cohomology-lemma-h1-extensions}
设 $(X, \mathcal{O}_X)$ 为环空间，$\mathcal{F}$ 为
$\mathcal{O}_X$-模层。存在典范双射
$$
\Ext^1_{\textit{Mod}(\mathcal{O}_X)}(\mathcal{O}_X, \mathcal{F})
\longrightarrow
H^1(X, \mathcal{F})
$$
，它把扩张
$$
0 \to \mathcal{F} \to \mathcal{E} \to \mathcal{O}_X \to 0
$$
映到 $1 \in \Gamma(X, \mathcal{O}_X)$ 在 $H^1(X, \mathcal{F})$ 中的像。
\end{lemma}

\begin{proof}
下面构造引理中映射的逆映射。设 $\xi \in H^1(X, \mathcal{F})$。
选取嵌入 $\mathcal{F} \subset \mathcal{I}$，其中 $\mathcal{I}$ 是
$\textit{Mod}(\mathcal{O}_X)$ 中的内射对象。置
$\mathcal{Q} = \mathcal{I}/\mathcal{F}$。由上同调长正合列，$\xi$ 是截面
$\tilde \xi \in \Gamma(X, \mathcal{Q}) =
\Hom_{\mathcal{O}_X}(\mathcal{O}_X, \mathcal{Q})$ 的像。现在取拉回
$$
\xymatrix{
0 \ar[r] &
\mathcal{F} \ar[r] \ar@{=}[d] &
\mathcal{E} \ar[r] \ar[d] &
\mathcal{O}_X \ar[r] \ar[d]^{\tilde \xi} &
0 \\
0 \ar[r] &
\mathcal{F} \ar[r] &
\mathcal{I} \ar[r] &
\mathcal{Q} \ar[r] &
0
}
$$
；见《同调》第 \ref{homology-section-extensions} 节。
\end{proof}








\section{第一上同调与可逆模层}
\label{cohomology-section-invertible-sheaves}

\noindent
环空间的 Picard 群定义见《模层》第 \ref{modules-section-invertible} 节。

\begin{lemma}
\label{cohomology-lemma-h1-invertible}
设 $(X, \mathcal{O}_X)$ 为环空间。若所有茎 $\mathcal{O}_{X, x}$ 都是
局部环，则存在阿贝尔群的典范同构
$$
H^1(X, \mathcal{O}_X^*) = \Pic(X).
$$
\end{lemma}

\begin{proof}
设 $\mathcal{L}$ 为可逆 $\mathcal{O}_X$-模。考虑由下述规则定义的预层
$\mathcal{L}^*$：
$$
U \longmapsto \{s \in \mathcal{L}(U)
\text{ 使得 } \mathcal{O}_U \xrightarrow{s \cdot -} \mathcal{L}_U
\text{ 是同构}\}
$$
该预层满足层条件。此外，若 $f \in \mathcal{O}_X^*(U)$ 且
$s \in \mathcal{L}^*(U)$，则显然 $fs \in \mathcal{L}^*(U)$。同理，若
$s, s' \in \mathcal{L}^*(U)$，则存在唯一的
$f \in \mathcal{O}_X^*(U)$ 使 $fs = s'$。再者，由《模层》引理
\ref{modules-lemma-invertible-is-locally-free-rank-1}，层
$\mathcal{L}^*$ 局部有截面。换言之，$\mathcal{L}^*$ 是
$\mathcal{O}_X^*$-挠子。因此得到映射
$$
\begin{matrix}
\text{环空间 }(X, \mathcal{O}_X)\text{ 上的可逆模层} \\
\text{不计同构}
\end{matrix}
\longrightarrow
\begin{matrix}
\mathcal{O}_X^*\text{-挠子} \\
\text{不计同构}
\end{matrix}
$$
关于它是阿贝尔群同态的验证从略。由引理 \ref{cohomology-lemma-torsors-h1}，右侧与
$H^1(X, \mathcal{O}_X^*)$ 典范双射。因此只需证明此映射既单又满。

\medskip\noindent
单射性。若挠子 $\mathcal{L}^*$ 平凡，则由引理
\ref{cohomology-lemma-trivial-torsor}，$\mathcal{L}^*$ 有整体截面。这恰好意味着
$\mathcal{L} \cong \mathcal{O}_X$ 是 $\Pic(X)$ 中的单位元。

\medskip\noindent
满射性。设 $\mathcal{F}$ 为 $\mathcal{O}_X^*$-挠子。考虑集合预层
$$
\mathcal{L}_1 : U \longmapsto
(\mathcal{F}(U) \times \mathcal{O}_X(U))/\mathcal{O}_X^*(U)
$$
，其中 $f \in \mathcal{O}_X^*(U)$ 在 $(s, g)$ 上的作用为
$(fs, f^{-1}g)$。规定
$(s, g) + (s', g') = (s, g + (s'/s)g')$，其中 $s'/s$ 是满足
$fs = s'$ 的 $\mathcal{O}_X^*$ 的局部截面 $f$；并对
$\mathcal{O}_X$ 的局部截面 $h$ 规定 $h(s, g) = (s, hg)$。如此
$\mathcal{L}_1$ 成为 $\mathcal{O}_X$-模预层。关于其层化
$\mathcal{L} = \mathcal{L}_1^\#$ 是可逆 $\mathcal{O}_X$-模，且相应的
$\mathcal{O}_X^*$-挠子 $\mathcal{L}^*$ 同构于 $\mathcal{F}$ 的验证从略。
\end{proof}













\section{上同调的局部性}
\label{cohomology-section-locality}

\noindent
下面的引理说明，在开集上定义层 $\mathcal{F}$ 的上同调不会产生歧义。

\begin{lemma}
\label{cohomology-lemma-cohomology-of-open}
设 $X$ 是环化空间。
设 $U \subset X$ 是开子空间。
\begin{enumerate}
\item 若 $\mathcal{I}$ 是内射 $\mathcal{O}_X$-模，
则 $\mathcal{I}|_U$ 是内射 $\mathcal{O}_U$-模。
\item 对任意 $\mathcal{O}_X$-模层 $\mathcal{F}$，有
$H^p(U, \mathcal{F}) = H^p(U, \mathcal{F}|_U)$.
\end{enumerate}
\end{lemma}

\begin{proof}
以 $j : U \to X$ 表示该开浸入。
回忆限制到 $U$ 的函子 $j^{-1}$ 是延零函子 $j_!$ 的右伴随，见
《层》引理 \ref{sheaves-lemma-j-shriek-modules}。
此外，$j_!$ 是正合的。因此 (1) 由
《同调代数》引理 \ref{homology-lemma-adjoint-preserve-injectives}
推出。

\medskip\noindent
按定义，
$H^p(U, \mathcal{F}) = H^p(\Gamma(U, \mathcal{I}^\bullet))$，
其中 $\mathcal{F} \to \mathcal{I}^\bullet$ 是
$\textit{Mod}(\mathcal{O}_X)$ 中的内射消解。
由上可知，$\mathcal{F}|_U \to \mathcal{I}^\bullet|_U$
是 $\textit{Mod}(\mathcal{O}_U)$ 中的内射消解。
因此 $H^p(U, \mathcal{F}|_U)$ 等于
$H^p(\Gamma(U, \mathcal{I}^\bullet|_U))$。
当然，对 $X$ 上任意层 $\mathcal{F}$，都有
$\Gamma(U, \mathcal{F}) = \Gamma(U, \mathcal{F}|_U)$。
故 (2) 中的等式成立。
\end{proof}

\noindent
设 $X$ 是环化空间，$\mathcal{F}$ 是 $\mathcal{O}_X$-模层，
并设 $U \subset V \subset X$ 是开子集。
则存在典范的{\it 限制映射}
\begin{equation}
\label{cohomology-equation-restriction-mapping}
H^n(V, \mathcal{F})
\longrightarrow
H^n(U, \mathcal{F}), \quad
\xi \longmapsto \xi|_U
\end{equation}
且它关于 $\mathcal{F}$ 是函子性的。事实上，任取内射消解
$\mathcal{F} \to \mathcal{I}^\bullet$。层 $\mathcal{I}^p$
的限制映射给出复形态射
$$
\Gamma(V, \mathcal{I}^\bullet)
\longrightarrow
\Gamma(U, \mathcal{I}^\bullet)
$$
左端是表示 $R\Gamma(V, \mathcal{F})$ 的复形，
右端是表示 $R\Gamma(U, \mathcal{F})$ 的复形。
施用函子 $H^n$ 即得到上同调群之间的映射。
如上所示，我们用记号 $\xi \mapsto \xi|_U$ 表示此映射。
因此规则 $U \mapsto H^n(U, \mathcal{F})$ 构成
$\mathcal{O}_X$-模预层。该预层通常记作
$\underline{H}^n(\mathcal{F})$。
我们将在引理 \ref{cohomology-lemma-include} 中给出它的另一种解释。

\begin{lemma}
\label{cohomology-lemma-kill-cohomology-class-on-covering}
设 $X$ 是环化空间，$\mathcal{F}$ 是 $\mathcal{O}_X$-模层，
$U \subset X$ 是开子空间。设 $n > 0$ 且
$\xi \in H^n(U, \mathcal{F})$。则存在开覆盖
$U = \bigcup_{i\in I} U_i$，使得对所有 $i \in I$，
均有 $\xi|_{U_i} = 0$。
\end{lemma}

\begin{proof}
设 $\mathcal{F} \to \mathcal{I}^\bullet$ 是内射消解。于是
$$
H^n(U, \mathcal{F}) =
\frac{\Ker(\mathcal{I}^n(U) \to \mathcal{I}^{n + 1}(U))}
{\Im(\mathcal{I}^{n - 1}(U) \to \mathcal{I}^n(U))}.
$$
取表示上述商中该上同调类的元素
$\tilde \xi \in \mathcal{I}^n(U)$。
由于 $\mathcal{I}^\bullet$ 是 $\mathcal{F}$ 的内射消解且 $n > 0$，
复形 $\mathcal{I}^\bullet$ 在次数 $n$ 处正合。因此，作为层有
$\Im(\mathcal{I}^{n - 1} \to \mathcal{I}^n) =
\Ker(\mathcal{I}^n \to \mathcal{I}^{n + 1})$。
由于 $\tilde \xi$ 是核层在 $U$ 上的截面，存在开覆盖
$U = \bigcup_{i \in I} U_i$，使得
$\tilde \xi|_{U_i}$ 是某个截面
$\xi_i \in \mathcal{I}^{n - 1}(U_i)$ 在 $d$ 下的像。
按限制 $\xi|_{U_i}$ 的定义，它对应于
$\tilde \xi|_{U_i}$ 的类，故结论成立。
\end{proof}

\begin{lemma}
\label{cohomology-lemma-describe-higher-direct-images}
设 $f : X \to Y$ 是环化空间的态射，
$\mathcal{F}$ 是 $\mathcal{O}_X$-模。
层 $R^if_*\mathcal{F}$ 是与预层
$$
V \longmapsto H^i(f^{-1}(V), \mathcal{F})
$$
相伴的层，其中限制映射如
Equation (\ref{cohomology-equation-restriction-mapping}) 所示。
对 $R^if_*$ 施用于下有界复形 $\mathcal{F}^\bullet$，也有类似陈述。
\end{lemma}

\begin{proof}
设 $\mathcal{F} \to \mathcal{I}^\bullet$ 是内射消解。
按定义，$R^if_*\mathcal{F}$ 是复形
$$
f_*\mathcal{I}^0 \to f_*\mathcal{I}^1 \to f_*\mathcal{I}^2 \to \ldots
$$
的第 $i$ 个上同调层。由 $\mathcal{O}_Y$-模上阿贝尔范畴结构的定义，
该上同调层是与预层
$$
V
\longmapsto
\frac{\Ker(f_*\mathcal{I}^i(V) \to f_*\mathcal{I}^{i + 1}(V))}
{\Im(f_*\mathcal{I}^{i - 1}(V) \to f_*\mathcal{I}^i(V))}
$$
相伴的层，而此商显然等于
$$
\frac{\Ker(\mathcal{I}^i(f^{-1}(V)) \to \mathcal{I}^{i + 1}(f^{-1}(V)))}
{\Im(\mathcal{I}^{i - 1}(f^{-1}(V)) \to \mathcal{I}^i(f^{-1}(V)))}
$$
后者等于 $H^i(f^{-1}(V), \mathcal{F})$，故结论成立。
\end{proof}

\begin{lemma}
\label{cohomology-lemma-localize-higher-direct-images}
设 $f : X \to Y$ 是环化空间的态射，
$\mathcal{F}$ 是 $\mathcal{O}_X$-模，且
$V \subset Y$ 是开子空间。以
$g : f^{-1}(V) \to V$ 表示 $f$ 的限制。则有
$$
R^pg_*(\mathcal{F}|_{f^{-1}(V)}) = (R^pf_*\mathcal{F})|_V
$$
对导出直像 $Rf_*\mathcal{F}^\bullet$ 也有类似陈述，
其中 $\mathcal{F}^\bullet$ 是 $\mathcal{O}_X$-模的下有界复形。
\end{lemma}

\begin{proof}
第一种证明。施用引理 \ref{cohomology-lemma-describe-higher-direct-images}
和 \ref{cohomology-lemma-cohomology-of-open} 即得所示等式。
第二种证明。取内射消解 $\mathcal{F} \to \mathcal{I}^\bullet$，
并利用
$\mathcal{F}|_{f^{-1}(V)} \to \mathcal{I}^\bullet|_{f^{-1}(V)}$
也是内射消解。
\end{proof}

\begin{remark}
\label{cohomology-remark-daniel}
下面给出上述引理 \ref{cohomology-lemma-kill-cohomology-class-on-covering}
和 \ref{cohomology-lemma-describe-higher-direct-images} 的另一种证明方法。
设 $(X, \mathcal{O}_X)$ 是环化空间，
$i_X : \textit{Mod}(\mathcal{O}_X) \to \textit{PMod}(\mathcal{O}_X)$
是包含函子，$\#$ 是层化函子。回忆 $i_X$ 左正合而 $\#$ 正合。
\begin{enumerate}
\item 先证明下面的引理 \ref{cohomology-lemma-include}；它说明 $i_X$ 的右导出函子为
$R^pi_X\mathcal{F} = \underline{H}^p(\mathcal{F})$.
也可如下证明：当 $p = 0$ 时等式显然成立。
$(R^pi_X)_{p \geq 0}$ 与 $(\underline{H}^p)_{p \geq 0}$
都是在内射对象上消失的 $\delta$-函子，故二者都是泛的，从而彼此同构。见
《同调代数》第 \ref{homology-section-cohomological-delta-functor} 节。
\item 引理 \ref{cohomology-lemma-kill-cohomology-class-on-covering}
可重述为：对任意 $\mathcal{O}_X$-模层 $\mathcal{F}$ 及 $p > 0$，
有 $(\underline{H}^p(\mathcal{F}))^\# = 0$。
为说明这一点，利用 ${}^\#$ 正合可得
$$
(\underline{H}^p(\mathcal{F}))^\# =
(R^pi_X\mathcal{F})^\# =
R^p(\# \circ i_X)(\mathcal{F}) = 0
$$
因为 $\# \circ i_X$ 是恒等函子。
\item 设 $f : X \to Y$ 是环化空间的态射，
$\mathcal{F}$ 是 $\mathcal{O}_X$-模。预层
$V \mapsto H^p(f^{-1}V, \mathcal{F})$ 等于
$R^p (i_Y \circ f_*)\mathcal{F}$。
与上面 (1) 的论证一样，注意二者都给出泛 $\delta$-函子即可证明此点。
因此引理 \ref{cohomology-lemma-describe-higher-direct-images} 说明
$R^p f_* \mathcal{F}= (R^p (i_Y \circ f_*)\mathcal{F})^\#$。
再次利用 $\#$ 正合且 $\# \circ i_Y$ 是恒等函子，得到
$$
R^p f_* \mathcal{F} =
R^p(\# \circ i_Y \circ f_*)\mathcal{F} =
(R^p (i_Y \circ f_*)\mathcal{F})^\#
$$
如所需。
\end{enumerate}
\end{remark}






\section{Mayer--Vietoris 序列}
\label{cohomology-section-mayer-vietoris}

\noindent
下文将构造从 {\v C}ech 上同调到层上同调的谱序列，见引理 \ref{cohomology-lemma-cech-spectral-sequence}。
该谱序列的一个特例是 Mayer--Vietoris 长正合序列。
由于它是谱序列的一种基础、实用而易于理解的变体，
我们在此单独讨论。

\begin{lemma}
\label{cohomology-lemma-injective-restriction-surjective}
\begin{slogan}
内射层是松弛层。
\end{slogan}
设 $X$ 是环化空间，$U' \subset U \subset X$ 是开子空间。
对任意内射 $\mathcal{O}_X$-模 $\mathcal{I}$，
限制映射 $\mathcal{I}(U) \to \mathcal{I}(U')$ 是满射。
\end{lemma}

\begin{proof}
设 $j : U \to X$ 和 $j' : U' \to X$ 为开浸入。
回忆 $j_!\mathcal{O}_U$ 是
$\mathcal{O}_U = \mathcal{O}_X|_U$ 的延零，见
《层》第 \ref{sheaves-section-open-immersions} 节。
由于 $j_!$ 是限制函子的左伴随，对任意 $\mathcal{O}_X$-模层
$\mathcal{F}$，有
$$
\Hom_{\mathcal{O}_X}(j_!\mathcal{O}_U, \mathcal{F})
=
\Hom_{\mathcal{O}_U}(\mathcal{O}_U, \mathcal{F}|_U)
=
\mathcal{F}(U)
$$
见 《层》引理 \ref{sheaves-lemma-j-shriek-modules}。
类似地，层 $j'_!\mathcal{O}_{U'}$ 表示函子
$\mathcal{F} \mapsto \mathcal{F}(U')$。
此外有一个显然的典范 $\mathcal{O}_X$-模映射
$$
j'_!\mathcal{O}_{U'} \longrightarrow j_!\mathcal{O}_U
$$
由 Yoneda 引理（《范畴》引理 \ref{categories-lemma-yoneda}），
它对应于限制映射
$\mathcal{F}(U) \to \mathcal{F}(U')$。
根据层 $j'_!\mathcal{O}_{U'}$ 和 $j_!\mathcal{O}_U$
之茎的描述可知，上述映射是单射（见上引引理）。
因此，若 $\mathcal{I}$ 是内射 $\mathcal{O}_X$-模，则映射
$$
\Hom_{\mathcal{O}_X}(j_!\mathcal{O}_U, \mathcal{I})
\longrightarrow
\Hom_{\mathcal{O}_X}(j'_!\mathcal{O}_{U'}, \mathcal{I})
$$
是满射，见
《同调代数》引理 \ref{homology-lemma-characterize-injectives}。
综合以上事实即得引理。
\end{proof}

\begin{lemma}[Mayer-Vietoris]
\label{cohomology-lemma-mayer-vietoris}
设 $X$ 是环化空间，并设 $X = U \cup V$ 是两个开子集之并。
对每个 $\mathcal{O}_X$-模 $\mathcal{F}$，存在上同调长正合序列
$$
0 \to
H^0(X, \mathcal{F}) \to
H^0(U, \mathcal{F}) \oplus H^0(V, \mathcal{F}) \to
H^0(U \cap V, \mathcal{F}) \to
H^1(X, \mathcal{F}) \to \ldots
$$
该长正合序列关于 $\mathcal{F}$ 是函子性的。
\end{lemma}

\begin{proof}
层条件说明，对任意阿贝尔层 $\mathcal{F}$，映射
$(1, -1) : \mathcal{F}(U) \oplus \mathcal{F}(V) \to \mathcal{F}(U \cap V)$
的核等于 $\mathcal{F}(X)$ 在第一个映射下的像。
上面的引理 \ref{cohomology-lemma-injective-restriction-surjective} 说明，
只要 $\mathcal{I}$ 是内射 $\mathcal{O}_X$-模，映射
$(1, -1) : \mathcal{I}(U) \oplus \mathcal{I}(V) \to \mathcal{I}(U \cap V)$
就是满射。因此，若
$\mathcal{F} \to \mathcal{I}^\bullet$ 是 $\mathcal{F}$ 的内射消解，
便得到复形的短正合序列
$$
0 \to
\mathcal{I}^\bullet(X) \to
\mathcal{I}^\bullet(U) \oplus \mathcal{I}^\bullet(V) \to
\mathcal{I}^\bullet(U \cap V) \to
0.
$$
取上同调即得结论（使用
《同调代数》引理 \ref{homology-lemma-long-exact-sequence-cochain}）。
该序列的函子性证明从略。
\end{proof}

\begin{lemma}[相对 Mayer--Vietoris]
\label{cohomology-lemma-relative-mayer-vietoris}
设 $f : X \to Y$ 是环化空间的态射，并设
$X = U \cup V$ 是两个开子集之并。记
$a = f|_U : U \to Y$、$b = f|_V : V \to Y$，以及
$c = f|_{U \cap V} : U \cap V \to Y$。
对每个 $\mathcal{O}_X$-模 $\mathcal{F}$，存在长正合序列
$$
0 \to
f_*\mathcal{F} \to
a_*(\mathcal{F}|_U) \oplus b_*(\mathcal{F}|_V) \to
c_*(\mathcal{F}|_{U \cap V}) \to
R^1f_*\mathcal{F} \to \ldots
$$
该长正合序列关于 $\mathcal{F}$ 是函子性的。
\end{lemma}

\begin{proof}
设 $\mathcal{F} \to \mathcal{I}^\bullet$ 是 $\mathcal{F}$ 的内射消解。
我们断言，存在复形的短正合序列
$$
0 \to
f_*\mathcal{I}^\bullet \to
a_*\mathcal{I}^\bullet|_U \oplus b_*\mathcal{I}^\bullet|_V \to
c_*\mathcal{I}^\bullet|_{U \cap V} \to
0.
$$
事实上，对任意开集 $W \subset Y$ 和任意 $n \geq 0$，
相应的 $W$ 上截面群序列
$$
0 \to
\mathcal{I}^n(f^{-1}(W)) \to
\mathcal{I}^n(U \cap f^{-1}(W))
\oplus \mathcal{I}^n(V \cap f^{-1}(W)) \to
\mathcal{I}^n(U \cap V \cap f^{-1}(W)) \to
0
$$
已在引理 \ref{cohomology-lemma-mayer-vietoris} 的证明中证为短正合。
取上同调层，并利用 $\mathcal{I}^\bullet|_U$
是 $\mathcal{F}|_U$ 的内射消解这一事实，即得引理。
对 $\mathcal{I}^\bullet|_V$、
$\mathcal{I}^\bullet|_{U \cap V}$ 亦然，见引理 \ref{cohomology-lemma-cohomology-of-open}。
\end{proof}




















\section{{\v C}ech 复形与 {\v C}ech 上同调}
\label{cohomology-section-cech}

\noindent
设 $X$ 是拓扑空间，
$\mathcal{U} : U = \bigcup_{i \in I} U_i$ 是开覆盖，见
《拓扑》基本概念（\ref{topology-item-covering}）。
按照惯例，以
$U_{i_0\ldots i_p} = U_{i_0} \cap \ldots \cap U_{i_p}$
表示 $\mathcal{U}$ 中成员的 $(p + 1)$ 重交。
设 $\mathcal{F}$ 是 $X$ 上的阿贝尔预层。置
$$
\check{\mathcal{C}}^p(\mathcal{U}, \mathcal{F})
=
\prod\nolimits_{(i_0, \ldots, i_p) \in I^{p + 1}}
\mathcal{F}(U_{i_0\ldots i_p}).
$$
这是一个阿贝尔群。对
$s \in \check{\mathcal{C}}^p(\mathcal{U}, \mathcal{F})$，
以 $s_{i_0\ldots i_p}$ 表示它在
$\mathcal{F}(U_{i_0\ldots i_p})$ 中的分量。
注意，若 $s \in \check{\mathcal{C}}^1(\mathcal{U}, \mathcal{F})$
且 $i, j \in I$，则 $s_{ij}$ 与 $s_{ji}$
都是 $\mathcal{F}(U_i \cap U_j)$ 的元素，但我们并未规定它们之间的关系。
换言之，这里{\it 不}使用交错上链（将在
第 \ref{cohomology-section-alternating-cech} 节 中定义）。定义
$$
d : \check{\mathcal{C}}^p(\mathcal{U}, \mathcal{F})
\longrightarrow
\check{\mathcal{C}}^{p + 1}(\mathcal{U}, \mathcal{F})
$$
为
\begin{equation}
\label{cohomology-equation-d-cech}
d(s)_{i_0\ldots i_{p + 1}}
=
\sum\nolimits_{j = 0}^{p + 1}
(-1)^j
s_{i_0\ldots \hat i_j \ldots i_{p + 1}}|_{U_{i_0\ldots i_{p + 1}}}
\end{equation}
直接可见 $d \circ d = 0$。换言之，
$\check{\mathcal{C}}^\bullet(\mathcal{U}, \mathcal{F})$ 是复形。

\begin{definition}
\label{cohomology-definition-cech-complex}
设 $X$ 是拓扑空间，
$\mathcal{U} : U = \bigcup_{i \in I} U_i$ 是开覆盖，
$\mathcal{F}$ 是 $X$ 上的阿贝尔预层。
复形 $\check{\mathcal{C}}^\bullet(\mathcal{U}, \mathcal{F})$
称为与 $\mathcal{F}$ 和开覆盖 $\mathcal{U}$ 相伴的
{\it {\v C}ech 复形}。其上同调群
$H^i(\check{\mathcal{C}}^\bullet(\mathcal{U}, \mathcal{F}))$
称为与 $\mathcal{F}$ 和覆盖 $\mathcal{U}$ 相伴的
{\it {\v C}ech 上同调群}，记作
$\check H^i(\mathcal{U}, \mathcal{F})$。
\end{definition}

\begin{lemma}
\label{cohomology-lemma-cech-h0}
设 $X$ 是拓扑空间，$\mathcal{F}$ 是 $X$ 上的阿贝尔预层。
下列条件等价：
\begin{enumerate}
\item $\mathcal{F}$ 是阿贝尔层；
\item 对每个开覆盖 $\mathcal{U} : U = \bigcup_{i \in I} U_i$，
自然映射
$$
\mathcal{F}(U) \to \check{H}^0(\mathcal{U}, \mathcal{F})
$$
是双射。
\end{enumerate}
\end{lemma}

\begin{proof}
这是因为层条件恰好说明，对每个开覆盖，
$\mathcal{F}(U) \to \check{H}^0(\mathcal{U}, \mathcal{F})$
都是双射。
\end{proof}

\begin{lemma}
\label{cohomology-lemma-cech-trivial}
设 $X$ 是拓扑空间，$\mathcal{F}$ 是 $X$ 上的阿贝尔预层，
$\mathcal{U} : U = \bigcup_{i \in I} U_i$ 是开覆盖。
若对某个 $i \in I$ 有 $U_i = U$，则增广 {\v C}ech 复形
$$
\mathcal{F}(U) \to \check{\mathcal{C}}^\bullet(\mathcal{U}, \mathcal{F})
$$
由将 $\mathcal{F}(U)$ 置于次数 $-1$ 并取微分为从
$\mathcal{F}(U)$ 到
$\check{\mathcal{C}}^0(\mathcal{U}, \mathcal{F})$
的典范映射而得；该复形链同伦等价于零复形。
\end{lemma}

\begin{proof}
固定满足 $U = U_i$ 的元素 $i \in I$。注意，当 $i_j = i$ 时，
$U_{i_0 \ldots i_p} = U_{i_0 \ldots \hat i_j \ldots i_p}$。
定义同伦
$$
h :
\prod\nolimits_{i_0 \ldots i_{p + 1}} \mathcal{F}(U_{i_0 \ldots i_{p + 1}})
\longrightarrow
\prod\nolimits_{i_0 \ldots i_p} \mathcal{F}(U_{i_0 \ldots i_p})
$$
如下：
$$
h(s)_{i_0 \ldots i_p} = s_{i i_0 \ldots i_p}
$$
换言之，
$h : \prod_{i_0} \mathcal{F}(U_{i_0}) \to \mathcal{F}(U)$
是到因子 $\mathcal{F}(U_i) = \mathcal{F}(U)$ 的投影；
一般地，映射 $h$ 等于到下列因子的投影：
$\mathcal{F}(U_{i i_1 \ldots i_{p + 1}}) =
\mathcal{F}(U_{i_1 \ldots i_{p + 1}})$.
计算得
\begin{align*}
(dh + hd)(s)_{i_0 \ldots i_p}
& =
\sum\nolimits_{j = 0}^p
(-1)^j
h(s)_{i_0 \ldots \hat i_j \ldots i_p}
+
d(s)_{i i_0 \ldots i_p}\\
& =
\sum\nolimits_{j = 0}^p
(-1)^j
s_{i i_0 \ldots \hat i_j \ldots i_p}
+
s_{i_0 \ldots i_p}
+
\sum\nolimits_{j = 0}^p
(-1)^{j + 1}
s_{i i_0 \ldots \hat i_j \ldots i_p} \\
& =
s_{i_0 \ldots i_p}
\end{align*}
这证明恒等映射同伦于零，如所需。
\end{proof}






\section{{\v C}ech 上同调作为预层上的函子}
\label{cohomology-section-cech-functor}

\noindent
警告：本节几乎全程处理预层与预层范畴；
在层与层范畴的情形，这些结论完全错误！

\medskip\noindent
设 $X$ 是环化空间，
$\mathcal{U} : U = \bigcup_{i \in I} U_i$ 是开覆盖，
$\mathcal{F}$ 是 $\mathcal{O}_X$-模预层。
只需将 $\mathcal{F}$ 看作阿贝尔群预层，便得到它的 {\v C}ech 复形
$\check{\mathcal{C}}^\bullet(\mathcal{U}, \mathcal{F})$
。不过，每一项
$\check{\mathcal{C}}^p(\mathcal{U}, \mathcal{F})$
都自然具有 $\mathcal{O}_X(U)$-模结构，且微分由
$\mathcal{O}_X(U)$-模映射给出。此外，构造
$$
\mathcal{F} \longmapsto \check{\mathcal{C}}^\bullet(\mathcal{U}, \mathcal{F})
$$
显然关于 $\mathcal{F}$ 是函子性的。事实上，它是函子
\begin{equation}
\label{cohomology-equation-cech-functor}
\check{\mathcal{C}}^\bullet(\mathcal{U}, -) :
\textit{PMod}(\mathcal{O}_X)
\longrightarrow
\text{Comp}^{+}(\text{Mod}_{\mathcal{O}_X(U)})
\end{equation}
；记号见
《导出范畴》定义 \ref{derived-definition-complexes-notation}。
回忆阿贝尔范畴中下有界复形所成的范畴仍是阿贝尔范畴，见
《同调代数》引理 \ref{homology-lemma-cat-cochain-abelian}。

\begin{lemma}
\label{cohomology-lemma-cech-exact-presheaves}
Equation (\ref{cohomology-equation-cech-functor}) 给出的函子是正合函子
（见 《同调代数》引理 \ref{homology-lemma-exact-functor}）。
\end{lemma}

\begin{proof}
对任意开集 $W \subset U$，函子
$\mathcal{F} \mapsto \mathcal{F}(W)$
是从 $\textit{PMod}(\mathcal{O}_X)$ 到
$\text{Mod}_{\mathcal{O}_X(U)}$ 的加性正合函子。
复形的各项
$\check{\mathcal{C}}^p(\mathcal{U}, \mathcal{F})$
是这些正合函子的积，因而正合。
此外，复形序列正合当且仅当它在每个给定次数的项所成序列正合。
故引理成立。
\end{proof}

\begin{lemma}
\label{cohomology-lemma-cech-cohomology-delta-functor-presheaves}
设 $X$ 是环化空间，
$\mathcal{U} : U = \bigcup_{i \in I} U_i$ 是开覆盖。
诸函子 $\mathcal{F} \mapsto \check{H}^n(\mathcal{U}, \mathcal{F})$
组成从 $\mathcal{O}_X$-模预层的阿贝尔范畴到
$\mathcal{O}_X(U)$-模范畴的 $\delta$-函子
（见 《同调代数》定义 \ref{homology-definition-cohomological-delta-functor}）。
\end{lemma}

\begin{proof}
由引理 \ref{cohomology-lemma-cech-exact-presheaves}，
$\mathcal{O}_X$-模预层的短正合序列
$0 \to \mathcal{F}_1 \to \mathcal{F}_2 \to \mathcal{F}_3 \to 0$
被送到 $\mathcal{O}_X(U)$-模复形的短正合序列。因此可用
《同调代数》引理 \ref{homology-lemma-long-exact-sequence-cochain}
得到连接映射
$\delta_{\mathcal{F}_1 \to \mathcal{F}_2 \to \mathcal{F}_3} :
\check{H}^n(\mathcal{U}, \mathcal{F}_3) \to
\check{H}^{n + 1}(\mathcal{U}, \mathcal{F}_1)$
及相应的长正合序列。
这些映射与预层短正合序列之间的映射相容，其验证从略。
\end{proof}


\noindent
在下面引理的表述中，我们使用相对于开浸入
$j : U \to X$ 的模预层延零函子 $j_{p!}$。
见 《层》第 \ref{sheaves-section-open-immersions} 节。
对任意开集 $W \subset X$ 及任意
$\mathcal{O}_X|_U$-模预层 $\mathcal{G}$，有
$$
(j_{p!}\mathcal{G})(W) =
\left\{
\begin{matrix}
\mathcal{G}(W) & \text{若 } W \subset U \\
0 & \text{else.}
\end{matrix}
\right.
$$
此外，函子 $j_{p!}$ 是限制函子的左伴随，见
《层》引理 \ref{sheaves-lemma-j-shriek-modules}。
特别地，有公式
$$
\Hom_{\mathcal{O}_X}(j_{p!}\mathcal{O}_U, \mathcal{F})
=
\Hom_{\mathcal{O}_U}(\mathcal{O}_U, \mathcal{F}|_U)
=
\mathcal{F}(U).
$$
由于函子 $\mathcal{F} \mapsto \mathcal{F}(U)$
在预层范畴上正合，预层 $j_{p!}\mathcal{O}_U$
是范畴 $\textit{PMod}(\mathcal{O}_X)$ 中的投射对象，见
《同调代数》引理 \ref{homology-lemma-characterize-projectives}。

\medskip\noindent
注意，若给定开子集 $U \subset V \subset X$
及相应的开浸入 $j_U,j_V$，便有典范映射
$(j_U)_{p!}\mathcal{O}_U \to (j_V)_{p!}\mathcal{O}_V$。
它在任意开集 $W \subset U$ 上的截面上为恒等映射，在其他情形为 $0$。
在恒等式
$\Hom_{\mathcal{O}_X}((j_U)_{p!}\mathcal{O}_U, (j_V)_{p!}\mathcal{O}_V) =
(j_V)_{p!}\mathcal{O}_V(U) = \mathcal{O}_V(U)$
之下，它对应于元素 $1 \in \mathcal{O}_V(U)$。

\begin{lemma}
\label{cohomology-lemma-cech-map-into}
设 $X$ 是环化空间，
$\mathcal{U} : U = \bigcup_{i \in I} U_i$ 是开覆盖。
以 $j_{i_0\ldots i_p} : U_{i_0 \ldots i_p} \to X$
表示开浸入。考虑 $\mathcal{O}_X$-模预层的链复形
$K(\mathcal{U})_\bullet$
$$
\ldots
\to
\bigoplus_{i_0i_1i_2} (j_{i_0i_1i_2})_{p!}\mathcal{O}_{U_{i_0i_1i_2}}
\to
\bigoplus_{i_0i_1} (j_{i_0i_1})_{p!}\mathcal{O}_{U_{i_0i_1}}
\to
\bigoplus_{i_0} (j_{i_0})_{p!}\mathcal{O}_{U_{i_0}}
\to 0 \to \ldots
$$
，其中最后一个非零项置于次数 $0$，而映射
$$
(j_{i_0\ldots i_{p + 1}})_{p!}\mathcal{O}_{U_{i_0\ldots i_{p + 1}}}
\longrightarrow
(j_{i_0\ldots \hat i_j \ldots i_{p + 1}})_{p!}
\mathcal{O}_{U_{i_0\ldots \hat i_j \ldots i_{p + 1}}}
$$
由典范映射的 $(-1)^j$ 倍给出。则存在同构
$$
\Hom_{\mathcal{O}_X}(K(\mathcal{U})_\bullet, \mathcal{F})
=
\check{\mathcal{C}}^\bullet(\mathcal{U}, \mathcal{F})
$$
且它关于
$\mathcal{F} \in \Ob(\textit{PMod}(\mathcal{O}_X))$ 是函子性的。
\end{lemma}

\begin{proof}
在引理之前的讨论中已经看到
$$
\Hom_{\mathcal{O}_X}(
(j_{i_0\ldots i_p})_{p!}\mathcal{O}_{U_{i_0\ldots i_p}},
\mathcal{F})
=
\mathcal{F}(U_{i_0\ldots i_p}).
$$
因此，直和
$$
\bigoplus\nolimits_{i_0 \ldots i_p}
(j_{i_0 \ldots i_p})_{p!}\mathcal{O}_{U_{i_0 \ldots i_p}}
$$
确实表示函子
$$
\mathcal{F}
\longmapsto
\prod\nolimits_{i_0\ldots i_p} \mathcal{F}(U_{i_0\ldots i_p}).
$$
由 《范畴》Yoneda 引理 \ref{categories-lemma-yoneda}，
可知存在具有上述各项的复形 $K(\mathcal{U})_\bullet$。
直接可见其映射正如引理所述。
\end{proof}

\begin{lemma}
\label{cohomology-lemma-homology-complex}
设 $X$ 是环化空间，
$\mathcal{U} : U = \bigcup_{i \in I} U_i$ 是开覆盖。
令 $\mathcal{O}_\mathcal{U} \subset \mathcal{O}_X$
为映射 $\bigoplus j_{p!}\mathcal{O}_{U_i} \to \mathcal{O}_X$
的像预层。上面引理 \ref{cohomology-lemma-cech-map-into}
中的预层链复形 $K(\mathcal{U})_\bullet$ 具有同调预层
$$
H_i(K(\mathcal{U})_\bullet) =
\left\{
\begin{matrix}
0 & \text{若} & i \not = 0 \\
\mathcal{O}_\mathcal{U} & \text{若} & i = 0
\end{matrix}
\right.
$$
\end{lemma}

\begin{proof}
考虑如下增广复形 $K^{ext}_\bullet$：将
$\mathcal{O}_\mathcal{U}$ 置于次数 $-1$，并取显然的映射
$K(\mathcal{U})_0 =
\bigoplus_{i_0} (j_{i_0})_{p!}\mathcal{O}_{U_{i_0}} \to
\mathcal{O}_\mathcal{U}$.
只需证明，对任意开集 $W \subset X$，取该增广复形在 $W$ 上的截面
所得复形无环。事实上，我们断言对每个 $W \subset X$，
复形 $K^{ext}_\bullet(W)$ 链同伦等价于零复形。
写作 $I = I_1 \amalg I_2$，其中
$W \subset U_i$ 当且仅当 $i \in I_1$。

\medskip\noindent
若 $I_1 = \emptyset$，则复形 $K^{ext}_\bullet(W) = 0$，无需证明。

\medskip\noindent
若 $I_1 \not = \emptyset$，则
$\mathcal{O}_\mathcal{U}(W) = \mathcal{O}_X(W)$，并且
$$
K^{ext}_p(W) =
\bigoplus\nolimits_{i_0 \ldots i_p \in I_1} \mathcal{O}_X(W).
$$
这是因为预层
$(j_{i_0 \ldots i_p})_{p!}\mathcal{O}_{U_{i_0 \ldots i_p}}$
有上述简单描述。此外，复形 $K^{ext}_\bullet(W)$ 的微分为
$$
d(s)_{i_0 \ldots i_p} =
\sum\nolimits_{j = 0, \ldots, p + 1} \sum\nolimits_{i \in I_1}
(-1)^j s_{i_0 \ldots i_{j - 1} i i_j \ldots i_p}.
$$
由于元素 $s$ 具有有限支撑，该和是有限的。
固定元素 $i_{\text{fix}} \in I_1$。定义映射
$$
h : K^{ext}_p(W) \longrightarrow K^{ext}_{p + 1}(W)
$$
如下：
$$
h(s)_{i_0 \ldots i_{p + 1}} =
\left\{
\begin{matrix}
0 & \text{若} & i_0 \not = i_{\text{fix}} \\
s_{i_1 \ldots i_{p + 1}} & \text{若} & i_0 = i_{\text{fix}}
\end{matrix}
\right.
$$
我们将使用简记
$h(s)_{i_0 \ldots i_{p + 1}} = (i_0 = i_{\text{fix}}) s_{i_1 \ldots i_p}$
。于是计算得
\begin{eqnarray*}
& & (dh + hd)(s)_{i_0 \ldots i_p} \\
& = &
\sum_j \sum_{i \in I_1} (-1)^j h(s)_{i_0 \ldots i_{j - 1} i i_j \ldots i_p}
+
(i = i_0) d(s)_{i_1 \ldots i_p} \\
& = &
s_{i_0 \ldots i_p} +
\sum_{j \geq 1}\sum_{i \in I_1}
(-1)^j (i_0 = i_{\text{fix}}) s_{i_1 \ldots i_{j - 1} i i_j \ldots i_p}
+
(i_0 = i_{\text{fix}}) d(s)_{i_1 \ldots i_p}
\end{eqnarray*}
它等于 $s_{i_0 \ldots i_p}$，如所需。
\end{proof}

\begin{lemma}
\label{cohomology-lemma-cech-cohomology-derived-presheaves}
设 $X$ 是环化空间，
$\mathcal{U} : U = \bigcup_{i \in I} U_i$
是 $U \subset X$ 的开覆盖。
{\v C}ech 上同调函子 $\check{H}^p(\mathcal{U}, -)$
作为 $\delta$-函子，典范同构于函子
$$
\check{H}^0(\mathcal{U}, -) :
\textit{PMod}(\mathcal{O}_X)
\longrightarrow
\text{Mod}_{\mathcal{O}_X(U)}.
$$
的右导出函子。此外，存在函子性拟同构
$$
\check{\mathcal{C}}^\bullet(\mathcal{U}, \mathcal{F})
\longrightarrow
R\check{H}^0(\mathcal{U}, \mathcal{F})
$$
，其中右端表示左正合函子
$$
R\check{H}^0(\mathcal{U}, -) :
D^{+}(\textit{PMod}(\mathcal{O}_X))
\longrightarrow
D^{+}(\mathcal{O}_X(U))
$$
$\check{H}^0(\mathcal{U}, -)$ 的右导出函子。
\end{lemma}

\begin{proof}
注意，$\mathcal{O}_X$-模预层的范畴有足够多的内射对象，见
《内射对象》Proposition \ref{injectives-proposition-presheaves-modules}。
又注意到 $\check{H}^0(\mathcal{U}, -)$
是从 $\mathcal{O}_X$-模预层范畴到
$\mathcal{O}_X(U)$-模范畴的左正合函子。
因此导出函子和右导出函子存在，见
《导出范畴》第 \ref{derived-section-right-derived-functor} 节。

\medskip\noindent
设 $\mathcal{I}$ 是内射 $\mathcal{O}_X$-模预层。
此时函子 $\Hom_{\mathcal{O}_X}(-, \mathcal{I})$
在 $\textit{PMod}(\mathcal{O}_X)$ 上正合。由引理 \ref{cohomology-lemma-cech-map-into}，有
$$
\Hom_{\mathcal{O}_X}(K(\mathcal{U})_\bullet, \mathcal{I})
=
\check{\mathcal{C}}^\bullet(\mathcal{U}, \mathcal{I}).
$$
由引理 \ref{cohomology-lemma-homology-complex}，
$K(\mathcal{U})_\bullet$ 拟同构于
$\mathcal{O}_\mathcal{U}[0]$。因此，由上述取到
$\mathcal{I}$ 的 Hom 函子的正合性可知，对所有 $i > 0$，
$\check{H}^i(\mathcal{U}, \mathcal{I}) = 0$。
故 $\delta$-函子 $(\check{H}^n, \delta)$
（见引理 \ref{cohomology-lemma-cech-cohomology-delta-functor-presheaves}）
满足
《同调代数》引理 \ref{homology-lemma-efface-implies-universal},
的假设，因而是泛 $\delta$-函子。

\medskip\noindent
由 《导出范畴》引理 \ref{derived-lemma-higher-derived-functors}，
序列 $R^i\check{H}^0(\mathcal{U}, -)$ 也组成泛 $\delta$-函子。
由泛 $\delta$-函子的唯一性，见
《同调代数》引理 \ref{homology-lemma-uniqueness-universal-delta-functor}
，得到
$R^i\check{H}^0(\mathcal{U}, -) = \check{H}^i(\mathcal{U}, -)$.
这对大多数应用已经足够，建议读者略过证明的其余部分。

\medskip\noindent
设 $\mathcal{F}$ 是任意 $\mathcal{O}_X$-模预层。
在范畴 $\textit{PMod}(\mathcal{O}_X)$ 中取内射消解
$\mathcal{F} \to \mathcal{I}^\bullet$。
考虑各项为
$\check{\mathcal{C}}^p(\mathcal{U}, \mathcal{I}^q)$
的双复形
$\check{\mathcal{C}}^\bullet(\mathcal{U}, \mathcal{I}^\bullet)$，
以及相伴的全复形
$\text{Tot}(\check{\mathcal{C}}^\bullet(\mathcal{U}, \mathcal{I}^\bullet))$,
见 《同调代数》定义 \ref{homology-definition-associated-simple-complex}。
由诸映射
$\check{\mathcal{C}}^p(\mathcal{U}, \mathcal{F})
\to \check{\mathcal{C}}^p(\mathcal{U}, \mathcal{I}^0)$
得到复形映射
$$
\check{\mathcal{C}}^\bullet(\mathcal{U}, \mathcal{F})
\longrightarrow
\text{Tot}(\check{\mathcal{C}}^\bullet(\mathcal{U}, \mathcal{I}^\bullet))
$$
，并且由诸映射
$\check{H}^0(\mathcal{U}, \mathcal{I}^q) \to
\check{\mathcal{C}}^0(\mathcal{U}, \mathcal{I}^q)$
得到复形映射
$$
\check{H}^0(\mathcal{U}, \mathcal{I}^\bullet)
\longrightarrow
\text{Tot}(\check{\mathcal{C}}^\bullet(\mathcal{U}, \mathcal{I}^\bullet))
$$
。施用 《同调代数》引理 \ref{homology-lemma-double-complex-gives-resolution}
可知，这两个映射都是拟同构。
事实上，由于作为函子的 {\v C}ech 复形是正合的
（引理 \ref{cohomology-lemma-cech-exact-presheaves}），
该双复形的各列在正次数处正合；
又由于刚才看到内射预层 $\mathcal{I}^q$
的高阶 {\v C}ech 上同调群为零，其各行在正次数处也正合。
拟同构在 $D^{+}(\mathcal{O}_X(U))$ 中可逆，
故得到引理中最后一个显示态射。
该态射的函子性验证从略。
\end{proof}





\section{{\v C}ech 上同调与层上同调}
\label{cohomology-section-cech-cohomology-cohomology}

\begin{lemma}
\label{cohomology-lemma-injective-trivial-cech}
设 $X$ 是环化空间，
$\mathcal{U} : U = \bigcup_{i \in I} U_i$ 是开覆盖，
$\mathcal{I}$ 是内射 $\mathcal{O}_X$-模。则
$$
\check{H}^p(\mathcal{U}, \mathcal{I}) =
\left\{
\begin{matrix}
\mathcal{I}(U) & \text{若} & p = 0 \\
0 & \text{若} & p > 0
\end{matrix}
\right.
$$
\end{lemma}

\begin{proof}
内射 $\mathcal{O}_X$-模作为范畴
$\textit{PMod}(\mathcal{O}_X)$ 中的对象仍是内射的
（例如，这是因为层化函子是包含函子的正合左伴随；使用
《同调代数》引理 \ref{homology-lemma-adjoint-preserve-injectives}）。
因此施用引理 \ref{cohomology-lemma-cech-cohomology-derived-presheaves}
（或其证明）即得结论。
\end{proof}

\begin{lemma}
\label{cohomology-lemma-cech-cohomology}
设 $X$ 是环化空间，
$\mathcal{U} : U = \bigcup_{i \in I} U_i$ 是开覆盖。
存在函子
$\textit{Mod}(\mathcal{O}_X) \to D^{+}(\mathcal{O}_X(U))$
之间的变换
$$
\check{\mathcal{C}}^\bullet(\mathcal{U}, -)
\longrightarrow
R\Gamma(U, -)
$$
。特别地，对遍历
$\textit{Mod}(\mathcal{O}_X)$ 的 $\mathcal{F}$，
它给出典范映射
$\check{H}^p(\mathcal{U}, \mathcal{F}) \to H^p(U, \mathcal{F})$。
\end{lemma}

\begin{proof}
设 $\mathcal{F}$ 是 $\mathcal{O}_X$-模，并取内射消解
$\mathcal{F} \to \mathcal{I}^\bullet$。
考虑各项为
$\check{\mathcal{C}}^p(\mathcal{U}, \mathcal{I}^q)$
的双复形
$\check{\mathcal{C}}^\bullet(\mathcal{U}, \mathcal{I}^\bullet)$。
由诸映射
$\mathcal{I}^q(U) \to \check{H}^0(\mathcal{U}, \mathcal{I}^q)$
得到复形映射
$$
\alpha :
\Gamma(U, \mathcal{I}^\bullet)
\longrightarrow
\text{Tot}(\check{\mathcal{C}}^\bullet(\mathcal{U}, \mathcal{I}^\bullet))
$$
，并由映射 $\mathcal{F} \to \mathcal{I}^0$
得到复形映射
$$
\beta :
\check{\mathcal{C}}^\bullet(\mathcal{U}, \mathcal{F})
\longrightarrow
\text{Tot}(\check{\mathcal{C}}^\bullet(\mathcal{U}, \mathcal{I}^\bullet))
$$
。施用
《同调代数》引理 \ref{homology-lemma-double-complex-gives-resolution}
可知 $\alpha$ 是拟同构。
事实上，引理 \ref{cohomology-lemma-injective-trivial-cech} 说明，
双复形
$\check{\mathcal{C}}^\bullet(\mathcal{U}, \mathcal{I}^\bullet)$
的第 $q$ 行是 $\Gamma(U, \mathcal{I}^q)$ 的消解。
故 $\alpha$ 在 $D^{+}(\mathcal{O}_X(U))$ 中可逆，
而引理中的变换是先施用 $\beta$ 再施用 $\alpha$ 的逆所得的复合。
其函子性验证从略。
\end{proof}

\begin{lemma}
\label{cohomology-lemma-cech-h1}
设 $X$ 是拓扑空间，$\mathcal{H}$ 是 $X$ 上的阿贝尔层，
$\mathcal{U} : X = \bigcup_{i \in I} U_i$ 是开覆盖。映射
$$
\check{H}^1(\mathcal{U}, \mathcal{H}) \longrightarrow H^1(X, \mathcal{H})
$$
是单射，并通过引理 \ref{cohomology-lemma-torsors-h1} 的双射，
将 $\check{H}^1(\mathcal{U}, \mathcal{H})$
等同于如下 $\mathcal{H}$-挠子的同构类集合：
这些挠子限制到每个 $U_i$ 上均为平凡挠子。
\end{lemma}

\begin{proof}
为此我们构造逆映射。设 $\mathcal{F}$ 是一个
$\mathcal{H}$-挠子，其在 $U_i$ 上的限制平凡。
由引理 \ref{cohomology-lemma-trivial-torsor}，这意味着存在截面
$s_i \in \mathcal{F}(U_i)$。在 $U_{i_0} \cap U_{i_1}$ 上，
存在唯一的 $\mathcal{H}$ 的截面 $s_{i_0i_1}$，使得
$s_{i_0i_1} \cdot s_{i_0}|_{U_{i_0} \cap U_{i_1}} =
s_{i_1}|_{U_{i_0} \cap U_{i_1}}$。
直接计算可见，$s_{i_0i_1}$ 是 {\v C}ech 上循环，
且其类是良定义的（即不依赖于截面 $s_i$ 的选取）。
逆映射把 $\mathcal{F}$ 的同构类送到上循环
$(s_{i_0i_1})$ 的上同调类。
该映射确为逆映射的验证从略。
\end{proof}

\begin{lemma}
\label{cohomology-lemma-include}
设 $X$ 是环化空间。考虑函子
$i : \textit{Mod}(\mathcal{O}_X) \to \textit{PMod}(\mathcal{O}_X)$.
这是左正合函子，其右导出函子为
$$
R^pi(\mathcal{F}) = \underline{H}^p(\mathcal{F}) :
U \longmapsto H^p(U, \mathcal{F})
$$
，见 第 \ref{cohomology-section-locality} 节 中的讨论。
\end{lemma}

\begin{proof}
显然 $i$ 左正合。取内射消解
$\mathcal{F} \to \mathcal{I}^\bullet$。
按定义，$R^pi$ 是复形 $\mathcal{I}^\bullet$ 的第 $p$ 个
上同调{\it 预层}。换言之，
$R^pi(\mathcal{F})$ 在开集 $U$ 上的截面为
$$
\frac{\Ker(\mathcal{I}^p(U) \to \mathcal{I}^{p + 1}(U))}
{\Im(\mathcal{I}^{p - 1}(U) \to \mathcal{I}^p(U))}.
$$
这正是 $H^p(U, \mathcal{F})$ 的定义。
\end{proof}

\begin{lemma}
\label{cohomology-lemma-cech-spectral-sequence}
设 $X$ 是环化空间，
$\mathcal{U} : U = \bigcup_{i \in I} U_i$ 是开覆盖。
对任意 $\mathcal{O}_X$-模层 $\mathcal{F}$，
存在谱序列 $(E_r, d_r)_{r \geq 0}$，满足
$$
E_2^{p, q} = \check{H}^p(\mathcal{U}, \underline{H}^q(\mathcal{F}))
$$
并收敛到 $H^{p + q}(U, \mathcal{F})$。
该谱序列关于 $\mathcal{F}$ 是函子性的。
\end{lemma}

\begin{proof}
这是函子
$$
i :  \textit{Mod}(\mathcal{O}_X) \to \textit{PMod}(\mathcal{O}_X)
\quad\text{且}\quad
\check{H}^0(\mathcal{U}, - ) : \textit{PMod}(\mathcal{O}_X)
\to \text{Mod}_{\mathcal{O}_X(U)}.
$$
的 Grothendieck 谱序列（见
《导出范畴》引理 \ref{derived-lemma-grothendieck-spectral-sequence})
）。事实上，由引理 \ref{cohomology-lemma-cech-h0}，
$\check{H}^0(\mathcal{U}, i(\mathcal{F})) = \mathcal{F}(U)$。
由引理 \ref{cohomology-lemma-injective-trivial-cech}，
$i(\mathcal{I})$ 是 {\v C}ech 无环的。
又由引理 \ref{cohomology-lemma-cech-cohomology-derived-presheaves}，
在 $\textit{PMod}(\mathcal{O}_X)$ 上有函子等式
$\check{H}^p(\mathcal{U}, -) = R^p\check{H}^0(\mathcal{U}, -)$。
综合这些事实即得引理。
\end{proof}

\begin{lemma}
\label{cohomology-lemma-cech-spectral-sequence-application}
设 $X$ 是环化空间，
$\mathcal{U} : U = \bigcup_{i \in I} U_i$ 是开覆盖，
$\mathcal{F}$ 是 $\mathcal{O}_X$-模。
假设对所有 $i > 0$、所有 $p \geq 0$
以及所有 $i_0, \ldots, i_p \in I$，均有
$H^i(U_{i_0 \ldots i_p}, \mathcal{F}) = 0$。
则作为 $\mathcal{O}_X(U)$-模，有
$\check{H}^p(\mathcal{U}, \mathcal{F}) = H^p(U, \mathcal{F})$。
\end{lemma}

\begin{proof}
使用引理 \ref{cohomology-lemma-cech-spectral-sequence} 的谱序列。
这些假设意味着，对所有满足 $q \not = 0$ 的 $(p,q)$，
均有 $E_2^{p, q} = 0$。因此谱序列在 $E_2$ 页退化，结论随之成立。
\end{proof}

\begin{lemma}
\label{cohomology-lemma-ses-cech-h1}
设 $X$ 是环化空间。设
$$
0 \to \mathcal{F} \to \mathcal{G} \to \mathcal{H} \to 0
$$
是 $\mathcal{O}_X$-模的短正合序列，
$U \subset X$ 是开子集。
若存在 $U$ 的开覆盖 $\mathcal{U}$ 所成的共尾系，
使得 $\check{H}^1(\mathcal{U}, \mathcal{F}) = 0$，
则映射 $\mathcal{G}(U) \to \mathcal{H}(U)$ 是满射。
\end{lemma}

\begin{proof}
取元素 $s \in \mathcal{H}(U)$。选择开覆盖
$\mathcal{U} : U = \bigcup_{i \in I} U_i$，使得
(a) $\check{H}^1(\mathcal{U}, \mathcal{F}) = 0$，且
(b) $s|_{U_i}$ 是某截面 $s_i \in \mathcal{G}(U_i)$ 的像。
显然能找到满足 (b) 的 $\mathcal{U}$；
由引理的假设，可找到同时满足 (a) 与 (b) 的 $\mathcal{U}$。
考虑截面
$$
s_{i_0i_1} = s_{i_1}|_{U_{i_0i_1}} - s_{i_0}|_{U_{i_0i_1}}.
$$
由于 $s_i$ 提升 $s$，可知
$s_{i_0i_1} \in \mathcal{F}(U_{i_0i_1})$。
由 $\check{H}^1(\mathcal{U}, \mathcal{F})$ 的消失性，
可找到截面 $t_i \in \mathcal{F}(U_i)$，使得
$$
s_{i_0i_1} = t_{i_1}|_{U_{i_0i_1}} - t_{i_0}|_{U_{i_0i_1}}.
$$
于是截面 $s_i-t_i$ 显然满足层条件，因而黏合成
$\mathcal{G}$ 在 $U$ 上的截面，且它映到 $s$。结论成立。
\end{proof}

\begin{lemma}
\label{cohomology-lemma-cech-vanish}
\begin{slogan}
若阿贝尔层对所有开覆盖的高阶 {\v C}ech 上同调均消失，
则其高阶层上同调消失。
\end{slogan}
设 $X$ 是环化空间，$\mathcal{F}$ 是 $\mathcal{O}_X$-模，且
$$
\check{H}^p(\mathcal{U}, \mathcal{F}) = 0
$$
对所有 $p > 0$ 以及 $X$ 中任意开集的任意开覆盖
$\mathcal{U} : U = \bigcup_{i \in I} U_i$ 均成立。
则对所有 $p > 0$ 及任意开集 $U \subset X$，有
$H^p(U, \mathcal{F}) = 0$。
\end{lemma}

\begin{proof}
设 $\mathcal{F}$ 是满足引理假设的层。
我们称之为“$\mathcal{F}$ 对任意开覆盖的高阶
{\v C}ech 上同调消失”。
选取到内射 $\mathcal{O}_X$-模的嵌入
$\mathcal{F} \to \mathcal{I}$。
由引理 \ref{cohomology-lemma-injective-trivial-cech}，
$\mathcal{I}$ 对任意开覆盖的高阶 {\v C}ech 上同调消失。
令 $\mathcal{Q} = \mathcal{I}/\mathcal{F}$，于是有短正合序列
$$
0 \to \mathcal{F} \to \mathcal{I} \to \mathcal{Q} \to 0.
$$
由引理 \ref{cohomology-lemma-ses-cech-h1} 及我们的假设，
该序列事实上作为预层序列也是正合的！
特别地，对任意开覆盖 $\mathcal{U}$，
都有 {\v C}ech 上同调群的长正合序列，例如见引理 \ref{cohomology-lemma-cech-cohomology-delta-functor-presheaves}
。这说明 $\mathcal{Q}$ 也是对所有开覆盖具有消失高阶
{\v C}ech 上同调的 $\mathcal{O}_X$-模。

\medskip\noindent
接着，对任意开集 $U \subset X$，考虑上同调长正合序列
$$
\xymatrix{
0 \ar[r] &
H^0(U, \mathcal{F}) \ar[r] &
H^0(U, \mathcal{I}) \ar[r] &
H^0(U, \mathcal{Q}) \ar[lld] \\
&
H^1(U, \mathcal{F}) \ar[r] &
H^1(U, \mathcal{I}) \ar[r] &
H^1(U, \mathcal{Q}) \ar[lld] \\
&
\ldots & \ldots & \ldots \\
}
$$
。由于 $\mathcal{I}$ 内射，对 $n > 0$ 有
$H^n(U, \mathcal{I}) = 0$（见
《导出范畴》引理 \ref{derived-lemma-higher-derived-functors}）。
由上可知，$H^0(U, \mathcal{I}) \to H^0(U, \mathcal{Q})$
是满射，因而 $H^1(U, \mathcal{F}) = 0$。
由于 $\mathcal{F}$ 是具有消失高阶 {\v C}ech 上同调的任意
$\mathcal{O}_X$-模，而 $\mathcal{Q}$ 也是这类层（见上），
同样得到 $H^1(U, \mathcal{Q}) = 0$。
长正合序列进而给出 $H^2(U, \mathcal{F}) = 0$，如此继续即可。
\end{proof}

\begin{lemma}
\label{cohomology-lemma-cech-vanish-basis}
（引理 \ref{cohomology-lemma-cech-vanish} 的变体。）
设 $X$ 是环化空间，$\mathcal{B}$ 是 $X$ 的拓扑基，
$\mathcal{F}$ 是 $\mathcal{O}_X$-模。
假设存在开覆盖的集合 $\text{Cov}$，满足下列性质：
\begin{enumerate}
\item 对每个 $\mathcal{U} \in \text{Cov}$，若
$\mathcal{U} : U = \bigcup_{i \in I} U_i$，则
$U,U_i \in \mathcal{B}$，且每个
$U_{i_0 \ldots i_p} \in \mathcal{B}$。
\item 对每个 $U \in \mathcal{B}$，$\text{Cov}$ 中出现的
$U$ 的开覆盖构成 $U$ 的开覆盖的共尾系。
\item 对每个 $\mathcal{U} \in \text{Cov}$ 及所有 $p > 0$，有
$\check{H}^p(\mathcal{U}, \mathcal{F}) = 0$。
\end{enumerate}
则对所有 $p > 0$ 及任意 $U \in \mathcal{B}$，有
$H^p(U, \mathcal{F}) = 0$。
\end{lemma}

\begin{proof}
设 $\mathcal{F}$ 与 $\text{Cov}$ 如引理所述。
我们称之为“$\mathcal{F}$ 对任意
$\mathcal{U} \in \text{Cov}$ 的高阶 {\v C}ech 上同调消失”。
选取到内射 $\mathcal{O}_X$-模的嵌入
$\mathcal{F} \to \mathcal{I}$。
由引理 \ref{cohomology-lemma-injective-trivial-cech}，
$\mathcal{I}$ 对任意 $\mathcal{U} \in \text{Cov}$
的高阶 {\v C}ech 上同调消失。
令 $\mathcal{Q} = \mathcal{I}/\mathcal{F}$，于是有短正合序列
$$
0 \to \mathcal{F} \to \mathcal{I} \to \mathcal{Q} \to 0.
$$
由引理 \ref{cohomology-lemma-ses-cech-h1} 及假设 (2)，
对每个 $U \in \mathcal{B}$，该序列给出正合序列
$$
0 \to \mathcal{F}(U) \to \mathcal{I}(U) \to \mathcal{Q}(U) \to 0.
$$
因此，对任意 $\mathcal{U} \in \text{Cov}$，
得到 {\v C}ech 复形的短正合序列
$$
0 \to
\check{\mathcal{C}}^\bullet(\mathcal{U}, \mathcal{F}) \to
\check{\mathcal{C}}^\bullet(\mathcal{U}, \mathcal{I}) \to
\check{\mathcal{C}}^\bullet(\mathcal{U}, \mathcal{Q}) \to 0
$$
，因为由假设 (1)，{\v C}ech 复形的每一项都是
$\mathcal{B}$ 中元素上取值的积。
特别地，对任意开覆盖 $\mathcal{U} \in \text{Cov}$，
都有 {\v C}ech 上同调群的长正合序列。
这说明 $\mathcal{Q}$ 也是对所有
$\mathcal{U} \in \text{Cov}$ 具有消失高阶 {\v C}ech 上同调的
$\mathcal{O}_X$-模。

\medskip\noindent
接着，对任意 $U \in \mathcal{B}$，考虑上同调长正合序列
$$
\xymatrix{
0 \ar[r] &
H^0(U, \mathcal{F}) \ar[r] &
H^0(U, \mathcal{I}) \ar[r] &
H^0(U, \mathcal{Q}) \ar[lld] \\
&
H^1(U, \mathcal{F}) \ar[r] &
H^1(U, \mathcal{I}) \ar[r] &
H^1(U, \mathcal{Q}) \ar[lld] \\
&
\ldots & \ldots & \ldots \\
}
$$
。由于 $\mathcal{I}$ 内射，对 $n > 0$ 有
$H^n(U, \mathcal{I}) = 0$（见
《导出范畴》引理 \ref{derived-lemma-higher-derived-functors}）。
由上可知，$H^0(U, \mathcal{I}) \to H^0(U, \mathcal{Q})$
是满射，因而 $H^1(U, \mathcal{F}) = 0$。
由于 $\mathcal{F}$ 是对所有
$\mathcal{U} \in \text{Cov}$ 具有消失高阶 {\v C}ech 上同调的任意
$\mathcal{O}_X$-模，而 $\mathcal{Q}$ 也是这类层（见上），
同样得到 $H^1(U, \mathcal{Q}) = 0$。
长正合序列进而给出 $H^2(U, \mathcal{F}) = 0$，如此继续即可。
\end{proof}

\begin{lemma}
\label{cohomology-lemma-pushforward-injective}
设 $f : X \to Y$ 是环化空间的态射，
$\mathcal{I}$ 是内射 $\mathcal{O}_X$-模。则
\begin{enumerate}
\item $\check{H}^p(\mathcal{V}, f_*\mathcal{I}) = 0$
对所有 $p > 0$ 及 $Y$ 的任意开覆盖
$\mathcal{V} : V = \bigcup_{j \in J} V_j$ 成立；
\item 对所有 $p > 0$ 及每个开集 $V \subset Y$，
有 $H^p(V, f_*\mathcal{I}) = 0$。
\end{enumerate}
换言之，对任意开集 $V \subset Y$，
$f_*\mathcal{I}$ 关于 $\Gamma(V, -)$ 是右无环的
（见 《导出范畴》定义 \ref{derived-definition-derived-functor}）。
\end{lemma}

\begin{proof}
置
$\mathcal{U} : f^{-1}(V) = \bigcup_{j \in J} f^{-1}(V_j)$。
这是 $X$ 的开覆盖，并且
$$
\check{\mathcal{C}}^\bullet(\mathcal{V}, f_*\mathcal{I}) =
\check{\mathcal{C}}^\bullet(\mathcal{U}, \mathcal{I}).
$$
这是因为
$$
f_*\mathcal{I}(V_{j_0 \ldots j_p})
= \mathcal{I}(f^{-1}(V_{j_0 \ldots j_p})) =
\mathcal{I}(f^{-1}(V_{j_0}) \cap \ldots \cap f^{-1}(V_{j_p}))
= \mathcal{I}(U_{j_0 \ldots j_p}).
$$
因此引理的第一条由引理 \ref{cohomology-lemma-injective-trivial-cech} 推出；
第二条由第一条及引理 \ref{cohomology-lemma-cech-vanish} 推出。
\end{proof}

\noindent
下面的引理特别说明，
$f_* : \textit{Ab}(X) \to \textit{Ab}(Y)$
把内射阿贝尔层送到内射阿贝尔层。

\begin{lemma}
\label{cohomology-lemma-pushforward-injective-flat}
设 $f : X \to Y$ 是环化空间的态射，并假设 $f$ 平坦。
则对任意内射 $\mathcal{O}_X$-模 $\mathcal{I}$，
$f_*\mathcal{I}$ 是内射 $\mathcal{O}_Y$-模。
\end{lemma}

\begin{proof}
此时函子 $f^*$ 把单射送到单射
（《模层》引理 \ref{modules-lemma-pullback-flat}）。
故结论由
《同调代数》引理 \ref{homology-lemma-adjoint-preserve-injectives}
推出。
\end{proof}

\begin{lemma}
\label{cohomology-lemma-cohomology-products}
设 $(X, \mathcal{O}_X)$ 是环化空间，$I$ 是集合。
对 $i \in I$，设 $\mathcal{F}_i$ 是 $\mathcal{O}_X$-模，
并设 $U \subset X$ 是开集。典范映射
$$
H^p(U, \prod\nolimits_{i \in I} \mathcal{F}_i)
\longrightarrow
\prod\nolimits_{i \in I} H^p(U, \mathcal{F}_i)
$$
当 $p = 0$ 时是同构，当 $p = 1$ 时是单射。
\end{lemma}

\begin{proof}
当 $p = 0$ 时的陈述成立，因为层的积等于其底层预层的积，见
《层》第 \ref{sheaves-section-limits-sheaves} 节。
下面证明 $p = 1$ 的情形。置
$\mathcal{F} = \prod \mathcal{F}_i$。
设 $\xi \in H^1(U, \mathcal{F})$ 在
$\prod H^1(U, \mathcal{F}_i)$ 中的像为零。
由上同调的局部性，见引理 \ref{cohomology-lemma-kill-cohomology-class-on-covering}，
存在开覆盖 $\mathcal{U} : U = \bigcup U_j$，
使得对所有 $j$ 有 $\xi|_{U_j} = 0$。
由引理 \ref{cohomology-lemma-cech-h1}，这意味着 $\xi$ 来自某元素
$\check \xi \in \check H^1(\mathcal{U}, \mathcal{F})$。
由于对所有 $i$，映射
$\check H^1(\mathcal{U}, \mathcal{F}_i) \to H^1(U, \mathcal{F}_i)$
均为单射（引理 \ref{cohomology-lemma-cech-h1}），
而 $\xi$ 在 $\prod H^1(U, \mathcal{F}_i)$ 中的像为零，
故其像 $\check \xi_i$ 在
$\check H^1(\mathcal{U}, \mathcal{F}_i)$ 中为零。
另一方面，由于 $\mathcal{F} = \prod \mathcal{F}_i$，
$\check{\mathcal{C}}^\bullet(\mathcal{U}, \mathcal{F})$
是诸复形
$\check{\mathcal{C}}^\bullet(\mathcal{U}, \mathcal{F}_i)$
的积。因此由
《同调代数》引理 \ref{homology-lemma-product-abelian-groups-exact}，
得到 $\check \xi = 0$，如所需。
\end{proof}







\section{松弛层}
\label{cohomology-section-flasque}

\noindent
定义如下。

\begin{definition}
\label{cohomology-definition-flasque}
设 $X$ 是拓扑空间。若对 $X$ 中任意开集
$U \subset V$，限制映射
$\mathcal{F}(V) \to \mathcal{F}(U)$ 都是满射，
则称集合预层 $\mathcal{F}$ 是{\it 松弛的}（亦称 flabby）。
\end{definition}

\noindent
对于阿贝尔层，以及当 $X$ 为环化空间时的模层，
我们也使用这一术语。
显然，只需假设对每个开集 $U \subset X$，
限制映射 $\mathcal{F}(X) \to \mathcal{F}(U)$ 是满射。

\begin{lemma}
\label{cohomology-lemma-injective-flasque}
设 $(X, \mathcal{O}_X)$ 是环化空间。
则任意内射 $\mathcal{O}_X$-模都是松弛的。
\end{lemma}

\begin{proof}
这是引理 \ref{cohomology-lemma-injective-restriction-surjective} 的改述。
\end{proof}

\begin{lemma}
\label{cohomology-lemma-flasque-acyclic}
设 $(X, \mathcal{O}_X)$ 是环化空间。
任意松弛 $\mathcal{O}_X$-模关于 $R\Gamma(X,-)$ 无环，
并且对 $X$ 的任意开集 $U$，关于 $R\Gamma(U,-)$ 也无环。
\end{lemma}

\begin{proof}
我们使用
《导出范畴》引理 \ref{derived-lemma-subcategory-right-acyclics}.
来证明。由于每个内射模都是松弛的，任意 $\mathcal{O}_X$-模
都可嵌入一个松弛模，见
《内射对象》引理 \ref{injectives-lemma-abelian-sheaves-space}。
因此只需证明：给定短正合序列
$$
0 \to \mathcal{F} \to \mathcal{G} \to \mathcal{H} \to 0
$$
，其中 $\mathcal{F}$、$\mathcal{G}$ 松弛，则 $\mathcal{H}$
松弛，且在 $X$ 的任意开集上取截面后，该序列仍短正合。
事实上，后一陈述蕴含前一陈述。
设 $U \subset X$ 是开子空间，$s \in \mathcal{H}(U)$。
我们证明 $s$ 可提升为 $\mathcal{G}$ 在 $U$ 上的截面。
考虑所有二元组 $(V,t)$ 所成的集合 $T$，其中
$V \subset U$ 开，且 $t \in \mathcal{G}(V)$ 是在
$\mathcal{H}$ 中映到 $s|_V$ 的截面。
规定
$(V,t) \leq (V',t')$ 当且仅当 $V \subset V'$ 且 $t'|_V=t$，
从而在 $T$ 上定义偏序。
若 $(V_\alpha,t_\alpha)$（$\alpha\in A$）是 $T$ 的全序子集，
则 $V=\bigcup V_\alpha$ 开；由 $\mathcal{G}$ 的层条件，
存在唯一的 $t\in\mathcal{G}(V)$，其在 $V_\alpha$ 上限制为
$t_\alpha$。故由 Zorn 引理，$T$ 中存在极大元 $(V,t)$。
我们证明 $V=U$，从而结束证明。
任取 $x\in U$。可找到 $x$ 的小开邻域 $W\subset U$
以及 $t'\in\mathcal{G}(W)$，使其在 $\mathcal{H}$ 中映到 $s|_W$。
于是 $t'|_{W\cap V}-t|_{W\cap V}$ 在 $\mathcal{H}$ 中映到零，
故来自某个截面 $r'\in\mathcal{F}(W\cap V)$。
由于 $\mathcal{F}$ 松弛，可找到截面
$r\in\mathcal{F}(W)$，其在 $W\cap V$ 上限制为 $r'$。
用 $r$ 的像修正 $t'$ 后，可假设 $t$ 与 $t'$
在 $W\cap V$ 上限制为同一截面。
由 $\mathcal{G}$ 的层条件，可找到 $\mathcal{G}$
在 $W\cup V$ 上的截面 $\tilde t$，其限制分别为 $t$ 和 $t'$。
由 $(V,t)$ 的极大性，$V\cup W=V$。故 $x\in V$，结论成立。
\end{proof}

\noindent
下面的引理对松弛预层并不成立。

\begin{lemma}
\label{cohomology-lemma-flasque-acyclic-cech}
设 $(X, \mathcal{O}_X)$ 是环化空间，
$\mathcal{F}$ 是 $\mathcal{O}_X$-模层，
$\mathcal{U}:U=\bigcup U_i$ 是开覆盖。
若 $\mathcal{F}$ 松弛，则对 $p>0$，
$\check{H}^p(\mathcal{U},\mathcal{F})=0$。
\end{lemma}

\begin{proof}
引理 \ref{cohomology-lemma-cech-spectral-sequence} 的陈述中所用的预层
$\underline{H}^q(\mathcal{F})$ 由引理 \ref{cohomology-lemma-flasque-acyclic} 可知为零。
再由引理 \ref{cohomology-lemma-flasque-acyclic}，得到
$\check{H}^p(U,\mathcal{F})=H^p(U,\mathcal{F})=0$。
\end{proof}

\begin{lemma}
\label{cohomology-lemma-flasque-acyclic-pushforward}
设 $f:(X,\mathcal{O}_X)\to(Y,\mathcal{O}_Y)$ 是环化空间的态射，
$\mathcal{F}$ 是 $\mathcal{O}_X$-模层。
若 $\mathcal{F}$ 松弛，则对 $p>0$，
$R^pf_*\mathcal{F}=0$。
\end{lemma}

\begin{proof}
这立即由引理 \ref{cohomology-lemma-describe-higher-direct-images}
及引理 \ref{cohomology-lemma-flasque-acyclic} 推出。
\end{proof}

\noindent
对于有限覆盖，下面的引理可用初等归纳法证明；
可与 \cite{FOAG} 中关于 {\v C}ech 上同调的讨论比较。

\begin{lemma}
\label{cohomology-lemma-vanishing-ravi}
设 $X$ 是拓扑空间，$\mathcal{F}$ 是 $X$ 上的阿贝尔层，
$\mathcal{U}:U=\bigcup_{i\in I}U_i$ 是开覆盖。
假设当 $U'$ 是形如 $U_{i_0\ldots i_p}$ 的开集的任意并时，
限制映射 $\mathcal{F}(U)\to\mathcal{F}(U')$ 都是满射。
则对 $p>0$，$\check{H}^p(\mathcal{U},\mathcal{F})$ 消失。
\end{lemma}

\begin{proof}
令 $Y$ 为 $I$ 的所有非空子集所成的集合。
我们用字母 $A,B,C,\ldots$ 表示 $Y$ 的元素，即 $I$ 的非空子集。
对有限非空子集 $J\subset I$，令
$$
V_J = \{A \in Y \mid J \subset A\}
$$
这意味着 $V_{\{i\}}=\{A\in Y\mid i\in A\}$，并且
$V_J = \bigcap_{j \in J} V_{\{j\}}$.
此时 $V_J\subset V_K$ 当且仅当 $J\supset K$。
$Y$ 上存在唯一的拓扑，使得诸子集 $V_J$ 构成拓扑基。
任意开集形如
$$
V = \bigcup\nolimits_{t \in T} V_{J_t}
$$
，其中 $(J_t)_{t\in T}$ 是有限子集族。
若 $J_t\subset J_{t'}$，则从该族中移去 $J_{t'}$ 不会改变 $V$。
因此可假设诸 $J_t$ 之间不存在包含关系。
此时 $V$ 的极小元素正是诸集合 $A=J_t$。
故可由开集 $V$ 恢复该族 $(J_t)_{t\in T}$。

\medskip\noindent
我们可以完全描述 $Y$ 中的开覆盖。首先，由于
$A\in Y$ 是 $I$ 的非空子集，有
$$
Y = \bigcup\nolimits_{i \in I} V_{\{i\}}
$$
为理解其他覆盖，设 $V$ 如上，并设开集 $V_s\subset Y$
对应于族 $(J_{s,t})_{t\in T_s}$。则
$$
V = \bigcup\nolimits_{s \in S} V_s
$$
当且仅当对每个 $t\in T$，存在 $s\in S$ 及
$t_s\in T_s$，使得 $J_t=J_{s,t_s}$。
事实上，由于族 $(J_t)_{t\in T}$ 是极小的，
极小元素 $A=J_t$ 必属于某个 $V_s$，故对某个
$t_s\in T_s$ 有 $A\in V_{J_{t_s}}$。
但 $A$ 在 $V_s$ 中也极小，故 $J_{t_s}=J_t$。

\medskip\noindent
接着，把 $Y$ 的开集集合映到 $X$ 的开集集合。
具体地，把 $Y$ 送到 $U$，并在开集 $V_J$ 上使用规则
$$
V_J \mapsto U_J = \bigcap\nolimits_{i \in J} U_i
$$
，再用规则
$$
V = \bigcup\nolimits_{t \in T} V_{J_t}
\mapsto
\bigcup\nolimits_{t \in T} U_{J_t}
$$
将其延拓到任意开集 $V$。
上述对 $Y$ 的开覆盖的分类说明，该规则把开覆盖送到开覆盖。
因此可如下得到 $Y$ 上的阿贝尔层 $\mathcal{G}$：
置 $\mathcal{G}(Y)=\mathcal{F}(U)$，并对
$V=\bigcup\nolimits_{t\in T}V_{J_t}$ 置
$$
\mathcal{G}(V) = \mathcal{F}\left(\bigcup\nolimits_{t \in T} U_{J_t}\right)
$$
，其限制映射取自 $\mathcal{F}$ 的限制映射。

\medskip\noindent
有了这些准备，即可证明引理。$Y$ 有开覆盖
$\mathcal{V} : Y = \bigcup_{i \in I} V_{\{i\}}$。
按构造，有 {\v C}ech 复形的等式
$$
\check{C}^\bullet(\mathcal{V}, \mathcal{G}) =
\check{C}^\bullet(\mathcal{U}, \mathcal{F})
$$
。由引理陈述中关于 $\mathcal{F}$ 的假设，
层 $\mathcal{G}$ 在 $Y$ 上松弛，故所需消失性由引理 \ref{cohomology-lemma-flasque-acyclic-cech} 推出。
\end{proof}




\section{Leray 谱序列}
\label{cohomology-section-Leray}

\begin{lemma}
\label{cohomology-lemma-before-Leray}
设 $f:X\to Y$ 是环化空间的态射。存在交换图
$$
\xymatrix{
D^{+}(X) \ar[rr]_-{R\Gamma(X, -)} \ar[d]_{Rf_*} & &
D^{+}(\mathcal{O}_X(X)) \ar[d]^{\text{restriction}} \\
D^{+}(Y) \ar[rr]^-{R\Gamma(Y, -)} & &
D^{+}(\mathcal{O}_Y(Y))
}
$$
更一般地，对任意开集 $V\subset Y$ 及 $U=f^{-1}(V)$，存在交换图
$$
\xymatrix{
D^{+}(X) \ar[rr]_-{R\Gamma(U, -)} \ar[d]_{Rf_*} & &
D^{+}(\mathcal{O}_X(U)) \ar[d]^{\text{restriction}} \\
D^{+}(Y) \ar[rr]^-{R\Gamma(V, -)} & &
D^{+}(\mathcal{O}_Y(V))
}
$$
更多说明见注记 \ref{cohomology-remark-elucidate-lemma}。
\end{lemma}

\begin{proof}
令
$\Gamma_{res} : \textit{Mod}(\mathcal{O}_X) \to \text{Mod}_{\mathcal{O}_Y(Y)}$
为如下函子：它把 $\mathcal{O}_X$-模 $\mathcal{F}$ 送到
$\mathcal{F}$ 的全局截面，并通过映射
$f^\sharp:\mathcal{O}_Y(Y)\to\mathcal{O}_X(X)$
将其视为 $\mathcal{O}_Y(Y)$-模。令
$restriction : \text{Mod}_{\mathcal{O}_X(X)} \to \text{Mod}_{\mathcal{O}_Y(Y)}$
为 $f^\sharp:\mathcal{O}_Y(Y)\to\mathcal{O}_X(X)$
诱导的限制标量函子。注意 $restriction$ 正合，故其右导出函子
只需逐项施用限制标量函子即可计算，见
《导出范畴》引理 \ref{derived-lemma-right-derived-exact-functor}.
显然
$$
\Gamma_{res}
=
restriction \circ \Gamma(X, -)
=
\Gamma(Y, -) \circ f_*
$$
我们断言
《导出范畴》引理 \ref{derived-lemma-compose-derived-functors}
适用于这两个复合。对第一个复合，这由上述讨论显然成立。
对第二个复合，它由引理 \ref{cohomology-lemma-pushforward-injective} 推出；该引理说明，
内射 $\mathcal{O}_X$-模被送到 $Y$ 上关于
$\Gamma(Y,-)$ 无环的层。
\end{proof}

\begin{remark}
\label{cohomology-remark-elucidate-lemma}
下面具体说明引理 \ref{cohomology-lemma-before-Leray} 的含义。
它说明，给定
$f : X \to Y$ 与 $\mathcal{F} \in \textit{Mod}(\mathcal{O}_X)$
以及内射消解 $\mathcal{F}\to\mathcal{I}^\bullet$，有
$$
\begin{matrix}
R\Gamma(X, \mathcal{F}) & \text{由下式表示} &
\Gamma(X, \mathcal{I}^\bullet) \\
Rf_*\mathcal{F} & \text{由下式表示} & f_*\mathcal{I}^\bullet \\
R\Gamma(Y, Rf_*\mathcal{F}) & \text{由下式表示} &
\Gamma(Y, f_*\mathcal{I}^\bullet)
\end{matrix}
$$
，最后一个事实来自 Leray 无环性引理
(《导出范畴》引理 \ref{derived-lemma-leray-acyclicity})
及引理 \ref{cohomology-lemma-pushforward-injective}。
最后，将此与显然的等式
$$
\Gamma(X, \mathcal{I}^\bullet)
=
\Gamma(Y, f_*\mathcal{I}^\bullet).
$$
结合，即得到引理中图表的交换性。
\end{remark}

\begin{lemma}
\label{cohomology-lemma-modules-abelian}
设 $X$ 是环化空间，$\mathcal{F}$ 是 $\mathcal{O}_X$-模。
\begin{enumerate}
\item 对开集 $U\subset X$，把 $\mathcal{F}$ 视为
$\mathcal{O}_X$-模或视为阿贝尔层所计算出的上同调群
$H^i(U,\mathcal{F})$ 相同。
\item 设 $f:X\to Y$ 是环化空间的态射。
把 $\mathcal{F}$ 视为 $\mathcal{O}_X$-模或视为阿贝尔层
所计算出的高阶直像 $R^if_*\mathcal{F}$ 相同。
\end{enumerate}
对 $\mathcal{O}_X$-模的下有界复形也有类似陈述。
\end{lemma}

\begin{proof}
考虑环化空间的态射
$(X, \mathcal{O}_X) \to (X, \underline{\mathbf{Z}}_X)$
；它在底层拓扑空间上为恒等映射，在结构层上由环层的唯一映射
$\underline{\mathbf{Z}}_X\to\mathcal{O}_X$ 给出。
设 $\mathcal{F}$ 是 $\mathcal{O}_X$-模。
以 $\mathcal{F}_{ab}$ 表示将同一层视为
$\underline{\mathbf{Z}}_X$-模，即视为阿贝尔群层。
设 $\mathcal{F}\to\mathcal{I}^\bullet$ 是内射消解。
由注记 \ref{cohomology-remark-elucidate-lemma}，
$\Gamma(X,\mathcal{I}^\bullet)$ 同时计算
$R\Gamma(X,\mathcal{F})$ 与
$R\Gamma(X,\mathcal{F}_{ab})$。这证明 (1)。

\medskip\noindent
为证明 (2)，使用 (1) 及引理 \ref{cohomology-lemma-describe-higher-direct-images}，结论立即成立。
\end{proof}

\begin{lemma}[Leray 谱序列]
\label{cohomology-lemma-Leray}
设 $f:X\to Y$ 是环化空间的态射，
$\mathcal{F}^\bullet$ 是 $\mathcal{O}_X$-模的下有界复形。
存在谱序列
$$
E_2^{p, q} = H^p(Y, R^qf_*(\mathcal{F}^\bullet))
$$
，它收敛到 $H^{p+q}(X,\mathcal{F}^\bullet)$。
\end{lemma}

\begin{proof}
这就是由函子复合
$\Gamma_{res}=\Gamma(Y,-)\circ f_*$ 所得的 Grothendieck 谱序列
《导出范畴》引理 \ref{derived-lemma-grothendieck-spectral-sequence}
，其中 $\Gamma_{res}$ 如引理 \ref{cohomology-lemma-before-Leray} 的证明中所定义。
关于
《导出范畴》引理 \ref{derived-lemma-grothendieck-spectral-sequence}
的假设为何满足，可参见引理 \ref{cohomology-lemma-before-Leray} 的证明或注记 \ref{cohomology-remark-elucidate-lemma}。
\end{proof}

\begin{remark}
\label{cohomology-remark-Leray-ss-more-structure}
按引理 \ref{cohomology-lemma-Leray} 中的证明方式，
Leray 谱序列是 $\Gamma(Y,\mathcal{O}_Y)$-模的谱序列。
不过很容易看出，它事实上是 $\Gamma(X,\mathcal{O}_X)$-模的谱序列。
例如，$f$ 给出环化空间的态射
$f' :  (X, \mathcal{O}_X) \to (Y, f_*\mathcal{O}_X)$.
由引理 \ref{cohomology-lemma-modules-abelian}，对 $\mathcal{O}_X$-模
$\mathcal{F}$ 而言，与 $f$ 相伴的 Leray 谱序列之各项
$E_r^{p,q}$ 至少在 $r\geq2$ 时，与 $f'$ 所给出的各项相同。
事实上，二者都等于把 $\mathcal{F}$ 视为阿贝尔层所得的
Leray 谱序列之各项。又因
$(f_*\mathcal{O}_X)(Y)=\mathcal{O}_X(X)$，结论成立。
Leray 谱序列往往还带有额外结构。
\end{remark}

\begin{lemma}
\label{cohomology-lemma-apply-Leray}
设 $f:X\to Y$ 是环化空间的态射，
$\mathcal{F}$ 是 $\mathcal{O}_X$-模。
\begin{enumerate}
\item 若对 $q>0$ 有 $R^qf_*\mathcal{F}=0$，则对所有 $p$，
$H^p(X,\mathcal{F})=H^p(Y,f_*\mathcal{F})$。
\item 若对所有 $q$ 及 $p>0$ 有
$H^p(Y,R^qf_*\mathcal{F})=0$，则对所有 $q$，
$H^q(X,\mathcal{F})=H^0(Y,R^qf_*\mathcal{F})$。
\end{enumerate}
\end{lemma}

\begin{proof}
这两个简单条件都迫使 Leray 谱序列在 $E_2$ 页退化。
也可以不使用谱序列直接证明这些事实；
这是层上同调方面的一个好练习。
\end{proof}

\begin{lemma}
\label{cohomology-lemma-higher-direct-images-compose}
\begin{slogan}
复合函子的全导出函子等于诸全导出函子的复合。
\end{slogan}
设 $f:X\to Y$ 及 $g:Y\to Z$ 是环化空间的态射。
此时，作为从 $D^{+}(X)$ 到 $D^{+}(Z)$ 的函子，有
$Rg_*\circ Rf_*=R(g\circ f)_*$。
\end{lemma}

\begin{proof}
施用
《导出范畴》引理 \ref{derived-lemma-compose-derived-functors}.
显然 $g_*\circ f_*=(g\circ f)_*$，见
《层》引理 \ref{sheaves-lemma-pushforward-composition}。
只需证明 $f_*\mathcal{I}$ 关于 $g_*$ 无环。
这由引理 \ref{cohomology-lemma-pushforward-injective}
及引理 \ref{cohomology-lemma-describe-higher-direct-images}
中对高阶直像 $R^ig_*$ 的描述推出。
\end{proof}

\begin{lemma}[相对 Leray 谱序列]
\label{cohomology-lemma-relative-Leray}
设 $f:X\to Y$ 及 $g:Y\to Z$ 是环化空间的态射，
$\mathcal{F}$ 是 $\mathcal{O}_X$-模。
存在谱序列，其
$$
E_2^{p, q} = R^pg_*(R^qf_*\mathcal{F})
$$
收敛到 $R^{p+q}(g\circ f)_*\mathcal{F}$。
该谱序列关于 $\mathcal{F}$ 是函子性的，
并且对 $\mathcal{O}_X$-模的下有界复形也有相应版本。
\end{lemma}

\begin{proof}
这是函子复合的 Grothendieck 谱序列，由引理 \ref{cohomology-lemma-higher-direct-images-compose} 及
《导出范畴》引理 \ref{derived-lemma-grothendieck-spectral-sequence} 推出。
\end{proof}














\section{上同调的函子性}
\label{cohomology-section-functoriality}

\begin{lemma}
\label{cohomology-lemma-functoriality}
设 $f:X\to Y$ 是环化空间的态射。
设 $\mathcal{G}^\bullet$、相应地 $\mathcal{F}^\bullet$
分别是 $\mathcal{O}_Y$-模、相应地 $\mathcal{O}_X$-模的下有界复形。
设
$\varphi : \mathcal{G}^\bullet \to f_*\mathcal{F}^\bullet$
是复形态射。则在 $D^{+}(Y)$ 中存在典范态射
$$
\mathcal{G}^\bullet
\longrightarrow
Rf_*(\mathcal{F}^\bullet)
$$
。此外，该构造关于三元组
$(\mathcal{G}^\bullet,\mathcal{F}^\bullet,\varphi)$ 是函子性的。
\end{lemma}

\begin{proof}
取内射消解
$\mathcal{F}^\bullet\to\mathcal{I}^\bullet$。
按定义，$Rf_*(\mathcal{F}^\bullet)$ 在
$K^{+}(\mathcal{O}_Y)$ 中由 $f_*\mathcal{I}^\bullet$ 表示。
复合
$$
\mathcal{G}^\bullet \to f_*\mathcal{F}^\bullet \to f_*\mathcal{I}^\bullet
$$
是 $K^{+}(Y)$ 中的态射；施用局部化函子
$j_Y:K^{+}(Y)\to D^{+}(Y)$ 后，它变为引理中的态射。
\end{proof}

\noindent
设 $f:X\to Y$ 是环化空间的态射，
$\mathcal{G}$ 是 $\mathcal{O}_Y$-模，
$\mathcal{F}$ 是 $\mathcal{O}_X$-模。
回忆，从 $\mathcal{G}$ 到 $\mathcal{F}$ 的 $f$-映射
$\varphi$ 是映射
$\varphi:\mathcal{G}\to f_*\mathcal{F}$；
等价地，它是映射
$\varphi:f^*\mathcal{G}\to\mathcal{F}$。见
《层》定义 \ref{sheaves-definition-f-map}。
这样的 $f$-映射给出 $D^{+}(\mathcal{O}_Y(Y))$ 中的复形态射
\begin{equation}
\label{cohomology-equation-functorial-derived}
\varphi :
R\Gamma(Y, \mathcal{G})
\longrightarrow
R\Gamma(X, \mathcal{F})
\end{equation}
。具体地，使用引理 \ref{cohomology-lemma-functoriality}
给出的 $D^{+}(Y)$ 中的态射
$\mathcal{G}\to Rf_*\mathcal{F}$，并施用 $R\Gamma(Y,-)$。
由引理 \ref{cohomology-lemma-before-Leray}，有
$R\Gamma(X, \mathcal{F}) = R\Gamma(Y, Rf_*\mathcal{F})$
，故得到所示箭头。下面的注记 \ref{cohomology-remark-explain-arrow} 将完整展开这一构造。
特别地，它给出上同调上的映射
\begin{equation}
\label{cohomology-equation-functorial}
\varphi : H^i(Y, \mathcal{G}) \longrightarrow H^i(X, \mathcal{F}).
\end{equation}

\begin{remark}
\label{cohomology-remark-explain-arrow}
设 $f:X\to Y$ 是环化空间的态射，
$\mathcal{G}$ 是 $\mathcal{O}_Y$-模，
$\mathcal{F}$ 是 $\mathcal{O}_X$-模，
$\varphi$ 是从 $\mathcal{G}$ 到 $\mathcal{F}$ 的 $f$-映射。
选取由内射 $\mathcal{O}_X$-模复形给出的消解
$\mathcal{F}\to\mathcal{I}^\bullet$。
再选取由内射 $\mathcal{O}_Y$-模复形给出的消解
$\mathcal{G}\to\mathcal{J}^\bullet$ 及
$f_*\mathcal{I}^\bullet\to(\mathcal{J}')^\bullet$。
由
《导出范畴》引理 \ref{derived-lemma-morphisms-lift}
，存在复形映射 $\beta$，使得图
\begin{equation}
\label{cohomology-equation-choice}
\xymatrix{
\mathcal{G} \ar[d] \ar[r] &
f_*\mathcal{F} \ar[r] &
f_*\mathcal{I}^\bullet \ar[d] \\
\mathcal{J}^\bullet \ar[rr]^\beta & &
(\mathcal{J}')^\bullet
}
\end{equation}
交换。施用全局截面函子，得到图
$$
\xymatrix{
 & & \Gamma(Y, f_*\mathcal{I}^\bullet) \ar[d]_{qis} \ar@{=}[r] &
\Gamma(X, \mathcal{I}^\bullet) \\
\Gamma(Y, \mathcal{J}^\bullet) \ar[rr]^\beta & &
\Gamma(Y, (\mathcal{J}')^\bullet) &
}
$$
。左下方的复形表示 $R\Gamma(Y,\mathcal{G})$，
右上方的复形表示 $R\Gamma(X,\mathcal{F})$。
由引理 \ref{cohomology-lemma-before-Leray}，竖直箭头是拟同构；
施用局部化函子
$K^{+}(\mathcal{O}_Y(Y))\to D^{+}(\mathcal{O}_Y(Y))$
后它变得可逆。
箭头 (\ref{cohomology-equation-functorial-derived})
由水平映射与竖直映射之逆的复合给出。
\end{remark}





\section{加细与 {\v C}ech 上同调}
\label{cohomology-section-refinements-cech}

\noindent
设 $(X,\mathcal{O}_X)$ 是环化空间，并设
$\mathcal{U} : X = \bigcup_{i \in I} U_i$ 及
$\mathcal{V} : X = \bigcup_{j \in J} V_j$ 是开覆盖。
假设 $\mathcal{U}$ 是 $\mathcal{V}$ 的加细。
选择映射 $c:I\to J$，使得对所有 $i\in I$，
$U_i\subset V_{c(i)}$。它诱导 {\v C}ech 复形之间的映射
$$
\gamma :
\check{\mathcal{C}}^\bullet(\mathcal{V}, \mathcal{F})
\longrightarrow
\check{\mathcal{C}}^\bullet(\mathcal{U}, \mathcal{F}),
\quad
(\xi_{j_0 \ldots j_p})
\longmapsto
(\xi_{c(i_0) \ldots c(i_p)}|_{U_{i_0 \ldots i_p}})
$$
，且关于 $\mathcal{O}_X$-模层 $\mathcal{F}$ 是函子性的。
假设 $c':I\to J$ 是另一个映射，使得对所有 $i\in I$，
$U_i\subset V_{c'(i)}$。则相应的映射 $\gamma$ 与 $\gamma'$
彼此同伦。具体地，
$\gamma - \gamma' = \text{d} \circ h + h \circ \text{d}$
，其中
$h : \check{\mathcal{C}}^{p + 1}(\mathcal{V}, \mathcal{F}) \to
\check{\mathcal{C}}^p(\mathcal{U}, \mathcal{F})$
由规则
$$
h(\alpha)_{i_0 \ldots i_p} =
\sum\nolimits_{a = 0}^{p}
(-1)^a
\alpha_{c(i_0)\ldots c(i_a) c'(i_a) \ldots c'(i_p)}
$$
给出。验证此公式的计算从略；更一般情形的证明见
(\ref{cohomology-equation-transformation}) 之后的讨论。
特别地，{\v C}ech 上同调群上的映射与 $c$ 的选择无关。
此外，若 $\mathcal{W}:X=\bigcup_{k\in K}W_k$ 是第三个开覆盖，
且 $\mathcal{V}$ 是 $\mathcal{W}$ 的加细，则显然映射之复合
$$
\check{\mathcal{C}}^\bullet(\mathcal{W}, \mathcal{F})
\longrightarrow
\check{\mathcal{C}}^\bullet(\mathcal{V}, \mathcal{F})
\longrightarrow
\check{\mathcal{C}}^\bullet(\mathcal{U}, \mathcal{F})
$$
（分别与映射 $I\to J$ 及 $J\to K$ 相伴）
等于与复合 $I\to K$ 相伴的映射。
特别地，可定义 {\v C}ech 上同调群
$$
\check{H}^p(X, \mathcal{F}) =
\colim_\mathcal{U} \check{H}^p(\mathcal{U}, \mathcal{F})
$$
，其中余极限取遍按加细预序排列的 $X$ 的所有开覆盖。

\medskip\noindent
上述映射 $\gamma$ 与到层上同调的映射相容。换言之，复合
$$
\check{H}^p(\mathcal{V}, \mathcal{F}) \to
\check{H}^p(\mathcal{U}, \mathcal{F})
\xrightarrow{\text{引理 \ref{cohomology-lemma-cech-cohomology}}}
H^p(X, \mathcal{F})
$$
正是引理 \ref{cohomology-lemma-cech-cohomology}
中从第一个群到层上同调的典范映射。
下面的引理将在稍一般的情形中证明这一点。
由此得到良定义映射
\begin{equation}
\label{cohomology-equation-cech-to-cohomology}
\check{H}^p(X, \mathcal{F}) =
\colim_\mathcal{U} \check{H}^p(\mathcal{U}, \mathcal{F})
\longrightarrow
H^p(X, \mathcal{F})
\end{equation}
，它从 {\v C}ech 上同调映到层上同调。

\begin{lemma}
\label{cohomology-lemma-functoriality-cech}
设 $f:X\to Y$ 是环化空间的态射，
$\varphi:f^*\mathcal{G}\to\mathcal{F}$
是从 $\mathcal{O}_Y$-模 $\mathcal{G}$ 到
$\mathcal{O}_X$-模 $\mathcal{F}$ 的 $f$-映射。
设 $\mathcal{U}:X=\bigcup_{i\in I}U_i$ 及
$\mathcal{V}:Y=\bigcup_{j\in J}V_j$ 是开覆盖。
假设 $\mathcal{U}$ 是
$f^{-1}\mathcal{V}:X=\bigcup_{j\in J}f^{-1}(V_j)$ 的加细。
此时在 $D^{+}(\mathcal{O}_X(X))$ 中存在交换图
$$
\xymatrix{
\check{\mathcal{C}}^\bullet(\mathcal{U}, \mathcal{F}) \ar[r] &
R\Gamma(X, \mathcal{F}) \\
\check{\mathcal{C}}^\bullet(\mathcal{V}, \mathcal{G}) \ar[r]
\ar[u]^\gamma &
R\Gamma(Y, \mathcal{G}) \ar[u]
}
$$
，其中水平箭头由引理 \ref{cohomology-lemma-cech-cohomology} 给出，
右侧竖直箭头由
(\ref{cohomology-equation-functorial-derived}) 给出。
特别地，得到上同调群的交换图
$$
\xymatrix{
\check{H}^p(\mathcal{U}, \mathcal{F}) \ar[r] &
H^p(X, \mathcal{F}) \\
\check{H}^p(\mathcal{V}, \mathcal{G}) \ar[r]
\ar[u]^\gamma &
H^p(Y, \mathcal{G}) \ar[u]
}
$$
，其中右侧竖直箭头是 (\ref{cohomology-equation-functorial})。
\end{lemma}

\begin{proof}
先定义左侧竖直箭头。选择映射 $c:I\to J$，使得对所有
$i\in I$，有 $U_i\subset f^{-1}(V_{c(i)})$。
在次数 $p$，用规则
$$
\gamma(s)_{i_0 \ldots i_p} = \varphi(s)_{c(i_0) \ldots c(i_p)}
$$
定义该映射。这是有意义的，因为由假设，
$\varphi$ 确实诱导映射
$\mathcal{G}(V_{c(i_0) \ldots c(i_p)}) \to \mathcal{F}(U_{i_0 \ldots i_p})$
。显然，这定义了复形态射。选取 $X$ 上的内射消解
$\mathcal{F} \to \mathcal{I}^\bullet$
以及 $Y$ 上的内射消解
$\mathcal{G}\to J^\bullet$。
按照引理 \ref{cohomology-lemma-cech-cohomology} 的证明，
引入双复形 $A^{\bullet,\bullet}$ 与 $B^{\bullet,\bullet}$，其各项为
$$
B^{p, q} = \check{\mathcal{C}}^p(\mathcal{V}, \mathcal{J}^q)
\quad
\text{以及}
\quad
A^{p, q} = \check{\mathcal{C}}^p(\mathcal{U}, \mathcal{I}^q).
$$
如上面的注记 \ref{cohomology-remark-explain-arrow}，还选择 $Y$ 上的内射消解
$f_*\mathcal{I} \to (\mathcal{J}')^\bullet$
及复形态射
$\beta:\mathcal{J}\to(\mathcal{J}')^\bullet$，
使 (\ref{cohomology-equation-choice}) 交换。
再引入双复形 $(B')^{\bullet,\bullet}$ 与
$(B'')^{\bullet,\bullet}$，其各项为
$$
(B')^{p, q} = \check{\mathcal{C}}^p(\mathcal{V}, (\mathcal{J}')^q)
\quad
\text{以及}
\quad
(B'')^{p, q} = \check{\mathcal{C}}^p(\mathcal{V}, f_*\mathcal{I}^q).
$$
注意，存在从 $f_*\mathcal{I}^\bullet$ 到
$\mathcal{I}^\bullet$ 的复形 $f$-映射。
因此显然，上述同一规则定义双复形态射
$$
\gamma : (B'')^{\bullet, \bullet} \longrightarrow A^{\bullet, \bullet}.
$$
考虑复形图
$$
\xymatrix{
\check{\mathcal{C}}^\bullet(\mathcal{U}, \mathcal{F})
\ar[r] &
\text{Tot}(A^{\bullet, \bullet}) & & &
\Gamma(X, \mathcal{I}^\bullet) \ar[lll]^{qis}
\ar@{=}[ddl]\\
\check{\mathcal{C}}^\bullet(\mathcal{V}, \mathcal{G})
\ar[r] \ar[u]^\gamma &
\text{Tot}(B^{\bullet, \bullet}) \ar[r]^\beta &
\text{Tot}((B')^{\bullet, \bullet}) &
\text{Tot}((B'')^{\bullet, \bullet}) \ar[l] \ar[llu]_{s\gamma} \\
& \Gamma(Y, \mathcal{J}^\bullet) \ar[u]^{qis} \ar[r]^\beta &
\Gamma(Y, (\mathcal{J}')^\bullet) \ar[u] &
\Gamma(Y, f_*\mathcal{I}^\bullet) \ar[u] \ar[l]_{qis}
}
$$
。以 $\text{Tot}(A^{\bullet,\bullet})$ 及
$\text{Tot}(B^{\bullet,\bullet})$ 为目标的两个水平箭头，
正是引理 \ref{cohomology-lemma-cech-cohomology} 中解释的箭头。
左上方的图形（五边形）交换，这直接来自
(\ref{cohomology-equation-choice}) 的交换性。
下面两个正方形显然交换。
由定义也立即可知，右上方的图形（正方形）交换。
现在由定义及下述事实得到引理：沿图的外边界，
无论从上方还是下方由
$\check{\mathcal{C}}^\bullet(\mathcal{V},\mathcal{G})$
走到 $\Gamma(X,\mathcal{I}^\bullet)$，所得结果相同
（途中须将所遇到的拟同构取逆）。
\end{proof}





\section{Hausdorff 拟紧空间上的上同调}
\label{cohomology-section-cohomology-LC}

\noindent
对这样的空间，{\v C}ech 上同调与层上同调一致。

\begin{lemma}
\label{cohomology-lemma-cech-always}
设 $X$ 是拓扑空间，$\mathcal{F}$ 是阿贝尔层。
则 (\ref{cohomology-equation-cech-to-cohomology}) 中定义的映射
$\check{H}^1(X,\mathcal{F})\to H^1(X,\mathcal{F})$ 是同构。
\end{lemma}

\begin{proof}
设 $\mathcal{U}$ 是 $X$ 的开覆盖。
由引理 \ref{cohomology-lemma-cech-spectral-sequence}，存在正合序列
$$
0 \to \check{H}^1(\mathcal{U}, \mathcal{F}) \to H^1(X, \mathcal{F})
\to \check{H}^0(\mathcal{U}, \underline{H}^1(\mathcal{F}))
$$
故该映射是单射。为证明满射性，只需证明：
用 $\mathcal{U}$ 的某个加细替换 $\mathcal{U}$ 后，
$\check{H}^0(\mathcal{U},\underline{H}^1(\mathcal{F}))$
中的任意元素都映到零。
这由定义及下述事实立即得到：
$\underline{H}^1(\mathcal{F})$ 是阿贝尔群预层，
由上同调的局部性，其层化为零；见引理 \ref{cohomology-lemma-kill-cohomology-class-on-covering}。
\end{proof}

\begin{lemma}
\label{cohomology-lemma-cech-Hausdorff-quasi-compact}
设 $X$ 是 Hausdorff 拟紧拓扑空间，
$\mathcal{F}$ 是 $X$ 上的阿贝尔层。
则对所有 $n$，(\ref{cohomology-equation-cech-to-cohomology}) 中定义的映射
$\check{H}^n(X,\mathcal{F})\to H^n(X,\mathcal{F})$ 是同构。
\end{lemma}

\begin{proof}
我们已经知道，当 $n=0,1$ 时，
$\check{H}^n(X,-)\to H^n(X,-)$ 是函子同构，见引理 \ref{cohomology-lemma-cech-always}。
诸函子 $H^n(X,-)$ 组成泛 $\delta$-函子，见
《导出范畴》引理 \ref{derived-lemma-higher-derived-functors}。
若证明 $\check{H}^n(X,-)$ 组成泛 $\delta$-函子，
且 $\check{H}^n(X,-)\to H^n(X,-)$ 与连接映射相容，
则由泛 $\delta$-函子的唯一性，该映射自动是同构，见
《同调代数》引理 \ref{homology-lemma-uniqueness-universal-delta-functor}。

\medskip\noindent
设
$0\to\mathcal{F}\to\mathcal{G}\to\mathcal{H}\to0$
是 $X$ 上阿贝尔层的短正合序列，
$\mathcal{U}:X=\bigcup_{i\in I}U_i$ 是开覆盖。
它给出复形的复形
$$
0 \to \check{\mathcal{C}}^\bullet(\mathcal{U}, \mathcal{F}) \to
\check{\mathcal{C}}^\bullet(\mathcal{U}, \mathcal{G}) \to
\check{\mathcal{C}}^\bullet(\mathcal{U}, \mathcal{H}) \to 0
$$
，该复形一般在右端不正合。此序列定义映射
$$
\check{H}^n(\mathcal{U}, \mathcal{F}) \to
\check{H}^n(\mathcal{U}, \mathcal{G}) \to
\check{H}^n(\mathcal{U}, \mathcal{H})
$$
，但不足以定义连接算子
$\delta : \check{H}^n(\mathcal{U}, \mathcal{H}) \to
\check{H}^{n + 1}(\mathcal{U}, \mathcal{F})$。
事实上，这样的算子一般不存在。不过，给定元素
$\overline{h}\in\check{H}^n(\mathcal{U},\mathcal{H})$，
它是上循环 $h=(h_{i_0\ldots i_n})$ 的上同调类；
我们可以选择开覆盖
$$
U_{i_0 \ldots i_n} = \bigcup W_{i_0 \ldots i_n, k}
$$
，使得
$h_{i_0\ldots i_n}|_{W_{i_0\ldots i_n,k}}$
可提升为 $\mathcal{G}$ 在 $W_{i_0\ldots i_n,k}$ 上的截面。
由 《拓扑》引理 \ref{topology-lemma-refine-covering}
（正是在此使用 $X$ 为 Hausdorff 拟紧空间的假设），
可选择开覆盖 $\mathcal{V}:X=\bigcup_{j\in J}V_j$
及 $\alpha:J\to I$，使得
$V_j\subset U_{\alpha(j)}$（故它是加细），并且对所有
$j_0,\ldots,j_n\in J$，存在 $k$ 使得
$V_{j_0 \ldots j_n} \subset W_{\alpha(j_0) \ldots \alpha(j_n), k}$.
于是得到复形之间的映射
$$
\xymatrix{
0 \ar[r] &
\check{\mathcal{C}}^\bullet(\mathcal{U}, \mathcal{F}) \ar[d] \ar[r] &
\check{\mathcal{C}}^\bullet(\mathcal{U}, \mathcal{G}) \ar[d] \ar[r] &
\check{\mathcal{C}}^\bullet(\mathcal{U}, \mathcal{H}) \ar[d] \ar[r] &
0 \\
0 \ar[r] &
\check{\mathcal{C}}^\bullet(\mathcal{V}, \mathcal{F}) \ar[r] &
\check{\mathcal{C}}^\bullet(\mathcal{V}, \mathcal{G}) \ar[r] &
\check{\mathcal{C}}^\bullet(\mathcal{V}, \mathcal{H}) \ar[r] &
0
}
$$
事实上，竖直箭头正是定义 {\v C}ech 上同调群之间转移映射
所用的复形映射。我们对加细的选择说明，可以选择
$$
g_{j_0 \ldots j_n} \in
\mathcal{G}(V_{j_0 \ldots j_n}),\quad
g_{j_0 \ldots j_n} \longmapsto
h_{\alpha(j_0) \ldots \alpha(j_n)}|_{V_{j_0 \ldots j_n}}
$$
上链 $g=(g_{j_0\ldots j_n})$ 一般不是上循环，
但其 {\v C}ech 边界 $\text{d}(g)$ 在
$\check{\mathcal{C}}^{n+1}(\mathcal{V},\mathcal{H})$
中映到零（由上面的交换图及 $h$ 为上循环这一事实）。
故 $\text{d}(g)$ 是
$\check{\mathcal{C}}^\bullet(\mathcal{V},\mathcal{F})$ 中的上循环。
这使我们可以定义
$$
\delta(\overline{h}) = \text{下式中 }\text{d}(g)\text{ 的类：}
\check{H}^{n + 1}(\mathcal{V}, \mathcal{F})
$$
现在，给定元素 $\xi\in\check{H}^n(X,\mathcal{G})$，
选择开覆盖 $\mathcal{U}$ 及元素
$\overline{h}\in\check{H}^n(\mathcal{U},\mathcal{G})$，
它在定义 {\v C}ech 上同调的余极限中映到 $\xi$。
再如上选择 $\mathcal{V}$ 与 $g$，并定义
$\delta(\xi)$ 为 $\delta(\overline{h})$
在 $\check{H}^n(X,\mathcal{F})$ 中的像。
此时需验证许多性质，但它们都是直接的。
例如，需验证该构造与
$\mathcal{U},\overline{h},\mathcal{V},\alpha:J\to I,g$
的选择无关。
$\text{d}(g)$ 的类与提升 $g_{i_0\ldots i_n}$ 的选择无关，
因为二者之差将是上边缘。
它与 $\alpha$ 的选择无关\footnote{这项验证很重要：
$\alpha$ 的非唯一性正是阻止我们对 $X$ 的所有开覆盖取
{\v C}ech 复形的余极限，从而得到计算 {\v C}ech 上同调的
复形短正合序列的唯一障碍。}，
因为 $\alpha$ 的不同选择在上图中确定彼此同伦的竖直复形映射，
见 第 \ref{cohomology-section-refinements-cech} 节。
对其他选择，我们使用如下事实：
给定 $X$ 的有限多个开覆盖，总能找到同时加细它们的开覆盖。
还需验证加性，方法相同。
最后，需验证映射
$\check{H}^n(X,-)\to H^n(X,-)$ 与连接映射相容。
为此选择内射消解
$$
\xymatrix{
0 \ar[r] &
\mathcal{F} \ar[r] \ar[d] &
\mathcal{G} \ar[r] \ar[d] &
\mathcal{H} \ar[r] \ar[d] &
0 \\
0 \ar[r] &
\mathcal{I}_1^\bullet \ar[r] &
\mathcal{I}_2^\bullet \ar[r] &
\mathcal{I}_3^\bullet \ar[r] &
0
}
$$
，如 《导出范畴》引理 \ref{derived-lemma-injective-resolution-ses} 中所述。
这给出交换图
$$
\xymatrix{
0 \ar[r] &
\check{\mathcal{C}}^\bullet(\mathcal{U}, \mathcal{F}) \ar[r] \ar[d] &
\check{\mathcal{C}}^\bullet(\mathcal{U}, \mathcal{F}) \ar[r] \ar[d] &
\check{\mathcal{C}}^\bullet(\mathcal{U}, \mathcal{F}) \ar[r] \ar[d] &
0 \\
0 \ar[r] &
\text{Tot}(\check{\mathcal{C}}^\bullet(\mathcal{U}, \mathcal{I}_1^\bullet))
\ar[r] &
\text{Tot}(\check{\mathcal{C}}^\bullet(\mathcal{U}, \mathcal{I}_2^\bullet))
\ar[r] &
\text{Tot}(\check{\mathcal{C}}^\bullet(\mathcal{U}, \mathcal{I}_3^\bullet))
\ar[r] &
0
}
$$
这里 $\mathcal{U}$ 是如上的开覆盖，
而竖直映射正是定义映射
$\check{H}^n(\mathcal{U},-)\to H^n(X,-)$ 时所用的映射，见引理 \ref{cohomology-lemma-cech-cohomology}。
下面的复形正合，因为内射对象的复形序列逐项分裂正合。
因此，上同调中的连接映射由下方正合序列的通常程序计算，见
《同调代数》引理 \ref{homology-lemma-long-exact-sequence-cochain}。
传递到定义 {\v C}ech 上同调连接映射的加细
$\mathcal{V}$ 后，同样如此。
故这些连接映射一致，因为它们使用相同的构造
（只要前者定义在某给定开覆盖的 {\v C}ech 上同调元素上）。
至此，我们完成了关于 {\v C}ech 上同调上
$\delta$-函子结构之构造，以及为何该结构与通常上同调上给定的
$\delta$-函子结构相容的讨论。

\medskip\noindent
最后，施用引理 \ref{cohomology-lemma-injective-trivial-cech}
可知内射层的高阶 {\v C}ech 上同调为零。
因此由
《同调代数》引理 \ref{homology-lemma-efface-implies-universal}，
{\v C}ech 上同调是泛 $\delta$-函子。
\end{proof}

\begin{lemma}
\label{cohomology-lemma-cohomology-of-closed}
\begin{reference}
\cite[Expose V bis, 4.1.3]{SGA4}
\end{reference}
设 $X$ 是拓扑空间，$Z\subset X$ 是拟紧子集，
且 $Z$ 中任意两点在 $X$ 中都有不交的开邻域。
对 $X$ 上每个阿贝尔层 $\mathcal{F}$，典范映射
$$
\colim H^p(U, \mathcal{F})
\longrightarrow
H^p(Z, \mathcal{F}|_Z)
$$
是同构，其中余极限取遍 $Z$ 在 $X$ 中的开邻域 $U$。
\end{lemma}

\begin{proof}
先证明 $p=0$ 的情形。单射性由 $\mathcal{F}|_Z$ 的定义推出，
并且一般都成立（对任意拓扑空间 $X$ 的任意子集）。
再设 $s\in H^0(Z,\mathcal{F}|_Z)$。
可找到开集 $U_i\subset X$，使得
$Z\subset\bigcup U_i$，且
$s|_{Z\cap U_i}$ 来自某个 $s_i\in\mathcal{F}(U_i)$。
于是存在开集 $W_{ij}\subset U_i\cap U_j$，满足
$W_{ij}\cap Z=U_i\cap U_j\cap Z$，且
$s_i|_{W_{ij}}=s_j|_{W_{ij}}$。
施用 《拓扑》引理
\ref{topology-lemma-lift-covering-of-quasi-compact-hausdorff-subset}，
可找到 $X$ 的开集 $V_i$，使得
$V_i\subset U_i$ 且 $V_i\cap V_j\subset W_{ij}$。
因此，诸 $s_i|_{V_i}$ 黏合成 $\mathcal{F}$
在 $Z$ 的开邻域 $\bigcup V_i$ 上的截面。

\medskip\noindent
为完成证明，只需说明：若 $\mathcal{I}$ 是 $X$ 上的内射阿贝尔层，
则对 $p>0$ 有 $H^p(Z,\mathcal{I}|_Z)=0$。
这可用短正合序列与维数转移推出；细节从略。
因此，设对某个 $p>0$，
$\overline{\xi}\in H^p(Z,\mathcal{I}|_Z)$。
由引理 \ref{cohomology-lemma-cech-Hausdorff-quasi-compact}，
元素 $\overline{\xi}$ 来自
$\check{H}^p(\mathcal{V},\mathcal{I}|_Z)$，
其中 $\mathcal{V}:Z=\bigcup V_i$ 是 $Z$ 的某个开覆盖。
设 $\overline{\xi}$ 是
$\check{\mathcal{C}}^p(\mathcal{V},\mathcal{I}|_Z)$
中上循环 $\xi=(\xi_{i_0\ldots i_p})$ 之类的像。

\medskip\noindent
令 $\mathcal{I}'\subset\mathcal{I}|_Z$ 为由下列规则定义的子预层：
$$
\mathcal{I}'(V) =
\{s \in \mathcal{I}|_Z(V) \mid
\exists (U, t),\ U \subset X\text{ 开},
\ t \in \mathcal{I}(U),\ V = Z \cap U,\ s = t|_{Z \cap U} \}
$$
则 $\mathcal{I}|_Z$ 是 $\mathcal{I}'$ 的层化。
因此，对每个 $(p+1)$-元组 $i_0\ldots i_p$，
可找到开覆盖
$V_{i_0\ldots i_p}=\bigcup W_{i_0\ldots i_p,k}$，
使得
$\xi_{i_0\ldots i_p}|_{W_{i_0\ldots i_p,k}}$
是 $\mathcal{I}'$ 的截面。
施用 《拓扑》引理 \ref{topology-lemma-refine-covering}，
在加细 $\mathcal{V}$ 后，可假设每个
$\xi_{i_0\ldots i_p}$ 都是预层 $\mathcal{I}'$ 的截面。

\medskip\noindent
将 $V_i$ 写成 $Z\cap U_i$，其中 $U_i\subset X$ 开。
由于 $\mathcal{I}$ 松弛
（引理 \ref{cohomology-lemma-injective-flasque}），
且对每个 $(p+1)$-元组 $i_0\ldots i_p$，
$\xi_{i_0\ldots i_p}$ 都是 $\mathcal{I}'$ 的截面，
可选择截面
$s_{i_0\ldots i_p}\in\mathcal{I}(U_{i_0\ldots i_p})$，
它在
$V_{i_0\ldots i_p}=Z\cap U_{i_0\ldots i_p}$
上限制为 $\xi_{i_0\ldots i_p}$。
（还可以再次施用
《拓扑》引理
\ref{topology-lemma-lift-covering-of-quasi-compact-hausdorff-subset}，
从而避免使用内射层松弛这一事实。）
令 $s=(s_{i_0\ldots i_p})$ 为开覆盖
$U=\bigcup U_i$ 的相应上链。
由于 $\text{d}(\xi)=0$，截面
$\text{d}(s)_{i_0\ldots i_{p+1}}$
在 $Z\cap U_{i_0\ldots i_{p+1}}$ 上限制为零。
故由证明开头的讨论，存在开子集
$W_{i_0\ldots i_{p+1}}\subset U_{i_0\ldots i_{p+1}}$，
满足
$Z\cap W_{i_0\ldots i_{p+1}}
=Z\cap U_{i_0\ldots i_{p+1}}$，且
$\text{d}(s)_{i_0\ldots i_{p+1}}|_{W_{i_0\ldots i_{p+1}}}=0$。
由 《拓扑》引理
\ref{topology-lemma-lift-covering-of-quasi-compact-hausdorff-subset}
，可找到 $U'_i\subset U_i$，使得
$Z\subset\bigcup U'_i$，且
$U'_{i_0\ldots i_{p+1}}\subset W_{i_0\ldots i_{p+1}}$。
令 $s'=(s'_{i_0\ldots i_p})$，其中
$s'_{i_0 \ldots i_p} = s_{i_0 \ldots i_p}|_{U'_{i_0 \ldots i_p}}$
。则 $s'$ 是 $\mathcal{I}$ 关于开覆盖
$U'=\bigcup U'_i$ 的上循环，其中 $U'$ 是 $Z$ 的开邻域。
由于 $\mathcal{I}$ 的高阶 {\v C}ech 上同调群为零
（引理 \ref{cohomology-lemma-injective-trivial-cech}），
$s'$ 是上边缘。
因此，$\xi$ 在开覆盖
$Z=\bigcup Z\cap U'_i$ 的 {\v C}ech 复形中的像是上边缘，证明完成。
\end{proof}






\section{基变换映射}
\label{cohomology-section-base-change-map}

\noindent
在若干情形下，我们需要知道如何构造基变换映射。由于尚未讨论
导出拉回，这里只讨论沿环空间的平坦态射作基变换的情形。
在陈述结果之前，先讨论导出范畴上的平坦拉回。具体而言，设
$g : X \to Y$ 为环空间的平坦态射。由《模层》引理
\ref{modules-lemma-pullback-flat}，函子
$g^* : \textit{Mod}(\mathcal{O}_Y) \to
\textit{Mod}(\mathcal{O}_X)$ 正合。因此它有导出函子
$$
g^* : D^{+}(Y) \to D^{+}(X)
$$
其计算只需拉回 $D^{+}(Y)$ 中给定对象的一个代表，见《导出范畴》
引理 \ref{derived-lemma-right-derived-exact-functor}。因此，如上所示，
我们将此函子记为 $g^*$ 而非 $Lg^*$。

\begin{lemma}
\label{cohomology-lemma-base-change-map-flat-case}
设
$$
\xymatrix{
X' \ar[r]_{g'} \ar[d]_{f'} &
X \ar[d]^f \\
S' \ar[r]^g &
S
}
$$
为环空间的交换图。设 $\mathcal{F}^\bullet$ 为
$\mathcal{O}_X$-模层的下有界复形。假设 $g$ 与 $g'$ 均平坦。
则在 $D^{+}(S')$ 中存在典范基变换映射
$$
g^*Rf_*\mathcal{F}^\bullet
\longrightarrow
R(f')_*(g')^*\mathcal{F}^\bullet
$$
\end{lemma}

\begin{proof}
取内射消解 $\mathcal{F}^\bullet \to \mathcal{I}^\bullet$
与 $(g')^*\mathcal{F}^\bullet \to \mathcal{J}^\bullet$。
由引理 \ref{cohomology-lemma-pushforward-injective-flat}，
$(g')_*\mathcal{J}^\bullet$ 是表示
$R(g')_*(g')^*\mathcal{F}^\bullet$ 的内射复形。因此，由《导出范畴》
引理 \ref{derived-lemma-morphisms-lift}
与 \ref{derived-lemma-morphisms-equal-up-to-homotopy}，图
$$
\xymatrix{
(g')_*(g')^*\mathcal{F}^\bullet \ar[r] &
(g')_*\mathcal{J}^\bullet \\
\mathcal{F}^\bullet \ar[u]^{adjunction} \ar[r] &
\mathcal{I}^\bullet \ar[u]_\beta
}
$$
中的箭头 $\beta$ 存在且在同伦意义下唯一。取到 $S$ 上的直像，得到
$$
f_*\beta :
f_*\mathcal{I}^\bullet
\longrightarrow
f_*(g')_*\mathcal{J}^\bullet
=
g_*(f')_*\mathcal{J}^\bullet
$$
由 $g^*$ 与 $g_*$ 的伴随性，得到复形态射
$g^*f_*\mathcal{I}^\bullet \to (f')_*\mathcal{J}^\bullet$。
注意，此映射在同伦意义下唯一，因为整个过程中唯一的选择是映射
$\beta$，而且所有构造均在复形层面完成。
\end{proof}

\begin{remark}
\label{cohomology-remark-correct-version-base-change-map}
基变换映射的“正确”版本是映射
$$
Lg^* Rf_* \mathcal{F}^\bullet
\longrightarrow
R(f')_* L(g')^*\mathcal{F}^\bullet.
$$
该映射的构造涉及无界复形，见评注 \ref{cohomology-remark-base-change}。
\end{remark}





\section{拓扑中的真基变换}
\label{cohomology-section-proper-base-change}

\noindent
本节证明拓扑中真基变换定理的一个非常一般的版本。它说明
高阶直像 $R^pf_*$ 的茎可在纤维上计算。

\begin{lemma}
\label{cohomology-lemma-proper-base-change}
设 $f : (X, \mathcal{O}_X) \to (Y, \mathcal{O}_Y)$ 为环空间态射，
且 $y \in Y$。假设
\begin{enumerate}
\item $f$ 是闭映射，
\item $f$ 是分离的，并且
\item $f^{-1}(y)$ 拟紧。
\end{enumerate}
则对 $D^+(\mathcal{O}_X)$ 中的 $E$，在
$D^+(\mathcal{O}_{Y, y})$ 中有
$(Rf_*E)_y = R\Gamma(f^{-1}(y), E|_{f^{-1}(y)})$。
\end{lemma}

\begin{proof}
引理 \ref{cohomology-lemma-base-change-map-flat-case} 的基变换映射给出典范映射
$(Rf_*E)_y \to R\Gamma(f^{-1}(y), E|_{f^{-1}(y)})$。
为证明此映射为同构，以内射对象的下有界复形
$\mathcal{I}^\bullet$ 表示 $E$。置 $Z = f^{-1}(\{y\})$。引理 \ref{cohomology-lemma-cohomology-of-closed} 的假设成立，见《拓扑》
引理 \ref{topology-lemma-separated}。因此各限制
$\mathcal{I}^n|_Z$ 对 $\Gamma(Z, -)$ 无环。于是
$R\Gamma(Z, E|_Z)$ 由复形
$\Gamma(Z, \mathcal{I}^\bullet|_Z)$ 表示，见《导出范畴》
引理 \ref{derived-lemma-leray-acyclicity}。换言之，我们须证明映射
$$
\colim_V \mathcal{I}^\bullet(f^{-1}(V))
\longrightarrow
\Gamma(Z, \mathcal{I}^\bullet|_Z)
$$
为同构。利用引理 \ref{cohomology-lemma-cohomology-of-closed} 可知，只需证明
$Z = f^{-1}(\{y\})$ 的开邻域 $f^{-1}(V)$ 所成的集合族在全部开邻域
组成的系统中共尾。若 $f^{-1}(\{y\}) \subset U$ 是一个开邻域，
则由于 $f$ 为闭映射，集合
$V = Y \setminus f(X \setminus U)$ 是 $y$ 的开邻域，且
$f^{-1}(V) \subset U$。这就证明了引理。
\end{proof}

\begin{theorem}[真基变换]
\label{cohomology-theorem-proper-base-change}
\begin{reference}
\cite[Expose V bis, 4.1.1]{SGA4}
\end{reference}
考虑拓扑空间的笛卡尔方块
$$
\xymatrix{
X' = Y' \times_Y X \ar[d]_{f'} \ar[r]_-{g'} & X \ar[d]^f \\
Y' \ar[r]^g & Y
}
$$
假设 $f$ 为真映射。设 $E$ 为 $D^+(X)$ 中的对象。则引理
\ref{cohomology-lemma-base-change-map-flat-case} 的基变换映射
$$
g^{-1}Rf_*E \longrightarrow Rf'_*(g')^{-1}E
$$
是 $D^+(Y')$ 中的同构。
\end{theorem}

\begin{proof}
设 $y' \in Y'$ 为像是 $y \in Y$ 的点。只需证明基变换映射在
$y'$ 处的茎上诱导同构。由于 $f$ 为真映射，可知 $f'$ 也是真映射，
$f$ 与 $f'$ 的纤维均拟紧，并且 $f$ 与 $f'$ 均为闭映射；见《拓扑》
定理 \ref{topology-theorem-characterize-proper} 与引理
\ref{topology-lemma-base-change-separated}。因而可两次应用引理
\ref{cohomology-lemma-proper-base-change}，得到
$$
(Rf'_*(g')^{-1}E)_{y'} = R\Gamma((f')^{-1}(y'), (g')^{-1}E|_{(f')^{-1}(y')})
$$
以及
$$
(Rf_*E)_y = R\Gamma(f^{-1}(y), E|_{f^{-1}(y)})
$$
纤维间的诱导映射 $(f')^{-1}(y') \to f^{-1}(y)$ 是拓扑空间的同胚，
而 $E|_{f^{-1}(y)}$ 的拉回为
$(g')^{-1}E|_{(f')^{-1}(y')}$。所需结论随即成立。
\end{proof}

\begin{lemma}[集合层的真基变换]
\label{cohomology-lemma-proper-base-change-sheaves-of-sets}
考虑拓扑空间的笛卡尔方块
$$
\xymatrix{
X' \ar[d]_{f'} \ar[r]_-{g'} & X \ar[d]^f \\
Y' \ar[r]^g & Y
}
$$
假设 $f$ 为真映射。则对 $X$ 上的任意集合层 $\mathcal{F}$，有
$g^{-1}f_*\mathcal{F} = f'_*(g')^{-1}\mathcal{F}$。
\end{lemma}

\begin{proof}
完全仿照定理 \ref{cohomology-theorem-proper-base-change} 的证明，可知只需证明
$(f_*\mathcal{F})_y = \Gamma(X_y, \mathcal{F}|_{X_y})$。
再仿照引理 \ref{cohomology-lemma-proper-base-change} 的论证，将此归结为集合层情形下引理 \ref{cohomology-lemma-cohomology-of-closed} 的 $p = 0$ 情形。引理
\ref{cohomology-lemma-cohomology-of-closed} 证明的第一部分同样适用于集合层，
从而完成证明。略去若干细节。
\end{proof}



\section{上同调与余极限}
\label{cohomology-section-limits}

\noindent
设 $X$ 为环空间。设 $(\mathcal{F}_i, \varphi_{ii'})$ 为由有向集
$I$ 指标的 $\mathcal{O}_X$-模层系统，见《范畴》第
\ref{categories-section-posets-limits} 节。由于对每个 $i$ 都有典范映射
$\mathcal{F}_i \to \colim_i \mathcal{F}_i$，我们对每个 $p \geq 0$
得到典范映射
$$
\colim_i H^p(X, \mathcal{F}_i)
\longrightarrow
H^p(X, \colim_i \mathcal{F}_i)
$$
对每个开集 $U \subset X$ 当然也有类似映射。一般而言，这些映射并非
同构，即使 $p = 0$ 也是如此。本节推广《层》引理
\ref{sheaves-lemma-directed-colimits-sections} 的结果。另见《模层》
引理 \ref{modules-lemma-finite-presentation-quasi-compact-colimit}
（取特殊情形 $\mathcal{G} = \mathcal{O}_X$）。

\begin{lemma}
\label{cohomology-lemma-quasi-separated-cohomology-colimit}
设 $X$ 为环空间。假设 $X$ 的底拓扑空间具有下列性质：
\begin{enumerate}
\item 存在由拟紧开子集组成的拓扑基，并且
\item 任意两个拟紧开集的交仍拟紧。
\end{enumerate}
则对 $\mathcal{O}_X$-模层的任意有向系统
$(\mathcal{F}_i, \varphi_{ii'})$ 以及任意拟紧开集 $U \subset X$，
典范映射
$$
\colim_i H^q(U, \mathcal{F}_i)
\longrightarrow
H^q(U, \colim_i \mathcal{F}_i)
$$
对每个 $q \geq 0$ 都是同构。
\end{lemma}

\begin{proof}
本证明的关键是同时对所有拟紧开集 $U \subset X$ 论证。由《层》引理
\ref{sheaves-lemma-directed-colimits-sections}（结合《拓扑》引理
\ref{topology-lemma-topology-quasi-separated-scheme}），当 $q = 0$
时结论对任意拟紧开集 $U \subset X$ 成立。假设已经对所有
$q \leq q_0$ 证明了结论；下面证明 $q = q_0 + 1$ 的情形。

\medskip\noindent
按我们关于有向系统的约定，指标集 $I$ 是有向的，而 $I$ 上的任意
$\mathcal{O}_X$-模层系统 $(\mathcal{F}_i, \varphi_{ii'})$ 都是有向系统。
由《内射对象》引理 \ref{injectives-lemma-sheaves-modules-space}，
$\mathcal{O}_X$-模层范畴有函子性内射嵌入。因此，对任意系统
$(\mathcal{F}_i, \varphi_{ii'})$，存在系统
$(\mathcal{I}_i, \varphi_{ii'})$，其中每个 $\mathcal{I}_i$ 都是内射
$\mathcal{O}_X$-模层，并且存在由内射 $\mathcal{O}_X$-模同态
$\mathcal{F}_i \to \mathcal{I}_i$ 给出的系统态射。记其余核为
$\mathcal{Q}_i$，于是有短正合列
$$
0 \to
\mathcal{F}_i \to
\mathcal{I}_i \to
\mathcal{Q}_i \to 0.
$$
我们断言序列
$$
0 \to
\colim_i \mathcal{F}_i \to
\colim_i \mathcal{I}_i \to
\colim_i \mathcal{Q}_i \to 0.
$$
也是 $\mathcal{O}_X$-模层的短正合列。可在茎上检验这一点。由《层》第
\ref{sheaves-section-limits-presheaves} 节与第
\ref{sheaves-section-limits-sheaves} 节，取茎与余极限可交换。由于
阿贝尔群短正合列的有向余极限仍短正合（见《代数》引理
\ref{algebra-lemma-directed-colimit-exact}），断言成立。我们还断言：
对每个拟紧开集 $U \subset X$ 及每个 $q \geq 1$，有
$H^q(U, \colim_i \mathcal{I}_i) = 0$。暂且接受这一断言，考虑图
$$
\xymatrix{
\colim_i H^{q_0}(U, \mathcal{I}_i) \ar[d] \ar[r] &
\colim_i H^{q_0}(U, \mathcal{Q}_i) \ar[d] \ar[r] &
\colim_i H^{q_0 + 1}(U, \mathcal{F}_i) \ar[d] \ar[r] &
0 \ar[d] \\
H^{q_0}(U, \colim_i \mathcal{I}_i) \ar[r] &
H^{q_0}(U, \colim_i \mathcal{Q}_i) \ar[r] &
H^{q_0 + 1}(U, \colim_i \mathcal{F}_i) \ar[r] &
0
}
$$
右下角的零来自上述断言，右上角的零则来自各层
$\mathcal{I}_i$ 的内射性。应用《代数》引理
\ref{algebra-lemma-directed-colimit-exact} 可知上行正合。因此由蛇形引理
得到 $q = q_0 + 1$ 的结论。

\medskip\noindent
只剩证明该断言。我们将使用引理 \ref{cohomology-lemma-cech-vanish-basis}。
令 $\mathcal{B}$ 为 $X$ 的全部拟紧开子集所成的集合。按假设，
这是 $X$ 的拓扑基。令 $\text{Cov}$ 为有限开覆盖
$\mathcal{U} : U = \bigcup_{j = 1, \ldots, m} U_j$ 所成的集合，其中
$U$ 与各 $U_j$ 均为 $X$ 中的拟紧开集。由 $q = 0$ 的结论，对
$\mathcal{U} \in \text{Cov}$ 有
$$
\check{\mathcal{C}}^\bullet(\mathcal{U}, \colim_i \mathcal{I}_i)
=
\colim_i \check{\mathcal{C}}^\bullet(\mathcal{U}, \mathcal{I}_i)
$$
因为所有多重交 $U_{j_0 \ldots j_p}$ 都拟紧。由引理
\ref{cohomology-lemma-injective-trivial-cech}，{\v C}ech 复形余极限中的每个复形
在次数 $\geq 1$ 都无环。因此由《代数》引理
\ref{algebra-lemma-directed-colimit-exact}，{\v C}ech 复形
$\check{\mathcal{C}}^\bullet(\mathcal{U}, \colim_i \mathcal{I}_i)$
在次数 $\geq 1$ 也无环。换言之，对所有 $p \geq 1$，有
$\check{H}^p(\mathcal{U},  \colim_i \mathcal{I}_i) = 0$。于是满足引理
\ref{cohomology-lemma-cech-vanish-basis} 的假设，断言得证。
\end{proof}

\begin{lemma}
\label{cohomology-lemma-higher-direct-image-colimit}
设 $f : X \to Y$ 为拓扑空间的连续映射。设
$(\mathcal{F}_i, \varphi_{ii'})$ 为 $X$ 上的阿贝尔层系统。置
$\mathcal{F} = \colim \mathcal{F}_i$。设 $p \geq 0$ 为整数。假设使得
$H^p(f^{-1}(V), \mathcal{F}) = \colim H^p(f^{-1}(V), \mathcal{F}_i)$
成立的开集 $V \subset Y$ 组成 $Y$ 的一个拓扑基。则
$R^pf_*\mathcal{F} = \colim R^pf_*\mathcal{F}_i$.
\end{lemma}

\begin{proof}
回忆 $R^pf_*\mathcal{F}$ 是预层 $\mathcal{G}$ 的层化，其中
$\mathcal{G}$ 将 $V$ 映为 $H^p(f^{-1}(V), \mathcal{F})$；见引理
\ref{cohomology-lemma-describe-higher-direct-images}。类似地，
$R^pf_*\mathcal{F}_i$ 是预层 $\mathcal{G}_i$ 的层化，其中
$\mathcal{G}_i$ 将 $V$ 映为 $H^p(f^{-1}(V), \mathcal{F}_i)$。
回忆层化函子是从层到预层的包含函子的左伴随，见《层》第
\ref{sheaves-section-sheafification} 节。因此层化与余极限可交换，
见《范畴》引理 \ref{categories-lemma-adjoint-exact}。故只需证明预层态射
（其中余极限在预层范畴中取）
$$
\colim \mathcal{G}_i \longrightarrow \mathcal{G}
$$
在层化后诱导同构。为此，只需证明预层 $\mathcal{G}$ 与
$\colim \mathcal{G}_i$ 在 $Y$ 的一个拓扑基上一致。事实上，此时两者
层化后的茎一致，因为这些茎可直接由预层在该基的元素上的取值计算。
所需的一致性恰为引理的假设。
\end{proof}

\noindent
下面给出《层》引理 \ref{sheaves-lemma-descend-opens} 的上同调版本。
设 $X$ 为谱拓扑空间，它写成谱拓扑空间 $X_i$ 的余滤过极限，相应图的
转移态射均为谱态射，如《拓扑》引理
\ref{topology-lemma-directed-inverse-limit-spectral-spaces} 所述。假设给定
\begin{enumerate}
\item 对每个 $i \in \Ob(\mathcal{I})$，给定 $X_i$ 上的阿贝尔层
$\mathcal{F}_i$；
\item 对每个 $a : j \to i$，给定阿贝尔层的一个 $f_a$-映射
$\varphi_a : \mathcal{F}_i \to \mathcal{F}_j$（见《层》定义
\ref{sheaves-definition-f-map}）。
\end{enumerate}
并且只要 $c = a \circ b$，就有
$\varphi_c = \varphi_b \circ \varphi_a$。在 $X$ 上置
$\mathcal{F} = \colim p_i^{-1}\mathcal{F}_i$。

\begin{lemma}
\label{cohomology-lemma-colimit}
在上述情形下，设 $i \in \Ob(\mathcal{I})$，并设
$U_i \subset X_i$ 为拟紧开集。则
$$
\colim_{a : j \to i} H^p(f_a^{-1}(U_i), \mathcal{F}_j) =
H^p(p_i^{-1}(U_i), \mathcal{F})
$$
对所有 $p \geq 0$ 成立。特别地，有
$H^p(X, \mathcal{F}) = \colim H^p(X_i, \mathcal{F}_i)$.
\end{lemma}

\begin{proof}
当 $p = 0$ 时，这就是《层》引理 \ref{sheaves-lemma-descend-opens}。

\medskip\noindent
本段证明可以找到系统态射
$(\gamma_i) : (\mathcal{F}_i, \varphi_a) \to (\mathcal{G}_i, \psi_a)$，
其中 $\mathcal{G}_i$ 为内射阿贝尔层，且 $\gamma_i$ 为单射。对每个
$i$，选取单射 $\mathcal{F}_i \to \mathcal{I}_i$，其中
$\mathcal{I}_i$ 是 $X_i$ 上的内射阿贝尔层。于是可考虑映射族
$$
\gamma_i :
\mathcal{F}_i
\longrightarrow
\prod\nolimits_{b : k \to i} f_{b, *}\mathcal{I}_k = \mathcal{G}_i
$$
其分量映射是映射
$f_b^{-1}\mathcal{F}_i \to \mathcal{F}_k \to \mathcal{I}_k$
的伴随映射。对 $\mathcal{I}$ 中的 $a : j \to i$，有典范映射
$$
\psi_a : f_a^{-1}\mathcal{G}_i \to \mathcal{G}_j
$$
其对应于 $b : k \to j$ 的分量是典范映射
$f_b^{-1}f_{a \circ b, *}\mathcal{I}_k \to f_{b, *}\mathcal{I}_k$。
因而得到系统的单射
$\{\gamma_i\} : \{\mathcal{F}_i, \varphi_a) \to (\mathcal{G}_i, \psi_a)$
。注意 $\mathcal{G}_i$ 是 $X_i$ 上的内射阿贝尔群层，见引理
\ref{cohomology-lemma-pushforward-injective-flat} 与《同调代数》引理
\ref{homology-lemma-product-injectives}。构造至此完成。

\medskip\noindent
完全仿照引理 \ref{cohomology-lemma-quasi-separated-cohomology-colimit} 的证明，
可知只需证明当 $p > 0$ 时
$H^p(X, \colim f_i^{-1}\mathcal{G}_i) = 0$。

\medskip\noindent
置 $\mathcal{G} = \colim f_i^{-1}\mathcal{G}_i$。为证明
$\mathcal{G}$ 在 $X$ 的每个拟紧开集上的上同调消失，只需证明：
对任意由有限个拟紧开集覆盖 $X$ 的一个拟紧开集的开覆盖
$\mathcal{U}$，$\mathcal{G}$ 的 {\v C}ech 上同调为零；见引理
\ref{cohomology-lemma-cech-vanish-basis}。由《拓扑》引理
\ref{topology-lemma-descend-opens}，这样的覆盖是某个 $X_i$ 上同类覆盖
$\mathcal{U}_i$ 在 $p_i$ 下的原像。由 $p = 0$ 的情形，有
$$
\check{\mathcal{C}}^\bullet(\mathcal{U}, \mathcal{G}) =
\colim_{a : j \to i}
\check{\mathcal{C}}^\bullet(f_a^{-1}(\mathcal{U}_i), \mathcal{G}_j)
$$
右端是复形的滤过余极限，而由引理
\ref{cohomology-lemma-injective-trivial-cech}，其中每个复形在正次数都无环。
因此由《代数》引理 \ref{algebra-lemma-directed-colimit-exact} 得到结论。
\end{proof}









\section{诺特拓扑空间上的消失}
\label{cohomology-section-vanishing-Noetherian}

\noindent
本节旨在证明 Grothendieck 的一个定理，即命题
\ref{cohomology-proposition-vanishing-Noetherian}。见 \cite{Tohoku}。

\begin{lemma}
\label{cohomology-lemma-cohomology-and-closed-immersions}
设 $i : Z \to X$ 为拓扑空间的闭浸入。
对 $Z$ 上任意阿贝尔层 $\mathcal{F}$，有
$H^p(Z, \mathcal{F}) = H^p(X, i_*\mathcal{F})$.
\end{lemma}

\begin{proof}
这是因为 $i_*$ 正合（见《模层》引理
\ref{modules-lemma-i-star-exact}），因而作为函子有
$R^pi_* = 0$（《导出范畴》引理
\ref{derived-lemma-right-derived-exact-functor}）。
故可应用引理 \ref{cohomology-lemma-apply-Leray}。
\end{proof}

\begin{lemma}
\label{cohomology-lemma-irreducible-constant-cohomology-zero}
设 $X$ 为不可约拓扑空间。则对任意阿贝尔群 $A$
及所有 $p > 0$，有 $H^p(X, \underline{A}) = 0$。
\end{lemma}

\begin{proof}
回忆 $\underline{A}$ 是《层》定义
\ref{sheaves-definition-constant-sheaf} 中所定义的常值层。
由于 $X$ 不可约，任一非空开集 $U$ 都不可约，因而尤其连通。
故对任意非空开集 $U \subset X$，有 $\underline{A}(U) = A$。
又有 $\underline{A}(\emptyset) = 0$。因此 $\underline{A}$
是 $X$ 上的松弛阿贝尔层，所述消失性由引理
\ref{cohomology-lemma-flasque-acyclic} 得到。
\end{proof}

\begin{lemma}
\label{cohomology-lemma-subsheaf-of-constant-sheaf}
\begin{reference}
\cite[Page 168]{Tohoku}.
\end{reference}
设 $X$ 为一个拓扑空间，且任意两个拟紧开集之交仍拟紧。设
$\mathcal{F} \subset \underline{\mathbf{Z}}$
为由拟紧开集上的有限多个截面生成的子层。则存在由阿贝尔子层组成的有限滤过
$$
(0) = \mathcal{F}_0 \subset \mathcal{F}_1 \subset \ldots \subset
\mathcal{F}_n = \mathcal{F}
$$
使得对每个 $0 < i \leq n$，都存在短正合列
$$
0 \to j'_!\underline{\mathbf{Z}}_V \to j_!\underline{\mathbf{Z}}_U \to
\mathcal{F}_i/\mathcal{F}_{i - 1} \to 0
$$
其中 $j : U \to X$ 与 $j' : V \to X$ 是拟紧开集到 $X$ 的包含映射。
\end{lemma}

\begin{proof}
设 $\mathcal{F}$ 由拟紧开集 $U_1, \ldots, U_t$ 上的截面
$s_1, \ldots, s_t$ 生成。由于 $U_i$ 拟紧且 $s_i$ 是到
$\mathbf{Z}$ 的局部常值函数，必要时可将 $U_i$ 换成其某个有限开闭分解的各部分，
从而不妨设 $s_i$ 为常值截面。
设 $s_i = n_i$，其中 $n_i \in \mathbf{Z}$。若 $n_i = 0$，
当然可从表中删去 $(U_i, n_i)$。必要时变号后，还可假设 $n_i > 0$。
其次，对任意子集 $I \subset \{1, \ldots, t\}$，可把
$\bigcap_{i \in I} U_i$ 与 $\gcd(n_i, i \in I)$ 加入表中。
这样处理后，表 $(U_1, n_1), \ldots, (U_t, n_t)$ 具有如下性质：
对 $x \in X$，令 $I_x = \{i \in \{1, \ldots, t\} \mid x \in U_i\}$。
则对某个 $i \in I_x$，$\gcd(n_i, i \in I_x)$ 等于 $n_i$。

\medskip\noindent
取 $\mathcal{F}_0 = (0)$，并令 $\mathcal{F}_n$ 由所有满足
$n_i \leq n$ 的 $i$ 所对应的 $U_i$ 上截面 $n_i$ 生成，以此作为所需滤过。
显然当 $n \gg 0$ 时 $\mathcal{F}_n = \mathcal{F}$。此外，商
$\mathcal{F}_n/\mathcal{F}_{n - 1}$ 由
$U = \bigcup_{n_i \leq n} U_i$ 上的截面 $n$ 生成，而映射
$j_!\underline{\mathbf{Z}}_U \to \mathcal{F}_n/\mathcal{F}_{n - 1}$
的核由 $V = \bigcup_{n_i \leq n - 1} U_i$ 上的截面 $n$ 生成。
因而得到引理陈述中的短正合列。
\end{proof}

\begin{lemma}
\label{cohomology-lemma-vanishing-generated-one-section}
\begin{reference}
这是 \cite[Proposition 3.6.1]{Tohoku} 的一个特例。
\end{reference}
设 $X$ 为拓扑空间，$d \geq 0$ 为整数。假设
\begin{enumerate}
\item $X$ 拟紧；
\item 拟紧开集构成 $X$ 的一组基；
\item 两个拟紧开集之交拟紧；
\item 对所有 $p > d$ 以及任意拟紧开集的包含映射
$j : U \to X$，有 $H^p(X, j_!\underline{\mathbf{Z}}_U) = 0$。
\end{enumerate}
则对 $X$ 上任意阿贝尔层 $\mathcal{F}$ 及所有 $p > d$，
有 $H^p(X, \mathcal{F}) = 0$。
\end{lemma}

\begin{proof}
令 $S = \coprod_{U \subset X} \mathcal{F}(U)$，其中 $U$ 遍历
$X$ 的拟紧开集。对任一有限子集
$A = \{s_1, \ldots, s_n\} \subset S$，令 $\mathcal{F}_A$
为由所有 $s_i$ 生成的 $\mathcal{F}$ 的子层（见《模层》定义
\ref{modules-definition-generated-by-local-sections}）。
注意，若 $A \subset A'$，则 $\mathcal{F}_A \subset \mathcal{F}_{A'}$。
因此 $\{\mathcal{F}_A\}$ 构成以 $S$ 的有限子集所成有向偏序集为指标的系统。
由《模层》引理 \ref{modules-lemma-generated-by-local-sections-stalk}，
观察各茎即知
$$
\colim_A \mathcal{F}_A = \mathcal{F}
$$
由引理 \ref{cohomology-lemma-quasi-separated-cohomology-colimit}，有
$$
H^p(X, \mathcal{F}) =
\colim_A H^p(X, \mathcal{F}_A)
$$
因此只需对阿贝尔层 $\mathcal{F}_A$ 证明所述消失性。换言之，
只需处理 $\mathcal{F}$ 由 $X$ 的拟紧开集上的有限多个局部截面生成的情形。

\medskip\noindent
设 $\mathcal{F}$ 由局部截面 $s_1, \ldots, s_n$ 生成，
并令 $\mathcal{F}' \subset \mathcal{F}$ 为由
$s_1, \ldots, s_{n - 1}$ 生成的子层。于是有短正合列
$$
0 \to \mathcal{F}' \to \mathcal{F} \to \mathcal{F}/\mathcal{F}' \to 0
$$
由上同调长正合列可知，只需对由少于 $n$ 个局部截面生成的阿贝尔层
$\mathcal{F}'$ 与 $\mathcal{F}/\mathcal{F}'$ 证明消失性。
故只需对至多由一个局部截面生成的层证明消失性。这样的层恰好是
$j_!\underline{\mathbf{Z}}_U$ 的商，其中 $U$ 是 $X$ 的拟紧开集。

\medskip\noindent
现设有短正合列
$$
0 \to \mathcal{K} \to j_!\underline{\mathbf{Z}}_U \to \mathcal{F} \to 0
$$
其中 $U$ 是 $X$ 的拟紧开集。只需证明当 $q \geq d + 1$ 时
$H^q(X, \mathcal{K})$ 为零。与上面一样，可将 $\mathcal{K}$
写成由拟紧开集上的有限多个截面生成的子层 $\mathcal{K}'$ 的滤过余极限。
于是 $\mathcal{F}$ 是各层
$j_!\underline{\mathbf{Z}}_U/\mathcal{K}'$ 的滤过余极限。
由此约化到 $\mathcal{K}$ 由拟紧开集上的有限多个截面生成的情形。
注意 $\mathcal{K}$ 是 $\underline{\mathbf{Z}}_X$ 的子层。
故由引理 \ref{cohomology-lemma-subsheaf-of-constant-sheaf}，$\mathcal{K}$
存在有限滤过，其逐次商 $\mathcal{Q}$ 嵌入短正合列
$$
0 \to j''_!\underline{\mathbf{Z}}_W \to
j'_!\underline{\mathbf{Z}}_V \to \mathcal{Q} \to 0
$$
其中 $j'' : W \to X$ 与 $j' : V \to X$ 是拟紧开集的包含映射。
因此，由本引理关于 $j''_!\underline{\mathbf{Z}}_W$
与 $j'_!\underline{\mathbf{Z}}_V$ 的上同调群消失的假设，
当 $p > d$ 时有 $H^p(X, \mathcal{Q}) = 0$。
回到 $\mathcal{K}$，沿上述滤过运用上同调长正合列归纳，
也得到所需的消失性。
\end{proof}

\begin{example}
\label{cohomology-example-datta}
令 $X = \mathbf{N}$，赋予它如下拓扑：开集为 $\emptyset$、$X$
以及 $n \geq 1$ 时的 $U_n = \{i \mid i \leq n\}$。
$X$ 上的阿贝尔层 $\mathcal{F}$ 等同于阿贝尔群逆系统
$A_n = \mathcal{F}(U_n)$，并且
$\Gamma(X, \mathcal{F}) = \lim A_n$。由于逆极限函子在逆系统范畴上不正合，
可知存在 $H^1$ 非零的阿贝尔层。
最后，读者可验证：若 $j : U = U_n \to X$ 为包含映射，则
$p \geq 1$ 时 $H^p(X, j_!\mathbf{Z}_U) = 0$。
因此，$X$ 是一个在 $d = 0$ 时满足引理
\ref{cohomology-lemma-vanishing-generated-one-section} 的条件 (2)、(3)、(4)
却不满足结论的空间。
\end{example}

\begin{lemma}
\label{cohomology-lemma-subsheaf-irreducible}
设 $X$ 为不可约拓扑空间，并设
$\mathcal{H} \subset \underline{\mathbf{Z}}$ 为该常值层的阿贝尔子层。
则存在非空开集 $U \subset X$，使得对某个 $d \in \mathbf{Z}$ 有
$\mathcal{H}|_U = \underline{d\mathbf{Z}}_U$。
\end{lemma}

\begin{proof}
回忆，对 $X$ 的任意非空开集 $V$，有
$\underline{\mathbf{Z}}(V) = \mathbf{Z}$（见引理
\ref{cohomology-lemma-irreducible-constant-cohomology-zero} 的证明）。
若 $\mathcal{H} = 0$，则取 $d = 0$ 即得引理。
若 $\mathcal{H} \not = 0$，则存在非空开集 $U \subset X$，
使得 $\mathcal{H}(U) \not = 0$。设对某个 $n \geq 1$ 有
$\mathcal{H}(U) = n\mathbf{Z}$。于是
$\underline{n\mathbf{Z}}_U
\subset \mathcal{H}|_U \subset
\underline{\mathbf{Z}}_U$。若第一个包含严格，则可找到非空开集
$U' \subset U$ 与整数 $1 \leq n' < n$，使得
$\underline{n'\mathbf{Z}}_{U'}
\subset \mathcal{H}|_{U'} \subset
\underline{\mathbf{Z}}_{U'}$。
此过程必在有限步后停止，因而得到引理。
\end{proof}

\begin{proposition}[Grothendieck]
\label{cohomology-proposition-vanishing-Noetherian}
\begin{reference}
\cite[Theorem 3.6.5]{Tohoku}.
\end{reference}
设 $X$ 为诺特拓扑空间。若 $\dim(X) \leq d$，则对 $X$ 上任意阿贝尔层
$\mathcal{F}$ 及所有 $p > d$，有 $H^p(X, \mathcal{F}) = 0$。
\end{proposition}

\begin{proof}
对 $d$ 归纳证明此引理。固定 $d$，并假设对所有维数 $< d$
的诺特拓扑空间，此引理均成立。

\medskip\noindent
设 $\mathcal{F}$ 为 $X$ 上的阿贝尔层。设 $U \subset X$ 为开集，
并以 $Z \subset X$ 表示其闭补集。以
$j : U \to X$ 与 $i : Z \to X$ 表示包含映射。于是有短正合列
$$
0 \to j_{!}j^*\mathcal{F} \to \mathcal{F} \to i_*i^*\mathcal{F} \to 0
$$
见《模层》引理 \ref{modules-lemma-canonical-exact-sequence}。
注意 $j_!j^*\mathcal{F}$ 支撑在 $U$ 的拓扑闭包 $Z'$ 上；
换言之，它形如 $i'_*\mathcal{F}'$，其中 $\mathcal{F}'$
是 $Z'$ 上的某个阿贝尔层，而 $i' : Z' \to X$ 是包含映射。

\medskip\noindent
利用这一点可约化到 $X$ 不可约的情形。事实上，由《拓扑》引理
\ref{topology-lemma-Noetherian}，$X$ 只有有限多个不可约分支。
若 $X$ 有多于一个不可约分支，取 $X$ 的一个不可约分支
$Z \subset X$，并令 $U = X \setminus Z$。由上述结论及上同调长正合列，
只需证明当 $p > d$ 时
$H^p(X, i_*i^*\mathcal{F})$ 与 $H^p(X, i'_*\mathcal{F}')$ 消失。
由引理 \ref{cohomology-lemma-cohomology-and-closed-immersions}，又只需证明当
$p > d$ 时 $H^p(Z, i^*\mathcal{F})$ 与
$H^p(Z', \mathcal{F}')$ 消失。由于 $Z'$ 和 $Z$
的不可约分支都更少，确实约化到了 $X$ 不可约的情形。

\medskip\noindent
若 $d = 0$ 且 $X$ 不可约，则 $X$ 是 $X$ 的唯一非空开子集。
因此每个层都是常值层，高阶上同调群消失（例如由引理
\ref{cohomology-lemma-irreducible-constant-cohomology-zero}）。

\medskip\noindent
设 $X$ 不可约且维数为 $d > 0$。由引理
\ref{cohomology-lemma-vanishing-generated-one-section}，可约化到
$\mathcal{F} = j_!\underline{\mathbf{Z}}_U$ 的情形，其中
$U \subset X$ 为某个开集。此时考察短正合列
$$
0 \to j_!(\underline{\mathbf{Z}}_U) \to
\underline{\mathbf{Z}}_X \to i_*\underline{\mathbf{Z}}_Z \to 0
$$
其中 $Z = X \setminus U$。由引理
\ref{cohomology-lemma-irreducible-constant-cohomology-zero}，对所有 $p \geq 1$，
$H^p(X, \underline{\mathbf{Z}}_X)$ 消失。由归纳假设，
$H^p(X, i_*\underline{\mathbf{Z}}_Z) = H^p(Z, \underline{\mathbf{Z}}_Z) = 0$
对 $p \geq d$ 成立。故由上同调长正合列即得结论。
\end{proof}





\section{支撑在闭子集上的上同调}
\label{cohomology-section-cohomology-support}

\noindent
本节只讨论最基本的内容；相关讨论将在第
\ref{cohomology-section-cohomology-support-bis} 节继续。

\medskip\noindent
设 $X$ 为拓扑空间，$Z \subset X$ 为闭子集，并设
$\mathcal{F}$ 为 $X$ 上的阿贝尔层。令
$$
\Gamma_Z(X, \mathcal{F}) =
\{s \in \mathcal{F}(X) \mid \text{Supp}(s) \subset Z\}
$$
为支集包含在 $Z$ 中的截面所成子集。截面的支集定义见
《模层》定义 \ref{modules-definition-support}。《模层》引理
\ref{modules-lemma-support-section-closed} 蕴含
$\Gamma_Z(X, \mathcal{F})$ 是 $\Gamma(X, \mathcal{F})$ 的子群。
同一引理还保证赋值
$\mathcal{F} \mapsto \Gamma_Z(X, \mathcal{F})$
关于 $\mathcal{F}$ 构成函子。此函子左正合，但一般并不正合。

\medskip\noindent
由于阿贝尔层范畴有足够多的内射对象（《内射对象》引理
\ref{injectives-lemma-abelian-sheaves-space}），可得到右导出函子
$$
R\Gamma_Z(X, -) : D^+(X) \longrightarrow D^+(\textit{Ab})
$$
见《导出范畴》引理
\ref{derived-lemma-enough-injectives-right-derived}。
为计算 $R\Gamma_Z(X, -)$ 在对象 $K$ 上的值，用内射阿贝尔层的下有界复形
$\mathcal{I}^\bullet$ 表示 $K$，再取
$\Gamma_Z(X, \mathcal{I}^\bullet)$；见《导出范畴》引理
\ref{derived-lemma-injective-acyclic}。阿贝尔层 $\mathcal{F}$
支撑在 $Z$ 上的上同调群定义为
$H^q_Z(X, \mathcal{F}) = R^q\Gamma_Z(X, \mathcal{F})$。

\medskip\noindent
设 $\mathcal{I}$ 为 $X$ 上的内射阿贝尔层，并令
$U = X \setminus Z$。则限制映射
$\mathcal{I}(X) \to \mathcal{I}(U)$ 是满射（引理
\ref{cohomology-lemma-injective-restriction-surjective}），其核为
$\Gamma_Z(X, \mathcal{I})$。由此立即得到：对 $K \in D^+(X)$，
$D^+(\textit{Ab})$ 中存在特异三角
$$
R\Gamma_Z(X, K) \to R\Gamma(X, K) \to R\Gamma(U, K) \to R\Gamma_Z(X, K)[1]
$$
因而得到上同调长正合列
$$
\ldots \to H^i_Z(X, K) \to H^i(X, K) \to H^i(U, K) \to
H^{i + 1}_Z(X, K) \to \ldots
$$
它对 $D^+(X)$ 中任意 $K$ 都成立。

\medskip\noindent
对 $X$ 上的阿贝尔层 $\mathcal{F}$，可考虑{\it 支撑在 $Z$ 上的截面子层}，
记为 $\mathcal{H}_Z(\mathcal{F})$，其定义为
$$
\mathcal{H}_Z(\mathcal{F})(U) =
\{s \in \mathcal{F}(U) \mid \text{Supp}(s) \subset U \cap Z\} =
\Gamma_{Z \cap U}(U, \mathcal{F}|_U)
$$
利用《模层》引理 \ref{modules-lemma-i-star-exact} 中的等价，
可将 $\mathcal{H}_Z(\mathcal{F})$ 看作 $Z$ 上的阿贝尔层；见《模层》注
\ref{modules-remark-sections-support-in-closed}。由此得到函子
$$
\textit{Ab}(X) \longrightarrow \textit{Ab}(Z),\quad
\mathcal{F} \longmapsto
\mathcal{H}_Z(\mathcal{F})\text{（看作 }Z\text{ 上的层）}
$$
此函子左正合，但一般并不正合。与上面完全相同，可得到右导出函子
$$
R\mathcal{H}_Z : D^+(X) \longrightarrow D^+(Z)
$$
并置 $\mathcal{H}^q_Z(\mathcal{F}) = R^q\mathcal{H}_Z(\mathcal{F})$，
于是 $\mathcal{H}^0_Z(\mathcal{F}) = \mathcal{H}_Z(\mathcal{F})$。

\medskip\noindent
注意，对任意阿贝尔层 $\mathcal{F}$，有
$\Gamma_Z(X, \mathcal{F}) = \Gamma(Z, \mathcal{H}_Z(\mathcal{F}))$
由下面的引理 \ref{cohomology-lemma-sections-with-support-acyclic}，
函子 $\mathcal{H}_Z$ 将内射阿贝尔层变为关于
$\Gamma(Z, -)$ 右无环的层。因此由《导出范畴》引理
\ref{derived-lemma-grothendieck-spectral-sequence}，得到收敛的
Grothendieck 谱序列
$$
E_2^{p, q} = H^p(Z, \mathcal{H}^q_Z(K)) \Rightarrow H^{p + q}_Z(X, K)
$$
，它关于 $D^+(X)$ 中的 $K$ 具有函子性。

\begin{lemma}
\label{cohomology-lemma-sections-with-support-acyclic}
设 $i : Z \to X$ 为闭子集的包含映射，并设 $\mathcal{I}$
为 $X$ 上的内射阿贝尔层。则
$\mathcal{H}_Z(\mathcal{I})$ 是 $Z$ 上的内射阿贝尔层。
\end{lemma}

\begin{proof}
由于 $\mathcal{H}_Z(-)$ 是正合函子 $i_*$ 的右伴随，
结论由《同调代数》引理
\ref{homology-lemma-adjoint-preserve-injectives} 得到。
另见《模层》引理 \ref{modules-lemma-i-star-exact} 与
\ref{modules-lemma-i-star-right-adjoint}。
\end{proof}




\section{谱拓扑空间上的上同调}
\label{cohomology-section-spectral}

\noindent
谱拓扑空间上同调的一个关键结果是引理 \ref{cohomology-lemma-colimit}。
粗略地说，它断言上同调与谱拓扑空间范畴中的余滤过极限可交换；
该范畴定义见《拓扑》定义 \ref{topology-definition-spectral-space}。
应用这一结果，可如下得到引理 \ref{cohomology-lemma-cohomology-of-closed}
与 \ref{cohomology-lemma-proper-base-change} 的类似结论。

\begin{lemma}
\label{cohomology-lemma-cohomology-of-neighbourhoods-of-closed}
设 $X$ 为谱拓扑空间，$\mathcal{F}$ 为 $X$ 上的阿贝尔层，
$E \subset X$ 为拟紧子集。令 $W \subset X$ 为 $X$
中所有特化到 $E$ 中某点的点所成集合。
\begin{enumerate}
\item $H^p(W, \mathcal{F}|_W) = \colim H^p(U, \mathcal{F})$
，其中余极限遍历 $E$ 的拟紧开邻域；
\item $H^p(W \setminus E, \mathcal{F}|_{W \setminus E}) =
\colim H^p(U \setminus E, \mathcal{F}|_{U \setminus E})$
，只要 $E$ 是可构造子集。
\end{enumerate}
\end{lemma}

\begin{proof}
由《拓扑》引理 \ref{topology-lemma-make-spectral-space}，
$W = \lim U$，其中极限遍历包含 $E$ 的拟紧开集。由《拓扑》引理
\ref{topology-lemma-spectral-sub}，每个 $U$ 都是谱拓扑空间。
故可应用引理 \ref{cohomology-lemma-colimit}，得到 (1)。
对 (2) 作相同证明，只需改用《拓扑》引理
\ref{topology-lemma-make-spectral-space-minus}。
\end{proof}

\begin{lemma}
\label{cohomology-lemma-proper-base-change-spectral}
设 $f : X \to Y$ 为谱拓扑空间之间的谱态射，$y \in Y$。
令 $E \subset Y$ 为所有特化到 $y$ 的点所成集合，并设
$\mathcal{F}$ 为 $X$ 上的阿贝尔层。则
$(R^pf_*\mathcal{F})_y = H^p(f^{-1}(E), \mathcal{F}|_{f^{-1}(E)})$。
\end{lemma}

\begin{proof}
注意 $E = \bigcap V$，其中 $V$ 遍历 $Y$ 中 $y$ 的拟紧开邻域。
故 $f^{-1}(E) = \bigcap f^{-1}(V)$，这又蕴含作为拓扑空间有
$f^{-1}(E) = \lim f^{-1}(V)$。由于 $f$ 是谱态射，每个
$f^{-1}(V)$ 也都是谱拓扑空间（《拓扑》引理
\ref{topology-lemma-spectral-sub}）。因此 $f^{-1}(E)$ 是谱拓扑空间，
并且由引理 \ref{cohomology-lemma-colimit} 有
$$
H^p(f^{-1}(E), \mathcal{F}|_{f^{-1}(E)}) =
\colim H^p(f^{-1}(V), \mathcal{F})
$$
另一方面，$R^pf_*\mathcal{F}$ 在 $y$ 处的茎正由右端的余极限给出。
\end{proof}

\begin{lemma}
\label{cohomology-lemma-vanishing-for-profinite}
设 $X$ 为预有限拓扑空间。则对任意阿贝尔层 $\mathcal{F}$
及所有 $q > 0$，有 $H^q(X, \mathcal{F}) = 0$。
\end{lemma}

\begin{proof}
$X$ 的任意开覆盖都可由一个以开集为各部分的有限不交并分解加细；
见《拓扑》引理
\ref{topology-lemma-profinite-refine-open-covering}。
因此，若 $\mathcal{F} \to \mathcal{G}$ 是 $X$ 上阿贝尔层的满射，
则 $\mathcal{F}(X) \to \mathcal{G}(X)$ 为满射。换言之，整体截面函子正合。
故其高阶导出函子均为零；见《导出范畴》引理
\ref{derived-lemma-right-derived-exact-functor}。
\end{proof}

\noindent
下面关于上同调消失的结果改进了 Grothendieck 的结果
（命题 \ref{cohomology-proposition-vanishing-Noetherian}），可见于
\cite{Scheiderer}。

\begin{proposition}
\label{cohomology-proposition-cohomological-dimension-spectral}
\begin{reference}
第 (1) 部分是 \cite{Scheiderer} 的主定理。
\end{reference}
设 $X$ 为 Krull 维数为 $d$ 的谱拓扑空间，$\mathcal{F}$
为 $X$ 上的阿贝尔层。
\begin{enumerate}
\item 当 $q > d$ 时，$H^q(X, \mathcal{F}) = 0$；
\item 对任意拟紧开集 $U \subset X$，映射
$H^d(X, \mathcal{F}) \to H^d(U, \mathcal{F})$ 为满射；
\item 对任意可构造闭子集 $Z \subset X$，当 $q > d$ 时，
$H^q_Z(X, \mathcal{F}) = 0$。
\end{enumerate}
\end{proposition}

\begin{proof}
对 $d$ 归纳证明。

\medskip\noindent
若 $d = 0$，则 $X$ 是预有限空间；见《拓扑》引理
\ref{topology-lemma-characterize-profinite-spectral}。
故 (1) 由引理 \ref{cohomology-lemma-vanishing-for-profinite} 成立。
若 $U \subset X$ 是拟紧开集，则作为 Hausdorff 空间的拟紧子集，
$U$ 也是闭集。因此作为拓扑空间有
$X = U \amalg (X \setminus U)$，从而 (2) 成立。
对 (3) 中的 $Z$，考察长正合列
$$
H^{q - 1}(X, \mathcal{F}) \to
H^{q - 1}(X \setminus Z, \mathcal{F}) \to
H^q_Z(X, \mathcal{F}) \to H^q(X, \mathcal{F})
$$
由于 $X$ 与 $U = X \setminus Z$ 都是预有限空间（确切地说，
因 $Z$ 可构造，$U$ 拟紧），再结合 (2) 与 (1)，便得到支撑在
$Z$ 上的上同调群所需的消失性。

\medskip\noindent
以下作归纳步骤。设 $d \geq 1$，并假设本命题对所有维数 $< d$
的谱拓扑空间成立。先证明 $X$ 的 (2)。设 $U$ 为拟紧开集，
$\xi \in H^d(U, \mathcal{F})$。置 $Z = X \setminus U$，
并令 $W \subset X$ 为所有特化到 $Z$ 的点所成集合。
由引理 \ref{cohomology-lemma-cohomology-of-neighbourhoods-of-closed}，有
$$
H^d(W \setminus Z, \mathcal{F}|_{W \setminus Z}) =
\colim_{Z \subset V} H^d(V \setminus Z, \mathcal{F})
$$
其中余极限遍历 $Z$ 在 $X$ 中的拟紧开邻域 $V$。
由《拓扑》引理 \ref{topology-lemma-make-spectral-space}，
$W \setminus Z$ 是谱拓扑空间。由于 $W$ 的每一点都特化到 $Z$
中的某点，$W \setminus Z$ 是 Krull 维数 $< d$ 的谱拓扑空间。
由归纳假设，$\xi$ 在
$H^d(W \setminus Z, \mathcal{F}|_{W \setminus Z})$ 中的像为零。
由上述显示式，存在拟紧开集 $Z \subset V \subset X$，使得
$\xi|_{V \setminus Z} = 0$。由于 $V \setminus Z = V \cap U$，
对覆盖 $X = U \cup V$ 应用 Mayer--Vietoris（引理
\ref{cohomology-lemma-mayer-vietoris}），可知存在
$\tilde \xi \in H^d(X, \mathcal{F})$，其在 $U$ 上限制为 $\xi$，
在 $V$ 上限制为零。换言之，(2) 成立。

\medskip\noindent
假设 (2) 成立，证明 (1)。选取内射消解
$\mathcal{F} \to \mathcal{I}^\bullet$，并置
$$
\mathcal{G} = \Im(\mathcal{I}^{d - 1} \to \mathcal{I}^d) =
\Ker(\mathcal{I}^d \to \mathcal{I}^{d + 1})
$$
对拟紧开集 $U \subset X$，有如下正合列之间的映射：
$$
\xymatrix{
\mathcal{I}^{d - 1}(X) \ar[r] \ar[d] &
\mathcal{G}(X) \ar[r] \ar[d] &
H^d(X, \mathcal{F}) \ar[d] \ar[r] & 0 \\
\mathcal{I}^{d - 1}(U) \ar[r] &
\mathcal{G}(U) \ar[r] &
H^d(U, \mathcal{F}) \ar[r] & 0
}
$$
由引理 \ref{cohomology-lemma-injective-flasque} 及 $d \geq 1$，
层 $\mathcal{I}^{d - 1}$ 松弛。由 (2)，右侧竖直箭头为满射。
作图追踪可知映射 $\mathcal{G}(X) \to \mathcal{G}(U)$ 为满射。
由引理 \ref{cohomology-lemma-vanishing-ravi}，对拟紧开集由拟紧开集组成的任意有限覆盖
$\mathcal{U} : U = U_1 \cup \ldots \cup U_n$ 及任意 $q > 0$，有
$\check{H}^q(\mathcal{U}, \mathcal{G}) = 0$。
应用引理 \ref{cohomology-lemma-cech-vanish-basis}，可知对 $X$
的所有拟紧开集 $U$ 及所有 $q > 0$，有
$H^q(U, \mathcal{G}) = 0$。由 Leray 无环性引理
（《导出范畴》引理 \ref{derived-lemma-leray-acyclicity}），得到
$$
H^q(X, \mathcal{F}) =
H^q\left(
\Gamma(X, \mathcal{I}^0) \to \ldots \to
\Gamma(X, \mathcal{I}^{d - 1}) \to \Gamma(X, \mathcal{G})
\right)
$$
特别地，当 $q > d$ 时该上同调群消失。

\medskip\noindent
最后证明 (3)。对 (3) 中的 $Z$，考察长正合列
$$
H^{q - 1}(X, \mathcal{F}) \to
H^{q - 1}(X \setminus Z, \mathcal{F}) \to
H^q_Z(X, \mathcal{F}) \to H^q(X, \mathcal{F})
$$
由于 $X$ 与 $U = X \setminus Z$ 都是维数 $\leq d$ 的谱拓扑空间
（《拓扑》引理 \ref{topology-lemma-spectral-sub}），再结合 (2) 与 (1)，
便得到所需的消失性。
\end{proof}













\section{交替 {\v C}ech 复形}
\label{cohomology-section-alternating-cech}

\noindent
本节比较 {\v C}ech 复形、交替 {\v C}ech 复形以及若干相关复形。

\medskip\noindent
设 $X$ 为拓扑空间，$\mathcal{U} : U = \bigcup_{i \in I} U_i$
为开覆盖。对 $p \geq 0$，置
$$
\check{\mathcal{C}}_{alt}^p(\mathcal{U}, \mathcal{F})
=
\left\{
\begin{matrix}
s \in  \check{\mathcal{C}}^p(\mathcal{U}, \mathcal{F})
\text{，满足 }
s_{i_0 \ldots i_p} = 0 \text{，若 } i_n = i_m \text{ 对某些 } n \not = m \text{ 成立}\\
\text{且}
s_{i_0\ldots i_n \ldots i_m \ldots i_p}
=
-s_{i_0\ldots i_m \ldots i_n \ldots i_p}
\text{总成立。}
\end{matrix}
\right\}
$$
我们略去如下验证：方程 (\ref{cohomology-equation-d-cech}) 中的微分 $d$ 将
$\check{\mathcal{C}}^p_{alt}(\mathcal{U}, \mathcal{F})$ 映入
$\check{\mathcal{C}}^{p + 1}_{alt}(\mathcal{U}, \mathcal{F})$。

\begin{definition}
\label{cohomology-definition-alternating-cech-complex}
设 $X$ 为拓扑空间，$\mathcal{U} : U = \bigcup_{i \in I} U_i$
为开覆盖，$\mathcal{F}$ 为 $X$ 上的阿贝尔预层。复形
$\check{\mathcal{C}}_{alt}^\bullet(\mathcal{U}, \mathcal{F})$
称为与 $\mathcal{F}$ 及开覆盖 $\mathcal{U}$ 相伴的
{\it 交替 {\v C}ech 复形}。
\end{definition}

\noindent
于是存在复形的典范态射
$$
\check{\mathcal{C}}_{alt}^\bullet(\mathcal{U}, \mathcal{F})
\longrightarrow
\check{\mathcal{C}}^\bullet(\mathcal{U}, \mathcal{F})
$$
即交替 {\v C}ech 复形到通常 {\v C}ech 复形的包含映射。

\medskip\noindent
设覆盖 $\mathcal{U} : U = \bigcup_{i \in I} U_i$ 的指标集 $I$
带有全序 $<$。此时置
$$
\check{\mathcal{C}}_{ord}^p(\mathcal{U}, \mathcal{F})
=
\prod\nolimits_{(i_0, \ldots, i_p) \in I^{p + 1}, i_0 < \ldots < i_p}
\mathcal{F}(U_{i_0\ldots i_p}).
$$
这是一个阿贝尔群。对
$s \in \check{\mathcal{C}}_{ord}^p(\mathcal{U}, \mathcal{F})$，
以 $s_{i_0\ldots i_p}$ 表示它在
$\mathcal{F}(U_{i_0\ldots i_p})$ 中的值。定义
$$
d : \check{\mathcal{C}}_{ord}^p(\mathcal{U}, \mathcal{F})
\longrightarrow
\check{\mathcal{C}}_{ord}^{p + 1}(\mathcal{U}, \mathcal{F})
$$
如下：
$$
d(s)_{i_0\ldots i_{p + 1}}
=
\sum\nolimits_{j = 0}^{p + 1}
(-1)^j
s_{i_0\ldots \hat i_j \ldots i_{p + 1}}|_{U_{i_0\ldots i_{p + 1}}}
$$
其中 $i_0 < \ldots < i_{p + 1}$。注意此公式与方程
(\ref{cohomology-equation-d-cech}) 完全相同。直接可见 $d \circ d = 0$；
换言之，$\check{\mathcal{C}}_{ord}^\bullet(\mathcal{U}, \mathcal{F})$
是复形。

\begin{definition}
\label{cohomology-definition-ordered-cech-complex}
设 $X$ 为拓扑空间，$\mathcal{U} : U = \bigcup_{i \in I} U_i$
为开覆盖，并在 $I$ 上给定全序。设 $\mathcal{F}$ 为 $X$ 上的阿贝尔预层。
复形 $\check{\mathcal{C}}_{ord}^\bullet(\mathcal{U}, \mathcal{F})$
称为与 $\mathcal{F}$、开覆盖 $\mathcal{U}$ 及 $I$ 上给定全序相伴的
{\it 有序 {\v C}ech 复形}。
\end{definition}

\noindent
此复形有时也称为交替 {\v C}ech 复形，因为有序 {\v C}ech 复形
与交替 {\v C}ech 复形之间存在明显的比较映射。具体地，考虑映射
$$
c :
\check{\mathcal{C}}_{ord}^\bullet(\mathcal{U}, \mathcal{F})
\longrightarrow
\check{\mathcal{C}}^\bullet(\mathcal{U}, \mathcal{F})
$$
，定义为
$$
c(s)_{i_0\ldots i_p} =
\left\{
\begin{matrix}
0 &
\text{若} &
i_n = i_m \text{ 对某些 } n \not = m \text{ 成立}\\
\text{sgn}(\sigma) s_{i_{\sigma(0)}\ldots i_{\sigma(p)}} &
\text{若} &
i_{\sigma(0)} < i_{\sigma(1)} < \ldots < i_{\sigma(p)}
\end{matrix}
\right.
$$
其中 $\sigma$ 表示 $\{0, \ldots, p\}$ 的一个置换，
$\text{sgn}(\sigma)$ 表示其符号。文献中常借映射 $c$
将交替与有序 {\v C}ech 复形等同起来。具体地，有如下简单引理。

\begin{lemma}
\label{cohomology-lemma-ordered-alternating}
设 $X$ 为拓扑空间，$\mathcal{U} : U = \bigcup_{i \in I} U_i$
为开覆盖，并设 $I$ 带有全序。映射 $c$ 是复形态射。
事实上，它诱导复形同构
$$
c : \check{\mathcal{C}}_{ord}^\bullet(\mathcal{U}, \mathcal{F})
\to \check{\mathcal{C}}_{alt}^\bullet(\mathcal{U}, \mathcal{F})
$$
\end{lemma}

\begin{proof}
略。
\end{proof}

\noindent
另有映射
$$
\pi :
\check{\mathcal{C}}^\bullet(\mathcal{U}, \mathcal{F})
\longrightarrow
\check{\mathcal{C}}_{ord}^\bullet(\mathcal{U}, \mathcal{F})
$$
，其定义为
$$
\pi(s)_{i_0\ldots i_p} = s_{i_0\ldots i_p}
$$
，其中 $i_0 < i_1 < \ldots < i_p$。

\begin{lemma}
\label{cohomology-lemma-project-to-ordered}
设 $X$ 为拓扑空间，$\mathcal{U} : U = \bigcup_{i \in I} U_i$
为开覆盖，并设 $I$ 带有全序。映射
$\pi : \check{\mathcal{C}}^\bullet(\mathcal{U}, \mathcal{F})
\to \check{\mathcal{C}}_{ord}^\bullet(\mathcal{U}, \mathcal{F})$
是复形态射。它诱导复形同构
$$
\pi : \check{\mathcal{C}}_{alt}^\bullet(\mathcal{U}, \mathcal{F})
\to \check{\mathcal{C}}_{ord}^\bullet(\mathcal{U}, \mathcal{F})
$$
，且该同构是态射 $c$ 的左逆。
\end{lemma}

\begin{proof}
略。
\end{proof}

\begin{remark}
\label{cohomology-remark-compared-ordered-complexes}
这意味着，若指标集 $I$ 上有两个全序 $<_1$ 与 $<_2$，则得到复形同构
$\tau = \pi_2 \circ c_1 :
\check{\mathcal{C}}_{ord\text{-}1}(\mathcal{U}, \mathcal{F}) \to
\check{\mathcal{C}}_{ord\text{-}2}(\mathcal{U}, \mathcal{F})$.
显然
$$
\tau(s)_{i_0 \ldots i_p} =
\text{符号}(\sigma) s_{i_{\sigma(0)} \ldots i_{\sigma(p)}}
$$
，其中 $i_0 <_1 i_1 <_1 \ldots <_1 i_p$，且
$i_{\sigma(0)} <_2 i_{\sigma(1)} <_2 \ldots <_2 i_{\sigma(p)}$。
正是在此意义下，有序 {\v C}ech 复形与所选全序无关。
\end{remark}

\begin{lemma}
\label{cohomology-lemma-alternating-usual}
设 $X$ 为拓扑空间，$\mathcal{U} : U = \bigcup_{i \in I} U_i$
为开覆盖，并设 $I$ 带有全序。映射 $c \circ \pi$
与 $\check{\mathcal{C}}^\bullet(\mathcal{U}, \mathcal{F})$ 上的恒等映射链同伦。
特别地，包含映射
$\check{\mathcal{C}}_{alt}^\bullet(\mathcal{U}, \mathcal{F}) \to
\check{\mathcal{C}}^\bullet(\mathcal{U}, \mathcal{F})$
是链同伦等价。
\end{lemma}

\begin{proof}
对任意多重指标 $(i_0, \ldots, i_p) \in I^{p + 1}$，存在唯一置换
$\sigma : \{0, \ldots, p\} \to \{0, \ldots, p\}$，使得
$$
i_{\sigma(0)} \leq i_{\sigma(1)} \leq \ldots \leq i_{\sigma(p)}
\quad
\text{且}
\quad
\sigma(j) < \sigma(j + 1)
\quad
\text{若}
\quad
i_{\sigma(j)} = i_{\sigma(j + 1)}.
$$
将此置换记为 $\sigma = \sigma^{i_0 \ldots i_p}$。

\medskip\noindent
对任意置换 $\sigma : \{0, \ldots, p\} \to \{0, \ldots, p\}$
及任意 $0 \leq a \leq p$，以 $\sigma_a$ 表示 $\{0, \ldots, p\}$
的唯一置换，使得当 $0 \leq j < a$ 时
$\sigma_a(j) = \sigma(j)$，并且
$\sigma_a(a) < \sigma_a(a + 1) < \ldots < \sigma_a(p)$。
例如，若 $p = 3$ 且 $\sigma,\tau$ 由下式给出：
$$
\begin{matrix}
\text{id} & 0 & 1 & 2 & 3 \\
\sigma & 3 & 2 & 1 & 0
\end{matrix}
\quad \text{以及} \quad
\begin{matrix}
\text{id} & 0 & 1 & 2 & 3 \\
\tau & 3 & 0 & 2 & 1
\end{matrix}
$$
则
$$
\begin{matrix}
\text{id} & 0 & 1 & 2 & 3 \\
\sigma_0 & 0 & 1 & 2 & 3 \\
\sigma_1 & 3 & 0 & 1 & 2 \\
\sigma_2 & 3 & 2 & 0 & 1 \\
\sigma_3 & 3 & 2 & 1 & 0 \\
\end{matrix}
\quad \text{以及} \quad
\begin{matrix}
\text{id} & 0 & 1 & 2 & 3 \\
\tau_0 & 0 & 1 & 2 & 3 \\
\tau_1 & 3 & 0 & 1 & 2 \\
\tau_2 & 3 & 0 & 1 & 2 \\
\tau_3 & 3 & 0 & 2 & 1 \\
\end{matrix}
$$
显然总有 $\sigma_0 = \text{id}$ 与 $\sigma_p = \sigma$。

\medskip\noindent
引入上述记号后，对
$s \in \check{\mathcal{C}}^{p + 1}(\mathcal{U}, \mathcal{F})$，
定义元素 $h(s) \in \check{\mathcal{C}}^p(\mathcal{U}, \mathcal{F})$，
其分量为
\begin{equation}
\label{cohomology-equation-first-homotopy}
h(s)_{i_0\ldots i_p} =
\sum\nolimits_{0 \leq a \leq p}
(-1)^a \text{符号}(\sigma_a)
s_{i_{\sigma(0)} \ldots i_{\sigma(a)} i_{\sigma_a(a)} \ldots i_{\sigma_a(p)}}
\end{equation}
其中 $\sigma = \sigma^{i_0 \ldots i_p}$。指标 $i_{\sigma(a)}$ 在
$i_{\sigma(0)} \ldots i_{\sigma(a)} i_{\sigma_a(a)} \ldots i_{\sigma_a(p)}$
中出现两次：一次出现在前一组 $a + 1$ 个指标中，另一次出现在后一组
$p - a + 1$ 个指标中。这是因为按 $\sigma_a$ 的定义，对某个
$j \geq a$ 有 $\sigma_a(j) = \sigma(a)$。因此该和有意义，因为每个元素
$s_{i_{\sigma(0)} \ldots i_{\sigma(a)} i_{\sigma_a(a)} \ldots i_{\sigma_a(p)}}$
都定义在开集 $U_{i_0 \ldots i_p}$ 上。还要注意，当 $a = 0$ 时得到
$s_{i_0 \ldots i_p}$，而当 $a = p$ 时得到
$(-1)^p \text{符号}(\sigma) s_{i_{\sigma(0)} \ldots i_{\sigma(p)}}$.

\medskip\noindent
我们断言
$$
(dh + hd)(s)_{i_0 \ldots i_p} =
s_{i_0 \ldots i_p} -
\text{符号}(\sigma) s_{i_{\sigma(0)} \ldots i_{\sigma(p)}}
$$
，其中 $\sigma = \sigma^{i_0 \ldots i_p}$。此断言的验证从略。
（stacks-project 的 scripts 子目录中有一个名为 first-homotopy.gp
的 PARI/gp 脚本，可用来检验此断言的有限多个实例。
我们编写该脚本是为了确认符号正确。）
记
$$
\kappa :
\check{\mathcal{C}}^\bullet(\mathcal{U}, \mathcal{F})
\longrightarrow
\check{\mathcal{C}}^\bullet(\mathcal{U}, \mathcal{F})
$$
为按下式定义的算子：
$$
\kappa(s)_{i_0 \ldots i_p} =
\text{符号}(\sigma^{i_0 \ldots i_p}) s_{i_{\sigma(0)} \ldots i_{\sigma(p)}}.
$$
上述断言蕴含 $\kappa$ 是复形态射，且 $\kappa$ 与 {\v C}ech
复形的恒等映射链同伦。这并不立即蕴含引理，因为算子 $\kappa$
的像并非交替子复形。具体地，$\kappa$ 的像是``半交替''复形
$\check{\mathcal{C}}_{semi\text{-}alt}^p(\mathcal{U}, \mathcal{F})$
，其中 $s$ 是此复形的 $p$-上链，当且仅当
$$
s_{i_0 \ldots i_p} = \text{符号}(\sigma) s_{i_{\sigma(0)} \ldots i_{\sigma(p)}}
$$
对任意 $(i_0, \ldots, i_p) \in I^{p + 1}$ 成立，其中
$\sigma = \sigma^{i_0 \ldots i_p}$。
再引入 {\v C}ech 复形的另一变体，即定义为下式的半有序
{\v C}ech 复形：
$$
\check{\mathcal{C}}_{semi\text{-}ord}^p(\mathcal{U}, \mathcal{F})
=
\prod\nolimits_{i_0 \leq i_1 \leq \ldots \leq i_p}
\mathcal{F}(U_{i_0 \ldots i_p})
$$
容易看出，方程 (\ref{cohomology-equation-d-cech}) 同样定义了微分，因而得到一个复形。
还显然有（类似于引理 \ref{cohomology-lemma-project-to-ordered}）投影映射
$$
\check{\mathcal{C}}_{semi\text{-}alt}^\bullet(\mathcal{U}, \mathcal{F})
\longrightarrow
\check{\mathcal{C}}_{semi\text{-}ord}^\bullet(\mathcal{U}, \mathcal{F})
$$
是复形同构。

\medskip\noindent
因此，只要证明明显的包含映射
$$
\check{\mathcal{C}}_{ord}^p(\mathcal{U}, \mathcal{F})
\longrightarrow
\check{\mathcal{C}}_{semi\text{-}ord}^p(\mathcal{U}, \mathcal{F})
$$
是链同伦等价，就得到引理。为此使用链同伦
\begin{equation}
\label{cohomology-equation-second-homotopy}
h(s)_{i_0 \ldots i_p} =
\left\{
\begin{matrix}
0 & \text{若} & i_0 < i_1 < \ldots < i_p \\
(-1)^a s_{i_0 \ldots i_{a - 1} i_a i_a i_{a + 1} \ldots i_p}
& \text{若} & i_0 < i_1 < \ldots < i_{a - 1} < i_a = i_{a + 1}
\end{matrix}
\right.
\end{equation}
我们断言
$$
(dh + hd)(s)_{i_0 \ldots i_p} =
\left\{
\begin{matrix}
0 & \text{若} & i_0 < i_1 < \ldots < i_p \\
s_{i_0 \ldots i_p}
& \text{否则} &
\end{matrix}
\right.
$$
验证从略。（stacks-project 的 scripts 子目录中有一个名为
second-homotopy.gp 的 PARI/gp 脚本，可用来检验此断言的有限多个实例。
我们编写该脚本是为了确认符号正确。）
该断言清楚表明，投影与自然包含的复合
$$
\check{\mathcal{C}}_{semi\text{-}ord}^\bullet(\mathcal{U}, \mathcal{F})
\longrightarrow
\check{\mathcal{C}}_{ord}^\bullet(\mathcal{U}, \mathcal{F})
\longrightarrow
\check{\mathcal{C}}_{semi\text{-}ord}^\bullet(\mathcal{U}, \mathcal{F})
$$
按所需与恒等映射链同伦。
\end{proof}

\begin{lemma}
\label{cohomology-lemma-alternating-cech-trivial}
设 $X$ 为拓扑空间，$\mathcal{F}$ 为 $X$ 上的阿贝尔预层，
$\mathcal{U} : U = \bigcup_{i \in I} U_i$ 为开覆盖。
若对某个 $i \in I$ 有 $U_i = U$，则扩展的交替 {\v C}ech 复形
$$
\mathcal{F}(U) \to \check{\mathcal{C}}_{alt}^\bullet(\mathcal{U}, \mathcal{F})
$$
是通过把 $\mathcal{F}(U)$ 置于次数 $-1$ 并以
$\mathcal{F}(U)$ 到 $\check{\mathcal{C}}^0(\mathcal{U}, \mathcal{F})$
的典范映射为微分所得；它链同伦等价于 $0$。类似地，对 $I$
上的任意全序，扩展的有序 {\v C}ech 复形
$$
\mathcal{F}(U) \to
\check{\mathcal{C}}_{ord}^\bullet(\mathcal{U}, \mathcal{F})
$$
链同伦等价于 $0$。
\end{lemma}

\begin{proof}[第一个证明]
结合引理 \ref{cohomology-lemma-cech-trivial} 与 \ref{cohomology-lemma-alternating-usual}。
\end{proof}

\begin{proof}[第二个证明]
由于交替与有序 {\v C}ech 复形同构，只需对有序复形证明。
采用标准记号：扩展的有序 {\v C}ech 复形中次数为 $p$ 的上链 $s$
形如 $s = (s_{i_0 \ldots i_p})$，其中
$s_{i_0 \ldots i_p} \in \mathcal{F}(U_{i_0 \ldots i_p})$，且
$i_0 < \ldots < i_p$。在此记号下，
$$
d(x)_{i_0 \ldots i_{p + 1}} =
\sum\nolimits_j (-1)^j x_{i_0 \ldots \hat i_j \ldots i_p}
$$
固定满足 $U = U_i$ 的指标 $i \in I$。作为链同伦，使用映射
$$
h : \text{次数为 }p + 1\text{ 的上链} \to \text{次数为 }p\text{ 的上链}
$$
，定义为
$$
h(s)_{i_0 \ldots i_p} = 0 \text{，若 } i \in \{i_0, \ldots, i_p\}
\text{；且}
h(s)_{i_0 \ldots i_p} = 
(-1)^j s_{i_0 \ldots i_j i i_{j + 1} \ldots i_p} \text{，否则}
$$
在第二种情形中，$j$ 是使得 $i_j < i < i_{j + 1}$ 的唯一指标；
又因 $U = U_i$，有等式
$$
\mathcal{F}(U_{i_0 \ldots i_p}) =
\mathcal{F}(U_{i_0 \ldots i_j i i_{j + 1} \ldots i_p})
$$
故可将
$(-1)^j s_{i_0 \ldots i_j i i_{j + 1} \ldots i_p}$
视为 $\mathcal{F}(U_{i_0 \ldots i_p})$ 的元素。
下面通过计算证明 $d h + h d$ 等于恒等映射的负值，从而完成证明。
为此固定次数为 $p$ 的上链 $s$，并取 $I$ 中的元素
$i_0 < \ldots < i_p$。

\medskip\noindent
情形 I：$i \in \{i_0, \ldots, i_p\}$。设 $i = i_t$。
于是 $h(d(s))_{i_0 \ldots i_p} = 0$。另一方面，
$$
d(h(s))_{i_0 \ldots i_p} =
\sum (-1)^j h(s)_{i_0 \ldots \hat i_j \ldots i_p} =
(-1)^t h(s)_{i_0 \ldots \hat i \ldots i_p} =
(-1)^t (-1)^{t - 1} s_{i_0 \ldots i_p}
$$
因此按所需有
$(dh + hd)(s)_{i_0 \ldots i_p} = -s_{i_0 \ldots i_p}$。

\medskip\noindent
情形 II：$i \not \in \{i_0, \ldots, i_p\}$。取 $j$ 使得
$i_j < i < i_{j + 1}$。于是
\begin{align*}
h(d(s))_{i_0 \ldots i_p}
& =
(-1)^j d(s)_{i_0 \ldots i_j i i_{j + 1} \ldots i_p} \\
& =
\sum\nolimits_{j' \leq j} (-1)^{j + j'}
s_{i_0 \ldots \hat i_{j'} \ldots i_j i i_{j + 1} \ldots i_p} -
s_{i_0 \ldots i_p} \\
&
+ \sum\nolimits_{j' > j} (-1)^{j + j' + 1}
s_{i_0 \ldots i_j i i_{j + 1} \ldots \hat i_{j'} \ldots i_p}
\end{align*}
另一方面，
\begin{align*}
d(h(s))_{i_0 \ldots i_p}
& =
\sum\nolimits_{j'} (-1)^{j'} h(s)_{i_0 \ldots \hat i_{j'} \ldots i_p} \\
& =
\sum\nolimits_{j' \leq j} (-1)^{j' + j - 1}
s_{i_0 \ldots \hat i_{j'} \ldots i_j i i_{j + 1} \ldots i_p} \\
& +
\sum\nolimits_{j' > j} (-1)^{j' + j}
s_{i_0 \ldots i_j i i_{j + 1} \ldots \hat i_{j'} \ldots i_p}
\end{align*}
两式相加得到
$(dh + hd)(s)_{i_0 \ldots i_p} = - s_{i_0 \ldots i_p}$
，正是所需结论。
\end{proof}





\section{{\v C}ech 复形的另一种看法}
\label{cohomology-section-locally-finite-cech}

\noindent
本节讨论建立 {\v C}ech 复形与上同调之间关系的另一种方法。

\begin{lemma}
\label{cohomology-lemma-covering-resolution}
设 $X$ 为环空间，$\mathcal{U} : X = \bigcup_{i \in I} U_i$
为 $X$ 的开覆盖，$\mathcal{F}$ 为 $\mathcal{O}_X$-模。
以 $\mathcal{F}_{i_0 \ldots i_p}$ 表示 $\mathcal{F}$
在 $U_{i_0 \ldots i_p}$ 上的限制。存在 $\mathcal{O}_X$-模复形
${\mathfrak C}^\bullet(\mathcal{U}, \mathcal{F})$
，满足
$$
{\mathfrak C}^p(\mathcal{U}, \mathcal{F}) =
\prod\nolimits_{i_0 \ldots i_p}
(j_{i_0 \ldots i_p})_* \mathcal{F}_{i_0 \ldots i_p}
$$
，且其微分
$d : {\mathfrak C}^p(\mathcal{U}, \mathcal{F})
\to {\mathfrak C}^{p + 1}(\mathcal{U}, \mathcal{F})$
如方程 (\ref{cohomology-equation-d-cech}) 所示。此外，存在典范映射
$$
\mathcal{F} \to {\mathfrak C}^\bullet(\mathcal{U}, \mathcal{F})
$$
，它是拟同构；换言之，
${\mathfrak C}^\bullet(\mathcal{U}, \mathcal{F})$ 是 $\mathcal{F}$ 的消解。
\end{lemma}

\begin{proof}
只需验证
$$
0 \to \mathcal{F} \to \mathfrak{C}^0(\mathcal{U}, \mathcal{F}) \to
\mathfrak{C}^1(\mathcal{U}, \mathcal{F}) \to  \ldots
$$
在茎上正合。取 $x \in X$，并选取满足
$x \in U_{i_{\text{fix}}}$ 的 $i_{\text{fix}} \in I$。定义
$$
h : \mathfrak{C}^p(\mathcal{U}, \mathcal{F})_x
\to \mathfrak{C}^{p - 1}(\mathcal{U}, \mathcal{F})_x
$$
如下：若 $s \in \mathfrak{C}^p(\mathcal{U}, \mathcal{F})_x$，取其在
$x$ 的某个邻域 $V$ 上定义的代表元
$$
\widetilde{s} \in
\mathfrak{C}^p(\mathcal{U}, \mathcal{F})(V) =
\prod\nolimits_{i_0 \ldots i_p}
\mathcal{F}(V \cap U_{i_0} \cap \ldots \cap U_{i_p})
$$
，并置
$$
h(s)_{i_0 \ldots i_{p - 1}} =
\widetilde{s}_{i_{\text{fix}} i_0 \ldots i_{p - 1}, x}.
$$
同一公式（取 $p = 0$）给出映射
$\mathfrak{C}^{0}(\mathcal{U},\mathcal{F})_x \to \mathcal{F}_x$。
形式计算如下：
\begin{align*}
(dh + hd)(s)_{i_0 \ldots i_p}
& =
\sum\nolimits_{j = 0}^p
(-1)^j
h(s)_{i_0 \ldots \hat i_j \ldots i_p}
+
d(s)_{i_{\text{fix}} i_0 \ldots i_p}\\
& =
\sum\nolimits_{j = 0}^p
(-1)^j
s_{i_{\text{fix}} i_0 \ldots \hat i_j \ldots i_p}
+
s_{i_0 \ldots i_p}
+
\sum\nolimits_{j = 0}^p
(-1)^{j + 1}
s_{i_{\text{fix}} i_0 \ldots \hat i_j \ldots i_p} \\
& =
s_{i_0 \ldots i_p}
\end{align*}
这表明 $h$ 给出扩展复形
$$
0 \to \mathcal{F}_x \to \mathfrak{C}^0(\mathcal{U}, \mathcal{F})_x
\to \mathfrak{C}^1(\mathcal{U}, \mathcal{F})_x \to \ldots
$$
的恒等映射与零映射之间的链同伦，结论得证。
\end{proof}

\noindent
利用此引理，很容易重新证明引理
\ref{cohomology-lemma-cech-spectral-sequence} 的 {\v C}ech 到上同调谱序列。
具体地，设 $X,\mathcal{U},\mathcal{F}$ 如引理
\ref{cohomology-lemma-covering-resolution}，并设
$\mathcal{F} \to \mathcal{I}^\bullet$ 为内射消解。于是可考虑双复形
$$
A^{\bullet, \bullet} =
\Gamma(X, {\mathfrak C}^\bullet(\mathcal{U}, \mathcal{I}^\bullet)).
$$
按构造有
$$
A^{p, q} = \prod\nolimits_{i_0 \ldots i_p} \mathcal{I}^q(U_{i_0 \ldots i_p})
$$
考虑《同调代数》第 \ref{homology-section-double-complex} 节中
与此双复形相伴的两个谱序列，特别见《同调代数》引理
\ref{homology-lemma-ss-double-complex}。
对谱序列 $({}'E_r, {}'d_r)_{r \geq 0}$，有
${}'E_2^{p, q} = \check{H}^p(\mathcal{U}, \underline{H}^q(\mathcal{F}))$
，因为取积正合（《同调代数》引理
\ref{homology-lemma-product-abelian-groups-exact}）。
对谱序列 $({}''E_r, {}''d_r)_{r \geq 0}$，若 $p > 0$，则
${}''E_2^{p, q} = 0$，且 ${}''E_2^{0, q} = H^q(X, \mathcal{F})$。
事实上，对固定的 $q$，层复形
${\mathfrak C}^\bullet(\mathcal{U}, \mathcal{I}^q)$
是内射层 $\mathcal{I}^q$ 的内射层消解（引理
\ref{cohomology-lemma-covering-resolution}；另由引理
\ref{cohomology-lemma-cohomology-of-open}、\ref{cohomology-lemma-pushforward-injective-flat}
及《同调代数》引理 \ref{homology-lemma-product-injectives}）。
因此
$\Gamma(X, {\mathfrak C}^\bullet(\mathcal{U}, \mathcal{I}^q))$
的上同调在正次数为零，在次数 $0$ 等于
$\Gamma(X, \mathcal{I}^q)$。再对下一个微分取上同调，
便得到关于谱序列 $({}''E_r, {}''d_r)_{r \geq 0}$ 的断言。
两个谱序列都收敛到 $A^{\bullet, \bullet}$ 相伴总复形的上同调，故得结论。

\begin{definition}
\label{cohomology-definition-covering-locally-finite}
设 $X$ 为拓扑空间。若对每个 $x \in X$ 都存在 $x$ 的开邻域 $W$，
使得 $\{i \in I \mid W \cap U_i \not = \emptyset\}$ 有限，
则称开覆盖 $X = \bigcup_{i \in I} U_i$ {\it 局部有限}。
\end{definition}

\begin{remark}
\label{cohomology-remark-locally-finite-sections}
设 $X = \bigcup_{i \in I} U_i$ 为局部有限开覆盖，
并以 $j_i : U_i \to X$ 表示包含映射。设对每个 $i$，
都给定 $U_i$ 上的阿贝尔层 $\mathcal{F}_i$。考虑阿贝尔层
$\mathcal{G} = \bigoplus_{i \in I} (j_i)_*\mathcal{F}_i$。
则对开集 $V \subset X$，事实上有
$$
\Gamma(V, \mathcal{G}) = \prod\nolimits_{i \in I} \mathcal{F}_i(V \cap U_i).
$$
换言之，
$$
\bigoplus\nolimits_{i \in I} (j_i)_*\mathcal{F}_i =
\prod\nolimits_{i \in I} (j_i)_*\mathcal{F}_i
$$
起初这似乎有些奇怪，但要注意，一族层的直和是其直观预期对象的层化。
相关讨论见《模层》第 \ref{modules-section-kernels} 节。
因此在此情形下，引理 \ref{cohomology-lemma-covering-resolution} 中复形的各项为
$$
{\mathfrak C}^p(\mathcal{U}, \mathcal{F}) =
\bigoplus\nolimits_{i_0 \ldots i_p}
(j_{i_0 \ldots i_p})_* \mathcal{F}_{i_0 \ldots i_p}
$$
，这有时很有用。
\end{remark}











\section{复形的 {\v C}ech 上同调}
\label{cohomology-section-cech-cohomology-of-complexes}

\noindent
一般而言，对 $X$ 上的阿贝尔群层 ${\mathcal F}$ 与 ${\mathcal G}$，
存在杯积映射
$$
H^i(X, {\mathcal F}) \times H^j(X, {\mathcal G})
\longrightarrow
H^{i + j}(X, {\mathcal F} \otimes_{\mathbf Z} {\mathcal G}).
$$
本节用 {\v C}ech 上循环和一个显式杯积公式定义此映射。
如果担心上同调未必等于 {\v C}ech 上同调，则可使用超覆盖，
同时仍采用上循环记号。这样还有一个优点：同一方法可定义任意拓扑斯上
超上同调的杯积（此处插入未来参考文献）。

\medskip\noindent
设 ${\mathcal F}^\bullet$ 为 $X$ 上阿贝尔群预层的下有界复形。
我们常可用 {\v C}ech 上循环计算 $H^n(X, {\mathcal F}^\bullet)$。
具体地，设 ${\mathcal U} : X = \bigcup_{i \in I} U_i$
为 $X$ 的开覆盖。由于 {\v C}ech 复形
$\check{\mathcal{C}}^\bullet(\mathcal{U}, \mathcal{F})$
（定义 \ref{cohomology-definition-cech-complex}）关于预层 $\mathcal{F}$ 具有函子性，
可得到双复形
$\check{\mathcal{C}}^\bullet(\mathcal{U}, \mathcal{F}^\bullet)$.
与 $\check{\mathcal{C}}^\bullet({\mathcal U}, {\mathcal F}^\bullet)$
相伴的总复形，其次数 $n$ 项为
$$
\text{Tot}^n(\check{\mathcal{C}}^\bullet({\mathcal U}, {\mathcal F}^\bullet))
=
\bigoplus\nolimits_{p + q = n}
\prod\nolimits_{i_0\ldots i_p} {\mathcal F}^q(U_{i_0\ldots i_p})
$$
；见《同调代数》定义
\ref{homology-definition-associated-simple-complex}。
以 $\alpha = \{\alpha_{i_0\ldots i_p}\}$ 表示 $\text{Tot}^n$
中的典型元素，其中
$\alpha_{i_0 \ldots i_p} \in \mathcal{F}^q(U_{i_0\ldots i_p})$。
换言之，$\alpha_{i_0\ldots i_p}$ 的 $\mathcal{F}$-次数为
$q = n - p$。此记号要求我们始终留意 $\alpha$ 所在的次数。
用公式
$\deg_{\mathcal F}(\alpha_{i_0\ldots i_p}) = q$.
表示这一情形。按照《同调代数》定义
\ref{homology-definition-associated-simple-complex} 中的约定，
次数为 $n$ 的元素 $\alpha$ 的微分为
$$
d(\alpha)_{i_0\ldots i_{p + 1}}
=
\sum\nolimits_{j = 0}^{p + 1}
(-1)^j \alpha_{i_0 \ldots \hat i_j \ldots i_{p + 1}} + 
(-1)^{p + 1}d_{{\mathcal F}}(\alpha_{i_0 \ldots i_{p + 1}})
$$
，其中 $d_\mathcal{F}$ 表示复形 $\mathcal{F}^\bullet$ 的微分。
表达式 $\alpha_{i_0 \ldots \hat i_j \ldots i_{p + 1}}$ 表示把
$\alpha_{i_0 \ldots \hat i_j \ldots i_{p + 1}}
\in {\mathcal F}(U_{i_0\ldots\hat i_j\ldots i_{p + 1}})$
限制到 $U_{i_0 \ldots i_{p + 1}}$ 所得的元素。

\medskip\noindent
构造
$\text{Tot}(\check{\mathcal{C}}^\bullet({\mathcal U}, {\mathcal F}^\bullet))$
关于 ${\mathcal F}^\bullet$ 具有函子性。此外，存在复形的函子性变换
\begin{equation}
\label{cohomology-equation-global-sections-to-cech}
\Gamma(X, {\mathcal F}^\bullet)
\longrightarrow
\text{Tot}(\check{\mathcal{C}}^\bullet({\mathcal U}, {\mathcal F}^\bullet))
\end{equation}
，定义如下：截面 $s\in \Gamma(X, {\mathcal F}^n)$
映为元素 $\alpha = \{\alpha_{i_0\ldots i_p}\}$，其中
$\alpha_{i_0} = s|_{U_{i_0}}$，且当 $p > 0$ 时
$\alpha_{i_0\ldots i_p} = 0$。

\medskip\noindent
加细。设 ${\mathcal V} = \{ V_j \}_{j\in J}$ 是
${\mathcal U}$ 的一个加细。这意味着存在映射 $t: J \to I$，
使得对所有 $j\in J$ 都有 $V_j \subset U_{t(j)}$。由此得到函子性变换
\begin{equation}
\label{cohomology-equation-transformation}
T_t :
\text{Tot}(\check{\mathcal{C}}^\bullet({\mathcal U}, {\mathcal F}^\bullet))
\longrightarrow
\text{Tot}(\check{\mathcal{C}}^\bullet({\mathcal V}, {\mathcal F}^\bullet))
\end{equation}
，其定义规则为
$$
T_t(\alpha)_{j_0\ldots j_p}
=
\alpha_{t(j_0)\ldots t(j_p)}|_{V_{j_0\ldots j_p}}.
$$
给定如上的两个映射 $t, t' : J \to I$，上面构造的映射
$T_t$ 与 $T_{t'}$ 链同伦。此链同伦由下式给出：
$$
h(\alpha)_{j_0\ldots j_p}
=
\sum\nolimits_{a = 0}^{p}
(-1)^a
\alpha_{t(j_0)\ldots t(j_a) t'(j_a) \ldots t'(j_p)}
$$
，其中 $\alpha$ 为次数 $n$ 的元素。下面的计算表明此式确实成立；
这里仍设 $\alpha$ 为次数 $n$ 的元素（因而 $d(\alpha)$
的次数为 $n + 1$，而 $h(\alpha)$ 的次数为 $n - 1$）：
\begin{align*}
(
d(h(\alpha)) + h(d(\alpha))
)_{j_0\ldots j_p}
= &
\sum\nolimits_{k = 0}^p
(-1)^k
h(\alpha)_{j_0 \ldots \hat j_k \ldots j_p}
+ \\
&
(-1)^p
d_{\mathcal F}(h(\alpha)_{j_0 \ldots j_p})
+ \\
&
\sum\nolimits_{a = 0}^p
(-1)^a
d(\alpha)_{t(j_0) \ldots t(j_a) t'(j_a) \ldots t'(j_p)}
\\
= &
\sum\nolimits_{k = 0}^p
\sum\nolimits_{a = 0}^{k - 1}
(-1)^{k + a}
\alpha_{t(j_0)\ldots t(j_a)t'(j_a)\ldots \hat{t'(j_k)}\ldots t'(j_p)}
+ \\
&
\sum\nolimits_{k = 0}^p
\sum\nolimits_{a = k + 1}^p
(-1)^{k + a - 1}
\alpha_{t(j_0)\ldots \hat{t(j_k)}\ldots t(j_a)t'(j_a)\ldots t'(j_p)}
+ \\
&
\sum\nolimits_{a = 0}^p
(-1)^{p + a}
d_{\mathcal F}(\alpha_{t(j_0)\ldots t(j_a) t'(j_a) \ldots t'(j_p)})
+ \\
&
\sum\nolimits_{a = 0}^p
\sum\nolimits_{k = 0}^a
(-1)^{a + k}
\alpha_{t(j_0)\ldots\hat{t(j_k)}\ldots t(j_a)t'(j_a)\ldots t'(j_p)}
+ \\
&
\sum\nolimits_{a = 0}^p
\sum\nolimits_{k = a}^p
(-1)^{a + k + 1}
\alpha_{t(j_0) \ldots t(j_a) t'(j_a) \ldots \hat{t'(j_k)} \ldots t'(j_p)}
+ \\
&
\sum\nolimits_{a = 0}^p
(-1)^{a + p + 1}
d_{\mathcal F}(\alpha_{t(j_0)\ldots t(j_a) t'(j_a) \ldots t'(j_p)})
\\
= &
\alpha_{t'(j_0)\ldots t'(j_p)} +
(-1)^{2p + 1}\alpha_{t(j_0)\ldots t(j_p)}
\\
= &
T_{t'}(\alpha)_{j_0\ldots j_p} - T_t(\alpha)_{j_0\ldots j_p}
\end{align*}
我们把各项相消的验证留给读者。（注意，第一、第二以及第四、第五个
和式中同时含有 $k$ 与 $a$ 的各项会相消；例外是 $a = k$ 的项，
它们只出现在第四和第五个和式中，且彼此相消，仅留下所需的两项。）
由此可知，诱导映射
$$
H^n(T_t) :
H^n(
\text{Tot}(\check{\mathcal{C}}^\bullet({\mathcal U}, {\mathcal F}^\bullet))
)
\to
H^n(
\text{Tot}(\check{\mathcal{C}}^\bullet({\mathcal V}, {\mathcal F}^\bullet))
)
$$
与 $t$ 的选取无关。我们把 {\it Cech 超上同调}定义为
经由映射 $H^\bullet(T_t)$ 对所有加细取 Cech 上同调群的余极限。

\medskip\noindent
在（遍历 $X$ 的所有开覆盖的）余极限中，下列引理给出一个
从 Cech 超上同调到上同调的映射；它常常是同构，而若使用超覆盖，
则它总是同构。

\begin{lemma}
\label{cohomology-lemma-cech-complex-complex}
设 $(X, \mathcal{O}_X)$ 是环空间，且
$\mathcal{U} : X = \bigcup_{i \in I} U_i$ 是开覆盖。
对 $\mathcal{O}_X$-模的有下界复形 $\mathcal{F}^\bullet$，
存在典范映射
$$
\text{Tot}(\check{\mathcal{C}}^\bullet(\mathcal{U}, \mathcal{F}^\bullet))
\longrightarrow
R\Gamma(X, \mathcal{F}^\bullet)
$$
，它关于 $\mathcal{F}^\bullet$ 具有函子性，并与
(\ref{cohomology-equation-global-sections-to-cech}) 和
(\ref{cohomology-equation-transformation}) 相容。存在谱序列
$(E_r, d_r)_{r \geq 0}$，满足
$$
E_2^{p, q} =
H^p(\text{Tot}(\check{\mathcal{C}}^\bullet(\mathcal{U},
\underline{H}^q(\mathcal{F}^\bullet)))
$$
且收敛到 $H^{p + q}(X, \mathcal{F}^\bullet)$。
\end{lemma}

\begin{proof}
设 ${\mathcal I}^\bullet$ 是内射对象的有下界复形。
对 $\mathcal{I}^\bullet$ 而言，
映射 (\ref{cohomology-equation-global-sections-to-cech}) 为
$\Gamma(X, {\mathcal I}^\bullet) \to
\text{Tot}(\check{\mathcal{C}}^\bullet({\mathcal U}, {\mathcal I}^\bullet))$.
把《同调代数》引理
\ref{homology-lemma-double-complex-gives-resolution}
应用于双复形
$\check{\mathcal{C}}^\bullet({\mathcal U}, {\mathcal I}^\bullet)$，
并使用引理 \ref{cohomology-lemma-injective-trivial-cech}，可知这是
阿贝尔群复形的拟同构。设
${\mathcal F}^\bullet \to {\mathcal I}^\bullet$ 是从
${\mathcal F}^\bullet$ 到内射对象的有下界复形的拟同构。
由于 $R\Gamma(X, {\mathcal F}^\bullet)$ 由复形
$\Gamma(X, {\mathcal I}^\bullet)$ 表示，利用映射
$$
\text{Tot}(\check{\mathcal{C}}^\bullet({\mathcal U}, {\mathcal F}^\bullet))
\longrightarrow
\text{Tot}(\check{\mathcal{C}}^\bullet({\mathcal U}, {\mathcal I}^\bullet)).
$$
便得到引理中的映射。我们略去函子性与相容性的验证。
为构造引理中的谱序列，选取 Cartan--Eilenberg 消解
$\mathcal{F}^\bullet \to \mathcal{I}^{\bullet, \bullet}$，见
《导出范畴》引理 \ref{derived-lemma-cartan-eilenberg}。此时
$\mathcal{F}^\bullet \to \text{Tot}(\mathcal{I}^{\bullet, \bullet})$
是内射消解，因而
$$
\text{Tot}(\check{\mathcal{C}}^\bullet({\mathcal U},
\text{Tot}({\mathcal I}^{\bullet, \bullet})))
$$
如上所述，此复形计算 $R\Gamma(X, \mathcal{F}^\bullet)$。
由《同调代数》注记 \ref{homology-remark-triple-complex}，
可把它看成三重复形
$\check{\mathcal{C}}^\bullet({\mathcal U},
{\mathcal I}^{\bullet, \bullet})$ 所伴随的总复形；因此再次使用同一注记，
还可把它看成双复形 $A^{\bullet, \bullet}$ 所伴随的总复形，其各项为
$$
A^{n, m} =
\bigoplus\nolimits_{p + q = n}
\check{\mathcal{C}}^p({\mathcal U}, \mathcal{I}^{q, m})
$$
由于 $\mathcal{I}^{q, \bullet}$ 是 $\mathcal{F}^q$ 的内射消解，
可应用 $A^{\bullet, \bullet}$ 所伴随的第一谱序列
（《同调代数》引理 \ref{homology-lemma-ss-double-complex}），
得到满足下式的谱序列：
$$
E_1^{n, m} =
\bigoplus\nolimits_{p + q = n}
\check{\mathcal{C}}^p(\mathcal{U}, \underline{H}^m(\mathcal{F}^q))
$$
这是复形
$\text{Tot}(\check{\mathcal{C}}^\bullet(\mathcal{U},
\underline{H}^m(\mathcal{F}^\bullet))$ 的第 $n$ 项。因此得到
引理中所述的 $E_2$ 项。收敛性由《同调代数》引理
\ref{homology-lemma-first-quadrant-ss} 给出。
\end{proof}

\begin{lemma}
\label{cohomology-lemma-cech-complex-complex-computes}
设 $(X, \mathcal{O}_X)$ 是环空间，
$\mathcal{U} : X = \bigcup_{i \in I} U_i$ 是开覆盖，且
$\mathcal{F}^\bullet$ 是 $\mathcal{O}_X$-模的有下界复形。
若对所有 $i > 0$ 及所有 $p, i_0, \ldots, i_p, q$ 都有
$H^i(U_{i_0 \ldots i_p}, \mathcal{F}^q) = 0$，则引理
\ref{cohomology-lemma-cech-complex-complex} 中的映射
$
\text{Tot}(\check{\mathcal{C}}^\bullet(\mathcal{U}, \mathcal{F}^\bullet))
\to
R\Gamma(X, \mathcal{F}^\bullet)
$
是同构。
\end{lemma}

\begin{proof}
这立即由引理 \ref{cohomology-lemma-cech-complex-complex} 的谱序列得到。
\end{proof}

\begin{remark}
\label{cohomology-remark-shift-complex-cech-complex}
设 $(X, \mathcal{O}_X)$ 是环空间，
$\mathcal{U} : X = \bigcup_{i \in I} U_i$ 是开覆盖，
$\mathcal{F}^\bullet$ 是 $\mathcal{O}_X$-模的有下界复形，
而 $b$ 是整数。我们断言在导出范畴中有交换图
$$
\xymatrix{
\text{Tot}(\check{\mathcal{C}}^\bullet(\mathcal{U}, \mathcal{F}^\bullet))[b]
\ar[r] \ar[d]_\gamma &
R\Gamma(X, \mathcal{F}^\bullet)[b] \ar[d] \\
\text{Tot}(\check{\mathcal{C}}^\bullet(\mathcal{U}, \mathcal{F}^\bullet[b]))
\ar[r] &
R\Gamma(X, \mathcal{F}^\bullet[b])
}
$$
，其中映射 $\gamma$ 是《同调代数》注记
\ref{homology-remark-shift-double-complex} 中构造的复形映射。
这是有意义的，因为双复形
$\check{\mathcal{C}}^\bullet(\mathcal{U}, \mathcal{F}^\bullet[b])$
显然就是《同调代数》注记
\ref{homology-remark-shift-double-complex} 中引入的双复形
$\check{\mathcal{C}}^\bullet(\mathcal{U}, \mathcal{F}^\bullet)[0, b]$
。为验证该图交换，可像引理
\ref{cohomology-lemma-cech-complex-complex} 的证明那样选取内射消解
$\mathcal{F}^\bullet \to \mathcal{I}^\bullet$。作图追踪可知，
只需验证把 $\mathcal{F}^\bullet$ 替换为
$\mathcal{I}^\bullet$ 后该图交换。于是考虑扩充图
$$
\xymatrix{
\Gamma(X, \mathcal{I}^\bullet)[b] \ar[r] \ar[d] &
\text{Tot}(\check{\mathcal{C}}^\bullet(\mathcal{U}, \mathcal{I}^\bullet))[b]
\ar[r] \ar[d]_\gamma &
R\Gamma(X, \mathcal{I}^\bullet)[b] \ar[d] \\
\Gamma(X, \mathcal{I}^\bullet[b]) \ar[r] &
\text{Tot}(\check{\mathcal{C}}^\bullet(\mathcal{U}, \mathcal{I}^\bullet[b]))
\ar[r] &
R\Gamma(X, \mathcal{I}^\bullet[b])
}
$$
，其中左侧水平箭头为
(\ref{cohomology-equation-global-sections-to-cech})。在此情形下，
水平箭头在导出范畴中均为同构（见引理
\ref{cohomology-lemma-cech-complex-complex} 的证明），故只需证明左方块交换。
这确实成立，因为映射 $\gamma$ 在直和项
$\check{\mathcal{C}}^0(\mathcal{U}, \mathcal{I}^{q + b})$
上所用的符号为 $1$；参见《同调代数》注记
\ref{homology-remark-shift-double-complex} 中的公式。
\end{remark}

\noindent
设 $X$ 是拓扑空间，$\mathcal{U} : X = \bigcup_{i \in I} U_i$
是开覆盖，而 $\mathcal{F}^\bullet$ 是阿贝尔群预层的有下界复形。
考虑映射
$\tau :
\text{Tot}(\check{\mathcal{C}}^\bullet({\mathcal U}, {\mathcal F}^\bullet))
\to
\text{Tot}(\check{\mathcal{C}}^\bullet({\mathcal U}, {\mathcal F}^\bullet))$
，其定义为
$$
\tau(\alpha)_{i_0 \ldots i_p} = (-1)^{p(p + 1)/2} \alpha_{i_p \ldots i_0}.
$$
于是对次数 $n$ 的元素 $\alpha$ 有
\begin{align*}
& d(\tau(\alpha))_{i_0 \ldots i_{p + 1}} \\
& =
\sum\nolimits_{j = 0}^{p + 1}
(-1)^j
\tau(\alpha)_{i_0 \ldots \hat i_j \ldots i_{p + 1}}
+
(-1)^{p + 1}
d_{\mathcal F}(\tau(\alpha)_{i_0 \ldots i_{p + 1}})
\\
& =
\sum\nolimits_{j = 0}^{p + 1}
(-1)^{j + \frac{p(p + 1)}{2}}
\alpha_{i_{p + 1} \ldots \hat i_j \ldots i_0}
+
(-1)^{p + 1 + \frac{(p + 1)(p + 2)}{2}}
d_{\mathcal F}(\alpha_{i_{p + 1} \ldots i_0})
\end{align*}
另一方面有
\begin{align*}
& \tau(d(\alpha))_{i_0\ldots i_{p + 1}} \\
& =
(-1)^{\frac{(p + 1)(p + 2)}{2}} d(\alpha)_{i_{p + 1} \ldots i_0}
\\
& =
(-1)^{\frac{(p + 1)(p + 2)}{2}}
\left(
\sum\nolimits_{j = 0}^{p + 1}
(-1)^j
\alpha_{i_{p + 1}\ldots \hat i_{p + 1 - j} \ldots i_0}
+
(-1)^{p + 1}
d_{\mathcal F}(\alpha_{i_{p + 1}\ldots i_0})
\right)
\end{align*}
由于
$p(p + 1)/2 \equiv (p + 1)(p + 2)/2 + p + 1 \bmod 2$，
可得 $d(\tau(\alpha)) = \tau(d(\alpha))$。换言之，
$\tau$ 是复形
$\text{Tot}(\check{\mathcal{C}}^\bullet({\mathcal U}, {\mathcal F}^\bullet))$
的自同态。注意，下图
$$
\begin{matrix}
\Gamma(X, {\mathcal F}^\bullet) &
\longrightarrow &
\text{Tot}(\check{\mathcal{C}}^\bullet({\mathcal U}, {\mathcal F}^\bullet)) \\
\downarrow \text{id} & & \downarrow \tau \\
\Gamma(X, {\mathcal F}^\bullet) &
\longrightarrow &
\text{Tot}(\check{\mathcal{C}}^\bullet({\mathcal U}, {\mathcal F}^\bullet))
\end{matrix}
$$
交换。此外，$\tau$ 显然与加细相容。这表明 $\tau$ 在
Cech 上同调上按恒等映射作用（即在余极限中如此——前提是
Cech 超上同调与超上同调一致；使用超覆盖时总是如此）。
我们断言，在总 Cech 复形
$\text{Tot}(\check{\mathcal{C}}^\bullet({\mathcal U}, {\mathcal F}^\bullet))$
上，$\tau$ 实际上与恒等映射链同伦。为证明这一点，取链同伦
$$
h(\alpha)_{i_0 \ldots i_p}
=
\sum\nolimits_{a = 0}^p
\epsilon_p(a)
\alpha_{i_0 \ldots i_a i_p \ldots i_a}
\quad\text{其中}\quad
\epsilon_p(a) = (-1)^{\frac{(p - a)(p - a - 1)}{2} + p}
$$
，其中 $\alpha$ 的次数为 $n$。与往常一样，我们省略
$|_{U_{i_0 \ldots i_p}}$。下面的计算表明此式成立；
这里仍设 $\alpha$ 为次数 $n$ 的元素：
\begin{align*}
(d(h(\alpha)) + h(d(\alpha)))_{i_0 \ldots i_p}
= &
\sum\nolimits_{k = 0}^p
(-1)^k
h(\alpha)_{i_0 \ldots \hat i_k \ldots i_p}
+ \\
&
(-1)^p
d_{\mathcal F}(h(\alpha)_{i_0 \ldots i_p})
+ \\
&
\sum\nolimits_{a = 0}^p
\epsilon_p(a)
d(\alpha)_{i_0 \ldots i_a i_p \ldots i_a}
\\
= &
\sum\nolimits_{k = 0}^p
\sum\nolimits_{a = 0}^{k - 1}
(-1)^k \epsilon_{p - 1}(a)
\alpha_{i_0 \ldots i_a i_p \ldots \hat{i_k} \ldots i_a}
+ \\
&
\sum\nolimits_{k = 0}^p
\sum\nolimits_{a = k + 1}^p
(-1)^k \epsilon_{p - 1}(a - 1)
\alpha_{i_0 \ldots \hat{i_k} \ldots i_a i_p \ldots i_a}
+ \\
&
\sum\nolimits_{a = 0}^p
(-1)^p \epsilon_p(a)
d_{\mathcal F}(\alpha_{i_0 \ldots i_a i_p \ldots i_a})
+ \\
&
\sum\nolimits_{a = 0}^p
\sum\nolimits_{k = 0}^a
\epsilon_p(a) (-1)^k
\alpha_{i_0 \ldots \hat{i_k} \ldots i_a i_p \ldots i_a}
+ \\
&
\sum\nolimits_{a = 0}^p
\sum\nolimits_{k = a}^p
\epsilon_p(a) (-1)^{p + a + 1 - k}
\alpha_{i_0 \ldots i_a i_p \ldots \hat{i_k} \ldots i_a}
+ \\
&
\sum\nolimits_{a = 0}^p
\epsilon_p(a) (-1)^{p + 1}
d_{\mathcal F}(\alpha_{i_0 \ldots i_a i_p \ldots i_a})
\\
= &
\epsilon_p(0) \alpha_{i_p \ldots i_0} +
\epsilon_p(p) (-1)^{p + 1} \alpha_{i_0 \ldots i_p} \\
= &
(-1)^{\frac{p(p + 1)}{2}}\alpha_{i_p \ldots i_0}
- \alpha_{i_0 \ldots i_p}
\end{align*}
各项相消是因为
$$
(-1)^k \epsilon_{p - 1}(a) + \epsilon_p(a)(-1)^{p + a + 1 - k} = 0
\quad\text{且}\quad
(-1)^k\epsilon_{p - 1}(a - 1) + \epsilon_p(a) (-1)^k = 0
$$
我们把相消的验证留给读者。

\medskip\noindent
设有阿贝尔层的两个有下界复形
${\mathcal F}^\bullet$ 与 ${\mathcal G}^\bullet$。我们把复形
$\text{Tot}({\mathcal F}^\bullet\otimes_{\mathbf Z} {\mathcal G}^\bullet)$
定义为各项为
$\bigoplus_{p + q = n} {\mathcal F}^p \otimes {\mathcal G}^q$
且微分遵循下列规则的复形：
\begin{equation}
\label{cohomology-equation-differential-tensor-product-complexes}
d(\alpha \otimes \beta) =
d(\alpha)\otimes \beta + (-1)^{\deg(\alpha)} \alpha \otimes d(\beta)
\end{equation}
，其中 $\alpha$ 与 $\beta$ 为齐次元素；见《同调代数》定义
\ref{homology-definition-associated-simple-complex}。

\medskip\noindent
设 $M^\bullet$ 与 $N^\bullet$ 是阿贝尔群的两个有下界复形。
若 $m$（相应地 $n$）是 $M^\bullet$（相应地 $N^\bullet$）
的上循环，则 $m \otimes n$ 显然是
$\text{Tot}(M^\bullet\otimes N^\bullet)$ 的上循环。因此有杯积
$$
H^i(M^\bullet) \times H^j(N^\bullet)
\longrightarrow
H^{i + j}(Tot(M^\bullet\otimes N^\bullet)).
$$
《代数详论》第 \ref{more-algebra-section-products-tor} 节也讨论了这一点。

\medskip\noindent
因此，复形超上同调中杯积的构造归结为构造复形映射
\begin{equation}
\label{cohomology-equation-needs-signs}
\text{Tot}\left(
\text{Tot}(\check{\mathcal{C}}^\bullet({\mathcal U}, {\mathcal F}^\bullet))
\otimes_{\mathbf Z}
\text{Tot}(\check{\mathcal{C}}^\bullet({\mathcal U}, {\mathcal G}^\bullet))
\right)
\longrightarrow
\text{Tot}(
\check{\mathcal{C}}^\bullet({\mathcal U},
\text{Tot}({\mathcal F}^\bullet\otimes {\mathcal G}^\bullet)
))
\end{equation}
此映射记为 $\cup$，并由下列规则给出：
$$
(\alpha \cup \beta)_{i_0 \ldots i_p}
=
\sum\nolimits_{r = 0}^p
\epsilon(n, m, p, r)
\alpha_{i_0 \ldots i_r} \otimes \beta_{i_r \ldots i_p}.
$$
，其中 $\alpha$ 的次数为 $n$，$\beta$ 的次数为 $m$，并且
$$
\epsilon(n, m, p, r) = (-1)^{(p + r)n + rp + r}.
$$
注意 $\epsilon(n, m, p, n) = 1$。因此，若
$\mathcal{F}^\bullet = \mathcal{F}[0]$ 是只由置于次数 $0$ 的
一个阿贝尔层 $\mathcal{F}$ 构成的复形，则 $\cup$ 的公式中没有符号
（因为此时除非 $r = n$，否则
$\alpha_{i_0 \ldots i_r} = 0$）。关于为何必须出现符号以及如何计算它，
见 Deligne 撰写的 \cite[Exposee XVII]{SGA4}。为验证
(\ref{cohomology-equation-needs-signs}) 是复形映射，必须证明
$$
d(\alpha \cup \beta) =
d(\alpha) \cup \beta +
(-1)^{\deg(\alpha)} \alpha \cup d(\beta)
$$
；这是由
$\text{Tot}(
\text{Tot}(\check{\mathcal{C}}^\bullet({\mathcal U}, {\mathcal F}^\bullet))
\otimes_{\mathbf Z}
\text{Tot}(\check{\mathcal{C}}^\bullet({\mathcal U}, {\mathcal G}^\bullet))
)$
上的微分之定义所要求的，该定义见《同调代数》定义
\ref{homology-definition-associated-simple-complex}。先计算
\begin{align*}
d(\alpha \cup \beta)_{i_0 \ldots i_{p + 1}}
= &
\sum\nolimits_{j = 0}^{p + 1}
(-1)^j
(\alpha \cup \beta)_{i_0 \ldots \hat i_j \ldots i_{p + 1}}
+
(-1)^{p + 1}
d_{{\mathcal F} \otimes {\mathcal G}}
((\alpha \cup \beta)_{i_0 \ldots i_{p + 1}})
\\
= &
\sum\nolimits_{j = 0}^{p + 1}
\sum\nolimits_{r = 0}^{j - 1}
(-1)^j \epsilon(n, m, p, r)
\alpha_{i_0 \ldots i_r} \otimes \beta_{i_r \ldots \hat i_j \ldots i_{p + 1}}
+ \\
&
\sum\nolimits_{j = 0}^{p + 1}
\sum\nolimits_{r = j + 1}^{p + 1}
(-1)^j \epsilon(n, m, p, r - 1)
\alpha_{i_0 \ldots \hat i_j \ldots i_r} \otimes \beta_{i_r \ldots i_{p + 1}}
+ \\
&
\sum\nolimits_{r = 0}^{p + 1}
(-1)^{p + 1} \epsilon(n, m, p + 1, r)
d_{{\mathcal F} \otimes {\mathcal G}}
(\alpha_{i_0 \ldots i_r} \otimes \beta_{i_r \ldots i_{p + 1}})
\end{align*}
并注意最后一项中的各和式项等于
$$
(-1)^{p + 1} \epsilon(n, m, p + 1, r)
\left(
d_{\mathcal F}(\alpha_{i_0 \ldots i_r}) \otimes 
\beta_{i_r \ldots i_{p + 1}} +
(-1)^{n - r}
\alpha_{i_0 \ldots i_r} \otimes d_{\mathcal G}(\beta_{i_r \ldots i_{p + 1}})
\right).
$$
，因为
$\deg_\mathcal{F}(\alpha_{i_0 \ldots i_r}) = n - r$。
另一方面，
\begin{align*}
(d(\alpha) \cup \beta)_{i_0\ldots i_{p + 1}}
= &
\sum\nolimits_{r = 0}^{p + 1}
\epsilon(n + 1, m, p + 1, r)
d(\alpha)_{i_0\ldots i_r} \otimes \beta_{i_r\ldots i_{p + 1}}
\\
= &
\sum\nolimits_{r = 0}^{p + 1}
\sum\nolimits_{j = 0}^{r}
\epsilon(n + 1, m, p + 1, r) (-1)^j
\alpha_{i_0\ldots\hat{i_j}\ldots i_r} \otimes \beta_{i_r\ldots i_{p + 1}}
+ \\
&
\sum\nolimits_{r = 0}^{p + 1}
\epsilon(n + 1, m, p + 1, r) (-1)^r
d_{\mathcal F}(\alpha_{i_0 \ldots i_r}) \otimes \beta_{i_r\ldots i_{p + 1}}
\end{align*}
且
\begin{align*}
(\alpha \cup d(\beta))_{i_0\ldots i_{p + 1}}
= &
\sum\nolimits_{r = 0}^{p + 1}
\epsilon(n, m + 1, p + 1, r)
\alpha_{i_0 \ldots i_r} \otimes d(\beta)_{i_r \ldots i_{p + 1}}
\\
= &
\sum\nolimits_{r = 0}^{p + 1}
\sum\nolimits_{j = r}^{p + 1}
\epsilon(n, m + 1, p + 1, r) (-1)^{j - r}
\alpha_{i_0 \ldots i_r} \otimes \beta_{i_r \ldots \hat{i_j}\ldots i_{p + 1}}
+ \\
&
\sum\nolimits_{r = 0}^{p + 1}
\epsilon(n, m + 1, p + 1, r) (-1)^{p + 1 - r}
\alpha_{i_0 \ldots i_r} \otimes d_{\mathcal G}(\beta_{i_r \ldots i_{p + 1}})
\end{align*}
若下列等式成立，则所需等式成立：
\begin{align*}
(-1)^{p + 1} \epsilon(n, m, p + 1, r)
& =
\epsilon(n + 1, m, p + 1, r) (-1)^r \\
(-1)^{p + 1} \epsilon(n, m, p + 1, r) (-1)^{n - r}
& =
(-1)^n \epsilon(n, m + 1, p + 1, r) (-1)^{p + 1 - r} \\
\epsilon(n + 1, m, p + 1, r) (-1)^r
& =
(-1)^{1 + n} \epsilon(n, m + 1, p + 1, r - 1) \\
(-1)^j \epsilon(n, m, p, r)
& =
(-1)^n \epsilon(n, m + 1, p + 1, r) (-1)^{j - r} \\
(-1)^j \epsilon(n, m, p, r - 1)
& =
\epsilon(n + 1, m, p + 1, r) (-1)^j
\end{align*}
（第三个等式是使 $d(\alpha) \cup \beta$ 与
$(-1)^n \alpha \cup d(\beta)$ 中 $r = j$ 的各项彼此相消所必需的。）
我们把验证留给读者。（也可检查 Stacks Project 的
$scripts$ 子目录中的脚本 $signs.gp$。）

\medskip\noindent
杯积的结合律。设 ${\mathcal F}^\bullet$、
${\mathcal G}^\bullet$ 与 ${\mathcal H}^\bullet$ 是 $X$ 上
阿贝尔群的有下界复形。不涉及任何符号的显然映射
$$
\text{Tot}(
\text{Tot}({\mathcal F}^\bullet \otimes_{\mathbf Z} {\mathcal G}^\bullet)
\otimes_{\mathbf Z} {\mathcal H}^\bullet
)
\longrightarrow
\text{Tot}(
{\mathcal F}^\bullet \otimes_{\mathbf Z}
\text{Tot}({\mathcal G}^\bullet \otimes_{\mathbf Z} {\mathcal H}^\bullet)
).
$$
是复形的同构。换言之，三重复形
${\mathcal F}^\bullet \otimes_{\mathbf Z} {\mathcal G}^\bullet
\otimes_{\mathbf Z} {\mathcal H}^\bullet$ 给出定义良好的总复形，
其微分满足
$$
d(\alpha \otimes \beta \otimes \gamma) =
d(\alpha) \otimes \beta \otimes \gamma +
(-1)^{\deg(\alpha)} \alpha \otimes d(\beta) \otimes \gamma +
(-1)^{\deg(\alpha) + \deg(\beta)} \alpha \otimes \beta \otimes d(\gamma)
$$
，其中各元素为齐次元素。利用此映射容易验证
$$
(\alpha \cup \beta) \cup \gamma = \alpha \cup ( \beta \cup \gamma)
$$
。事实上，若 $\alpha$ 的次数为 $a$，$\beta$ 的次数为 $b$，
而 $\gamma$ 的次数为 $c$，则
\begin{align*}
((\alpha \cup \beta) \cup \gamma)_{i_0 \ldots i_p}
= &
\sum\nolimits_{r = 0}^p
\epsilon(a + b, c, p, r)
(\alpha \cup \beta)_{i_0 \ldots i_r} \otimes \gamma_{i_r \ldots i_p}
\\
= &
\sum\nolimits_{r = 0}^p
\sum\nolimits_{s = 0}^r
\epsilon(a + b, c, p, r) \epsilon(a, b, r, s)
\alpha_{i_0 \ldots i_s} \otimes
\beta_{i_s \ldots i_r} \otimes
\gamma_{i_r \ldots i_p}
\end{align*}
并且
\begin{align*}
(\alpha \cup (\beta \cup \gamma)_{i_0\ldots i_p}
= &
\sum\nolimits_{s = 0}^p
\epsilon(a, b + c, p, s)
\alpha_{i_0 \ldots i_s} \otimes (\beta \cup \gamma)_{i_s \ldots i_p}
\\
= &
\sum\nolimits_{s = 0}^p
\sum\nolimits_{r = s}^p
\epsilon(a, b + c, p, s) \epsilon(b, c, p - s, r - s)
\alpha_{i_0 \ldots i_s} \otimes \beta_{i_s \ldots i_r} \otimes
\gamma_{i_r \ldots i_p}
\end{align*}
一个直接的模 $2$ 计算表明这些符号相符。
（也可检查 Stacks Project 的 $scripts$ 子目录中的脚本 $signs.gp$。）

\medskip\noindent
最后说明杯积为何至少在上同调层面保持分次交换结构。为此使用上面引入的
算子 $\tau$。设 ${\mathcal F}^\bullet$ 是阿贝尔群的有下界复形，
并假设给定分次交换乘法
$$
\wedge^\bullet :
\text{Tot}({\mathcal F}^\bullet\otimes {\mathcal F}^\bullet)
\longrightarrow
{\mathcal F}^\bullet.
$$
这意味着：若 $s$ 是 ${\mathcal F}^a$ 的局部截面，而 $t$ 是
${\mathcal F}^b$ 的局部截面，则 $s \wedge t$ 是
${\mathcal F}^{a + b}$ 的局部截面。分次交换是指
$s \wedge t = (-1)^{ab} t \wedge s$。由于 $\wedge$ 是复形映射，故有
$d(s\wedge t) = d(s) \wedge t + (-1)^a s \wedge d(t)$.
复合
$$
\text{Tot}(
\text{Tot}(\check{\mathcal{C}}^\bullet({\mathcal U}, {\mathcal F}^\bullet))
\otimes
\text{Tot}(\check{\mathcal{C}}^\bullet({\mathcal U}, {\mathcal F}^\bullet))
)
\to
\text{Tot}(
\check{\mathcal{C}}^\bullet({\mathcal U},
\text{Tot}({\mathcal F}^\bullet\otimes_{\mathbf Z}{\mathcal F}^\bullet))
)
\to
\text{Tot}(\check{\mathcal{C}}^\bullet({\mathcal U}, {\mathcal F}^\bullet))
$$
在上同调上诱导杯积
$$
H^n(
\text{Tot}(\check{\mathcal{C}}^\bullet({\mathcal U}, {\mathcal F}^\bullet))
)
\times
H^m(
\text{Tot}(\check{\mathcal{C}}^\bullet({\mathcal U}, {\mathcal F}^\bullet))
)
\longrightarrow
H^{n + m}(
\text{Tot}(\check{\mathcal{C}}^\bullet({\mathcal U}, {\mathcal F}^\bullet))
)
$$
，从而在余极限中也诱导 Cech 上同调上的乘法，并因而（必要时使用
超覆盖）诱导 ${\mathcal F}^\bullet$ 的上同调上的乘法。
我们断言这个（上同调上的）乘法也分次交换。为证明这一点，先取
$\text{Tot}(\check{\mathcal{C}}^\bullet({\mathcal U}, {\mathcal F}^\bullet))$
中次数为 $n$ 的元素 $\alpha$，以及
$\text{Tot}(\check{\mathcal{C}}^\bullet({\mathcal U}, {\mathcal F}^\bullet))$
中次数为 $m$ 的元素 $\beta$，并计算
\begin{align*}
\wedge^\bullet(\alpha \cup \beta)_{i_0 \ldots i_p}
= &
\sum\nolimits_{r = 0}^p
\epsilon(n, m, p, r)
\alpha_{i_0 \ldots i_r} \wedge \beta_{i_r \ldots i_p} \\
= &
\sum\nolimits_{r = 0}^p
\epsilon(n, m, p, r)
(-1)^{\deg(\alpha_{i_0 \ldots i_r})\deg(\beta_{i_r \ldots i_p})}
\beta_{i_r \ldots i_p} \wedge \alpha_{i_0 \ldots i_r}
\end{align*}
，这里使用了 $\wedge$ 的分次交换性。另一方面有
\begin{align*}
\tau(\wedge^\bullet(\tau(\beta) \cup \tau(\alpha)))_{i_0 \ldots i_p}
= &
\chi(p)
\sum\nolimits_{r = 0}^p
\epsilon(m, n, p, r)
\tau(\beta)_{i_p \ldots i_{p - r}} \wedge \tau(\alpha)_{i_{p - r} \ldots i_0}
\\
= &
\chi(p)
\sum\nolimits_{r = 0}^p
\epsilon(m, n, p, r) \chi(r) \chi(p - r)
\beta_{i_{p - r} \ldots i_p} \wedge \alpha_{i_0 \ldots i_{p - r}}
\\
= &
\chi(p)
\sum\nolimits_{r = 0}^p
\epsilon(m, n, p, p - r) \chi(r) \chi(p - r)
\beta_{i_r \ldots i_p} \wedge \alpha_{i_0 \ldots i_r}
\end{align*}
，其中 $\chi(t) = (-1)^{\frac{t(t + 1)}{2}}$。前面已经证明
$\tau$ 在上同调上按恒等映射作用，因此只需验证
$$
\epsilon(n, m, p, r)
(-1)^{(n - r)(m - (p - r))}
=
(-1)^{nm} \chi(p)\epsilon(m, n, p, p - r) \chi(r) \chi(p - r)
$$
一个直接的模 $2$ 计算表明这些符号相符。
（也可检查 Stacks Project 的 $scripts$ 子目录中的脚本 $signs.gp$。）

\medskip\noindent
最后研究杯积与连接映射的相容性。设
$$
0
\to
{\mathcal F}_1^\bullet
\to
{\mathcal F}_2^\bullet
\to
{\mathcal F}_3^\bullet
\to
0
\quad\text{且}\quad
0
\leftarrow
{\mathcal G}_1^\bullet
\leftarrow
{\mathcal G}_2^\bullet
\leftarrow
{\mathcal G}_3^\bullet
\leftarrow
0
$$
是 $X$ 上阿贝尔层的有下界复形的短正合列。再设
${\mathcal H}^\bullet$ 是阿贝尔层的另一个有下界复形，并假设有复形映射
$$
\gamma_i :
\text{Tot}({\mathcal F}_i^\bullet \otimes_{\mathbf Z} {\mathcal G}_i^\bullet)
\longrightarrow
{\mathcal H}^\bullet
$$
，它们与这些复形之间的映射相容，即下列各图
$$
\xymatrix{
\text{Tot}({\mathcal F}_1^\bullet \otimes_{\mathbf Z} {\mathcal G}_1^\bullet)
\ar[d]_{\gamma_1}
&
\text{Tot}({\mathcal F}_1^\bullet \otimes_{\mathbf Z} {\mathcal G}_2^\bullet)
\ar[l] \ar[d]
\\
\mathcal{H}^\bullet &
\text{Tot}({\mathcal F}_2^\bullet \otimes_{\mathbf Z} {\mathcal G}_2^\bullet)
\ar[l]_-{\gamma_2}
}
$$
且
$$
\xymatrix{
\text{Tot}({\mathcal F}_2^\bullet \otimes_{\mathbf Z} {\mathcal G}_2^\bullet)
\ar[d]_{\gamma_2}
&
\text{Tot}({\mathcal F}_2^\bullet \otimes_{\mathbf Z} {\mathcal G}_3^\bullet)
\ar[l] \ar[d]
\\
\mathcal{H}^\bullet &
\text{Tot}({\mathcal F}_3^\bullet \otimes_{\mathbf Z} {\mathcal G}_3^\bullet)
\ar[l]_-{\gamma_3}
}
$$
交换。

\begin{lemma}
\label{cohomology-lemma-compute-sign-cup-product-boundaries}
在上述情形中，假设对各层 $\mathcal{F}_i^p$ 与
$\mathcal{G}_j^q$，Cech 上同调与上同调一致。设
$a_3 \in H^n(X, \mathcal{F}_3^\bullet)$ 且
$b_1 \in H^m(X, \mathcal{G}_1^\bullet)$。则有
$$
\gamma_1( \partial a_3 \cup b_1) =
(-1)^{n + 1} \gamma_3( a_3 \cup \partial b_1)
$$
在 $H^{n + m + 1}(X, \mathcal{H}^\bullet)$ 中成立，其中
$\partial$ 表示由上述复形短正合列所诱导的上同调连接映射。
\end{lemma}

\begin{proof}
采用下列约定与记号。把 ${\mathcal F}_1^p$ 看成
${\mathcal F}_2^p$ 的子层，并把 ${\mathcal G}_3^q$ 看成
${\mathcal G}_2^q$ 的子层。因此，若 $s$ 是
${\mathcal F}_1^p$ 的局部截面，也用 $s$ 表示
${\mathcal F}_2^p$ 的相应截面；对 ${\mathcal G}_3^q$ 的局部截面
也作同样约定。此外，若 $s$ 是 ${\mathcal F}_2^p$ 的局部截面，
则用 $\bar s$ 表示它在 ${\mathcal F}_3^p$ 中的像。
对映射 ${\mathcal G}_2^q \to {\mathcal G}^q_1$ 也作同样约定。
特别地，若 $s$ 是 ${\mathcal F}_2^p$ 的局部截面且
$\bar s = 0$，则 $s$ 是 ${\mathcal F}_1^p$ 的局部截面。
上述图的交换性蕴含：对 ${\mathcal F}_2^p$ 的局部截面 $s$
和 ${\mathcal G}_3^q$ 的局部截面 $t$，作为
${\mathcal H}^{p + q}$ 的截面有
$\gamma_2(s \otimes t) = \gamma_3(\bar s \otimes t)$。

\medskip\noindent
设 ${\mathcal U} : X =  \bigcup_{i \in I} U_i$ 是 $X$ 的开覆盖。
假设 $\alpha_3$（相应地 $\beta_1$）是
$\text{Tot}(
\check{\mathcal{C}}^\bullet({\mathcal U}, {\mathcal F}_3^\bullet))$
（相应地
$\text{Tot}(
\check{\mathcal{C}}^\bullet({\mathcal U}, {\mathcal G}_1^\bullet))$）
中次数为 $n$（相应地 $m$）的上循环，代表
$a_3$（相应地 $b_1$）。必要时加细 $\mathcal{U}$ 后，可以找到
$\text{Tot}(
\check{\mathcal{C}}^\bullet({\mathcal U}, {\mathcal F}_2^\bullet))$
中次数为 $n$ 的上链 $\alpha_2$，以及
$\text{Tot}(
\check{\mathcal{C}}^\bullet({\mathcal U}, {\mathcal G}_2^\bullet))$
中次数为 $m$ 的上链 $\beta_2$，它们分别映到
$\alpha_3$ 与 $\beta_1$。于是
$$
\overline{d(\alpha_2)} = d(\bar \alpha_2) = 0
\quad\text{且}\quad
\overline{d(\beta_2)} = d(\bar \beta_2) = 0.
$$
这意味着 $\alpha_1 = d(\alpha_2)$ 是
$\text{Tot}(\check{\mathcal{C}}^\bullet({\mathcal U}, {\mathcal F}_1^\bullet))$
中次数为 $n + 1$ 的上循环，代表 $\partial a_3$。类似地，
$\beta_3 = d(\beta_2)$ 是
$\text{Tot}(\check{\mathcal{C}}^\bullet({\mathcal U}, {\mathcal G}_3^\bullet))$
中次数为 $m + 1$ 的上循环，代表 $\partial b_1$。
因此可作如下计算：
\begin{align*}
d(\gamma_2(\alpha_2 \cup \beta_2))
& =
\gamma_2(d(\alpha_2 \cup \beta_2))
\\
& =
\gamma_2(d(\alpha_2) \cup \beta_2 + (-1)^n \alpha_2 \cup d(\beta_2) )
\\
& =
\gamma_2( \alpha_1 \cup \beta_2)  + (-1)^n \gamma_2( \alpha_2 \cup \beta_3)
\\
& =
\gamma_1(\alpha_1 \cup \beta_1) + (-1)^n \gamma_3(\alpha_3 \cup \beta_3)
\end{align*}
这甚至表明符号确如引理所述为 $(-1)^{n + 1}$
\footnote{该符号取决于与 $D(X)$ 中三角所对应的上同调长正合列的
符号约定。Stacks Project 的约定是：（a）特异三角对应于逐项分裂的
正合列；（b）长正合列中的连接映射由蛇引理中的映射给出，不另加符号。
见《导出范畴》第 \ref{derived-section-homotopy-triangulated} 节。}。
\end{proof}

\begin{lemma}
\label{cohomology-lemma-boundary-derivation-over-cup-product}
设 $X$ 是拓扑空间。设 $\mathcal{O}' \to \mathcal{O}$ 是环层的满射，
其核 $\mathcal{I} \subset \mathcal{O}'$ 的平方为零。则
$M = H^1(X, \mathcal{I})$ 是 $R = H^0(X, \mathcal{O})$-模，
并且由短正合列
$$
0 \to \mathcal{I} \to \mathcal{O}' \to \mathcal{O} \to 0
$$
所诱导的连接映射 $\partial : R \to M$ 是导子
（《代数》定义 \ref{algebra-definition-derivation}）。
\end{lemma}

\begin{proof}
由假设 $\mathcal{I} \cdot \mathcal{I} = 0$，映射
$\mathcal{O}' \to \SheafHom(\mathcal{I}, \mathcal{I})$
通过 $\mathcal{O}$ 分解。因此 $\mathcal{I}$ 是
$\mathcal{O}$-模层，这定义了 $M$ 上的 $R$-模结构。
连接映射是可加的，故只需证明 Leibniz 法则。设 $f \in R$。
选取开覆盖 $\mathcal{U} : X = \bigcup U_i$，使得存在
$f_i \in \mathcal{O}'(U_i)$ 提升
$f|_{U_i} \in \mathcal{O}(U_i)$。注意
$f_i - f_j$ 是 $\mathcal{I}(U_i \cap U_j)$ 的元素。
于是 $\partial(f)$ 对应于上循环 $\alpha$ 的 Cech 上同调类，
其中 $\alpha_{i_0i_1} = f_{i_0} - f_{i_1}$。
（注意，由引理 \ref{cohomology-lemma-cech-h1}，关于 $\mathcal{U}$ 的第一
Cech 上同调群是 $M$ 的子模。）
接着取另一个元素 $g \in R$，并假设（必要时加细开覆盖后）
$g_i \in \mathcal{O}'(U_i)$ 提升
$g|_{U_i} \in \mathcal{O}(U_i)$。于是
$\partial(g)$ 由上循环 $\beta$ 给出，其中
$\beta_{i_0i_1} = g_{i_0} - g_{i_1}$。由于
$f_ig_i \in \mathcal{O}'(U_i)$ 提升 $fg|_{U_i}$，
可知 $\partial(fg)$ 由上循环 $\gamma$ 给出，其中
$$
\gamma_{i_0i_1} = f_{i_0}g_{i_0} - f_{i_1}g_{i_1} =
(f_{i_0} - f_{i_1})g_{i_0} + f_{i_1}(g_{i_0} - g_{i_1}) =
\alpha_{i_0i_1}g + f\beta_{i_0i_1}
$$
，这里使用了我们对 $\mathcal{I}$ 上 $\mathcal{O}$-模结构的定义。
这就证明了 Leibniz 法则，证明完毕。
\end{proof}









\section{平坦消解}
\label{cohomology-section-flat}

\noindent
本节材料可参考 \cite{Spaltenstein}。设
$(X, \mathcal{O}_X)$ 是环空间。由《模层》引理
\ref{modules-lemma-module-quotient-flat}，任意 $\mathcal{O}_X$-模
都是某个平坦 $\mathcal{O}_X$-模的商。由《导出范畴》引理
\ref{derived-lemma-subcategory-left-resolution}，任意
$\mathcal{O}_X$-模的有上界复形都具有由平坦
$\mathcal{O}_X$-模的有上界复形给出的左消解。然而，对无界复形而言，
事实表明平坦消解还不够好。

\begin{lemma}
\label{cohomology-lemma-derived-tor-exact}
设 $(X, \mathcal{O}_X)$ 是环空间，而
$\mathcal{G}^\bullet$ 是 $\mathcal{O}_X$-模复形。函子
$$
K(\textit{Mod}(\mathcal{O}_X))
\longrightarrow
K(\textit{Mod}(\mathcal{O}_X)),
\quad
\mathcal{F}^\bullet \longmapsto
\text{Tot}(\mathcal{G}^\bullet \otimes_{\mathcal{O}_X} \mathcal{F}^\bullet)
$$
以及
$$
K(\textit{Mod}(\mathcal{O}_X))
\longrightarrow
K(\textit{Mod}(\mathcal{O}_X)),
\quad
\mathcal{F}^\bullet \longmapsto
\text{Tot}(\mathcal{F}^\bullet \otimes_{\mathcal{O}_X} \mathcal{G}^\bullet)
$$
是三角范畴之间的正合函子。
\end{lemma}

\begin{proof}
这由《导出范畴》注记
\ref{derived-remark-double-complex-as-tensor-product-of} 得到。
\end{proof}

\begin{definition}
\label{cohomology-definition-K-flat}
设 $(X, \mathcal{O}_X)$ 是环空间。称 $\mathcal{O}_X$-模复形
$\mathcal{K}^\bullet$ 为 {\it K-平坦的}，如果对任意
$\mathcal{O}_X$-模的无环复形 $\mathcal{F}^\bullet$，复形
$$
\text{Tot}(\mathcal{F}^\bullet \otimes_{\mathcal{O}_X} \mathcal{K}^\bullet)
$$
都是无环的。
\end{definition}

\begin{lemma}
\label{cohomology-lemma-K-flat-quasi-isomorphism}
设 $(X, \mathcal{O}_X)$ 是环空间，而
$\mathcal{K}^\bullet$ 是 K-平坦复形。则函子
$$
K(\textit{Mod}(\mathcal{O}_X))
\longrightarrow
K(\textit{Mod}(\mathcal{O}_X)), \quad
\mathcal{F}^\bullet
\longmapsto
\text{Tot}(\mathcal{F}^\bullet \otimes_{\mathcal{O}_X} \mathcal{K}^\bullet)
$$
将拟同构映为拟同构。
\end{lemma}

\begin{proof}
这由引理 \ref{cohomology-lemma-derived-tor-exact} 以及拟同构可由其锥无环刻画这一事实得到。
\end{proof}

\begin{lemma}
\label{cohomology-lemma-check-K-flat-stalks}
设 $(X, \mathcal{O}_X)$ 是环空间，且 $\mathcal{K}^\bullet$
是 $\mathcal{O}_X$-模复形。则 $\mathcal{K}^\bullet$ 是
K-平坦的，当且仅当对所有 $x \in X$，
$\mathcal{O}_{X, x}$-模复形 $\mathcal{K}_x^\bullet$ 都是
K-平坦的（《代数详论》定义 \ref{more-algebra-definition-K-flat}）。
\end{lemma}

\begin{proof}
若对所有 $x \in X$，$\mathcal{K}_x^\bullet$ 都是 K-平坦的，
则 $\mathcal{K}^\bullet$ 是 K-平坦的；这是因为 $\otimes$ 与直和
都和取茎交换，并且正合性可在茎上检验，见《模层》引理
\ref{modules-lemma-abelian}。反过来，假设
$\mathcal{K}^\bullet$ 是 K-平坦的。取 $x \in X$，并设
$M^\bullet$ 是 $\mathcal{O}_{X, x}$-模的无环复形。于是
$i_{x, *}M^\bullet$ 是 $\mathcal{O}_X$-模的无环复形，因而
$\text{Tot}(i_{x, *}M^\bullet \otimes_{\mathcal{O}_X}
\mathcal{K}^\bullet)$ 无环。在 $x$ 处取茎可知
$\text{Tot}(M^\bullet \otimes_{\mathcal{O}_{X, x}} \mathcal{K}_x^\bullet)$
无环。
\end{proof}

\begin{lemma}
\label{cohomology-lemma-tensor-product-K-flat}
设 $(X, \mathcal{O}_X)$ 是环空间。若
$\mathcal{K}^\bullet$、$\mathcal{L}^\bullet$ 是
$\mathcal{O}_X$-模的 K-平坦复形，则
$\text{Tot}(\mathcal{K}^\bullet \otimes_{\mathcal{O}_X} \mathcal{L}^\bullet)$
是 $\mathcal{O}_X$-模的 K-平坦复形。
\end{lemma}

\begin{proof}
这由同构
$$
\text{Tot}(\mathcal{M}^\bullet \otimes_{\mathcal{O}_X}
\text{Tot}(\mathcal{K}^\bullet \otimes_{\mathcal{O}_X} \mathcal{L}^\bullet))
=
\text{Tot}(\text{Tot}(\mathcal{M}^\bullet \otimes_{\mathcal{O}_X}
\mathcal{K}^\bullet) \otimes_{\mathcal{O}_X} \mathcal{L}^\bullet)
$$
和定义得到。
\end{proof}

\begin{lemma}
\label{cohomology-lemma-K-flat-two-out-of-three}
设 $(X, \mathcal{O}_X)$ 是环空间。设
$(\mathcal{K}_1^\bullet, \mathcal{K}_2^\bullet,
\mathcal{K}_3^\bullet)$ 是
$K(\textit{Mod}(\mathcal{O}_X))$ 中的特异三角。
若三个 $\mathcal{K}_i^\bullet$ 中有两个是 K-平坦的，
则第三个也是 K-平坦的。
\end{lemma}

\begin{proof}
这由引理 \ref{cohomology-lemma-derived-tor-exact} 以及下述事实得到：
在 $K(\textit{Mod}(\mathcal{O}_X))$ 的特异三角中，
若三者中有两个无环，则第三个也无环。
\end{proof}

\begin{lemma}
\label{cohomology-lemma-K-flat-two-out-of-three-ses}
设 $(X, \mathcal{O}_X)$ 是环空间。设
$0 \to \mathcal{K}_1^\bullet \to \mathcal{K}_2^\bullet \to
\mathcal{K}_3^\bullet \to 0$ 是复形的短正合列，并且
$\mathcal{K}_3^\bullet$ 的各项都是平坦
$\mathcal{O}_X$-模。若三个 $\mathcal{K}_i^\bullet$ 中有两个
是 K-平坦的，则第三个也是 K-平坦的。
\end{lemma}

\begin{proof}
由《模层》引理 \ref{modules-lemma-flat-tor-zero}，对每个复形
$\mathcal{L}^\bullet$ 都得到复形的短正合列
$$
0 \to
\text{Tot}(\mathcal{L}^\bullet \otimes_{\mathcal{O}_X} \mathcal{K}_1^\bullet)
\to
\text{Tot}(\mathcal{L}^\bullet \otimes_{\mathcal{O}_X} \mathcal{K}_1^\bullet)
\to
\text{Tot}(\mathcal{L}^\bullet \otimes_{\mathcal{O}_X} \mathcal{K}_1^\bullet)
\to 0
$$
。因此该引理由上同调层长正合列和 K-平坦复形的定义得到。
\end{proof}

\begin{lemma}
\label{cohomology-lemma-pullback-K-flat}
设 $f : (X, \mathcal{O}_X) \to (Y, \mathcal{O}_Y)$ 是环空间的态射。
$\mathcal{O}_Y$-模的 K-平坦复形的拉回是
$\mathcal{O}_X$-模的 K-平坦复形。
\end{lemma}

\begin{proof}
可在茎上检验，见引理 \ref{cohomology-lemma-check-K-flat-stalks}。
因此这由《层论》引理 \ref{sheaves-lemma-stalk-pullback-modules}
和《代数详论》引理 \ref{more-algebra-lemma-base-change-K-flat} 得到。
\end{proof}

\begin{lemma}
\label{cohomology-lemma-bounded-flat-K-flat}
设 $(X, \mathcal{O}_X)$ 是环空间。平坦
$\mathcal{O}_X$-模的有上界复形是 K-平坦的。
\end{lemma}

\begin{proof}
可在茎上检验，见引理 \ref{cohomology-lemma-check-K-flat-stalks}。
故该引理由《模层》引理 \ref{modules-lemma-flat-stalks-flat}
和《代数详论》引理
\ref{more-algebra-lemma-derived-tor-quasi-isomorphism} 得到。
\end{proof}

\noindent
在下列引理中，复形系统的余极限是指逐项余极限。

\begin{lemma}
\label{cohomology-lemma-colimit-K-flat}
设 $(X, \mathcal{O}_X)$ 是环空间，且
$\mathcal{K}_1^\bullet \to \mathcal{K}_2^\bullet \to \ldots$
是 K-平坦复形系统。则
$\colim_i \mathcal{K}_i^\bullet$ 是 K-平坦的。
\end{lemma}

\begin{proof}
由于取的是逐项余极限，显然
$$
\colim_i \text{Tot}(
\mathcal{F}^\bullet \otimes_{\mathcal{O}_X} \mathcal{K}_i^\bullet)
=
\text{Tot}(\mathcal{F}^\bullet \otimes_{\mathcal{O}_X}
\colim_i \mathcal{K}_i^\bullet)
$$
因此该引理由滤过余极限为正合这一事实得到。
\end{proof}

\begin{lemma}
\label{cohomology-lemma-resolution-by-direct-sums-extensions-by-zero}
设 $(X, \mathcal{O}_X)$ 是环空间。对任意
$\mathcal{O}_X$-模复形 $\mathcal{G}^\bullet$，存在
$\mathcal{O}_X$-模复形的交换图
$$
\xymatrix{
\mathcal{K}_1^\bullet \ar[d] \ar[r] &
\mathcal{K}_2^\bullet \ar[d] \ar[r] & \ldots \\
\tau_{\leq 1}\mathcal{G}^\bullet \ar[r] &
\tau_{\leq 2}\mathcal{G}^\bullet \ar[r] & \ldots
}
$$
，满足下列性质：（1）竖直箭头是拟同构且逐项满射；
（2）每个 $\mathcal{K}_n^\bullet$ 都是有上界复形，其各项为形如
$j_{U!}\mathcal{O}_U$ 的 $\mathcal{O}_X$-模的直和；
（3）映射
$\mathcal{K}_n^\bullet \to \mathcal{K}_{n + 1}^\bullet$
是逐项分裂单射，其余核为形如 $j_{U!}\mathcal{O}_U$ 的
$\mathcal{O}_X$-模的直和。此外，映射
$\colim \mathcal{K}_n^\bullet \to \mathcal{G}^\bullet$ 是拟同构。
\end{lemma}

\begin{proof}
该图以及性质（1）、（2）、（3）的存在性立即由《模层》引理
\ref{modules-lemma-module-quotient-flat} 和《导出范畴》引理
\ref{derived-lemma-special-direct-system} 得到。由于滤过余极限正合，
诱导映射
$\colim \mathcal{K}_n^\bullet \to \mathcal{G}^\bullet$ 是拟同构。
\end{proof}

\begin{lemma}
\label{cohomology-lemma-K-flat-resolution}
设 $(X, \mathcal{O}_X)$ 是环空间。对任意复形
$\mathcal{G}^\bullet$，存在各项为平坦 $\mathcal{O}_X$-模的
$K$-平坦复形 $\mathcal{K}^\bullet$，以及逐项满射的拟同构
$\mathcal{K}^\bullet \to \mathcal{G}^\bullet$。
\end{lemma}

\begin{proof}
选取引理 \ref{cohomology-lemma-resolution-by-direct-sums-extensions-by-zero}
中的图。每个复形 $\mathcal{K}_n^\bullet$ 都是平坦模的有上界复形，
见《模层》引理 \ref{modules-lemma-j-shriek-flat}。故由引理
\ref{cohomology-lemma-bounded-flat-K-flat}，$\mathcal{K}_n^\bullet$ 是
K-平坦的；再由引理 \ref{cohomology-lemma-colimit-K-flat}，
$\colim \mathcal{K}_n^\bullet$ 是 K-平坦的。按构造，诱导映射
$\colim \mathcal{K}_n^\bullet \to \mathcal{G}^\bullet$
是拟同构且逐项满射。引理
\ref{cohomology-lemma-resolution-by-direct-sums-extensions-by-zero} 的性质（3）
表明 $\colim \mathcal{K}_n^m$ 是平坦模的直和，因而平坦；
这就证明了最后一个断言。
\end{proof}

\begin{lemma}
\label{cohomology-lemma-derived-tor-quasi-isomorphism-other-side}
设 $(X, \mathcal{O}_X)$ 是环空间。设
$\alpha : \mathcal{P}^\bullet \to \mathcal{Q}^\bullet$ 是
$\mathcal{O}_X$-模的 K-平坦复形之间的拟同构。对任意
$\mathcal{O}_X$-模复形 $\mathcal{F}^\bullet$，诱导映射
$$
\text{Tot}(\text{id}_{\mathcal{F}^\bullet} \otimes \alpha) :
\text{Tot}(\mathcal{F}^\bullet \otimes_{\mathcal{O}_X} \mathcal{P}^\bullet)
\longrightarrow
\text{Tot}(\mathcal{F}^\bullet \otimes_{\mathcal{O}_X} \mathcal{Q}^\bullet)
$$
是拟同构。
\end{lemma}

\begin{proof}
选取拟同构 $\mathcal{K}^\bullet \to \mathcal{F}^\bullet$，
其中 $\mathcal{K}^\bullet$ 是 K-平坦复形，见引理
\ref{cohomology-lemma-K-flat-resolution}。考虑交换图
$$
\xymatrix{
\text{Tot}(\mathcal{K}^\bullet
\otimes_{\mathcal{O}_X} \mathcal{P}^\bullet) \ar[r] \ar[d] &
\text{Tot}(\mathcal{K}^\bullet
\otimes_{\mathcal{O}_X} \mathcal{Q}^\bullet) \ar[d] \\
\text{Tot}(\mathcal{F}^\bullet
\otimes_{\mathcal{O}_X} \mathcal{P}^\bullet) \ar[r] &
\text{Tot}(\mathcal{F}^\bullet
\otimes_{\mathcal{O}_X} \mathcal{Q}^\bullet)
}
$$
由引理 \ref{cohomology-lemma-K-flat-quasi-isomorphism}，竖直箭头和上方水平箭头
都是拟同构，故结论成立。
\end{proof}

\noindent
设 $(X, \mathcal{O}_X)$ 是环空间，且
$\mathcal{F}^\bullet$ 是 $D(\mathcal{O}_X)$ 的对象。
选取 K-平坦消解
$\mathcal{K}^\bullet \to \mathcal{F}^\bullet$，见引理
\ref{cohomology-lemma-K-flat-resolution}。由引理
\ref{cohomology-lemma-derived-tor-exact} 得到三角范畴的正合函子
$$
K(\mathcal{O}_X)
\longrightarrow
K(\mathcal{O}_X),
\quad
\mathcal{G}^\bullet
\longmapsto
\text{Tot}(\mathcal{G}^\bullet \otimes_{\mathcal{O}_X} \mathcal{K}^\bullet)
$$
由引理 \ref{cohomology-lemma-K-flat-quasi-isomorphism}，此函子诱导函子
$D(\mathcal{O}_X) \to D(\mathcal{O}_X)$；这是因为
$D(\mathcal{O}_X)$ 是 $K(\mathcal{O}_X)$ 关于拟同构的局部化。
由于 $\mathcal{F}^\bullet$ 的 $K$-平坦消解所成的范畴是余滤过的，
再结合引理 \ref{cohomology-lemma-derived-tor-quasi-isomorphism-other-side}，
所得函子（在同构意义下）与 $\mathcal{K}^\bullet$ 的选取无关。

\begin{definition}
\label{cohomology-definition-derived-tor}
设 $(X, \mathcal{O}_X)$ 是环空间，而
$\mathcal{F}^\bullet$ 是 $D(\mathcal{O}_X)$ 的对象。{\it 导出张量积}
$$
- \otimes_{\mathcal{O}_X}^{\mathbf{L}} \mathcal{F}^\bullet :
D(\mathcal{O}_X)
\longrightarrow
D(\mathcal{O}_X)
$$
是上面所述的三角范畴正合函子。
\end{definition}

\noindent
由上述显式构造，显然有典范同构
$$
\mathcal{F}^\bullet \otimes_{\mathcal{O}_X}^{\mathbf{L}} \mathcal{G}^\bullet
\cong
\mathcal{G}^\bullet \otimes_{\mathcal{O}_X}^{\mathbf{L}} \mathcal{F}^\bullet
$$
，其中 $\mathcal{G}^\bullet$ 与 $\mathcal{F}^\bullet$ 属于
$D(\mathcal{O}_X)$。这里使用《代数详论》第
\ref{more-algebra-section-sign-rules} 节给出的符号规则。因此写
$\mathcal{F}^\bullet \otimes_{\mathcal{O}_X}^{\mathbf{L}}
\mathcal{G}^\bullet$ 时，我们通常不区分使用哪一个变量来定义导出张量积。

\begin{definition}
\label{cohomology-definition-tor}
设 $(X, \mathcal{O}_X)$ 是环空间，而
$\mathcal{F}$、$\mathcal{G}$ 是 $\mathcal{O}_X$-模。
$\mathcal{F}$ 与 $\mathcal{G}$ 的 {\it Tor} 由公式
$$
\text{Tor}_p^{\mathcal{O}_X}(\mathcal{F}, \mathcal{G}) =
H^{-p}(\mathcal{F} \otimes_{\mathcal{O}_X}^\mathbf{L} \mathcal{G})
$$
定义，其中导出张量积如上定义。
\end{definition}

\noindent
此定义蕴含：对 $\mathcal{O}_X$-模的任意短正合列
$0 \to \mathcal{F}_1 \to \mathcal{F}_2 \to
\mathcal{F}_3 \to 0$，都有上同调长正合列
$$
\xymatrix{
\mathcal{F}_1 \otimes_{\mathcal{O}_X} \mathcal{G} \ar[r] &
\mathcal{F}_2 \otimes_{\mathcal{O}_X} \mathcal{G} \ar[r] &
\mathcal{F}_3 \otimes_{\mathcal{O}_X} \mathcal{G} \ar[r] & 0 \\
\text{Tor}_1^{\mathcal{O}_X}(\mathcal{F}_1, \mathcal{G}) \ar[r] &
\text{Tor}_1^{\mathcal{O}_X}(\mathcal{F}_2, \mathcal{G}) \ar[r] &
\text{Tor}_1^{\mathcal{O}_X}(\mathcal{F}_3, \mathcal{G}) \ar[ull]
}
$$
，其中 $\mathcal{G}$ 是任意 $\mathcal{O}_X$-模。称其为与该情形
相伴的 $\text{Tor}$ 长正合列。

\begin{lemma}
\label{cohomology-lemma-flat-tor-zero}
\begin{slogan}
Tor 衡量偏离平坦性的程度。
\end{slogan}
设 $(X, \mathcal{O}_X)$ 是环空间，而
$\mathcal{F}$ 是 $\mathcal{O}_X$-模。下列条件等价：
\begin{enumerate}
\item $\mathcal{F}$ 是平坦 $\mathcal{O}_X$-模；
\item $\text{Tor}_1^{\mathcal{O}_X}(\mathcal{F}, \mathcal{G}) = 0$
对每个 $\mathcal{O}_X$-模 $\mathcal{G}$ 都成立。
\end{enumerate}
\end{lemma}

\begin{proof}
若 $\mathcal{F}$ 平坦，则
$\mathcal{F} \otimes_{\mathcal{O}_X} -$ 是正合函子，且卫星函子消失。
反过来，假设（2）成立。若
$\mathcal{G} \to \mathcal{H}$ 是余核为 $\mathcal{Q}$ 的单射，
则 $\text{Tor}$ 长正合列表明映射
$\mathcal{F} \otimes_{\mathcal{O}_X} \mathcal{G} \to
\mathcal{F} \otimes_{\mathcal{O}_X} \mathcal{H}$
的核是
$\text{Tor}_1^{\mathcal{O}_X}(\mathcal{F}, \mathcal{Q})$
的商；后者依假设为零。因此 $\mathcal{F}$ 平坦。
\end{proof}

\begin{lemma}
\label{cohomology-lemma-factor-through-K-flat}
设 $(X, \mathcal{O}_X)$ 是环空间。设
$a : \mathcal{K}^\bullet \to \mathcal{L}^\bullet$ 是
$\mathcal{O}_X$-模复形的映射。若 $\mathcal{K}^\bullet$
是 K-平坦的，则存在复形 $\mathcal{N}^\bullet$ 和复形映射
$b : \mathcal{K}^\bullet \to \mathcal{N}^\bullet$、
$c : \mathcal{N}^\bullet \to \mathcal{L}^\bullet$，使得
\begin{enumerate}
\item $\mathcal{N}^\bullet$ 是 K-平坦的；
\item $c$ 是拟同构；
\item $a$ 与 $c \circ b$ 链同伦。
\end{enumerate}
若 $\mathcal{K}^\bullet$ 的各项平坦，则可选取
$\mathcal{N}^\bullet$、$b$、$c$，使得
$\mathcal{N}^\bullet$ 的各项也平坦。
\end{lemma}

\begin{proof}
我们将使用同伦范畴 $K(\textit{Mod}(\mathcal{O}_X))$ 为三角范畴这一事实，
见《导出范畴》命题
\ref{derived-proposition-homotopy-category-triangulated}。
选取特异三角
$\mathcal{K}^\bullet \to \mathcal{L}^\bullet \to
\mathcal{C}^\bullet \to \mathcal{K}^\bullet[1]$.
选取拟同构 $\mathcal{M}^\bullet \to \mathcal{C}^\bullet$，
其中 $\mathcal{M}^\bullet$ 为各项平坦的 K-平坦复形，见引理
\ref{cohomology-lemma-K-flat-resolution}。由三角范畴的公理，可把复合
$\mathcal{M}^\bullet \to \mathcal{C}^\bullet \to \mathcal{K}^\bullet[1]$
嵌入特异三角
$\mathcal{K}^\bullet \to \mathcal{N}^\bullet \to
\mathcal{M}^\bullet \to \mathcal{K}^\bullet[1]$.
由引理 \ref{cohomology-lemma-K-flat-two-out-of-three}，
$\mathcal{N}^\bullet$ 是 K-平坦的。再次使用三角范畴的公理，
可选取映射 $\mathcal{N}^\bullet \to \mathcal{L}^\bullet$，
把它嵌入下列特异三角的态射：
$$
\xymatrix{
\mathcal{K}^\bullet \ar[r] \ar[d] &
\mathcal{N}^\bullet \ar[r] \ar[d] &
\mathcal{M}^\bullet \ar[r] \ar[d] &
\mathcal{K}^\bullet[1] \ar[d] \\
\mathcal{K}^\bullet \ar[r] &
\mathcal{L}^\bullet \ar[r] &
\mathcal{C}^\bullet \ar[r] &
\mathcal{K}^\bullet[1]
}
$$
由于三个箭头中有两个是拟同构，由这些特异三角所伴随的上同调长正合列，
第三个箭头 $\mathcal{N}^\bullet \to \mathcal{L}^\bullet$ 也是拟同构
（也可考察此图在 $D(\mathcal{O}_X)$ 中的像，并使用《导出范畴》引理
\ref{derived-lemma-third-isomorphism-triangle}）。这就证明了
（1）、（2）、（3）。为证明最后一个断言，可以选取
$\mathcal{N}^\bullet$ 使得
$\mathcal{N}^n \cong \mathcal{M}^n \oplus \mathcal{K}^n$，见《导出范畴》引理
\ref{derived-lemma-improve-distinguished-triangle-homotopy}。
因此，若 $\mathcal{K}^\bullet$ 的各项平坦，就得到所需的平坦性。
\end{proof}











\section{导出拉回}
\label{cohomology-section-derived-pullback}

\noindent
设 $f : (X, \mathcal{O}_X) \to (Y, \mathcal{O}_Y)$
是环空间的态射。可以用 K-平坦消解定义导出拉回函子
$$
Lf^* : D(\mathcal{O}_Y) \to D(\mathcal{O}_X)
$$
。具体地，对任意 $\mathcal{O}_Y$-模复形
$\mathcal{G}^\bullet$，可选取 K-平坦消解
$\mathcal{K}^\bullet \to \mathcal{G}^\bullet$，并令
$Lf^*\mathcal{G}^\bullet = f^*\mathcal{K}^\bullet$。
利用引理 \ref{cohomology-lemma-pullback-K-flat}、
\ref{cohomology-lemma-K-flat-resolution} 和
\ref{cohomology-lemma-derived-tor-quasi-isomorphism-other-side} 可知此定义良好。
不过，为把所有细节完整处理，使用一些一般理论或许更为方便。

\begin{lemma}
\label{cohomology-lemma-derived-base-change}
上述构造与选择无关，并定义三角范畴的正合函子
$Lf^* : D(\mathcal{O}_Y) \to D(\mathcal{O}_X)$.
\end{lemma}

\begin{proof}
为此使用《导出范畴》第 \ref{derived-section-derived-functors} 节
发展的一般理论。令
$\mathcal{D} = K(\mathcal{O}_Y)$ 及
$\mathcal{D}' = D(\mathcal{O}_X)$。用
$F : \mathcal{D} \to \mathcal{D}'$ 表示由规则
$F(\mathcal{G}^\bullet) = f^*\mathcal{G}^\bullet$ 定义的
三角范畴正合函子。令 $S$ 为
$\mathcal{D} = K(\mathcal{O}_Y)$ 中所有拟同构所成的集合。
这给出《导出范畴》情形
\ref{derived-situation-derived-functor} 中的情形，因而可应用
《导出范畴》定义
\ref{derived-definition-right-derived-functor-defined}。
我们断言 $LF$ 处处有定义。取
$\mathcal{P} \subset \Ob(\mathcal{D})$ 为所有 $K$-平坦复形所成的族，
这由《导出范畴》引理
\ref{derived-lemma-find-existence-computes} 得到：
（1）由引理 \ref{cohomology-lemma-K-flat-resolution} 得到；为验证（2），
须证明对 $\mathcal{O}_Y$-模的 K-平坦复形之间的拟同构
$\mathcal{K}_1^\bullet  \to \mathcal{K}_2^\bullet$，映射
$f^*\mathcal{K}_1^\bullet  \to f^*\mathcal{K}_2^\bullet$
也是拟同构。把该映射写成
$$
f^{-1}\mathcal{K}_1^\bullet \otimes_{f^{-1}\mathcal{O}_Y} \mathcal{O}_X
\longrightarrow
f^{-1}\mathcal{K}_2^\bullet \otimes_{f^{-1}\mathcal{O}_Y} \mathcal{O}_X
$$
。函子 $f^{-1}$ 正合，故映射
$f^{-1}\mathcal{K}_1^\bullet \to
f^{-1}\mathcal{K}_2^\bullet$ 是拟同构。把引理
\ref{cohomology-lemma-pullback-K-flat} 应用于态射
$(X, f^{-1}\mathcal{O}_Y) \to (Y, \mathcal{O}_Y)$，可知复形
$f^{-1}\mathcal{K}_1^\bullet$ 与
$f^{-1}\mathcal{K}_2^\bullet$ 是
$f^{-1}\mathcal{O}_Y$-模的 K-平坦复形。因此引理
\ref{cohomology-lemma-derived-tor-quasi-isomorphism-other-side}
保证所显示的映射是拟同构。于是得到导出函子
$$
LF :
D(\mathcal{O}_Y) = S^{-1}\mathcal{D}
\longrightarrow
\mathcal{D}' = D(\mathcal{O}_X)
$$
；见《导出范畴》公式 (\ref{derived-equation-everywhere})。最后，
《导出范畴》引理 \ref{derived-lemma-find-existence-computes} 还保证
$LF(\mathcal{K}^\bullet) = F(\mathcal{K}^\bullet) = f^*\mathcal{K}^\bullet$
在 $\mathcal{K}^\bullet$ 为 K-平坦时成立；即
$Lf^* = LF$ 确实按上述方式计算。
\end{proof}

\begin{lemma}
\label{cohomology-lemma-derived-pullback-composition}
设 $f : X \to Y$ 与 $g : Y \to Z$ 是环空间的态射。
则作为函子 $D(\mathcal{O}_Z) \to D(\mathcal{O}_X)$，有
$Lf^* \circ Lg^* = L(g \circ f)^*$。
\end{lemma}

\begin{proof}
设 $E$ 是 $D(\mathcal{O}_Z)$ 的对象。按构造，为计算
$Lg^*E$，选取在 $Z$ 上表示 $E$ 的 K-平坦复形
$\mathcal{K}^\bullet$，并令
$Lg^*E = g^*\mathcal{K}^\bullet$。由引理
\ref{cohomology-lemma-pullback-K-flat}，$g^*\mathcal{K}^\bullet$ 在
$Y$ 上 K-平坦。于是 $Lf^*Lg^*E$ 由
$f^*g^*\mathcal{K}^\bullet = (g \circ f)^*\mathcal{K}^\bullet$
给出，而后者也表示 $L(g \circ f)^*E$。
\end{proof}

\begin{lemma}
\label{cohomology-lemma-pullback-tensor-product}
设 $f : (X, \mathcal{O}_X) \to (Y, \mathcal{O}_Y)$
是环空间的态射。存在典范的双函子性同构
$$
Lf^*(
\mathcal{F}^\bullet \otimes_{\mathcal{O}_Y}^{\mathbf{L}} \mathcal{G}^\bullet
) =
Lf^*\mathcal{F}^\bullet 
\otimes_{\mathcal{O}_X}^{\mathbf{L}}
Lf^*\mathcal{G}^\bullet 
$$
，其中
$\mathcal{F}^\bullet, \mathcal{G}^\bullet \in
\Ob(D(\mathcal{O}_Y))$。
\end{lemma}

\begin{proof}
可以假设 $\mathcal{F}^\bullet$ 与 $\mathcal{G}^\bullet$
都是 K-平坦复形。此时
$\mathcal{F}^\bullet \otimes_{\mathcal{O}_Y}^{\mathbf{L}} \mathcal{G}^\bullet$
就是双复形
$\mathcal{F}^\bullet \otimes_{\mathcal{O}_Y}
\mathcal{G}^\bullet$ 所伴随的总复形。由引理
\ref{cohomology-lemma-tensor-product-K-flat}，
$\text{Tot}(\mathcal{F}^\bullet \otimes_{\mathcal{O}_Y} \mathcal{G}^\bullet)$
也是 K-平坦的。因此引理中的同构来自同构
$$
\text{Tot}(f^*\mathcal{F}^\bullet \otimes_{\mathcal{O}_X}
f^*\mathcal{G}^\bullet)
\longrightarrow
f^*\text{Tot}(\mathcal{F}^\bullet \otimes_{\mathcal{O}_Y} \mathcal{G}^\bullet)
$$
，其各分量为《模层》引理
\ref{modules-lemma-tensor-product-pullback} 中的同构
$f^*\mathcal{F}^p \otimes_{\mathcal{O}_X} f^*\mathcal{G}^q \to
 f^*(\mathcal{F}^p \otimes_{\mathcal{O}_Y} \mathcal{G}^q)$。
\end{proof}

\begin{lemma}
\label{cohomology-lemma-variant-derived-pullback}
设 $f : (X, \mathcal{O}_X) \to (Y, \mathcal{O}_Y)$
是环空间的态射。存在典范的双函子性同构
$$
\mathcal{F}^\bullet
\otimes_{\mathcal{O}_X}^{\mathbf{L}}
Lf^*\mathcal{G}^\bullet
=
\mathcal{F}^\bullet 
\otimes_{f^{-1}\mathcal{O}_Y}^{\mathbf{L}}
f^{-1}\mathcal{G}^\bullet 
$$
，其中 $\mathcal{F}^\bullet$ 属于 $D(\mathcal{O}_X)$，
而 $\mathcal{G}^\bullet$ 属于 $D(\mathcal{O}_Y)$。
\end{lemma}

\begin{proof}
设 $\mathcal{F}$ 是 $\mathcal{O}_X$-模，而
$\mathcal{G}$ 是 $\mathcal{O}_Y$-模。则
$\mathcal{F} \otimes_{\mathcal{O}_X} f^*\mathcal{G} =
\mathcal{F} \otimes_{f^{-1}\mathcal{O}_Y} f^{-1}\mathcal{G}$
，这是因为
$f^*\mathcal{G} =
\mathcal{O}_X \otimes_{f^{-1}\mathcal{O}_Y} f^{-1}\mathcal{G}$.
引理由此及各定义得到。
\end{proof}

\begin{lemma}
\label{cohomology-lemma-tensor-pull-compatibility}
设 $f : (X, \mathcal{O}_X) \to (Y, \mathcal{O}_Y)$ 是环空间的态射。
设 $\mathcal{K}^\bullet$ 与 $\mathcal{M}^\bullet$ 是
$\mathcal{O}_Y$-模复形。下图
$$
\xymatrix{
Lf^*(\mathcal{K}^\bullet
\otimes_{\mathcal{O}_Y}^\mathbf{L}
\mathcal{M}^\bullet) \ar[r] \ar[d] &
Lf^*\text{Tot}(\mathcal{K}^\bullet
\otimes_{\mathcal{O}_Y}
\mathcal{M}^\bullet) \ar[d] \\
Lf^*\mathcal{K}^\bullet \otimes_{\mathcal{O}_X}^\mathbf{L}
Lf^*\mathcal{M}^\bullet \ar[d] &
f^*\text{Tot}(\mathcal{K}^\bullet
\otimes_{\mathcal{O}_Y}
\mathcal{M}^\bullet) \ar[d] \\
f^*\mathcal{K}^\bullet \otimes_{\mathcal{O}_X}^\mathbf{L}
f^*\mathcal{M}^\bullet \ar[r] &
\text{Tot}(f^*\mathcal{K}^\bullet \otimes_{\mathcal{O}_X}
f^*\mathcal{M}^\bullet)
}
$$
交换。
\end{lemma}

\begin{proof}
我们将使用引理 \ref{cohomology-lemma-pullback-K-flat} 中所述的 K-平坦消解的存在性。
若选取这样的消解
$\mathcal{P}^\bullet \to \mathcal{K}^\bullet$ 与
$\mathcal{Q}^\bullet \to \mathcal{M}^\bullet$，则可见
$$
\xymatrix{
Lf^*\text{Tot}(\mathcal{P}^\bullet
\otimes_{\mathcal{O}_Y}
\mathcal{Q}^\bullet) \ar[r] \ar[d] &
Lf^*\text{Tot}(\mathcal{K}^\bullet
\otimes_{\mathcal{O}_Y}
\mathcal{M}^\bullet) \ar[d] \\
f^*\text{Tot}(\mathcal{P}^\bullet
\otimes_{\mathcal{O}_Y}
\mathcal{Q}^\bullet) \ar[d] \ar[r] &
f^*\text{Tot}(\mathcal{K}^\bullet
\otimes_{\mathcal{O}_Y}
\mathcal{M}^\bullet) \ar[d] \\
\text{Tot}(f^*\mathcal{P}^\bullet \otimes_{\mathcal{O}_X}
f^*\mathcal{Q}^\bullet) \ar[r] &
\text{Tot}(f^*\mathcal{K}^\bullet \otimes_{\mathcal{O}_X}
f^*\mathcal{M}^\bullet)
}
$$
交换。不过，按我们对 $\mathcal{P}^\bullet$ 与
$\mathcal{Q}^\bullet$ 的选取以及引理
\ref{cohomology-lemma-tensor-product-K-flat}，此图的左边正是原图的左边。
\end{proof}









\section{无界复形的上同调}
\label{cohomology-section-unbounded}

\noindent
设 $(X, \mathcal{O}_X)$ 是环空间。范畴
$\textit{Mod}(\mathcal{O}_X)$ 是 Grothendieck 阿贝尔范畴：
它具有所有余极限，滤过余极限正合，并且具有生成元，即
$$
\bigoplus\nolimits_{U \subset X\text{开}} j_{U!}\mathcal{O}_U,
$$
；见《模层》第 \ref{modules-section-kernels} 节及引理
\ref{modules-lemma-j-shriek-flat}、
\ref{modules-lemma-module-quotient-flat}。由《内射对象》定理
\ref{injectives-theorem-K-injective-embedding-grothendieck}
，对 $\mathcal{O}_X$-模的每个复形 $\mathcal{F}^\bullet$，
都存在到 K-内射 $\mathcal{O}_X$-模复形的单射拟同构
$\mathcal{F}^\bullet \to \mathcal{I}^\bullet$，其中
$\mathcal{I}^\bullet$ 的各项均为内射 $\mathcal{O}_X$-模；
此外，可使此嵌入关于复形 $\mathcal{F}^\bullet$ 具有函子性。
由《导出范畴》引理
\ref{derived-lemma-enough-K-injectives-implies} 可得：
\begin{enumerate}
\item 任意到三角范畴 $\mathcal{D}$ 的正合函子
$F : K(\textit{Mod}(\mathcal{O}_X)) \to \mathcal{D}$
都具有右导出函子 $RF : D(\mathcal{O}_X) \to \mathcal{D}$；
\item 对任意到阿贝尔范畴 $\mathcal{A}$ 的加性函子
$F : \textit{Mod}(\mathcal{O}_X) \to \mathcal{A}$
，考虑由 $F$ 诱导的正合函子
$F : K(\textit{Mod}(\mathcal{O}_X)) \to K(\mathcal{A})$，
可得到右导出函子
$RF : D(\mathcal{O}_X) \to D(\mathcal{A})$。
\end{enumerate}
按构造有
$RF(\mathcal{F}^\bullet) = F(\mathcal{I}^\bullet)$，其中
$\mathcal{F}^\bullet \to \mathcal{I}^\bullet$ 如上。

\medskip\noindent
下面给出上述构造的一些例子：
\begin{enumerate}
\item 函子 $\Gamma(X, -) : \textit{Mod}(\mathcal{O}_X) \to
\text{Mod}_{\Gamma(X, \mathcal{O}_X)}$ 给出
$$
R\Gamma(X, -) : D(\mathcal{O}_X) \to D(\Gamma(X, \mathcal{O}_X))
$$
我们将用记号 $H^i(X, K) = H^i(R\Gamma(X, K))$ 表示上同调。
\item 对开集 $U \subset X$，考虑函子
$\Gamma(U, -) : \textit{Mod}(\mathcal{O}_X) \to
\text{Mod}_{\Gamma(U, \mathcal{O}_X)}$。它给出
$$
R\Gamma(U, -) : D(\mathcal{O}_X) \to D(\Gamma(U, \mathcal{O}_X))
$$
我们将用记号 $H^i(U, K) = H^i(R\Gamma(U, K))$ 表示上同调。
\item 对环空间态射
$f : (X, \mathcal{O}_X) \to (Y, \mathcal{O}_Y)$，考虑函子
$f_* : \textit{Mod}(\mathcal{O}_X) \to \textit{Mod}(\mathcal{O}_Y)$
，它在无界导出范畴上给出总直像
$$
Rf_* : D(\mathcal{O}_X) \longrightarrow D(\mathcal{O}_Y)
$$
\end{enumerate}

\begin{lemma}
\label{cohomology-lemma-adjoint}
设 $f : (X, \mathcal{O}_X) \to (Y, \mathcal{O}_Y)$ 是环空间的态射。
上面定义的函子 $Rf_*$ 与引理
\ref{cohomology-lemma-derived-base-change} 中定义的函子 $Lf^*$ 互为伴随：
$$
\Hom_{D(\mathcal{O}_X)}(Lf^*\mathcal{G}^\bullet, \mathcal{F}^\bullet)
=
\Hom_{D(\mathcal{O}_Y)}(\mathcal{G}^\bullet, Rf_*\mathcal{F}^\bullet)
$$
，它关于
$\mathcal{F}^\bullet \in \Ob(D(\mathcal{O}_X))$ 与
$\mathcal{G}^\bullet \in \Ob(D(\mathcal{O}_Y))$ 具有双函子性。
\end{lemma}

\begin{proof}
这形式地由 $Rf_*$ 与 $Lf^*$ 的存在性得到，见《导出范畴》引理
\ref{derived-lemma-derived-adjoint-functors}。
\end{proof}

\begin{lemma}
\label{cohomology-lemma-derived-pushforward-composition}
设 $f : X \to Y$ 与 $g : Y \to Z$ 是环空间的态射。
则作为函子 $D(\mathcal{O}_X) \to D(\mathcal{O}_Z)$，有
$Rg_* \circ Rf_* = R(g \circ f)_*$。
\end{lemma}

\begin{proof}
由引理 \ref{cohomology-lemma-adjoint}，$Rg_* \circ Rf_*$ 伴随于
$Lf^* \circ Lg^*$。由引理
\ref{cohomology-lemma-derived-pullback-composition} 有
$Lf^* \circ Lg^* = L(g \circ f)^*$，故由伴随函子的唯一性，
$Rg_* \circ Rf_* = R(g \circ f)_*$。
\end{proof}

\begin{remark}
\label{cohomology-remark-base-change}
无界导出函子 $Lf^*$ 与 $Rf_*$ 的构造使我们能在完全一般的情形下
构造基变换映射。具体地，假设
$$
\xymatrix{
X' \ar[r]_{g'} \ar[d]_{f'} &
X \ar[d]^f \\
S' \ar[r]^g &
S
}
$$
是环空间的交换图。设 $K$ 是 $D(\mathcal{O}_X)$ 的对象。
则在 $D(\mathcal{O}_{S'})$ 中存在典范基变换映射
$$
Lg^*Rf_*K \longrightarrow R(f')_*L(g')^*K
$$
。具体地，此映射伴随于映射
$L(f')^*Lg^*Rf_*K \to L(g')^*K$。
由于 $L(f')^*Lg^* = L(g')^*Lf^*$，这等同于映射
$L(g')^*Lf^*Rf_*K \to L(g')^*K$
，可取它为伴随映射 $Lf^*Rf_*K \to K$ 在 $L(g')^*$ 下的像。
\end{remark}

\begin{remark}
\label{cohomology-remark-compose-base-change}
考虑环空间的交换图
$$
\xymatrix{
X' \ar[r]_k \ar[d]_{f'} & X \ar[d]^f \\
Y' \ar[r]^l \ar[d]_{g'} & Y \ar[d]^g \\
Z' \ar[r]^m & Z
}
$$
。则注记 \ref{cohomology-remark-base-change} 中对应于两个方块的基变换映射
复合后给出对应于外部矩形的基变换映射。更确切地，复合
\begin{align*}
Lm^* \circ R(g \circ f)_*
& =
Lm^* \circ Rg_* \circ Rf_* \\
& \to Rg'_* \circ Ll^* \circ Rf_* \\
& \to Rg'_* \circ Rf'_* \circ Lk^* \\
& = R(g' \circ f')_* \circ Lk^*
\end{align*}
是该矩形的基变换映射。我们略去验证。
\end{remark}

\begin{remark}
\label{cohomology-remark-compose-base-change-horizontal}
考虑环空间的交换图
$$
\xymatrix{
X'' \ar[r]_{g'} \ar[d]_{f''} & X' \ar[r]_g \ar[d]_{f'} & X \ar[d]^f \\
Y'' \ar[r]^{h'} & Y' \ar[r]^h & Y
}
$$
。则注记 \ref{cohomology-remark-base-change} 中对应于两个方块的基变换映射
复合后给出对应于外部矩形的基变换映射。更确切地，复合
\begin{align*}
L(h \circ h')^* \circ Rf_*
& =
L(h')^* \circ Lh_* \circ Rf_* \\
& \to L(h')^* \circ Rf'_* \circ Lg^* \\
& \to Rf''_* \circ L(g')^* \circ Lg^* \\
& = Rf''_* \circ L(g \circ g')^*
\end{align*}
是该矩形的基变换映射。我们略去验证。
\end{remark}

\begin{lemma}
\label{cohomology-lemma-adjoints-push-pull-compatibility}
设 $f : (X, \mathcal{O}_X) \to (Y, \mathcal{O}_Y)$ 是环空间的态射。
设 $\mathcal{K}^\bullet$ 是 $\mathcal{O}_X$-模复形。图
$$
\xymatrix{
Lf^*f_*\mathcal{K}^\bullet \ar[r] \ar[d] &
f^*f_*\mathcal{K}^\bullet \ar[d] \\
Lf^*Rf_*\mathcal{K}^\bullet \ar[r] &
\mathcal{K}^\bullet
}
$$
来自复形上的 $Lf^* \to f^*$、$f_* \to Rf_*$ 以及伴随
$Lf^* \circ Rf_* \to \text{id}$，并在 $D(\mathcal{O}_X)$ 中交换。
\end{lemma}

\begin{proof}
我们将使用 K-平坦消解与 K-内射消解的存在性，见引理
\ref{cohomology-lemma-pullback-K-flat} 及上面的讨论。选取拟同构
$\mathcal{K}^\bullet \to \mathcal{I}^\bullet$，其中
$\mathcal{I}^\bullet$ 作为 $\mathcal{O}_X$-模复形是 K-内射的。
选取拟同构
$\mathcal{Q}^\bullet \to f_*\mathcal{I}^\bullet$，其中
$\mathcal{Q}^\bullet$ 作为 $\mathcal{O}_Y$-模复形是 K-平坦的。
可以选取 $\mathcal{O}_Y$-模的 K-平坦复形
$\mathcal{P}^\bullet$ 以及复形态射图
$$
\xymatrix{
\mathcal{P}^\bullet \ar[r] \ar[d] &
f_*\mathcal{K}^\bullet \ar[d] \\
\mathcal{Q}^\bullet \ar[r] & f_*\mathcal{I}^\bullet
}
$$
，它在同伦意义下交换，且上方水平箭头为拟同构。事实上，
由于拟同构在复形的同伦范畴中构成乘法系统，可以先对某个复形
$\mathcal{P}^\bullet$ 选取这样的图，然后用 K-平坦复形替换
$\mathcal{P}^\bullet$。取拉回后得到复形态射图
$$
\xymatrix{
f^*\mathcal{P}^\bullet \ar[r] \ar[d] &
f^*f_*\mathcal{K}^\bullet \ar[d] \ar[r] &
\mathcal{K}^\bullet \ar[d] \\
f^*\mathcal{Q}^\bullet \ar[r] &
f^*f_*\mathcal{I}^\bullet \ar[r] &
\mathcal{I}^\bullet
}
$$
，它在同伦意义下交换。外部矩形证明了引理中的断言。
\end{proof}

\begin{remark}
\label{cohomology-remark-cup-product}
设 $f : (X, \mathcal{O}_X) \to (Y, \mathcal{O}_Y)$ 是环空间的态射。
$Lf^*$ 与 $Rf_*$ 的伴随性使我们能构造相对杯积
$$
Rf_*K \otimes_{\mathcal{O}_Y}^\mathbf{L} Rf_*L
\longrightarrow
Rf_*(K \otimes_{\mathcal{O}_X}^\mathbf{L} L)
$$
；它对所有 $D(\mathcal{O}_X)$ 中的 $K, L$ 都是
$D(\mathcal{O}_Y)$ 中的映射。具体地，此映射伴随于映射
$Lf^*(Rf_*K \otimes_{\mathcal{O}_Y}^\mathbf{L} Rf_*L) \to
 K \otimes_{\mathcal{O}_X}^\mathbf{L} L$；可以把它取为同构
$Lf^*(Rf_*K \otimes_{\mathcal{O}_Y}^\mathbf{L} Rf_*L) =
Lf^*Rf_*K \otimes_{\mathcal{O}_X}^\mathbf{L} Lf^*Rf_*L$
（引理 \ref{cohomology-lemma-pullback-tensor-product}）与映射
$Lf^*Rf_*K \otimes_{\mathcal{O}_X}^\mathbf{L} Lf^*Rf_*L
\to K \otimes_{\mathcal{O}_X}^\mathbf{L} L$
的复合，其中后一个映射来自余单位
$Lf^* \circ Rf_* \to \text{id}$。
\end{remark}






\section{滤过复形的上同调}
\label{cohomology-section-cohomology-filtered-object}

\noindent
研究代数簇的上同调时，层的滤过复形经常自然出现；例如 de Rham
复形带有 Hodge 滤过。本节使用非常一般的《内射对象》引理
\ref{injectives-lemma-K-injective-embedding-filtration}
构造上同调上的谱序列，并把它们与先前构造的谱序列联系起来。

\begin{lemma}
\label{cohomology-lemma-spectral-sequence-filtered-object}
设 $(X, \mathcal{O}_X)$ 是环空间，而
$\mathcal{F}^\bullet$ 是 $\mathcal{O}_X$-模的滤过复形。
存在双分次 $\Gamma(X, \mathcal{O}_X)$-模的典范谱序列
$(E_r, \text{d}_r)_{r \geq 1}$，其中 $d_r$ 的双次数为
$(r, -r + 1)$，且
$$
E_1^{p, q} = H^{p + q}(X, \text{gr}^p\mathcal{F}^\bullet)
$$
若对每个 $n$ 都有
$$
H^n(X, F^p\mathcal{F}^\bullet) = 0\text{，当 }p \gg 0
\quad\text{且}\quad
H^n(X, F^p\mathcal{F}^\bullet) = H^n(X, \mathcal{F}^\bullet)\text{，当 }p \ll 0
$$
，则该谱序列有界并收敛到
$H^*(X, \mathcal{F}^\bullet)$。
\end{lemma}

\begin{proof}
（当该复形为具有有限滤过的模的有下界复形时，证明见下面的注记。）
选取《内射对象》引理
\ref{injectives-lemma-K-injective-embedding-filtration}
中的滤过复形映射
$j : \mathcal{F}^\bullet \to \mathcal{J}^\bullet$。
该谱序列是《同调代数》第
\ref{homology-section-filtered-complex} 节中与滤过复形
$$
\Gamma(X, \mathcal{J}^\bullet)
\quad\text{其中}\quad
F^p\Gamma(X, \mathcal{J}^\bullet) = \Gamma(X, F^p\mathcal{J}^\bullet)
$$
相伴的谱序列。由于上同调通过在 K-内射代表元上取值来计算，
可知 $E_1$ 页如引理所述。在所给条件下，收敛性与有界性由
《同调代数》引理 \ref{homology-lemma-ss-converges-trivial} 得到。
\end{proof}

\begin{remark}
\label{cohomology-remark-spectral-sequence-filtered-object}
设 $(X, \mathcal{O}_X)$ 是环空间，而
$\mathcal{F}^\bullet$ 是 $\mathcal{O}_X$-模的滤过复形。
若 $\mathcal{F}^\bullet$ 有下界，且对每个 $n$，
$\mathcal{F}^n$ 上的滤过有限，则可构造引理
\ref{cohomology-lemma-spectral-sequence-filtered-object} 中的谱序列，而无需使用《内射对象》引理
\ref{injectives-lemma-K-injective-embedding-filtration}。
事实上，由《导出范畴》引理
\ref{derived-lemma-right-resolution-by-filtered-injectives}，
存在滤过复形的滤过拟同构
$i : \mathcal{F}^\bullet \to \mathcal{I}^\bullet$，
其中 $\mathcal{I}^\bullet$ 有下界，对所有 $n$，
$\mathcal{I}^n$ 上的滤过有限，并且每个
$\text{gr}^p\mathcal{I}^n$ 都是内射 $\mathcal{O}_X$-模。
然后取与
$$
\Gamma(X, \mathcal{I}^\bullet)
\quad\text{其中}\quad
F^p\Gamma(X, \mathcal{I}^\bullet) = \Gamma(X, F^p\mathcal{I}^\bullet)
$$
相伴的谱序列。由于上同调可通过在内射对象的有下界复形上取值来计算，
可知 $E_1$ 页如引理所述。在所给条件下，收敛性与有界性由
《同调代数》引理 \ref{homology-lemma-biregular-ss-converges} 得到。
事实上，这是《导出范畴》引理
\ref{derived-lemma-ss-filtered-derived} 中谱序列的特例。
\end{remark}

\begin{example}
\label{cohomology-example-spectral-sequence}
设 $(X, \mathcal{O}_X)$ 是环空间，而
$\mathcal{F}^\bullet$ 是 $\mathcal{O}_X$-模复形。令
$F^p\mathcal{F}^\bullet =
\tau_{\leq -p}\mathcal{F}^\bullet$，即可应用引理
\ref{cohomology-lemma-spectral-sequence-filtered-object}。
（若 $\mathcal{F}^\bullet$ 有下界，可使用注记
\ref{cohomology-remark-spectral-sequence-filtered-object}。）于是得到谱序列
$$
E_1^{p, q} = H^{p + q}(X, H^{-p}(\mathcal{F}^\bullet)[p]) =
H^{2p + q}(X, H^{-p}(\mathcal{F}^\bullet))
$$
作重新编号 $p = -j$、$q = i + 2j$ 后可知，对任意
$K \in D(\mathcal{O}_X)$，存在双分次模的谱序列
$(E'_r, d'_r)_{r \geq 2}$，其中 $d'_r$ 的双次数为
$(r, -r + 1)$，并且
$$
(E'_2)^{i, j} = H^i(X, H^j(K))
$$
例如，若 $K$ 有下界，则此谱序列有界，并收敛到
$H^{i + j}(X, K)$。在有下界的情形下，此谱序列是《导出范畴》引理
\ref{derived-lemma-two-ss-complex-functor} 的第二谱序列
（用 Cartan--Eilenberg 消解构造）的一个例子。
\end{example}

\begin{example}
\label{cohomology-example-spectral-sequence-bis}
设 $(X, \mathcal{O}_X)$ 是环空间，而
$\mathcal{F}^\bullet$ 是 $\mathcal{O}_X$-模复形。令
$F^p\mathcal{F}^\bullet =
\sigma_{\geq p}\mathcal{F}^\bullet$，即可应用引理
\ref{cohomology-lemma-spectral-sequence-filtered-object}。于是得到谱序列
$$
E_1^{p, q} = H^{p + q}(X, \mathcal{F}^p[-p]) = H^q(X, \mathcal{F}^p)
$$
若 $\mathcal{F}^\bullet$ 有下界，则
\begin{enumerate}
\item 可以用注记 \ref{cohomology-remark-spectral-sequence-filtered-object}
构造此谱序列；
\item 此谱序列有界，并收敛到
$H^{i + j}(X, \mathcal{F}^\bullet)$；
\item 此谱序列等于《导出范畴》引理
\ref{derived-lemma-two-ss-complex-functor} 的第一谱序列
（用 Cartan--Eilenberg 消解构造）。
\end{enumerate}
\end{example}

\begin{lemma}
\label{cohomology-lemma-relative-spectral-sequence-filtered-object}
设 $f : (X, \mathcal{O}_X) \to (Y, \mathcal{O}_Y)$ 是环空间的态射。
设 $\mathcal{F}^\bullet$ 是 $\mathcal{O}_X$-模的滤过复形。
存在双分次 $\mathcal{O}_Y$-模的典范谱序列
$(E_r, \text{d}_r)_{r \geq 1}$，其中 $d_r$ 的双次数为
$(r, -r + 1)$，且
$$
E_1^{p, q} = R^{p + q}f_*\text{gr}^p\mathcal{F}^\bullet
$$
若对每个 $n$ 都有
$$
R^nf_*F^p\mathcal{F}^\bullet = 0 \text{，当 }p \gg 0
\quad\text{且}\quad
R^nf_*F^p\mathcal{F}^\bullet = R^nf_*\mathcal{F}^\bullet \text{，当 }p \ll 0
$$
，则该谱序列有界并收敛到
$Rf_*\mathcal{F}^\bullet$。
\end{lemma}

\begin{proof}
证明与引理 \ref{cohomology-lemma-spectral-sequence-filtered-object} 的证明完全相同。
\end{proof}





\section{Godement 消解}
\label{cohomology-section-godement}

\noindent
参考文献为 \cite{Godement}。

\medskip\noindent
设 $(X, \mathcal{O}_X)$ 是环空间。以 $X_{disc}$ 表示与 $X$
具有相同点集的离散拓扑空间，以 $f : X_{disc} \to X$
表示显然的连续映射。令
$\mathcal{O}_{X_{disc}} = f^{-1}\mathcal{O}_X$。
则 $f : (X_{disc}, \mathcal{O}_{X_{disc}}) \to (X, \mathcal{O}_X)$
是环空间的平坦态射。可将《单纯方法》第
\ref{simplicial-section-standard} 节中材料的
{\it 对偶}应用于模层上的伴随函子对 $f^*, f_*$。由此得到增广余单纯对象
$$
\xymatrix{
\text{id} \ar[r] &
f_*f^* \ar@<1ex>[r] \ar@<-1ex>[r] &
f_*f^*f_*f^* \ar@<0ex>[l] \ar@<-2ex>[r] \ar@<0ex>[r] \ar@<2ex>[r] &
f_*f^*f_*f^*f_*f^* \ar@<1ex>[l] \ar@<-1ex>[l]
}
$$
，它位于从 $\textit{Mod}(\mathcal{O}_X)$ 到自身的函子范畴中；
参见《单纯方法》引理 \ref{simplicial-lemma-standard-simplicial}。
此外，增广
$$
\xymatrix{
f^* \ar[r] &
f^*f_*f^* \ar@<1ex>[r] \ar@<-1ex>[r] &
f^*f_*f^*f_*f^* \ar@<0ex>[l] \ar@<-2ex>[r] \ar@<0ex>[r] \ar@<2ex>[r] &
f^*f_*f^*f_*f^*f_*f^* \ar@<1ex>[l] \ar@<-1ex>[l]
}
$$
是同伦等价；参见《单纯方法》引理
\ref{simplicial-lemma-standard-simplicial-homotopy}。

\begin{lemma}
\label{cohomology-lemma-godement-resolution}
设 $(X, \mathcal{O}_X)$ 是环空间。对每个
$\mathcal{O}_X$-模层 $\mathcal{F}$，存在消解
$$
0 \to
\mathcal{F} \to
f_*f^*\mathcal{F} \to
f_*f^*f_*f^*\mathcal{F} \to
f_*f^*f_*f^*f_*f^*\mathcal{F} \to \ldots
$$
，它关于 $\mathcal{F}$ 具有函子性，并且每一项
$f_*f^* \ldots f_*f^*\mathcal{F}$ 都是松弛
$\mathcal{O}_X$-模；此外，对所有 $x \in X$，映射
$$
\mathcal{F}_x[0] \to \Big(
(f_*f^*\mathcal{F})_x \to
(f_*f^*f_*f^*\mathcal{F})_x \to
(f_*f^*f_*f^*f_*f^*\mathcal{F})_x \to \ldots
\Big)
$$
是 $\mathcal{O}_{X, x}$-模复形范畴中的链同伦等价。
\end{lemma}

\begin{proof}
复形 $f_*f^*\mathcal{F} \to  f_*f^*f_*f^*\mathcal{F} \to
f_*f^*f_*f^*f_*f^*\mathcal{F} \to \ldots$ 是与上述余单纯对象相伴的复形，
其各项为
$f_*f^*\mathcal{F}, f_*f^*f_*f^*\mathcal{F},
f_*f^*f_*f^*f_*f^*\mathcal{F}, \ldots$；参见《单纯方法》第
\ref{simplicial-section-dold-kan-cosimplicial} 节。
增广给出所示的映射 $\mathcal{F} \to f_*f^*\mathcal{F}$。
对 $X_{disc}$ 上任意阿贝尔层 $\mathcal{H}$，推前
$f_*\mathcal{H}$ 都是松弛的，因为 $X_{disc}$ 是离散空间，
而松弛层的推前仍是松弛层。因此该复形的各项都是松弛
$\mathcal{O}_X$-模。

\medskip\noindent
若 $x \in X_{disc} = X$ 是一点，则对任意
$\mathcal{O}_X$-模 $\mathcal{G}$ 都有
$(f^*\mathcal{G})_x = \mathcal{G}_x$。因此 $f^*$ 是正合函子，并且
$\mathcal{O}_X$-模复形
$\mathcal{G}_1 \to \mathcal{G}_2 \to \mathcal{G}_3$
正合，当且仅当
$f^*\mathcal{G}_1 \to f^*\mathcal{G}_2 \to f^*\mathcal{G}_3$
正合（参见《模》引理 \ref{modules-lemma-abelian}）。
本节引言中提到的结果表明，应用 $f^*$ 后，常值余单纯对象
$f^*\mathcal{F}$ 到各项为
$f_*f^*\mathcal{F}, f_*f^*f_*f^*\mathcal{F},
f_*f^*f_*f^*f_*f^*\mathcal{F}, \ldots$
的余单纯对象之间得到一个同伦等价。由《单纯方法》引理
\ref{simplicial-lemma-homotopy-equivalence-s-Q} 可知
$$
f^*\mathcal{F}[0] \to \Big(
f^*f_*f^*\mathcal{F} \to
f^*f_*f^*f_*f^*\mathcal{F} \to
f^*f_*f^*f_*f^*f_*f^*\mathcal{F} \to \ldots
\Big)
$$
是链同伦等价。这立即推出引理其余两个断言。
\end{proof}

\begin{lemma}
\label{cohomology-lemma-godement-resolution-bounded-below}
设 $(X, \mathcal{O}_X)$ 是环空间，而
$\mathcal{F}^\bullet$ 是 $\mathcal{O}_X$-模的有下界复形。
存在拟同构
$\mathcal{F}^\bullet \to \mathcal{G}^\bullet$
，其中 $\mathcal{G}^\bullet$ 是松弛 $\mathcal{O}_X$-模的有下界复形；
并且对所有 $x \in X$，映射
$\mathcal{F}^\bullet_x \to \mathcal{G}^\bullet_x$
是 $\mathcal{O}_{X, x}$-模复形范畴中的链同伦等价。
\end{lemma}

\begin{proof}
令 $\mathcal{A}$ 为 $\mathcal{O}_X$-模复形的范畴，
并令 $\mathcal{B}$ 为 $\mathcal{O}_X$-模复形的范畴。
于是可将上述讨论应用于 $\mathcal{A}$ 与 $\mathcal{B}$
之间的伴随函子 $f^*$ 和 $f_*$。完全仿照引理
\ref{cohomology-lemma-godement-resolution} 的证明，可得消解
$$
0 \to
\mathcal{F}^\bullet \to
f_*f^*\mathcal{F}^\bullet \to
f_*f^*f_*f^*\mathcal{F}^\bullet \to
f_*f^*f_*f^*f_*f^*\mathcal{F}^\bullet \to \ldots
$$
，它位于阿贝尔范畴 $\mathcal{A}$ 中，并且每个
$f_*f^*\ldots f_*f^*\mathcal{F}^\bullet$ 的每一项都是松弛
$\mathcal{O}_X$-模；此外，对所有 $x \in X$，映射
$$
\mathcal{F}^\bullet_x[0] \to \Big(
(f_*f^*\mathcal{F}^\bullet)_x \to
(f_*f^*f_*f^*\mathcal{F}^\bullet)_x \to
(f_*f^*f_*f^*f_*f^*\mathcal{F}^\bullet)_x \to \ldots
\Big)
$$
是 $\mathcal{O}_{X, x}$-模的复形之复形范畴中的链同伦等价。
由于复形的复形就是双复形，可以考虑所诱导的映射
$$
\mathcal{F}^\bullet \to
\mathcal{G}^\bullet =
\text{Tot}(
f_*f^*\mathcal{F}^\bullet \to
f_*f^*f_*f^*\mathcal{F}^\bullet \to
f_*f^*f_*f^*f_*f^*\mathcal{F}^\bullet \to \ldots
)
$$
。由于复形 $\mathcal{F}^\bullet$ 有下界，
$\mathcal{G}^\bullet$ 也有下界。事实上，
$\mathcal{G}^\bullet$ 的每一项都是复形
$f_*f^*\ldots f_*f^*\mathcal{F}^\bullet$ 各项的有限直和，
因而是松弛的。引理的最后一个断言现由《同调代数》引理
\ref{homology-lemma-homotopy-complex-complexes} 推出。
这尤其表明
$\mathcal{F}^\bullet \to \mathcal{G}^\bullet$
是拟同构，故证明完成。
\end{proof}














\section{杯积}
\label{cohomology-section-cup-product}

\noindent
设 $(X, \mathcal{O}_X)$ 是环空间，而 $K, M$ 是
$D(\mathcal{O}_X)$ 的对象。令 $A = \Gamma(X, \mathcal{O}_X)$。
在此情形中，（整体）杯积是 $D(A)$ 中的映射
$$
\mu :
R\Gamma(X, K) \otimes_A^\mathbf{L} R\Gamma(X, M)
\longrightarrow
R\Gamma(X, K \otimes_{\mathcal{O}_X}^\mathbf{L} M)
$$
。通过 $D(pt, A) = D(A)$，按照注记
\ref{cohomology-remark-cup-product}，把它定义为环空间态射
$f : (X, \mathcal{O}_X) \to (pt, A)$ 的相对杯积。
此映射特别给出配对
$$
\cup :
H^i(X, K) \times H^j(X, M)
\longrightarrow
H^{i + j}(X, K \otimes_{\mathcal{O}_X}^\mathbf{L} M)
$$
具体地，给定 $\xi \in H^i(X, K) = H^i(R\Gamma(X, K))$ 与
$\eta \in H^j(X, M) = H^j(R\Gamma(X, M))$，可以先将二者“张量”，
得到
$H^{i + j}(R\Gamma(X, K) \otimes_A^\mathbf{L} R\Gamma(X, M))$
中的元素 $\xi \otimes \eta$；参见《更多代数》第
\ref{more-algebra-section-products-tor} 节。然后应用 $\mu$，
得到 $H^{i + j}(X, K \otimes_{\mathcal{O}_X}^\mathbf{L} M)$
中的所需元素
$\xi \cup \eta = \mu(\xi \otimes \eta)$。

\medskip\noindent
还可以用另一种方式理解 $\xi$ 与 $\eta$ 的杯积。具体地，可以写成
$$
H^i(R\Gamma(X, K)) = \Hom_{D(\mathcal{O}_X)}(\mathcal{O}_X[-i], K)
\quad\text{且}\quad
H^j(R\Gamma(X, M)) = \Hom_{D(\mathcal{O}_X)}(\mathcal{O}_X[-j], M)
$$
，因为 $\Hom(\mathcal{O}_X, -) = \Gamma(X, -)$。
因此，$\xi$ 与 $\eta$ 分别与映射
$$
\tilde \xi : \mathcal{O}_X[-i] \to K
\quad\text{及}\quad
\tilde \eta : \mathcal{O}_X[-j] \to M
$$
是“同一”回事。结合导出张量积的函子性，可得
$$
\mathcal{O}_X[-i - j] =
\mathcal{O}_X[-i] \otimes_{\mathcal{O}_X}^\mathbf{L} \mathcal{O}_X[-j]
\xrightarrow{\tilde \xi \otimes \tilde \eta}
K \otimes_{\mathcal{O}_X}^\mathbf{L} M
$$
；由与上面相同的理由，这给出
$H^{i + j}(X, K \otimes_{\mathcal{O}_X}^\mathbf{L} M)$ 的一个元素。

\begin{lemma}
\label{cohomology-lemma-second-cup-equals-first}
此构造给出杯积。
\end{lemma}

\begin{proof}
取上述 $f : (X, \mathcal{O}_X) \to (pt, A)$，则
$Rf_*(-) = R\Gamma(X, -)$，且映射 $\mu$ 与下列映射伴随：
$$
Lf^*(Rf_*K \otimes_A^\mathbf{L} Rf_*M) =
Lf^*Rf_*K \otimes_{\mathcal{O}_X}^\mathbf{L} Lf^*Rf_*M
\xrightarrow{\epsilon_K \otimes \epsilon_M}
K \otimes_{\mathcal{O}_X}^\mathbf{L} M
$$
，其中 $\epsilon$ 是 $Lf^*$ 与 $Rf_*$ 之间伴随的余单位。
若把 $\xi$ 与 $\eta$ 看作映射
$\xi : A[-i] \to R\Gamma(X, K)$ 和
$\eta : A[-j] \to R\Gamma(X, M)$，则张量
$\xi \otimes \eta$ 对应于映射\footnote{这里隐藏着一个符号。
具体地，该等式由复合
$$
A[-i - j] \to (A \otimes_A^\mathbf{L} A)[-i - j] \to
A[-i] \otimes_A^\mathbf{L} A[-j]
$$
定义；第二步使用《更多代数》第
(\ref{more-algebra-item-shift-tensor}) 项的等同，该等同原则上带有符号。
不过依照我们的约定，此处符号为 $+1$；即使它不是 $+1$ 也无妨，
因为在等同
$\mathcal{O}_X[-i - j] =
\mathcal{O}_X[-i] \otimes_{\mathcal{O}_X}^\mathbf{L} \mathcal{O}_X[-j]$
中使用了相同的符号。}
$$
A[-i - j] = A[-i] \otimes_A^\mathbf{L} A[-j]
\xrightarrow{\xi \otimes \eta}
R\Gamma(X, K) \otimes_A^\mathbf{L} R\Gamma(X, M)
$$
按照定义，杯积 $\xi \cup \eta$ 是映射
$A[-i - j] \to R\Gamma(X, K \otimes_{\mathcal{O}_X}^\mathbf{L} M)$
，它与下列映射伴随：
$$
(\epsilon_K \otimes \epsilon_M) \circ Lf^*(\xi \otimes \eta) =
(\epsilon_K \circ Lf^*\xi) \otimes (\epsilon_M \circ Lf^*\eta)
$$
然而，容易看出
$\epsilon_K \circ Lf^*\xi = \tilde \xi$，且
$\epsilon_M \circ Lf^*\eta = \tilde \eta$.
由此得到
$\widetilde{\xi \cup \eta} = \tilde \xi \otimes \tilde \eta$，
这正是所需的一致性。
\end{proof}

\begin{remark}
\label{cohomology-remark-cup-with-element-map-total-cohomology}
设 $(X, \mathcal{O}_X)$ 是环空间，而 $K, M$ 是
$D(\mathcal{O}_X)$ 的对象。令 $A = \Gamma(X, \mathcal{O}_X)$。
给定 $\xi \in H^i(X, K)$，可得相伴映射
$$
\xi = \text{“}\xi \cup -\text{”} :
R\Gamma(X, M)[-i]
\to
R\Gamma(X, K \otimes_{\mathcal{O}_X}^\mathbf{L} M)
$$
：先如引理 \ref{cohomology-lemma-second-cup-equals-first} 的证明中那样，
把 $\xi$ 表示为映射 $\xi : A[-i] \to R\Gamma(X, K)$，
再使用复合
$$
R\Gamma(X, M)[-i] = A[-i] \otimes_A^\mathbf{L} R\Gamma(X, M)
\xrightarrow{\xi \otimes 1}
R\Gamma(X, K) \otimes_A^\mathbf{L} R\Gamma(X, M)
\to
R\Gamma(X, K \otimes_{\mathcal{O}_X}^\mathbf{L} M)
$$
，其中第二个箭头是上述整体杯积 $\mu$。
由引理 \ref{cohomology-lemma-second-cup-equals-first} 及其证明可知，
在上同调上，这恢复了与 $\xi$ 作杯积。
\end{remark}

\noindent
下面陈述并证明相对杯积的一种自然相容性。设有环空间态射
$f : (X, \mathcal{O}_X) \to (Y, \mathcal{O}_Y)$，
并设 $\mathcal{K}^\bullet$ 与 $\mathcal{M}^\bullet$
为 $\mathcal{O}_X$-模复形。存在朴素杯积
$$
\text{Tot}(
f_*\mathcal{K}^\bullet
\otimes_{\mathcal{O}_Y}
f_*\mathcal{M}^\bullet)
\longrightarrow
f_*\text{Tot}(\mathcal{K}^\bullet
\otimes_{\mathcal{O}_X}
\mathcal{M}^\bullet)
$$
我们断言它与相对杯积相关。

\begin{lemma}
\label{cohomology-lemma-cup-compatible-with-naive}
在上述情形中，下图交换：
$$
\xymatrix{
f_*\mathcal{K}^\bullet
\otimes_{\mathcal{O}_Y}^\mathbf{L}
f_*\mathcal{M}^\bullet \ar[r] \ar[d]
&
Rf_*\mathcal{K}^\bullet
\otimes_{\mathcal{O}_Y}^\mathbf{L}
Rf_*\mathcal{M}^\bullet \ar[d]^{\text{注记 \ref{cohomology-remark-cup-product}}} \\
\text{Tot}(
f_*\mathcal{K}^\bullet
\otimes_{\mathcal{O}_Y}
f_*\mathcal{M}^\bullet) \ar[d]_{\text{朴素杯积}} &
Rf_*(\mathcal{K}^\bullet
\otimes_{\mathcal{O}_X}^\mathbf{L}
\mathcal{M}^\bullet) \ar[d] \\
f_*\text{Tot}(\mathcal{K}^\bullet
\otimes_{\mathcal{O}_X}
\mathcal{M}^\bullet) \ar[r] &
Rf_*\text{Tot}(\mathcal{K}^\bullet
\otimes_{\mathcal{O}_X}
\mathcal{M}^\bullet)
}
$$
\end{lemma}

\begin{proof}
由注记 \ref{cohomology-remark-cup-product} 中的构造可知，沿图顺时针绕行所得映射
$$
f_*\mathcal{K}^\bullet
\otimes_{\mathcal{O}_Y}^\mathbf{L}
f_*\mathcal{M}^\bullet 
\longrightarrow
Rf_*\text{Tot}(\mathcal{K}^\bullet
\otimes_{\mathcal{O}_X}
\mathcal{M}^\bullet)
$$
与下列映射伴随：
\begin{align*}
Lf^*(f_*\mathcal{K}^\bullet
\otimes_{\mathcal{O}_Y}^\mathbf{L}
f_*\mathcal{M}^\bullet)
& =
Lf^*f_*\mathcal{K}^\bullet
\otimes_{\mathcal{O}_Y}^\mathbf{L}
Lf^*f_*\mathcal{M}^\bullet \\
& \to
Lf^*Rf_*\mathcal{K}^\bullet
\otimes_{\mathcal{O}_Y}^\mathbf{L}
Lf^*Rf_*\mathcal{M}^\bullet \\
& \to
\mathcal{K}^\bullet
\otimes_{\mathcal{O}_Y}^\mathbf{L}
\mathcal{M}^\bullet \\
& \to
\text{Tot}(\mathcal{K}^\bullet
\otimes_{\mathcal{O}_X}
\mathcal{M}^\bullet)
\end{align*}
由引理 \ref{cohomology-lemma-adjoints-push-pull-compatibility}，这也等于
\begin{align*}
Lf^*(f_*\mathcal{K}^\bullet
\otimes_{\mathcal{O}_Y}^\mathbf{L}
f_*\mathcal{M}^\bullet)
& =
Lf^*f_*\mathcal{K}^\bullet
\otimes_{\mathcal{O}_Y}^\mathbf{L}
Lf^*f_*\mathcal{M}^\bullet \\
& \to
f^*f_*\mathcal{K}^\bullet
\otimes_{\mathcal{O}_Y}^\mathbf{L}
f^*f_*\mathcal{M}^\bullet \\
& \to
\mathcal{K}^\bullet
\otimes_{\mathcal{O}_Y}^\mathbf{L}
\mathcal{M}^\bullet \\
& \to
\text{Tot}(\mathcal{K}^\bullet
\otimes_{\mathcal{O}_X}
\mathcal{M}^\bullet)
\end{align*}
沿图逆时针绕行，得到与下列映射伴随的映射：
\begin{align*}
Lf^*(f_*\mathcal{K}^\bullet
\otimes_{\mathcal{O}_Y}^\mathbf{L}
f_*\mathcal{M}^\bullet)
& \to
Lf^*\text{Tot}(
f_*\mathcal{K}^\bullet
\otimes_{\mathcal{O}_Y}
f_*\mathcal{M}^\bullet) \\
& \to
Lf^*f_*\text{Tot}(\mathcal{K}^\bullet
\otimes_{\mathcal{O}_X}
\mathcal{M}^\bullet) \\
& \to
Lf^*Rf_*\text{Tot}(\mathcal{K}^\bullet
\otimes_{\mathcal{O}_X}
\mathcal{M}^\bullet) \\
& \to
\text{Tot}(\mathcal{K}^\bullet
\otimes_{\mathcal{O}_X}
\mathcal{M}^\bullet)
\end{align*}
由引理 \ref{cohomology-lemma-adjoints-push-pull-compatibility}，这也等于
\begin{align*}
Lf^*(f_*\mathcal{K}^\bullet
\otimes_{\mathcal{O}_Y}^\mathbf{L}
f_*\mathcal{M}^\bullet)
& \to
Lf^*\text{Tot}(
f_*\mathcal{K}^\bullet
\otimes_{\mathcal{O}_Y}
f_*\mathcal{M}^\bullet) \\
& \to
Lf^*f_*\text{Tot}(\mathcal{K}^\bullet
\otimes_{\mathcal{O}_X}
\mathcal{M}^\bullet) \\
& \to
f^*f_*\text{Tot}(\mathcal{K}^\bullet
\otimes_{\mathcal{O}_X}
\mathcal{M}^\bullet) \\
& \to
\text{Tot}(\mathcal{K}^\bullet
\otimes_{\mathcal{O}_X}
\mathcal{M}^\bullet)
\end{align*}
考察下图即可完成证明：
$$
\xymatrix{
Lf^*(f_*\mathcal{K}^\bullet
\otimes_{\mathcal{O}_Y}^\mathbf{L}
f_*\mathcal{M}^\bullet) \ar[d] \ar[rr] & &
Lf^*f_*\mathcal{K}^\bullet \otimes_{\mathcal{O}_X}^\mathbf{L}
Lf^*f_*\mathcal{M}^\bullet \ar[d] \\
Lf^*\text{Tot}(
f_*\mathcal{K}^\bullet
\otimes_{\mathcal{O}_Y}
f_*\mathcal{M}^\bullet) \ar[d]_{\text{朴素}} \ar[r] &
f^*\text{Tot}(
f_*\mathcal{K}^\bullet
\otimes_{\mathcal{O}_Y}
f_*\mathcal{M}^\bullet) \ar[ldd]^{\text{朴素}} \ar[dd] &
f^*f_*\mathcal{K}^\bullet \otimes_{\mathcal{O}_X}^\mathbf{L}
f^*f_*\mathcal{M}^\bullet \ar[dd] \ar[ldd] \\
Lf^*f_*\text{Tot}(\mathcal{K}^\bullet
\otimes_{\mathcal{O}_X}
\mathcal{M}^\bullet) \ar[d] \\
f^*f_*\text{Tot}(\mathcal{K}^\bullet \otimes_{\mathcal{O}_X}
\mathcal{M}^\bullet) \ar[rd] &
\text{Tot}(f^*f_*\mathcal{K}^\bullet \otimes_{\mathcal{O}_X}
f^*f_*\mathcal{M}^\bullet) \ar[d] &
\mathcal{K}^\bullet \otimes_{\mathcal{O}_X}^\mathbf{L}
\mathcal{M}^\bullet \ar[ld] \\
& \text{Tot}(\mathcal{K}^\bullet
\otimes_{\mathcal{O}_X}
\mathcal{M}^\bullet)
}
$$
此图中的所有多边形都交换。最上方的多边形由引理
\ref{cohomology-lemma-tensor-pull-compatibility} 交换。
含有两个朴素杯积的方块交换，因为
$Lf^* \to f^*$ 关于模复形具有函子性。
含有下列两个映射的方块也同样交换：
$\mathcal{A}^\bullet \otimes^\mathbf{L} \mathcal{B}^\bullet \to
\text{Tot}(\mathcal{A}^\bullet \otimes \mathcal{B}^\bullet)$.
最后，剩余方块在复形层次上交换，并且可视为朴素杯积的定义
（利用 $f^*$ 与 $f_*$ 的伴随性）。沿图的外边界绕行所得的
正是上面给出的两个映射，故证明完成。
\end{proof}

\noindent
设 $(X, \mathcal{O}_X)$ 是环空间，而 $\mathcal{K}^\bullet$ 与
$\mathcal{M}^\bullet$ 是 $\mathcal{O}_X$-模复形。
则有“朴素”杯积
$$
\mu' :
\text{Tot}(
\Gamma(X, \mathcal{K}^\bullet) \otimes_A \Gamma(X, \mathcal{M}^\bullet))
\longrightarrow
\Gamma(X, \text{Tot}(
\mathcal{K}^\bullet \otimes_{\mathcal{O}_X} \mathcal{M}^\bullet))
$$
将引理 \ref{cohomology-lemma-cup-compatible-with-naive}
应用于态射 $(X, \mathcal{O}_X) \to (pt, A)$，可知此朴素杯积与
本节第一段定义的杯积 $\mu$ 由下列交换图联系：
$$
\xymatrix{
\Gamma(X, \mathcal{K}^\bullet)
\otimes_A^\mathbf{L}
\Gamma(X, \mathcal{M}^\bullet) \ar[d] \ar[r] &
R\Gamma(X, \mathcal{K}^\bullet)
\otimes_A^\mathbf{L}
R\Gamma(X, \mathcal{M}^\bullet) \ar[d]^-\mu \\
\text{Tot}(\Gamma(X, \mathcal{K}^\bullet)
\otimes_A
\Gamma(X, \mathcal{M}^\bullet)) \ar[d]_-{\mu'} &
R\Gamma(X, \mathcal{K}^\bullet
\otimes_{\mathcal{O}_X}^\mathbf{L}
\mathcal{M}^\bullet) \ar[d] \\
\Gamma(X, \text{Tot}(\mathcal{K}^\bullet
\otimes_{\mathcal{O}_X}
\mathcal{M}^\bullet)) \ar[r] &
R\Gamma(X, \text{Tot}(\mathcal{K}^\bullet
\otimes_{\mathcal{O}_X}
\mathcal{M}^\bullet))
}
$$
该图位于 $D(A)$ 中。在上同调上得到交换图
$$
\xymatrix{
H^i(\Gamma(X, \mathcal{K}^\bullet)) \times
H^j(\Gamma(X, \mathcal{M}^\bullet)) \ar[d] \ar[r] &
H^{i + j}(X,
\text{Tot}(\mathcal{K}^\bullet \otimes_{\mathcal{O}_X} \mathcal{M}^\bullet)) \\
H^i(X, \mathcal{K}^\bullet) \times
H^j(X, \mathcal{M}^\bullet) \ar[r]^\cup &
H^{i + j}(X, \mathcal{K}^\bullet \otimes_{\mathcal{O}_X}^\mathbf{L}
\mathcal{M}^\bullet) \ar[u]
}
$$
，它把朴素杯积与真正的杯积联系起来。

\begin{lemma}
\label{cohomology-lemma-diagrams-commute}
设 $(X, \mathcal{O}_X)$ 是环空间。设
$\mathcal{K}^\bullet$ 与 $\mathcal{M}^\bullet$
是 $\mathcal{O}_X$-模的有下界复形，并设
$\mathcal{U} : X = \bigcup_{i \in I} U_i$ 是开覆盖。则
$$
\xymatrix{
\text{Tot}(\check{\mathcal{C}}^\bullet(\mathcal{U}, \mathcal{K}^\bullet))
\otimes_A^\mathbf{L}
\text{Tot}(\check{\mathcal{C}}^\bullet(\mathcal{U}, \mathcal{M}^\bullet))
\ar[d] \ar[r] &
R\Gamma(X, \mathcal{K}^\bullet)
\otimes_A^\mathbf{L}
R\Gamma(X, \mathcal{M}^\bullet) \ar[d]^\mu \\
\text{Tot}(
\text{Tot}(\check{\mathcal{C}}^\bullet(\mathcal{U}, \mathcal{K}^\bullet))
\otimes_A
\text{Tot}(\check{\mathcal{C}}^\bullet(\mathcal{U}, \mathcal{M}^\bullet)))
\ar[d]^{(\ref{cohomology-equation-needs-signs})} &
R\Gamma(X,
\mathcal{K}^\bullet \otimes_{\mathcal{O}_X}^\mathbf{L} \mathcal{M}^\bullet)
\ar[d] \\
\text{Tot}(
\check{\mathcal{C}}^\bullet({\mathcal U},
\text{Tot}(\mathcal{K}^\bullet \otimes_{\mathcal{O}_X} \mathcal{M}^\bullet)
)) \ar[r] &
R\Gamma(X,
\text{Tot}(\mathcal{K}^\bullet \otimes_{\mathcal{O}_X} \mathcal{M}^\bullet))
}
$$
在 $D(A)$ 中交换，其中水平箭头是引理
\ref{cohomology-lemma-cech-complex-complex} 中的箭头。
\end{lemma}

\begin{proof}
如引理 \ref{cohomology-lemma-godement-resolution-bounded-below} 中那样，
选取复形的拟同构
$a : \mathcal{K}^\bullet \to \mathcal{K}_1^\bullet$ 与
$b : \mathcal{M}^\bullet \to \mathcal{M}_1^\bullet$。
由于映射 $a$ 和 $b$ 在茎上都是链同伦等价，所诱导的映射
$$
\text{Tot}(\mathcal{K}^\bullet \otimes_{\mathcal{O}_X} \mathcal{M}^\bullet)
\to
\text{Tot}(\mathcal{K}_1^\bullet \otimes_{\mathcal{O}_X} \mathcal{M}_1^\bullet)
$$
在茎上也为链同伦等价（《更多代数》引理
\ref{more-algebra-lemma-derived-tor-homotopy}），因而是拟同构。因此两个图的目标
$$
R\Gamma(X,
\text{Tot}(\mathcal{K}^\bullet
\otimes_{\mathcal{O}_X} \mathcal{M}^\bullet)) =
R\Gamma(X,
\text{Tot}(\mathcal{K}_1^\bullet
\otimes_{\mathcal{O}_X} \mathcal{M}_1^\bullet))
$$
在 $D(A)$ 中相同。因此只需证明把 $\mathcal{K}$ 与 $\mathcal{M}$
分别替换为 $\mathcal{K}_1$ 与 $\mathcal{M}_1$ 后该图交换。
这就归结为下一段讨论的情形。

\medskip\noindent
假设 $\mathcal{K}^\bullet$ 与 $\mathcal{M}^\bullet$
是松弛 $\mathcal{O}_X$-模的有下界复形，并考察把杯积与
{\v C}ech 复形上的杯积 (\ref{cohomology-equation-needs-signs}) 联系起来的图。
此时可考察交换图
$$
\xymatrix{
\Gamma(X, \mathcal{K}^\bullet)
\otimes_A^\mathbf{L}
\Gamma(X, \mathcal{M}^\bullet) \ar[d] \ar[r] &
\text{Tot}(\check{\mathcal{C}}^\bullet(\mathcal{U}, \mathcal{K}^\bullet))
\otimes_A^\mathbf{L}
\text{Tot}(\check{\mathcal{C}}^\bullet(\mathcal{U}, \mathcal{M}^\bullet))
\ar[d] \\
\text{Tot}(\Gamma(X, \mathcal{K}^\bullet)
\otimes_A
\Gamma(X, \mathcal{M}^\bullet)) \ar[d] \ar[r] &
\text{Tot}(
\text{Tot}(\check{\mathcal{C}}^\bullet(\mathcal{U}, \mathcal{K}^\bullet))
\otimes_A
\text{Tot}(\check{\mathcal{C}}^\bullet(\mathcal{U}, \mathcal{M}^\bullet)))
\ar[d]^{(\ref{cohomology-equation-needs-signs})} \\
\Gamma(X, \text{Tot}(\mathcal{K}^\bullet
\otimes_{\mathcal{O}_X}
\mathcal{M}^\bullet)) \ar[r] &
\text{Tot}(
\check{\mathcal{C}}^\bullet({\mathcal U},
\text{Tot}(\mathcal{K}^\bullet \otimes_{\mathcal{O}_X} \mathcal{M}^\bullet)
))
}
$$
此图中的水平箭头是 $D(A)$ 中的同构，因为对
$\mathcal{K}^\bullet$ 这样的松弛模有下界复形，有
$$
\Gamma(X, \mathcal{K}^\bullet) =
\text{Tot}(\check{\mathcal{C}}^\bullet(\mathcal{U}, \mathcal{K}^\bullet)) =
R\Gamma(X, \mathcal{K}^\bullet)
$$
，此等式位于 $D(A)$ 中。它由引理 \ref{cohomology-lemma-flasque-acyclic}、
《导出范畴》引理 \ref{derived-lemma-leray-acyclicity} 以及引理 \ref{cohomology-lemma-cech-complex-complex-computes} 得到。
因此，引理之图中涉及 (\ref{cohomology-equation-needs-signs}) 的部分交换性，
由已经证明的引理 \ref{cohomology-lemma-cup-compatible-with-naive}
在 $f$ 为到一点的态射时的交换性得到（见引理
\ref{cohomology-lemma-cup-compatible-with-naive} 后的讨论）。
\end{proof}

\begin{lemma}
\label{cohomology-lemma-cup-product-associative}
设 $f : (X, \mathcal{O}_X) \to (Y, \mathcal{O}_Y)$
是环空间的态射。注记 \ref{cohomology-remark-cup-product} 的相对杯积具有结合性，
意即图
$$
\xymatrix{
Rf_*K \otimes_{\mathcal{O}_Y}^\mathbf{L}
Rf_*L \otimes_{\mathcal{O}_Y}^\mathbf{L}
Rf_*M \ar[r] \ar[d] &
Rf_*(K \otimes_{\mathcal{O}_X}^\mathbf{L} L)
\otimes_{\mathcal{O}_Y}^\mathbf{L} Rf_*M \ar[d] \\
Rf_*K \otimes_{\mathcal{O}_Y}^\mathbf{L}
Rf_*(L \otimes_{\mathcal{O}_X}^\mathbf{L} M) \ar[r] &
Rf_*(K \otimes_{\mathcal{O}_X}^\mathbf{L} 
L \otimes_{\mathcal{O}_X}^\mathbf{L} M)
}
$$
对 $D(\mathcal{O}_X)$ 中所有 $K, L, M$，都在
$D(\mathcal{O}_Y)$ 中交换。
\end{lemma}

\begin{proof}
沿任一侧绕行，都会得到与下列显然映射伴随的映射：
\begin{align*}
Lf^*(Rf_*K \otimes_{\mathcal{O}_Y}^\mathbf{L}
Rf_*L \otimes_{\mathcal{O}_Y}^\mathbf{L}
Rf_*M) & =
Lf^*(Rf_*K) \otimes_{\mathcal{O}_X}^\mathbf{L}
Lf^*(Rf_*L) \otimes_{\mathcal{O}_X}^\mathbf{L}
Lf^*(Rf_*M) \\
& \to
K \otimes_{\mathcal{O}_X}^\mathbf{L} 
L \otimes_{\mathcal{O}_X}^\mathbf{L} M
\end{align*}
该映射位于 $D(\mathcal{O}_X)$ 中。
\end{proof}

\begin{lemma}
\label{cohomology-lemma-cup-product-commutative}
设 $f : (X, \mathcal{O}_X) \to (Y, \mathcal{O}_Y)$
是环空间的态射。注记 \ref{cohomology-remark-cup-product} 的相对杯积具有交换性，
意即图
$$
\xymatrix{
Rf_*K \otimes_{\mathcal{O}_Y}^\mathbf{L} Rf_*L \ar[r] \ar[d]_\psi &
Rf_*(K \otimes_{\mathcal{O}_X}^\mathbf{L} L) \ar[d]^{Rf_*\psi} \\
Rf_*L \otimes_{\mathcal{O}_Y}^\mathbf{L} Rf_*K \ar[r] &
Rf_*(L \otimes_{\mathcal{O}_X}^\mathbf{L} K)
}
$$
对 $D(\mathcal{O}_X)$ 中所有 $K, L$，都在
$D(\mathcal{O}_Y)$ 中交换。这里 $\psi$ 是导出范畴上的交换约束
（引理 \ref{cohomology-lemma-symmetric-monoidal-derived}）。
\end{lemma}

\begin{proof}
从略。
\end{proof}

\begin{lemma}
\label{cohomology-lemma-compose-cup-product}
设 $f : (X, \mathcal{O}_X) \to (Y, \mathcal{O}_Y)$ 与
$g : (Y, \mathcal{O}_Y) \to (Z, \mathcal{O}_Z)$
是环空间的态射。注记 \ref{cohomology-remark-cup-product} 的相对杯积与复合相容，
意即图
$$
\xymatrix{
R(g \circ f)_*K \otimes_{\mathcal{O}_Z}^\mathbf{L} R(g \circ f)_*L
\ar@{=}[rr] \ar[d] & &
Rg_*Rf_*K \otimes_{\mathcal{O}_Z}^\mathbf{L} Rg_*Rf_*L \ar[d] \\
R(g \circ f)_*(K \otimes_{\mathcal{O}_X}^\mathbf{L} L) \ar@{=}[r] &
Rg_*Rf_*(K \otimes_{\mathcal{O}_X}^\mathbf{L} L) &
Rg_*(Rf_*K \otimes_{\mathcal{O}_Y}^\mathbf{L}  Rf_*L) \ar[l]
}
$$
对 $D(\mathcal{O}_X)$ 中所有 $K, L$，都在
$D(\mathcal{O}_Z)$ 中交换。
\end{lemma}

\begin{proof}
这是因为无论沿图以哪种方式绕行，都会得到与下列映射伴随的映射：
\begin{align*}
& L(g \circ f)^*\left(R(g \circ f)_*K
\otimes_{\mathcal{O}_Z}^\mathbf{L}
R(g \circ f)_*L\right) \\
& =
L(g \circ f)^*R(g \circ f)_*K
\otimes_{\mathcal{O}_X}^\mathbf{L}
L(g \circ f)^*R(g \circ f)_*L) \\
& \to
K \otimes_{\mathcal{O}_X}^\mathbf{L} L
\end{align*}
该映射位于 $D(\mathcal{O}_X)$ 中。为此要用到如下余单位复合：
$$
L(g \circ f)^*R(g \circ f)_* =
Lf^* Lg^* Rg_* Rf_*  \to
Lf^* Rf_* \to \text{id}
$$
它是 $L(g \circ f)^*$ 与 $R(g \circ f)_*$ 的余单位；参见《范畴》引理
\ref{categories-lemma-compose-counits}。
\end{proof}

\begin{lemma}
\label{cohomology-lemma-base-change-cup-product}
考察环空间的交换方块
$$
\xymatrix{
(X', \mathcal{O}_{X'}) \ar[r]_{g'} \ar[d]_{f'} &
(X, \mathcal{O}_X) \ar[d]^f \\
(Y', \mathcal{O}_{Y'}) \ar[r]^g & (Y, \mathcal{O}_Y) 
}
$$
。设 $K, L$ 属于 $D(\mathcal{O}_X)$。相对杯积与此方块相容，
意即图
$$
\xymatrix{
Lg^*(Rf_*K \otimes_{\mathcal{O}_Y}^\mathbf{L} Rf_*L)
\ar[r] \ar@{=}[d] &
Lg^*(Rf_*(K \otimes_{\mathcal{O}_X}^\mathbf{L} L)) \ar[d] \\
Lg^*Rf_*K \otimes_{\mathcal{O}_{Y'}}^\mathbf{L} Lg^*Rf_*L \ar[d] &
R(f')_*L(g')^*(K \otimes_{\mathcal{O}_X}^\mathbf{L} L) \ar@{=}[d] \\
R(f')_*(L(g')^*K \otimes_{\mathcal{O}_{Y'}} R(f')_*(L(g')^*L \ar[r] &
R(f')_*(L(g')^*K \otimes_{\mathcal{O}_{X'}}^\mathbf{L} L(g')^*L)
}
$$
在 $D(\mathcal{O}_{Y'})$ 中交换。水平箭头由相对杯积给出
（注记 \ref{cohomology-remark-cup-product}），竖直箭头由基变换映射
（注记 \ref{cohomology-remark-base-change}）与引理
\ref{cohomology-lemma-pullback-tensor-product} 给出。
\end{lemma}

\begin{proof}
从略。
\end{proof}














\section{K-内射复形的一些性质}
\label{cohomology-section-properties-K-injective}

\noindent
设 $(X, \mathcal{O}_X)$ 是环空间，$U \subset X$ 是开子集，并以
$j : (U, \mathcal{O}_U) \to (X, \mathcal{O}_X)$
表示相应的开浸入。拉回函子 $j^*$ 就是限制函子，因而正合。
所以在任意复形上，导出拉回 $Lj^*$ 只需限制该复形即可计算。
我们通常把相应函子简记为
$$
D(\mathcal{O}_X) \to D(\mathcal{O}_U),
\quad
E \mapsto j^*E = E|_U
$$
类似地，零延拓
$j_! : \textit{Mod}(\mathcal{O}_U) \to \textit{Mod}(\mathcal{O}_X)$
是正合函子（参见《层》第 \ref{sheaves-section-open-immersions} 节及
《模》引理 \ref{modules-lemma-j-shriek-exact}）。因此它诱导函子
$$
j_! : D(\mathcal{O}_U) \to D(\mathcal{O}_X),\quad
F \mapsto j_!F
$$
：只需对表示对象 $F$ 的任意复形逐项应用 $j_!$。

\begin{lemma}
\label{cohomology-lemma-restrict-K-injective-to-open}
设 $X$ 是环空间，而 $U \subset X$ 是开子空间。
$\mathcal{O}_X$-模的 K-内射复形限制到 $U$ 后，
是 $\mathcal{O}_U$-模的 K-内射复形。
\end{lemma}

\begin{proof}
这由《导出范畴》引理
\ref{derived-lemma-adjoint-preserve-K-injectives}
以及限制函子具有正合左伴随 $j_!$ 这一事实推出。
$j_!$ 的构造见《层》第 \ref{sheaves-section-open-immersions} 节，
其正合性见《模》引理 \ref{modules-lemma-j-shriek-exact}。
\end{proof}

\begin{lemma}
\label{cohomology-lemma-unbounded-cohomology-of-open}
设 $X$ 是环空间，而 $U \subset X$ 是开子空间。
对 $D(\mathcal{O}_X)$ 中的 $K$，有
$H^p(U, K) = H^p(U, K|_U)$.
\end{lemma}

\begin{proof}
取表示 $K$ 的 $\mathcal{O}_X$-模 K-内射复形
$\mathcal{I}^\bullet$。则
$$
H^q(U, K) = H^q(\Gamma(U, \mathcal{I}^\bullet)) =
H^q(\Gamma(U, \mathcal{I}^\bullet|_U))
$$
。这是由上同调的构造得到的。由引理
\ref{cohomology-lemma-restrict-K-injective-to-open}，
复形 $\mathcal{I}^\bullet|_U$ 是表示 $K|_U$ 的 K-内射复形，
故引理成立。
\end{proof}

\begin{lemma}
\label{cohomology-lemma-sheafification-cohomology}
设 $(X, \mathcal{O}_X)$ 是环空间，而 $K$ 是
$D(\mathcal{O}_X)$ 的对象。预层
$$
U \mapsto H^q(U, K) = H^q(U, K|_U)
$$
的层化是 $K$ 的第 $q$ 个上同调层 $H^q(K)$。
\end{lemma}

\begin{proof}
等式 $H^q(U, K) = H^q(U, K|_U)$ 由引理
\ref{cohomology-lemma-unbounded-cohomology-of-open} 成立。
选取表示 $K$ 的 K-内射复形 $\mathcal{I}^\bullet$。则
$$
H^q(U, K) =
\frac{\Ker(\mathcal{I}^q(U) \to \mathcal{I}^{q + 1}(U))}
{\Im(\mathcal{I}^{q - 1}(U) \to \mathcal{I}^q(U))}.
$$
此式由我们的上同调构造得到。又因为
$H^q(K) = \Ker(\mathcal{I}^q \to \mathcal{I}^{q + 1})/
\Im(\mathcal{I}^{q - 1} \to \mathcal{I}^q)$，结论显然。
\end{proof}

\begin{lemma}
\label{cohomology-lemma-restrict-direct-image-open}
设 $f : (X, \mathcal{O}_X) \to (Y, \mathcal{O}_Y)$ 是环空间的态射。
给定开子空间 $V \subset Y$，令 $U = f^{-1}(V)$，并以
$g : U \to V$ 表示所诱导的态射。则对 $D(\mathcal{O}_X)$ 中的 $E$，
有 $(Rf_*E)|_V = Rg_*(E|_U)$。
\end{lemma}

\begin{proof}
用 $\mathcal{O}_X$-模 K-内射复形
$\mathcal{I}^\bullet$ 表示 $E$。则
$Rf_*(E) = f_*\mathcal{I}^\bullet$，并且由引理
\ref{cohomology-lemma-restrict-K-injective-to-open}，
$Rg_*(E|_U) = g_*(\mathcal{I}^\bullet|_U)$。
对 $X$ 上任意层 $\mathcal{F}$，显然有
$(f_*\mathcal{F})|_V = g_*(\mathcal{F}|_U)$，故结论成立。
\end{proof}

\begin{lemma}
\label{cohomology-lemma-Leray-unbounded}
设 $f : X \to Y$ 是环空间的态射。作为函子
$D(\mathcal{O}_X) \to D(\Gamma(Y, \mathcal{O}_Y))$，有
$R\Gamma(Y, -) \circ Rf_* = R\Gamma(X, -)$。
更一般地，若 $V \subset Y$ 开且 $U = f^{-1}(V)$，则
$R\Gamma(U, -) = R\Gamma(V, -) \circ Rf_*$。
\end{lemma}

\begin{proof}
令 $Z$ 为底空间是单点且满足
$\Gamma(Z, \mathcal{O}_Z) = \Gamma(Y, \mathcal{O}_Y)$ 的环空间。
存在环空间的典范态射 $Y \to Z$，它在结构层的整体截面上诱导上述等同。
于是 $D(\mathcal{O}_Z) = D(\Gamma(Y, \mathcal{O}_Y))$。
因此，把引理 \ref{cohomology-lemma-derived-pushforward-composition}
应用于 $X \to Y \to Z$，即得
$R\Gamma(Y, -) \circ Rf_* = R\Gamma(X, -)$。

\medskip\noindent
将引理 \ref{cohomology-lemma-restrict-direct-image-open} 应用于第一个断言，
即得第二个（更一般的）断言。
\end{proof}

\begin{lemma}
\label{cohomology-lemma-unbounded-describe-higher-direct-images}
设 $f : (X, \mathcal{O}_X) \to (Y, \mathcal{O}_Y)$ 是环空间的态射，
而 $K$ 属于 $D(\mathcal{O}_X)$。则 $H^i(Rf_*K)$ 是与下列预层相伴的层：
$$
V \mapsto H^i(f^{-1}(V), K) = H^i(V, Rf_*K)
$$
\end{lemma}

\begin{proof}
对引理 \ref{cohomology-lemma-Leray-unbounded} 的第二个断言取上同调，
即得等式 $H^i(f^{-1}(V), K) = H^i(V, Rf_*K)$。
层化断言于是由引理 \ref{cohomology-lemma-sheafification-cohomology} 得到。
\end{proof}

\begin{lemma}
\label{cohomology-lemma-modules-abelian-unbounded}
设 $X$ 是环空间，$K$ 是 $D(\mathcal{O}_X)$ 的对象，并以
$K_{ab}$ 表示它在 $D(\underline{\mathbf{Z}}_X)$ 中的像。
\begin{enumerate}
\item 对任意开集 $U \subset X$，存在典范映射
$R\Gamma(U, K) \to R\Gamma(U, K_{ab})$
，它是 $D(\textit{Ab})$ 中的同构。
\item 设 $f : X \to Y$ 是环空间的态射。存在典范映射
$Rf_*K \to Rf_*(K_{ab})$，它是
$D(\underline{\mathbf{Z}}_Y)$ 中的同构。
\end{enumerate}
\end{lemma}

\begin{proof}
映射构造如下。选取表示 $K$ 的 K-内射复形
$\mathcal{I}^\bullet$，并选取拟同构
$\mathcal{I}^\bullet \to \mathcal{J}^\bullet$，其中
$\mathcal{J}^\bullet$ 是阿贝尔群的 K-内射复形。则 (1) 中的映射为
$\Gamma(U, \mathcal{I}^\bullet) \to \Gamma(U, \mathcal{J}^\bullet)$，
而 (2) 中的映射为
$f_*\mathcal{I}^\bullet \to f_*\mathcal{J}^\bullet$。
要证明这些映射是同构，只需证明它们在上同调群和上同调层上诱导同构。
由引理 \ref{cohomology-lemma-unbounded-cohomology-of-open} 与
\ref{cohomology-lemma-unbounded-describe-higher-direct-images}，只需证明映射
$$
H^0(X, K) \longrightarrow H^0(X, K_{ab})
$$
是同构。注意
$$
H^0(X, K) = \Hom_{D(\mathcal{O}_X)}(\mathcal{O}_X, K)
$$
，另一个群也有类似表达式。选取表示 $K$ 的任意
$\mathcal{O}_X$-模复形 $\mathcal{K}^\bullet$。
由导出范畴作为局部化的构造，有
$$
\Hom_{D(\mathcal{O}_X)}(\mathcal{O}_X, K) =
\colim_{s : \mathcal{F}^\bullet \to \mathcal{O}_X}
\Hom_{K(\mathcal{O}_X)}(\mathcal{F}^\bullet, \mathcal{K}^\bullet)
$$
，其中余极限取遍 $\mathcal{O}_X$-模复形的拟同构 $s$。类似地，有
$$
\Hom_{D(\underline{\mathbf{Z}}_X)}(\underline{\mathbf{Z}}_X, K) =
\colim_{s : \mathcal{G}^\bullet \to \underline{\mathbf{Z}}_X}
\Hom_{K(\underline{\mathbf{Z}}_X)}(\mathcal{G}^\bullet, \mathcal{K}^\bullet)
$$
接着注意，拟同构
$s : \mathcal{G}^\bullet \to \underline{\mathbf{Z}}_X$
中满足 $\mathcal{G}^\bullet$ 是平坦
$\underline{\mathbf{Z}}_X$-模的有上界复形者，在该系统中余终。
（这由《模》引理 \ref{modules-lemma-module-quotient-flat}
与《导出范畴》引理 \ref{derived-lemma-subcategory-left-resolution} 得到；
参见第 \ref{cohomology-section-flat} 节的讨论。）因此可以构造映射
$H^0(X, K) \longrightarrow H^0(X, K_{ab})$
的逆映射：把元素 $\xi \in H^0(X, K_{ab})$ 表示为一对
$$
(s : \mathcal{G}^\bullet \to \underline{\mathbf{Z}}_X,
a : \mathcal{G}^\bullet \to \mathcal{K}^\bullet)
$$
，其中 $\mathcal{G}^\bullet$ 是平坦
$\underline{\mathbf{Z}}_X$-模的有上界复形，并将其映到
$$
(\mathcal{G}^\bullet \otimes_{\underline{\mathbf{Z}}_X} \mathcal{O}_X
\to \mathcal{O}_X,
\mathcal{G}^\bullet  \otimes_{\underline{\mathbf{Z}}_X} \mathcal{O}_X
\to \mathcal{K}^\bullet)
$$
。这里唯一需要注意的是，由引理
\ref{cohomology-lemma-derived-tor-quasi-isomorphism-other-side} 与
\ref{cohomology-lemma-bounded-flat-K-flat}，第一个箭头是拟同构。
关于此构造确为逆映射的详细验证从略。
\end{proof}

\begin{lemma}
\label{cohomology-lemma-adjoint-lower-shriek-restrict}
设 $(X, \mathcal{O}_X)$ 是环空间，$U \subset X$ 是开子集，并以
$j : (U, \mathcal{O}_U) \to (X, \mathcal{O}_X)$
表示相应的开浸入。限制函子
$D(\mathcal{O}_X) \to D(\mathcal{O}_U)$ 是零延拓
$j_! : D(\mathcal{O}_U) \to D(\mathcal{O}_X)$ 的右伴随。
\end{lemma}

\begin{proof}
这由 $j_!$ 与 $j^*$ 互为伴随且均正合这一事实形式地推出
（因而 $Lj_! = j_!$ 与 $Rj^* = j^*$ 存在）；参见《导出范畴》引理
\ref{derived-lemma-derived-adjoint-functors}。
\end{proof}

\begin{lemma}
\label{cohomology-lemma-K-injective-flat}
设 $f : X \to Y$ 是环空间的平坦态射。
若 $\mathcal{I}^\bullet$ 是 $\mathcal{O}_X$-模的 K-内射复形，
则 $f_*\mathcal{I}^\bullet$ 作为 $\mathcal{O}_Y$-模复形是 K-内射的。
\end{lemma}

\begin{proof}
这是因为
$$
\Hom_{K(\mathcal{O}_Y)}(\mathcal{F}^\bullet, f_*\mathcal{I}^\bullet)
=
\Hom_{K(\mathcal{O}_X)}(f^*\mathcal{F}^\bullet, \mathcal{I}^\bullet)
$$
；此等式由《层》引理
\ref{sheaves-lemma-adjoint-pullback-pushforward-modules}
以及 $f$ 平坦从而 $f^*$ 正合这一事实得到。
\end{proof}





\section{无界 Mayer--Vietoris}
\label{cohomology-section-unbounded-mayer-vietoris}

\noindent
无界上同调也有 Mayer--Vietoris 序列。

\begin{lemma}
\label{cohomology-lemma-exact-sequence-lower-shriek}
设 $(X, \mathcal{O}_X)$ 是环空间，且
$X = U \cup V$ 是两个开子空间之并。对 $D(\mathcal{O}_X)$
的任意对象 $E$，有特异三角
$$
j_{U \cap V!}E|_{U \cap V} \to
j_{U!}E|_U \oplus j_{V!}E|_V \to E \to 
j_{U \cap V!}E|_{U \cap V}[1]
$$
，它位于 $D(\mathcal{O}_X)$ 中。
\end{lemma}

\begin{proof}
我们已在第 \ref{cohomology-section-properties-K-injective} 节看到，
限制函子与零延拓函子在任意复形上都只需直接应用相应函子即可计算。
设 $\mathcal{E}^\bullet$ 是表示 $E$ 的
$\mathcal{O}_X$-模复形。引理的特异三角，是与下列
$\mathcal{O}_X$-模复形短正合列相伴的特异三角
（见《导出范畴》第 \ref{derived-section-canonical-delta-functor} 节，
尤其是引理 \ref{derived-lemma-derived-canonical-delta-functor}）：
$$
0 \to j_{U \cap V!}\mathcal{E}^\bullet|_{U \cap V} \to
j_{U!}\mathcal{E}^\bullet|_U \oplus j_{V!}\mathcal{E}^\bullet|_V
\to \mathcal{E}^\bullet \to 0
$$
利用《层》引理 \ref{sheaves-lemma-j-shriek-modules}
在茎上检验即可看出此序列正合（计算从略）。
\end{proof}

\begin{lemma}
\label{cohomology-lemma-exact-sequence-j-star}
设 $(X, \mathcal{O}_X)$ 是环空间，且
$X = U \cup V$ 是两个开子空间之并。对 $D(\mathcal{O}_X)$
的任意对象 $E$，有特异三角
$$
E \to 
Rj_{U, *}E|_U \oplus Rj_{V, *}E|_V \to
Rj_{U \cap V, *}E|_{U \cap V} \to
E[1]
$$
，它位于 $D(\mathcal{O}_X)$ 中。
\end{lemma}

\begin{proof}
选取表示 $E$ 的 K-内射复形 $\mathcal{I}^\bullet$，
使其各项 $\mathcal{I}^n$ 都是
$\textit{Mod}(\mathcal{O}_X)$ 的内射对象；参见《内射对象》定理
\ref{injectives-theorem-K-injective-embedding-grothendieck}。
我们已经看到 $\mathcal{I}^\bullet|U$ 也是 K-内射复形
（引理 \ref{cohomology-lemma-restrict-K-injective-to-open}）。因此
$Rj_{U, *}E|_U$ 由 $j_{U, *}\mathcal{I}^\bullet|_U$ 表示；
$V$ 与 $U \cap V$ 的情形类似。于是引理的特异三角，是与下列
复形短正合列相伴的特异三角
（见《导出范畴》第 \ref{derived-section-canonical-delta-functor} 节，
尤其是引理 \ref{derived-lemma-derived-canonical-delta-functor}）：
$$
0 \to
\mathcal{I}^\bullet \to
j_{U, *}\mathcal{I}^\bullet|_U \oplus j_{V, *}\mathcal{I}^\bullet|_V \to
j_{U \cap V, *}\mathcal{I}^\bullet|_{U \cap V} \to
0.
$$
此序列正合，因为对任意开集 $W \subset X$ 及任意 $n$，序列
$$
0 \to
\mathcal{I}^n(W) \to
\mathcal{I}^n(W \cap U) \oplus \mathcal{I}^n(W \cap V) \to
\mathcal{I}^n(W \cap U \cap V) \to
0
$$
正合（见引理 \ref{cohomology-lemma-mayer-vietoris} 的证明）。
\end{proof}

\begin{lemma}
\label{cohomology-lemma-mayer-vietoris-hom}
设 $(X, \mathcal{O}_X)$ 是环空间，且 $X = U \cup V$
是 $X$ 的两个开子空间之并。对 $D(\mathcal{O}_X)$ 的对象
$E,F$，有 Mayer--Vietoris 序列
$$
\xymatrix{
& \ldots \ar[r] & \Ext^{-1}(E_{U \cap V}, F_{U \cap V}) \ar[lld] \\
\Hom(E, F) \ar[r] &
\Hom(E_U, F_U) \oplus
\Hom(E_V, F_V) \ar[r] &
\Hom(E_{U \cap V}, F_{U \cap V})
}
$$
，其中下标表示限制到相应开集，而 $\Hom$ 与 $\Ext$
均取自相应的导出范畴。
\end{lemma}

\begin{proof}
使用引理 \ref{cohomology-lemma-exact-sequence-lower-shriek} 的特异三角，
得到 $\Hom$ 的长正合列
（由《导出范畴》引理 \ref{derived-lemma-representable-homological}），
再使用
$$
\Hom_{D(\mathcal{O}_X)}(j_{U!}E|_U, F) =
\Hom_{D(\mathcal{O}_U)}(E|_U, F|_U)
$$
；此等式由引理 \ref{cohomology-lemma-adjoint-lower-shriek-restrict} 得到。
\end{proof}

\begin{lemma}
\label{cohomology-lemma-unbounded-mayer-vietoris}
设 $(X, \mathcal{O}_X)$ 是环空间，并假设
$X = U \cup V$ 是两个开子集之并。对 $D(\mathcal{O}_X)$
的对象 $E$，有特异三角
$$
R\Gamma(X, E) \to R\Gamma(U, E) \oplus R\Gamma(V, E) \to
R\Gamma(U \cap V, E) \to R\Gamma(X, E)[1]
$$
，从而特别有上同调长正合列
$$
\ldots \to
H^n(X, E) \to
H^n(U, E) \oplus H^0(V, E) \to
H^n(U \cap V, E) \to
H^{n + 1}(X, E) \to \ldots
$$
特异三角与长正合列的构造关于 $E$ 具有函子性。
\end{lemma}

\begin{proof}
选取表示 $E$ 的 K-内射复形 $\mathcal{I}^\bullet$。
可假设对所有 $n$，$\mathcal{I}^n$ 都是
$\textit{Mod}(\mathcal{O}_X)$ 的内射对象；参见《内射对象》定理
\ref{injectives-theorem-K-injective-embedding-grothendieck}。
于是 $R\Gamma(X, E)$ 由 $\Gamma(X, \mathcal{I}^\bullet)$ 计算。
由引理 \ref{cohomology-lemma-restrict-K-injective-to-open}，
$U$、$V$ 及 $U \cap V$ 的情形也类似。因此，引理的特异三角
是与下列复形短正合列相伴的特异三角
（见《导出范畴》第 \ref{derived-section-canonical-delta-functor} 节，
尤其是引理 \ref{derived-lemma-derived-canonical-delta-functor}）：
$$
0 \to
\mathcal{I}^\bullet(X) \to
\mathcal{I}^\bullet(U) \oplus \mathcal{I}^\bullet(V) \to
\mathcal{I}^\bullet(U \cap V) \to
0.
$$
我们已在引理 \ref{cohomology-lemma-mayer-vietoris} 的证明中看到这是短正合列。
最后一个断言由《内射对象》定理
\ref{injectives-theorem-K-injective-embedding-grothendieck}
中构造的函子性得到。
\end{proof}

\begin{lemma}
\label{cohomology-lemma-unbounded-relative-mayer-vietoris}
设 $f : X \to Y$ 是环空间的态射，并假设
$X = U \cup V$ 是两个开子集之并。以
$a = f|_U : U \to Y$、$b = f|_V : V \to Y$ 及
$c = f|_{U \cap V} : U \cap V \to Y$
表示相应的限制。对 $D(\mathcal{O}_X)$ 的每个对象 $E$，
存在特异三角
$$
Rf_*E \to
Ra_*(E|_U) \oplus Rb_*(E|_V) \to
Rc_*(E|_{U \cap V}) \to
Rf_*E[1]
$$
此三角关于 $E$ 具有函子性。
\end{lemma}

\begin{proof}
选取表示 $E$ 的 K-内射复形 $\mathcal{I}^\bullet$。
可假设对所有 $n$，$\mathcal{I}^n$ 都是
$\textit{Mod}(\mathcal{O}_X)$ 的内射对象；参见《内射对象》定理
\ref{injectives-theorem-K-injective-embedding-grothendieck}。
于是 $Rf_*E$ 由 $f_*\mathcal{I}^\bullet$ 计算。
由引理 \ref{cohomology-lemma-restrict-K-injective-to-open}，
$U$、$V$ 及 $U \cap V$ 的情形也类似。因此，引理的特异三角
是与下列复形短正合列相伴的特异三角
（见《导出范畴》第 \ref{derived-section-canonical-delta-functor} 节，
尤其是引理 \ref{derived-lemma-derived-canonical-delta-functor}）：
$$
0 \to
f_*\mathcal{I}^\bullet \to
a_*\mathcal{I}^\bullet|_U \oplus b_*\mathcal{I}^\bullet|_V \to
c_*\mathcal{I}^\bullet|_{U \cap V} \to
0.
$$
这是复形的短正合列；这由引理
\ref{cohomology-lemma-relative-mayer-vietoris} 以及下述事实得到：
对 $\textit{Mod}(\mathcal{O}_X)$ 的内射对象 $\mathcal{I}$，
$R^1f_*\mathcal{I} = 0$。
最后一个断言由《内射对象》定理
\ref{injectives-theorem-K-injective-embedding-grothendieck}
中构造的函子性得到。
\end{proof}

\begin{lemma}
\label{cohomology-lemma-pushforward-restriction}
设 $(X, \mathcal{O}_X)$ 是环空间。设 $j : U \to X$ 是开子空间，
而 $T \subset X$ 是包含于 $U$ 的闭子集。
\begin{enumerate}
\item 若 $E$ 是 $D(\mathcal{O}_X)$ 的对象，且其上同调层支撑在
$T$ 上，则 $E \to Rj_*(E|_U)$ 是同构。
\item 若 $F$ 是 $D(\mathcal{O}_U)$ 的对象，且其上同调层支撑在
$T$ 上，则 $j_!F \to Rj_*F$ 是同构。
\end{enumerate}
\end{lemma}

\begin{proof}
令 $V = X \setminus T$ 且 $W = U \cap V$。注意
$X = U \cup V$ 是 $X$ 的开覆盖，并以
$j_W : W \to V$ 表示开浸入。
设 $E$ 是 $D(\mathcal{O}_X)$ 的对象，且其上同调层支撑在 $T$ 上。
由引理 \ref{cohomology-lemma-restrict-direct-image-open}，有
$(Rj_*E|_U)|_V = Rj_{W, *}(E|_W) = 0$，因为根据假设 $E|_W = 0$。
另一方面，$Rj_*(E|_U)|_U = E|_U$。故 (1) 显然。
设 $F$ 是 $D(\mathcal{O}_U)$ 的对象，且其上同调层支撑在 $T$ 上。
由引理 \ref{cohomology-lemma-restrict-direct-image-open}，有
$(Rj_*F)|_V = Rj_{W, *}(F|_W) = 0$，因为根据假设 $F|_W = 0$。
还有 $(j_!F)|_V = j_{W!}(F|_W) = 0$
（第一个等式直接来自零延拓的定义）。由于
$(Rj_*F)|_U = F$ 且 $(j_!F)|_U = F$，可知 (2) 成立。
\end{proof}

\begin{lemma}
\label{cohomology-lemma-mayer-vietoris-cup}
设 $(X, \mathcal{O}_X)$ 是环空间，并令
$A = \Gamma(X, \mathcal{O}_X)$。假设 $X = U \cup V$
是两个开子集之并。对 $D(\mathcal{O}_X)$ 的对象 $K$ 与 $M$，
有特异三角的态射
$$
\xymatrix{
R\Gamma(X, K) \otimes_A^\mathbf{L} R\Gamma(X, M) \ar[r] \ar[d] &
R\Gamma(X, K \otimes_{\mathcal{O}_X}^\mathbf{L} M) \ar[d] \\
R\Gamma(X, K) \otimes_A^\mathbf{L}
(R\Gamma(U, M) \oplus R\Gamma(V, M)) \ar[r] \ar[d] &
R\Gamma(U, K \otimes_{\mathcal{O}_X}^\mathbf{L} M)
\oplus R\Gamma(V, K \otimes_{\mathcal{O}_X}^\mathbf{L} M)) \ar[d] \\
R\Gamma(X, K) \otimes_A^\mathbf{L} R\Gamma(U \cap V, M) \ar[r] \ar[d] &
R\Gamma(U \cap V, K \otimes_{\mathcal{O}_X}^\mathbf{L} M) \ar[d] \\
R\Gamma(X, K) \otimes_A^\mathbf{L} R\Gamma(X, M)[1] \ar[r] &
R\Gamma(X, K \otimes_{\mathcal{O}_X}^\mathbf{L} M)[1]
}
$$
，其中
\begin{enumerate}
\item 水平箭头由杯积给出；
\item 右侧是引理 \ref{cohomology-lemma-unbounded-mayer-vietoris}
对 $K \otimes_{\mathcal{O}_X}^\mathbf{L} M$ 所给出的特异三角；
\item 左侧是将正合函子
$R\Gamma(X, K) \otimes_A^\mathbf{L} - $
应用于引理 \ref{cohomology-lemma-unbounded-mayer-vietoris} 对 $M$
所给出的特异三角。
\end{enumerate}
\end{lemma}

\begin{proof}
选取表示 $R\Gamma(X, K)$ 的平坦 $A$-模 K-平坦复形
$T^\bullet$；参见《更多代数》引理
\ref{more-algebra-lemma-K-flat-resolution}。
以 $T^\bullet \otimes_A \mathcal{O}_X$ 表示 $T^\bullet$
沿环空间态射 $(X, \mathcal{O}_X) \to (pt, A)$ 的拉回。
在 $D(\mathcal{O}_X)$ 中有自然的伴随映射
$\epsilon : T^\bullet \otimes_A \mathcal{O}_X \to K$。
注意 $T^\bullet \otimes_A \mathcal{O}_X$
是逐项平坦的 $\mathcal{O}_X$-模 K-平坦复形；参见引理
\ref{cohomology-lemma-pullback-K-flat} 与《模》引理
\ref{modules-lemma-pullback-flat}。
由引理 \ref{cohomology-lemma-factor-through-K-flat}，可找到复形态射
$$
T^\bullet \otimes_A \mathcal{O}_X \longrightarrow \mathcal{K}^\bullet
$$
，它是表示 $\epsilon$ 的 $\mathcal{O}_X$-模复形态射，
且 $\mathcal{K}^\bullet$ 是逐项平坦的 K-平坦复形。
具体地，由 $D(\mathcal{O}_X)$ 的构造，可先用某个表示
$\epsilon$ 的 $\mathcal{O}_X$-模复形态射
$e : T^\bullet \otimes_A \mathcal{O}_X \to \mathcal{L}^\bullet$
表示 $\epsilon$，再把该引理应用于 $e$。
选取表示 $M$ 的 K-内射复形 $\mathcal{I}^\bullet$，
其各项都是内射 $\mathcal{O}_X$-模。最后选取拟同构
$$
\text{Tot}(\mathcal{K}^\bullet \otimes_\mathcal{O} \mathcal{I}^\bullet)
\longrightarrow
\mathcal{J}^\bullet
$$
，其目标为逐项是内射 $\mathcal{O}_X$-模的 K-内射复形。
注意，此箭头的源与目标都在 $D(\mathcal{O}_X)$ 中表示
$K \otimes_{\mathcal{O}_X}^\mathbf{L} M$。
此时，对任意开集 $W \subset X$，得到复形态射
$$
\text{Tot}(T^\bullet \otimes_A \mathcal{I}^\bullet(W))
\to
\text{Tot}(\mathcal{K}^\bullet(W) \otimes_A \mathcal{I}^\bullet(W))
\to
\mathcal{J}^\bullet(W)
$$
；它是 $A$-模复形的态射，其复合表示映射
$$
R\Gamma(X, K) \otimes_A^\mathbf{L} R\Gamma(W, M)
\longrightarrow
R\Gamma(W, K \otimes_{\mathcal{O}_X}^\mathbf{L} M)
$$
，该映射位于 $D(A)$ 中。显然，这些映射与限制映射相容。
现在可以考察下列 $A$-模复形的交换图：
$$
\xymatrix{
0 \ar[d] & 0 \ar[d] \\
\text{Tot}(T^\bullet \otimes_A \mathcal{I}^\bullet(X)) \ar[d] \ar[r] &
\mathcal{J}^\bullet(X) \ar[d] \\
\text{Tot}(T^\bullet \otimes_A
(\mathcal{I}^\bullet(U) \oplus \mathcal{I}^\bullet(V)) \ar[d] \ar[r] &
\mathcal{J}^\bullet(U) \oplus \mathcal{J}^\bullet(V) \ar[d] \\
\text{Tot}(T^\bullet \otimes_A \mathcal{I}^\bullet(U \cap V)) \ar[r] \ar[d] &
\mathcal{J}^\bullet(U \cap V) \ar[d] \\
0 & 0
}
$$
由引理 \ref{cohomology-lemma-mayer-vietoris} 的证明，各列都是
$A$-模复形的正合列（这里还用到：
由于 $T^\bullet$ 的各项是平坦 $A$-模，
$\text{Tot}(T^\bullet \otimes_A -)$
把 $A$-模复形的短正合列变为短正合列）。引理 \ref{cohomology-lemma-unbounded-mayer-vietoris} 的特异三角
正是与这些复形短正合列相伴的特异三角。因此，所需结论由
《导出范畴》第 \ref{derived-section-canonical-delta-functor} 节所讨论的
“取相伴特异三角”这一操作的函子性得到。
\end{proof}












\section{支撑在闭子集上的上同调，II}
\label{cohomology-section-cohomology-support-bis}

\noindent
继续第 \ref{cohomology-section-cohomology-support} 节开始的讨论。

\medskip\noindent
设 $(X, \mathcal{O}_X)$ 是环空间，而 $Z \subset X$ 是闭子集。
在此情形中，可以考察函子
$\textit{Mod}(\mathcal{O}_X) \to \textit{Mod}(\mathcal{O}_X(X))$
，它由 $\mathcal{F} \mapsto \Gamma_Z(X, \mathcal{F})$ 给出；
参见《模》定义 \ref{modules-definition-support} 与引理
\ref{modules-lemma-support-section-closed}。
使用 K-内射消解（见第 \ref{cohomology-section-unbounded} 节），得到右导出函子
$$
R\Gamma_Z(X, - ) : D(\mathcal{O}_X) \to D(\mathcal{O}_X(X))
$$
给定 $D(\mathcal{O}_X)$ 中的对象 $K$，以
$H^q_Z(X, K) = H^q(R\Gamma_Z(X, K))$
表示支撑在 $Z$ 上的上同调模。稍后将看到
（引理 \ref{cohomology-lemma-sections-support-abelian-unbounded}），
它与第 \ref{cohomology-section-cohomology-support} 节的构造一致。

\medskip\noindent
对 $\mathcal{O}_X$-模 $\mathcal{F}$，可以考察{\it 支撑在 $Z$
上的截面子层}，记为 $\mathcal{H}_Z(\mathcal{F})$，其定义为
$$
\mathcal{H}_Z(\mathcal{F})(U) =
\{s \in \mathcal{F}(U) \mid \text{Supp}(s) \subset U \cap Z\} =
\Gamma_{Z \cap U}(U, \mathcal{F}|_U)
$$
如《模》注记
\ref{modules-remark-sections-support-in-closed-modules} 所述，
可把 $\mathcal{H}_Z(\mathcal{F})$ 看作 $Z$ 上的
$\mathcal{O}_X|_Z$-模，并得到函子
$$
\textit{Mod}(\mathcal{O}_X) \longrightarrow \textit{Mod}(\mathcal{O}_X|_Z),
\quad
\mathcal{F} \longmapsto \mathcal{H}_Z(\mathcal{F})
\text{（看作 }\mathcal{O}_X|_Z\text{-模，定义在 }Z\text{ 上）}
$$
此函子左正合，但一般并不正合。与上面完全相同，可得右导出函子
$$
R\mathcal{H}_Z : D(\mathcal{O}_X) \longrightarrow D(\mathcal{O}_X|_Z)
$$
令 $\mathcal{H}^q_Z(K) = H^q(R\mathcal{H}_Z(K))$，于是
$\mathcal{H}^0_Z(\mathcal{F}) = \mathcal{H}_Z(\mathcal{F})$
对任意 $\mathcal{O}_X$-模层 $\mathcal{F}$ 都成立。

\begin{lemma}
\label{cohomology-lemma-cohomology-with-support-sheaf-on-support}
设 $(X, \mathcal{O}_X)$ 是环化空间。设 $i : Z \to X$ 是闭子集的嵌入。
\begin{enumerate}
\item $R\mathcal{H}_Z : D(\mathcal{O}_X) \to D(\mathcal{O}_X|_Z)$
是 $i_* : D(\mathcal{O}_X|_Z) \to D(\mathcal{O}_X)$ 的右伴随。
\item 对 $D(\mathcal{O}_X|_Z)$ 中的 $K$，有
$R\mathcal{H}_Z(i_*K) = K$。
\item 设 $\mathcal{G}$ 是 $Z$ 上的 $\mathcal{O}_X|_Z$-模层。则当
$p > 0$ 时，$\mathcal{H}^p_Z(i_*\mathcal{G}) = 0$。
\end{enumerate}
\end{lemma}

\begin{proof}
函子 $i_*$ 正合，故 $i_* = Ri_* = Li_*$。因此，本引理的 (1) 由
《模层》引理 \ref{modules-lemma-adjoint-section-with-support}
与
《导出范畴》引理 \ref{derived-lemma-derived-adjoint-functors}
推出。设 $K$ 如 (2) 中所述。可用一个由 $\mathcal{O}_X|_Z$-模组成的
K-内射复形 $\mathcal{I}^\bullet$ 表示 $K$。由引理 \ref{cohomology-lemma-K-injective-flat}，
表示 $i_*K$ 的复形 $i_*\mathcal{I}^\bullet$
是由 $\mathcal{O}_X$-模组成的 K-内射复形。因此，
$R\mathcal{H}_Z(i_*K)$ 由
$\mathcal{H}_Z(i_*\mathcal{I}^\bullet) = \mathcal{I}^\bullet$
计算，这就证明了 (2)。(3) 是 (2) 的特例。
\end{proof}

\noindent
设 $(X, \mathcal{O}_X)$ 是环化空间，且 $Z \subset X$ 是闭子集。
支撑包含于 $Z$ 的 $\mathcal{O}_X$-模所成的范畴，是所有
$\mathcal{O}_X$-模所成范畴的一个 Serre 子范畴；参见
《同调代数》定义 \ref{homology-definition-serre-subcategory}
以及
《模层》引理 \ref{modules-lemma-support-section-closed}。
记 $D_Z(\mathcal{O}_X)$ 为 $D(\mathcal{O}_X)$ 的严格满的饱和三角子范畴，
其对象为上同调层支撑在 $Z$ 上的复形；参见
《导出范畴》第 \ref{derived-section-triangulated-sub} 节。

\begin{lemma}
\label{cohomology-lemma-complexes-with-support-on-closed}
设 $(X, \mathcal{O}_X)$ 是环化空间。设 $i : Z \to X$ 是闭子集的嵌入。
\begin{enumerate}
\item 对 $D(\mathcal{O}_X|_Z)$ 中的 $K$，有
$i_*K \in D_Z(\mathcal{O}_X)$。
\item 函子 $i_* : D(\mathcal{O}_X|_Z) \to D_Z(\mathcal{O}_X)$
是等价，其拟逆为
$i^{-1}|_{D_Z(\mathcal{O}_X)} = R\mathcal{H}_Z|_{D_Z(\mathcal{O}_X)}$。
\item 函子
$i_* \circ R\mathcal{H}_Z : D(\mathcal{O}_X) \to D_Z(\mathcal{O}_X)$
是嵌入函子
$D_Z(\mathcal{O}_X) \to D(\mathcal{O}_X)$ 的右伴随。
\end{enumerate}
\end{lemma}

\begin{proof}
(1) 由定义立即得到。(3) 是 (2) 与引理 \ref{cohomology-lemma-cohomology-with-support-sheaf-on-support}
的形式推论。以下证明 (2)。

\medskip\noindent
将 $i$ 看作环化空间的态射
$i : (Z, \mathcal{O}_X|_Z) \to (X, \mathcal{O}_X)$。
回忆 $i^*$ 与 $i_*$ 构成一对伴随函子。由于 $i$ 是闭浸入，
$i_*$ 正合。又因 $i^{-1}\mathcal{O}_X = \mathcal{O}_X|_Z$
是 $(Z, \mathcal{O}_X|_Z)$ 的结构层，可见 $i^* = i^{-1}$ 正合，
且 $i^*i_* = i^{-1}i_*$ 同构于恒等函子。参见
《模层》引理 \ref{modules-lemma-exactness-pushforward-pullback} 和
\ref{modules-lemma-i-star-exact}。因此
$i_* : D(\mathcal{O}_X|_Z) \to D_Z(\mathcal{O}_X)$
是全忠实的，而 $i^{-1}$ 给出其左逆。另一方面，设 $K$ 是
$D_Z(\mathcal{O}_X)$ 的对象，并考虑伴随态射
$K \to i_*i^{-1}K$。
由 $i_*$ 与 $i^{-1}$ 的正合性，它在上同调层上诱导伴随态射
$H^n(K) \to i_*i^{-1}H^n(K)$。由于这些上同调层支撑在 $Z$ 上，
这些伴随态射都是同构，从而
$i_* : D(\mathcal{O}_X|_Z) \to D_Z(\mathcal{O}_X)$ 是等价。

\medskip\noindent
最后，只需证明：若 $K$ 是 $D_Z(\mathcal{O}_X)$ 的对象，则
$R\mathcal{H}_Z(K) = i^{-1}K$。为此，可使用刚刚证明的
$K = i_*i^{-1}K$。于是引理 \ref{cohomology-lemma-cohomology-with-support-sheaf-on-support}
给出所需结论。
\end{proof}

\begin{lemma}
\label{cohomology-lemma-sections-with-support-K-injective}
设 $(X, \mathcal{O}_X)$ 是环化空间。设 $i : Z \to X$ 是闭子集的嵌入。
若 $\mathcal{I}^\bullet$ 是由 $\mathcal{O}_X$-模组成的 K-内射复形，
则 $\mathcal{H}_Z(\mathcal{I}^\bullet)$
是由 $\mathcal{O}_X|_Z$-模组成的 K-内射复形。
\end{lemma}

\begin{proof}
由于 $i_* : \textit{Mod}(\mathcal{O}_X|_Z) \to
\textit{Mod}(\mathcal{O}_X)$ 正合，并且是 $\mathcal{H}_Z$ 的左伴随
（《模层》引理 \ref{modules-lemma-adjoint-section-with-support}），
结论由
《导出范畴》引理 \ref{derived-lemma-adjoint-preserve-K-injectives}
推出。
\end{proof}

\begin{lemma}
\label{cohomology-lemma-local-to-global-sections-with-support}
设 $(X, \mathcal{O}_X)$ 是环化空间。设 $i : Z \to X$ 是闭子集的嵌入。则
$R\Gamma(Z, - ) \circ R\mathcal{H}_Z = R\Gamma_Z(X, - )$
作为从 $D(\mathcal{O}_X)$ 到 $D(\mathcal{O}_X(X))$ 的函子成立。
\end{lemma}

\begin{proof}
这由利用 K-内射消解构造右导出函子、引理 \ref{cohomology-lemma-sections-with-support-K-injective}
以及等式 $\Gamma_Z(X, -) = \Gamma(Z, -) \circ \mathcal{H}_Z$ 得到。
\end{proof}

\begin{lemma}
\label{cohomology-lemma-triangle-sections-with-support}
设 $(X, \mathcal{O}_X)$ 是环化空间。设 $i : Z \to X$ 是闭子集的嵌入。
令 $U = X \setminus Z$。存在特异三角
$$
R\Gamma_Z(X, K) \to R\Gamma(X, K) \to R\Gamma(U, K) \to
R\Gamma_Z(X, K)[1]
$$
位于 $D(\mathcal{O}_X(X))$ 中，并关于 $D(\mathcal{O}_X)$ 中的
$K$ 具有函子性。
\end{lemma}

\begin{proof}
选择一个表示 $K$ 的 K-内射复形 $\mathcal{I}^\bullet$，其每一项都是
内射 $\mathcal{O}_X$-模；参见 第 \ref{cohomology-section-unbounded} 节。
回忆 $\mathcal{I}^\bullet|_U$ 是由 $\mathcal{O}_U$-模组成的
K-内射复形；参见引理 \ref{cohomology-lemma-restrict-K-injective-to-open}。
因此，特异三角中的每个导出函子，都可通过将相应原函子作用于
$\mathcal{I}^\bullet$ 来得到。故只需证明：对任意内射
$\mathcal{O}_X$-模 $\mathcal{I}$，都有短正合列
$$
0 \to \Gamma_Z(X, \mathcal{I}) \to \Gamma(X, \mathcal{I})
\to \Gamma(U, \mathcal{I}) \to 0
$$
这由引理 \ref{cohomology-lemma-injective-restriction-surjective} 和定义得到。
\end{proof}

\begin{lemma}
\label{cohomology-lemma-triangle-sections-with-support-sheaves}
设 $(X, \mathcal{O}_X)$ 是环化空间。设 $i : Z \to X$ 是闭子集的嵌入。
记 $j : U = X \setminus Z \to X$ 为其补集的嵌入。存在特异三角
$$
i_*R\mathcal{H}_Z(K) \to K \to Rj_*(K|_U) \to
i_*R\mathcal{H}_Z(K)[1]
$$
位于 $D(\mathcal{O}_X)$ 中，并关于 $D(\mathcal{O}_X)$ 中的
$K$ 具有函子性。
\end{lemma}

\begin{proof}
选择一个表示 $K$ 的 K-内射复形 $\mathcal{I}^\bullet$，其每一项都是
内射 $\mathcal{O}_X$-模；参见 第 \ref{cohomology-section-unbounded} 节。
回忆 $\mathcal{I}^\bullet|_U$ 是由 $\mathcal{O}_U$-模组成的
K-内射复形；参见引理 \ref{cohomology-lemma-restrict-K-injective-to-open}。
因此，特异三角中的每个导出函子，都可通过将相应原函子作用于
$\mathcal{I}^\bullet$ 来得到。故只需证明：对任意内射
$\mathcal{O}_X$-模 $\mathcal{I}$，都有短正合列
$$
0 \to i_*\mathcal{H}_Z(\mathcal{I}) \to \mathcal{I}
\to j_*(\mathcal{I}|_U) \to 0
$$
这由引理 \ref{cohomology-lemma-injective-restriction-surjective} 和定义得到。
\end{proof}

\begin{lemma}
\label{cohomology-lemma-sections-support-in-closed-disjoint-open}
设 $(X, \mathcal{O}_X)$ 是环化空间。设 $Z \subset X$ 是闭子集。
设 $j : U \to X$ 是满足 $U \cap Z = \emptyset$ 的开子集的嵌入。
则对 $D(\mathcal{O}_U)$ 中的所有 $K$，有
$R\mathcal{H}_Z(Rj_*K) = 0$。
\end{lemma}

\begin{proof}
选择一个由 $\mathcal{O}_U$-模组成且表示 $K$ 的 K-内射复形
$\mathcal{I}^\bullet$。于是 $j_*\mathcal{I}^\bullet$ 表示 $Rj_*K$。
由引理 \ref{cohomology-lemma-K-injective-flat}，复形
$j_*\mathcal{I}^\bullet$ 是由 $\mathcal{O}_X$-模组成的
K-内射复形。因此，$\mathcal{H}_Z(j_*\mathcal{I}^\bullet)$
表示 $R\mathcal{H}_Z(Rj_*K)$。故只需证明：对 $U$ 上任意阿贝尔层
$\mathcal{G}$，都有 $\mathcal{H}_Z(j_*\mathcal{G}) = 0$。
换言之，必须证明：若 $j_*\mathcal{G}$ 在某个开集 $W$ 上的截面
$s$ 支撑在 $W \cap Z$ 上，则 $s$ 为零。支撑条件意味着
$s|_{W \setminus W \cap Z} = 0$。由于
$j_*\mathcal{G}(W) =
\mathcal{G}(U \cap W) = j_*\mathcal{G}(W \setminus W \cap Z)$
，故 $s$ 确如所需为零。
\end{proof}

\begin{lemma}
\label{cohomology-lemma-sections-support-abelian-unbounded}
设 $(X, \mathcal{O}_X)$ 是环化空间。设 $Z \subset X$ 是闭子集。
设 $K$ 是 $D(\mathcal{O}_X)$ 的对象，并记 $K_{ab}$ 为其在
$D(\underline{\mathbf{Z}}_X)$ 中的像。
\begin{enumerate}
\item 存在典范态射 $R\Gamma_Z(X, K) \to R\Gamma_Z(X, K_{ab})$，
它是 $D(\textit{Ab})$ 中的同构。
\item 存在典范态射
$R\mathcal{H}_Z(K) \to R\mathcal{H}_Z(K_{ab})$
，它是 $D(\underline{\mathbf{Z}}_Z)$ 中的同构。
\end{enumerate}
\end{lemma}

\begin{proof}
证明 (1)。此态射构造如下。选择一个由 $\mathcal{O}_X$-模组成、表示
$K$ 的 K-内射复形 $\mathcal{I}^\bullet$。选择拟同构
$\mathcal{I}^\bullet \to \mathcal{J}^\bullet$，其中
$\mathcal{J}^\bullet$ 是由阿贝尔群组成的 K-内射复形。则 (1) 中的态射为
$$
\Gamma_Z(X, \mathcal{I}^\bullet) \to \Gamma_Z(X, \mathcal{J}^\bullet)
$$
；它由 $\Gamma_Z$ 是阿贝尔层上的函子这一事实确定。容易验证，
所得态射与引理 \ref{cohomology-lemma-modules-abelian-unbounded} 的典范态射
共同构成如下特异三角的态射：
$$
\xymatrix{
R\Gamma_Z(X, K) \ar[r] \ar[d] &
R\Gamma(X, K) \ar[r] \ar[d] &
R\Gamma(U, K) \ar[d] \\
R\Gamma_Z(X, K_{ab}) \ar[r] &
R\Gamma(X, K_{ab}) \ar[r] &
R\Gamma(U, K_{ab})
}
$$
这里使用引理 \ref{cohomology-lemma-triangle-sections-with-support}。
由所引引理，三个箭头中的两个是同构，故可由 《导出范畴》引理
\ref{derived-lemma-third-isomorphism-triangle}.
得出结论。

\medskip\noindent
省略 (2) 的证明。提示：将同一论证用于引理 \ref{cohomology-lemma-triangle-sections-with-support-sheaves}
的特异三角。
\end{proof}

\begin{remark}
\label{cohomology-remark-support-cup-product}
设 $(X, \mathcal{O}_X)$ 是环化空间。设 $i : Z \to X$ 是闭子集的嵌入。
给定 $D(\mathcal{O}_X)$ 中的 $K$ 与 $M$，存在典范态射
$$
K|_Z \otimes_{\mathcal{O}_X|_Z}^\mathbf{L} R\mathcal{H}_Z(M)
\longrightarrow
R\mathcal{H}_Z(K \otimes_{\mathcal{O}_X}^\mathbf{L} M)
$$
位于 $D(\mathcal{O}_X|_Z)$ 中。这里 $K|_Z = i^{-1}K$ 是 $K$ 在
$Z$ 上的限制，并看作 $D(\mathcal{O}_X|_Z)$ 的对象。由引理 \ref{cohomology-lemma-cohomology-with-support-sheaf-on-support}
中的 $i_*$ 与 $R\mathcal{H}_Z$ 的伴随性，为构造此态射，只需给出典范态射
$$
i_*\left(K|_Z \otimes_{\mathcal{O}_X|_Z}^\mathbf{L} R\mathcal{H}_Z(M)\right)
\longrightarrow
K \otimes_{\mathcal{O}_X}^\mathbf{L} M
$$
为构造此态射，选择一个表示 $M$、由 $\mathcal{O}_X$-模组成的
K-内射复形 $\mathcal{I}^\bullet$，以及一个表示 $K$、由
$\mathcal{O}_X$-模组成的 K-平坦复形 $\mathcal{K}^\bullet$。
注意，$\mathcal{K}^\bullet|_Z$ 是由 $\mathcal{O}_X|_Z$-模组成、
表示 $K|_Z$ 的 K-平坦复形；参见引理 \ref{cohomology-lemma-pullback-K-flat}。因此，需要构造复形态射
$$
i_*\text{Tot}\left(
\mathcal{K}^\bullet|_Z \otimes_{\mathcal{O}_X|_Z}
\mathcal{H}_Z(\mathcal{I}^\bullet)\right)
\longrightarrow
\text{Tot}(\mathcal{K}^\bullet \otimes_{\mathcal{O}_X} \mathcal{I}^\bullet)
$$
，它们由 $\mathcal{O}_X$-模组成。为此，只需给出态射
$$
i_*(\mathcal{K}^a|_Z \otimes_{\mathcal{O}_X|_Z}
\mathcal{H}_Z(\mathcal{I}^b))
\longrightarrow
\mathcal{K}^a \otimes_{\mathcal{O}_X} \mathcal{I}^b
$$
例如考察茎，可见此式左端等于
$\mathcal{K}^a \otimes_{\mathcal{O}_X} i_*\mathcal{H}_Z(\mathcal{I}^b)$
，再利用嵌入 $\mathcal{H}_Z(\mathcal{I}^b) \to \mathcal{I}^b$
即可得到所需态射。
\end{remark}

\begin{remark}
\label{cohomology-remark-support-cup-product-global}
采用注记 \ref{cohomology-remark-support-cup-product} 中的记号，可得到典范杯积
\begin{align*}
H^a(X, K) \times H^b_Z(X, M)
& =
H^a(X, K) \times H^b(Z, R\mathcal{H}_Z(M)) \\
& \to
H^a(Z, K|_Z) \times H^b(Z, R\mathcal{H}_Z(M)) \\
& \to
H^{a + b}(Z, K|_Z \otimes_{\mathcal{O}_X|_Z}^\mathbf{L} R\mathcal{H}_Z(M)) \\
& \to
H^{a + b}(Z, R\mathcal{H}_Z(K \otimes_{\mathcal{O}_X}^\mathbf{L} M)) \\
& =
H^{a + b}_Z(X, K \otimes_{\mathcal{O}_X}^\mathbf{L} M)
\end{align*}
这里的等号由引理 \ref{cohomology-lemma-local-to-global-sections-with-support}
给出；第一个箭头是在 $Z$ 上的限制，第二个箭头是杯积
（第 \ref{cohomology-section-cup-product} 节），第三个箭头则是注记 \ref{cohomology-remark-support-cup-product} 中的态射。
\end{remark}

\begin{lemma}
\label{cohomology-lemma-support-cup-product}
采用注记 \ref{cohomology-remark-support-cup-product} 中的记号，下图交换：
$$
\xymatrix{
H^i(X, K) \times H^j_Z(X, M) \ar[r] \ar[d] &
H^{i + j}_Z(X, K \otimes_{\mathcal{O}_X}^\mathbf{L} M) \ar[d] \\
H^i(X, K) \times H^j(X, M) \ar[r] &
H^{i + j}(X, K \otimes_{\mathcal{O}_X}^\mathbf{L} M)
}
$$
其中上方水平箭头是注记 \ref{cohomology-remark-support-cup-product-global}
中的杯积。
\end{lemma}

\begin{proof}
从略。
\end{proof}

\begin{remark}
\label{cohomology-remark-support-functorial}
设 $f : (X', \mathcal{O}_{X'}) \to (X, \mathcal{O}_X)$ 是环化空间的态射。
设 $Z \subset X$ 是闭子集，并令 $Z' = f^{-1}(Z)$。记
$f|_{Z'} : (Z', \mathcal{O}_{X'}|_{Z'}) \to (Z, \mathcal{O}_X|Z)$
为诱导的环化空间态射。对 $D(\mathcal{O}_X)$ 中任意 $K$，存在典范态射
$$
L(f|_{Z'})^*R\mathcal{H}_Z(K) \longrightarrow R\mathcal{H}_{Z'}(Lf^*K)
$$
位于 $D(\mathcal{O}_{X'}|_{Z'})$ 中。记 $i : Z \to X$ 与
$i' : Z' \to X'$ 为嵌入态射。将引理 \ref{cohomology-lemma-complexes-with-support-on-closed} 的 (2) 用于 $i'$，
则给出上述态射等价于给出态射
$$
i'_* L(f|_{Z'})^* R\mathcal{H}_Z(K)
\longrightarrow
i'_*R\mathcal{H}_{Z'}(Lf^*K)
$$
位于 $D_{Z'}(\mathcal{O}_{X'})$ 中。注记 \ref{cohomology-remark-base-change}
中的函子态射
$Lf^* \circ i_* \to i'_* \circ L(f|_{Z'})^*$
在此情形下是同构（可在茎上检验；这里使用 $i_*$ 与 $i'_*$ 的正合性、
$\mathcal{O}_{Z, z} = \mathcal{O}_{X, z}$，以及 $Z'$ 的类似等式）。
因此，只需构造下图上方的水平箭头
$$
\xymatrix{
Lf^* i_* R\mathcal{H}_Z(K) \ar[rr] \ar[rd] & &
i'_* R\mathcal{H}_{Z'}(Lf^*K) \ar[ld] \\
& Lf^*K
}
$$
复形 $Lf^* i_* R\mathcal{H}_Z(K)$ 支撑在 $Z'$ 上。东南方向的箭头来自
伴随态射 $i_*R\mathcal{H}_Z(K) \to K$
（引理 \ref{cohomology-lemma-cohomology-with-support-sheaf-on-support}）。
由引理 \ref{cohomology-lemma-complexes-with-support-on-closed} 的 (3)，伴随态射
$i'_* R\mathcal{H}_{Z'}(Lf^*K) \to Lf^*K$ 具有泛性质。因此，
东南方向的箭头唯一地经由西南方向的箭头分解，从而得到所需箭头。
\end{remark}

\begin{lemma}
\label{cohomology-lemma-support-functorial}
采用注记 \ref{cohomology-remark-support-functorial} 中的记号与假设，下图交换：
$$
\xymatrix{
H^p_Z(X, K) \ar[r] \ar[d] & H^p_{Z'}(X', Lf^*K) \ar[d] \\
H^p(X, K) \ar[r] & H^p(X', Lf^*K)
}
$$
其中上方水平箭头来自如下等同：
$H^p_Z(X, K) = H^p(Z, R\mathcal{H}_Z(K))$ 与
$H^p_{Z'}(X', Lf^*K) = H^p(Z', R\mathcal{H}_{Z'}(K'))$、拉回态射
$H^p(Z, R\mathcal{H}_Z(K)) \to H^p(Z', L(f|_{Z'})^*R\mathcal{H}_Z(K))$,
以及注记 \ref{cohomology-remark-support-functorial} 中构造的态射。
\end{lemma}

\begin{proof}
从略。提示：利用
$H^p(Z, R\mathcal{H}_Z(K)) = H^p(X, i_*R\mathcal{H}_Z(K))$
以及 $R\mathcal{H}_{Z'}(Lf^*K)$ 的类似等式，结论可由拉回态射的函子性
和定义注记 \ref{cohomology-remark-support-functorial} 中态射时所用的交换图得到。
\end{proof}






\section{逆系统与上同调，I}
\label{cohomology-section-inverse-systems}

\noindent
设 $A$ 是环，且 $I \subset A$ 是理想。我们证明若干关于
$A/I^n$-模层的逆系统的结果。

\begin{lemma}
\label{cohomology-lemma-ML-general}
设 $I$ 是环 $A$ 的理想。设 $X$ 是拓扑空间。设
$$
\ldots \to \mathcal{F}_3 \to \mathcal{F}_2 \to \mathcal{F}_1
$$
是 $X$ 上 $A$-模层的逆系统，并满足
$\mathcal{F}_n = \mathcal{F}_{n + 1}/I^n\mathcal{F}_{n + 1}$。
设 $p \geq 0$。假设
$$
\bigoplus\nolimits_{n \geq 0} H^{p + 1}(X, I^n\mathcal{F}_{n + 1})
$$
作为分次 $\bigoplus_{n \geq 0} I^n/I^{n + 1}$-模满足升链条件。
则逆系统 $M_n = H^p(X, \mathcal{F}_n)$ 满足 Mittag--Leffler
条件\footnote{事实上，存在 $c \geq 0$，使得对所有 $n \geq c$，
$\Im(M_n \to M_{n - c})$ 都是稳定像。}。
\end{lemma}

\begin{proof}
令 $N_n = H^{p + 1}(X, I^n\mathcal{F}_{n + 1})$，并令
$\delta_n : M_n \to N_n$ 为短正合列
$0 \to I^n\mathcal{F}_{n + 1} \to \mathcal{F}_{n + 1}
\to \mathcal{F}_n \to 0$
在上同调上诱导的边界态射。则
$\bigoplus \Im(\delta_n) \subset \bigoplus N_n$ 是分次子模。
事实上，若 $s \in M_n$ 且 $f \in I^m$，则有交换图
$$
\xymatrix{
0 \ar[r] &
I^n\mathcal{F}_{n + 1} \ar[d]_f \ar[r] &
\mathcal{F}_{n + 1} \ar[d]_f \ar[r] &
\mathcal{F}_n \ar[d]_f \ar[r] & 0 \\
0 \ar[r] &
I^{n + m}\mathcal{F}_{n + m + 1} \ar[r] &
\mathcal{F}_{n + m + 1} \ar[r] &
\mathcal{F}_{n + m} \ar[r] & 0
}
$$
中间的竖直态射由如下方式给出：先将 $\mathcal{F}_{n + 1}$ 的局部截面
提升为 $\mathcal{F}_{n + m + 1}$ 的截面，再乘以 $f$；其余竖直箭头亦然。
由此得到 $\delta_{n + m}(fs) = f \delta_n(s)$。根据假设，可找到
$s_j \in M_{n_j}$，$j = 1, \ldots, N$，使得
$\delta_{n_j}(s_j)$ 作为分次模生成
$\bigoplus \Im(\delta_n)$。设 $n > c = \max(n_j)$，并取 $s \in M_n$。
则可找到 $f_j \in I^{n - n_j}$，使得
$\delta_n(s) = \sum f_j \delta_{n_j}(s_j)$。因此
$\delta(s - \sum f_j s_j) = 0$；换言之，可找到 $s' \in M_{n + 1}$，
其在 $M_n$ 中的像为 $s - \sum f_js_j$。于是
$$
\Im(M_{n + 1} \to M_{n - c}) = \Im(M_n \to M_{n - c})
$$
这是因为各元素 $f_js_j$ 在 $M_{n - c}$ 中的像为零。引理得证。
\end{proof}

\begin{lemma}
\label{cohomology-lemma-ML-general-better}
设 $I$ 是环 $A$ 的理想。设 $X$ 是拓扑空间。设
$$
\ldots \to \mathcal{F}_3 \to \mathcal{F}_2 \to \mathcal{F}_1
$$
是 $X$ 上 $A$-模的逆系统，并满足
$\mathcal{F}_n = \mathcal{F}_{n + 1}/I^n\mathcal{F}_{n + 1}$。
设 $p \geq 0$。对给定的 $n$，定义
$$
N_n =
\bigcap\nolimits_{m \geq n}
\Im\left(
H^{p + 1}(X, I^n\mathcal{F}_{m + 1}) \to H^{p + 1}(X, I^n\mathcal{F}_{n + 1})
\right)
$$
若 $\bigoplus N_n$ 作为分次
$\bigoplus_{n \geq 0} I^n/I^{n + 1}$-模满足升链条件，则逆系统
$M_n = H^p(X, \mathcal{F}_n)$ 满足 Mittag--Leffler
条件\footnote{事实上，存在 $c \geq 0$，使得对所有 $n \geq c$，
$\Im(M_n \to M_{n - c})$ 都是稳定像。}。
\end{lemma}

\begin{proof}
证明与引理 \ref{cohomology-lemma-ML-general} 的证明完全相同。事实上，只要证明
$\bigoplus N_n$ 是
$\bigoplus H^{p + 1}(X, I^n\mathcal{F}_{n + 1})$
的分次 $\bigoplus_{n \geq 0} I^n/I^{n + 1}$-子模，并且边界态射
$\delta_n : M_n \to H^{p + 1}(X, I^n\mathcal{F}_{n + 1})$
的像包含于 $N_n$，结论便由其中的论证推出。

\medskip\noindent
设 $\xi \in N_n$ 且 $f \in I^k$。选择 $m \gg n + k$，并选择
$\xi' \in H^{p + 1}(X, I^n\mathcal{F}_{m + 1})$ 作为 $\xi$ 的提升。
考虑交换图
$$
\xymatrix{
0 \ar[r] &
I^n\mathcal{F}_{m + 1} \ar[d]_f \ar[r] &
\mathcal{F}_{m + 1} \ar[d]_f \ar[r] &
\mathcal{F}_n \ar[d]_f \ar[r] & 0 \\
0 \ar[r] &
I^{n + k}\mathcal{F}_{m + 1} \ar[r] &
\mathcal{F}_{m + 1} \ar[r] &
\mathcal{F}_{n + k} \ar[r] & 0
}
$$
其构造如引理 \ref{cohomology-lemma-ML-general} 的证明所示。它在上同调上诱导态射，
并可见 $f \xi' \in H^{p + 1}(X, I^{n + k}\mathcal{F}_{m + 1})$
映到 $f \xi$。由于这对所有 $m \gg n + k$ 都成立，故如所需，
$f\xi$ 属于 $N_{n + k}$。

\medskip\noindent
为证明边界态射 $\delta_n$ 的像包含于 $N_n$，考虑各图
$$
\xymatrix{
0 \ar[r] &
I^n\mathcal{F}_{m + 1} \ar[d] \ar[r] &
\mathcal{F}_{m + 1} \ar[d] \ar[r] &
\mathcal{F}_n \ar[d] \ar[r] & 0 \\
0 \ar[r] &
I^n\mathcal{F}_{n + 1} \ar[r] &
\mathcal{F}_{n + 1} \ar[r] &
\mathcal{F}_n \ar[r] & 0
}
$$
其中 $m \geq n$。考察诱导的上同调态射即可得出结论。
\end{proof}

\begin{lemma}
\label{cohomology-lemma-topology-I-adic-general}
设 $I$ 是环 $A$ 的理想。设 $X$ 是拓扑空间。设
$$
\ldots \to \mathcal{F}_3 \to \mathcal{F}_2 \to \mathcal{F}_1
$$
是 $X$ 上 $A$-模层的逆系统，并满足
$\mathcal{F}_n = \mathcal{F}_{n + 1}/I^n\mathcal{F}_{n + 1}$。
设 $p \geq 0$。假设
$$
\bigoplus\nolimits_{n \geq 0} H^p(X, I^n\mathcal{F}_{n + 1})
$$
作为分次 $\bigoplus_{n \geq 0} I^n/I^{n + 1}$-模满足升链条件。
则 $M = \lim H^p(X, \mathcal{F}_n)$ 上的极限拓扑就是 $I$-进拓扑。
\end{lemma}

\begin{proof}
对 $n \geq 1$，令 $F^n = \Ker(M \to H^p(X, \mathcal{F}_n))$，并令
$F^0 = M$。注意 $I F^n \subset F^{n + 1}$；特别地，
$I^n M \subset F^n$。因此，$I$-进拓扑比极限拓扑更细。反过来，
我们将证明：给定 $n$，存在 $m \geq n$，使得
$F^m \subset I^nM$\footnote{事实上，存在 $c \geq 0$，使得对所有
$n$，都有 $F^{n + c} \subset I^nM$。}。存在单射
$$
F^n/F^{n + 1} \longrightarrow H^p(X, \mathcal{F}_{n + 1})
$$
其像包含于态射
$H^p(X, I^n\mathcal{F}_{n + 1}) \to H^p(X, \mathcal{F}_{n + 1})$
的像中。记
$$
E_n \subset H^p(X, I^n\mathcal{F}_{n + 1})
$$
为 $F^n/F^{n + 1}$ 的逆像。则 $\bigoplus E_n$ 是
$\bigoplus H^p(X, I^n\mathcal{F}_{n + 1})$ 的分次
$\bigoplus I^n/I^{n + 1}$-子模，而
$\bigoplus E_n \to \bigoplus F^n/F^{n + 1}$ 是分次模同态；
细节从略。根据假设，$\bigoplus E_n$ 由有限多个齐次元素作为
$\bigoplus I^n/I^{n + 1}$-模生成。由于
$E_n \to F^n/F^{n + 1}$ 满射，可见
$\bigoplus F^n/F^{n + 1}$ 亦具有相同性质。因此，可找到 $r$、
$c_1, \ldots, c_r \geq 0$ 以及 $a_i \in F^{c_i}$，使其在
$\bigoplus F^n/F^{n + 1}$ 中的像生成该模。令 $c = \max(c_i)$。

\medskip\noindent
对 $n \geq c$，我们断言 $I F^n = F^{n + 1}$。由此断言可得
$F^{n + c} = I^nF^c \subset I^nM$，正合所需。为证明断言，设
$a \in F^{n + 1}$。$a$ 在 $F^{n + 1}/F^{n + 2}$ 中的像是各
$a_i$ 的线性组合。因此，对某些 $f_i \in I^{n + 1 - c_i}$，有
$a  - \sum f_i a_i \in F^{n + 2}$。由于 $n \geq c_i$ 时
$I^{n + 1 - c_i} = I \cdot I^{n - c_i}$，可写
$f_i = \sum g_{i, j} h_{i, j}$，其中 $g_{i, j} \in I$ 且
$h_{i, j}a_i \in F^n$。于是
$F^{n + 1} = F^{n + 2} + IF^n$.
简单归纳即得：对所有 $e > 0$，
$F^{n + 1} = F^{n + e} + IF^n$。因此，$IF^n$ 在
$F^{n + 1}$ 中稠密。选择 $I$ 的生成元
$k_1, \ldots, k_r$，并考虑连续态射
$$
u : (F^n)^{\oplus r} \longrightarrow F^{n + 1},\quad
(x_1, \ldots, x_r) \mapsto \sum k_i x_i
$$
（取极限拓扑）。由上可知，对所有 $m \geq n$，
$(F^m)^{\oplus r}$ 在 $u$ 下的像都在 $F^{m + 1}$ 中稠密。
由开映射引理（《更多代数》引理
\ref{more-algebra-lemma-open-mapping}），$u$ 是开映射，因而是满射。
故当 $n \geq c$ 时，$IF^n = F^{n + 1}$。证明完毕。
\end{proof}

\begin{lemma}
\label{cohomology-lemma-topology-I-adic-general-better}
设 $I$ 是环 $A$ 的理想。设 $X$ 是拓扑空间。设
$$
\ldots \to \mathcal{F}_3 \to \mathcal{F}_2 \to \mathcal{F}_1
$$
是 $X$ 上 $A$-模层的逆系统，并满足
$\mathcal{F}_n = \mathcal{F}_{n + 1}/I^n\mathcal{F}_{n + 1}$。
设 $p \geq 0$。对给定的 $n$，定义
$$
N_n =
\bigcap\nolimits_{m \geq n}
\Im\left(
H^p(X, I^n\mathcal{F}_{m + 1}) \to H^p(X, I^n\mathcal{F}_{n + 1})
\right)
$$
若 $\bigoplus N_n$ 作为分次
$\bigoplus_{n \geq 0} I^n/I^{n + 1}$-模满足升链条件，则
$M = \lim H^p(X, \mathcal{F}_n)$ 上的极限拓扑就是 $I$-进拓扑。
\end{lemma}

\begin{proof}
证明与引理 \ref{cohomology-lemma-topology-I-adic-general} 的证明完全相同。
事实上，只要证明 $\bigoplus N_n$ 是
$\bigoplus H^{p + 1}(X, I^n\mathcal{F}_{n + 1})$
的分次 $\bigoplus_{n \geq 0} I^n/I^{n + 1}$-子模，并且
$F^n/F^{n + 1} \subset H^p(X, \mathcal{F}_{n + 1})$
包含于 $N_n \to H^p(X, \mathcal{F}_{n + 1})$ 的像中，
结论便由其中的论证推出。在引理 \ref{cohomology-lemma-ML-general-better}
的证明中，我们已经看到了关于模结构的陈述。

\medskip\noindent
设 $t \in F^n$。选择元素
$s \in H^p(X, I^n\mathcal{F}_{n + 1})$，使其映到 $t$ 在
$H^p(X, \mathcal{F}_{n + 1})$ 中的像。我们必须证明 $s \in N_n$。
由于 $F^n$ 是态射 $M \to H^p(X, \mathcal{F}_n)$ 的核，故对所有
$m \geq n$，可将 $t$ 映到元素
$t_m \in H^p(X, \mathcal{F}_{m + 1})$，且它在
$H^p(X, \mathcal{F}_n)$ 中的像为零。考虑上同调序列
$$
H^{p - 1}(X, \mathcal{F}_n) \to
H^p(X, I^n\mathcal{F}_{m + 1}) \to
H^p(X, \mathcal{F}_{m + 1}) \to
H^p(X, \mathcal{F}_n)
$$
它由短正合列
$0 \to I^n\mathcal{F}_{m + 1} \to \mathcal{F}_{m + 1}
\to \mathcal{F}_n \to 0$
诱导。
可选择 $s_m \in H^p(X, I^n\mathcal{F}_{m + 1})$，使其映到 $t_m$。
将上述序列与 $m = n$ 时的序列比较，可见 $s_m$ 的像与 $s$
相差一个属于如下态射之像的元素：
$H^{p - 1}(X, \mathcal{F}_n) \to
H^p(X, I^n\mathcal{F}_{n + 1})$。
然而，此态射经由态射
$H^p(X, I^n\mathcal{F}_{m + 1}) \to H^p(X, I^n\mathcal{F}_{n + 1})$
分解，因而可见 $s$ 如所需属于其像。
\end{proof}





\section{逆系统与上同调，II}
\label{cohomology-section-inverse-systems-bis}

\noindent
本节在理想为主理想的情形下，继续
第 \ref{cohomology-section-inverse-systems} 节 中的讨论。

\begin{lemma}
\label{cohomology-lemma-equivalent-f-good}
设 $(X, \mathcal{O}_X)$ 是环化空间。设
$f \in \Gamma(X, \mathcal{O}_X)$。设
$$
\ldots \to \mathcal{F}_3 \to \mathcal{F}_2 \to \mathcal{F}_1
$$
是 $\mathcal{O}_X$-模的逆系统。考虑下列条件：
\begin{enumerate}
\item 对所有 $n \geq 1$，态射
$f : \mathcal{F}_{n + 1} \to \mathcal{F}_{n + 1}$
经由 $\mathcal{F}_{n + 1} \to \mathcal{F}_n$ 分解，并给出短正合列
$0 \to \mathcal{F}_n \to \mathcal{F}_{n + 1} \to \mathcal{F}_1 \to 0$,
\item 对所有 $n \geq 1$，态射
$f^n : \mathcal{F}_{n + 1} \to \mathcal{F}_{n + 1}$
经由 $\mathcal{F}_{n + 1} \to \mathcal{F}_1$ 分解，并给出短正合列
$0 \to \mathcal{F}_1 \to \mathcal{F}_{n + 1} \to \mathcal{F}_n \to 0$
\item 存在 $f$-可除的 $\mathcal{O}_X$-模 $\mathcal{G}$，使得
$\mathcal{F}_n = \mathcal{G}[f^n]$；并且
\item 存在无 $f$-挠的 $\mathcal{O}_X$-模 $\mathcal{F}$，使得
$\mathcal{F}_n = \mathcal{F}/f^n\mathcal{F}$.
\end{enumerate}
则 (4) $\Rightarrow$ (3) $\Leftrightarrow$ (2) $\Leftrightarrow$ (1)。
\end{lemma}

\begin{proof}
省略 (1) 与 (2) 等价的证明，也省略 (3) 推出 (1) 的证明。
给定 (1) 中的 $\mathcal{F}_n$，为证明 (3)，令
$\mathcal{G} = \colim \mathcal{F}_n$，其中各态射
$\mathcal{F}_1 \to \mathcal{F}_2 \to \mathcal{F}_3 \to \ldots$
如 (1) 中所述。态射 $f : \mathcal{G} \to \mathcal{G}$ 满射；
这是因为由 (1) 的短正合列，
$\mathcal{F}_{n + 1} \subset \mathcal{G}$ 的像是
$\mathcal{F}_n \subset \mathcal{G}$。因此，$\mathcal{G}$ 是
$f$-可除的 $\mathcal{O}_X$-模，并且
$\mathcal{F}_n = \mathcal{G}[f^n]$。

\medskip\noindent
设给定 (4) 中的 $\mathcal{F}$。态射
$\mathcal{F}/f^{n + 1}\mathcal{F} \to \mathcal{F}/f^n\mathcal{F}$
总是满射，其核是态射
$\mathcal{F}/f\mathcal{F} \to \mathcal{F}/f^{n + 1}\mathcal{F}$
的像；该态射由乘以 $f^n$ 诱导。为验证 (2)，只需说明态射
$f^n : \mathcal{F} \to \mathcal{F}/f^{n + 1}\mathcal{F}$
的核为 $f\mathcal{F}$。为此，只需证明：给定开集 $U \subset X$
上 $\mathcal{F}$ 的截面 $s,t$，若 $f^ns = f^{n + 1}t$，则
$s = ft$。这很清楚，因为 $\mathcal{F}$ 无 $f$-挠，故
$f : \mathcal{F} \to \mathcal{F}$ 是单射。
\end{proof}

\begin{lemma}
\label{cohomology-lemma-topology-I-adic-f}
设 $X$、$f$、$(\mathcal{F}_n)$ 满足引理 \ref{cohomology-lemma-equivalent-f-good} 的条件 (1)。设 $p \geq 0$，并令
$H^p = \lim H^p(X, \mathcal{F}_n)$.
则对所有 $c \geq 1$，$f^cH^p$ 是
$H^p \to H^p(X, \mathcal{F}_c)$ 的核。因此，$H^p$ 上的极限拓扑
就是 $f$-进拓扑。
\end{lemma}

\begin{proof}
设 $c \geq 1$。显然，$f^c H^p$ 在 $H^p(X, \mathcal{F}_c)$ 中映到零。
若 $\xi = (\xi_n) \in H^p$ 在极限拓扑中足够小，则 $\xi_c = 0$，
因而当 $n \geq c$ 时，$\xi_n$ 在
$H^p(X, \mathcal{F}_c)$ 中映到零。考虑短正合列的逆系统
$$
0 \to \mathcal{F}_{n - c} \xrightarrow{f^c} \mathcal{F}_n \to
\mathcal{F}_c \to 0
$$
以及相应的上同调长正合列的逆系统
$$
H^{p - 1}(X, \mathcal{F}_c) \to
H^p(X, \mathcal{F}_{n - c}) \to
H^p(X, \mathcal{F}_n) \to
H^p(X, \mathcal{F}_c)
$$
由于项 $H^{p - 1}(X, \mathcal{F}_c)$ 与 $n$ 无关，可以选择相容的元素序列
$\xi'_n \in H^p(X, \mathcal{F}_{n - c})$
作为 $\xi_n$ 的提升。令 $\xi' = (\xi'_n)$，便得到
$\xi = f^c \xi'$，正合所需。
\end{proof}

\begin{lemma}
\label{cohomology-lemma-limit-finite}
设 $A$ 是关于主理想 $(f)$ 完备的 Noether 环。设 $X$ 是拓扑空间。设
$$
\ldots \to \mathcal{F}_3 \to \mathcal{F}_2 \to \mathcal{F}_1
$$
是 $X$ 上 $A$-模层的逆系统。假设
\begin{enumerate}
\item $\Gamma(X, \mathcal{F}_1)$ 是有限 $A$-模；
\item $X$、$f$、$(\mathcal{F}_n)$ 满足引理 \ref{cohomology-lemma-equivalent-f-good} 的条件 (1)。
\end{enumerate}
则
$$
M = \lim \Gamma(X, \mathcal{F}_n)
$$
是有限 $A$-模，$f$ 是 $M$ 上的非零因子，而且 $M/fM$ 是
$M$ 在 $\Gamma(X, \mathcal{F}_1)$ 中的像。
\end{lemma}

\begin{proof}
由引理 \ref{cohomology-lemma-topology-I-adic-f}，有
$M/fM \subset H^0(X, \mathcal{F}_1)$。由 (1) 与 $A$ 的 Noether
性质，$M/fM$ 是有限 $A$-模。注意
$\bigcap f^nM = 0$，因为 $f^nM$ 在
$H^0(X, \mathcal{F}_n)$ 中映到零。由
《代数》引理 \ref{algebra-lemma-finite-over-complete-ring}
可得 $M$ 是有限 $A$-模。最后，设
$s = (s_n) \in M = \lim H^0(X, \mathcal{F}_n)$ 满足 $fs = 0$。
由引理 \ref{cohomology-lemma-equivalent-f-good} 的条件 (1)，
$s_{n + 1}$ 属于 $\mathcal{F}_{n + 1} \to \mathcal{F}_n$ 的核。
因此 $s_n = 0$。由于 $n$ 任意，故 $s = 0$。于是 $f$ 是 $M$
上的非零因子。
\end{proof}

\begin{lemma}
\label{cohomology-lemma-ML}
设 $A$ 是环，$f \in A$，且 $X$ 是拓扑空间。设
$$
\ldots \to \mathcal{F}_3 \to \mathcal{F}_2 \to \mathcal{F}_1
$$
是 $X$ 上 $A$-模层的逆系统。设 $p \geq 0$。假设
\begin{enumerate}
\item 或者 $H^{p + 1}(X, \mathcal{F}_1)$ 是有限长度的 $A$-模，
或者 $A$ 是 Noether 环且 $H^{p + 1}(X, \mathcal{F}_1)$ 是有限 $A$-模；
\item $X$、$f$、$(\mathcal{F}_n)$ 满足引理 \ref{cohomology-lemma-equivalent-f-good} 的条件 (1)。
\end{enumerate}
则逆系统 $M_n = H^p(X, \mathcal{F}_n)$ 满足 Mittag--Leffler 条件。
\end{lemma}

\begin{proof}
令 $I = (f)$。我们将使用引理 \ref{cohomology-lemma-ML-general} 的判别法。
注意，对所有 $n \geq 0$，
$f^n : \mathcal{F}_1 \to I^n\mathcal{F}_{n + 1}$ 都是同构。
因此，只需证明
$$
\bigoplus\nolimits_{n \geq 1} H^{p + 1}(X, \mathcal{F}_1) \cdot f^{n + 1}
$$
作为分次 $S = \bigoplus_{n \geq 0} A/(f) \cdot f^n$-模满足升链条件。
若 $A$ 不是 Noether 环，则 $H^{p + 1}(X, \mathcal{F}_1)$ 长度有限，
结论成立。若 $A$ 是 Noether 环，则 $S$ 是 Noether 环；由所假设的
$H^{p + 1}(X, \mathcal{F}_1)$ 的有限性，该模是有限 $S$-模，
故结论亦成立。省略若干细节。
\end{proof}

\begin{lemma}
\label{cohomology-lemma-ML-better}
设 $A$ 是环，$f \in A$，且 $X$ 是拓扑空间。设
$$
\ldots \to \mathcal{F}_3 \to \mathcal{F}_2 \to \mathcal{F}_1
$$
是 $X$ 上 $A$-模层的逆系统。设 $p \geq 0$。假设
\begin{enumerate}
\item 或者存在 $m \geq 1$，使得
$H^{p + 1}(X, \mathcal{F}_m) \to H^{p + 1}(X, \mathcal{F}_1)$
的像是有限长度的 $A$-模；或者 $A$ 是 Noether 环，且各态射
$H^{p + 1}(X, \mathcal{F}_m) \to H^{p + 1}(X, \mathcal{F}_1)$
之像的交是有限 $A$-模；
\item $X$、$f$、$(\mathcal{F}_n)$ 满足引理 \ref{cohomology-lemma-equivalent-f-good} 的条件 (1)。
\end{enumerate}
则逆系统 $M_n = H^p(X, \mathcal{F}_n)$ 满足 Mittag--Leffler 条件。
\end{lemma}

\begin{proof}
令 $I = (f)$。我们将使用引理 \ref{cohomology-lemma-ML-general-better}
中涉及模 $N_n$ 的判别法。当 $m \geq n$ 时，有
$I^n\mathcal{F}_{m + 1} = \mathcal{F}_{m + 1 - n}$。因此
$$
N_n = \bigcap\nolimits_{m \geq 1} \Im\left(
H^{p + 1}(X, \mathcal{F}_m) \to H^{p + 1}(X, \mathcal{F}_1)
\right)
$$
与 $n$ 无关，并且
$\bigoplus N_n = \bigoplus N_1 \cdot f^{n + 1}$。
于是，可完全仿照引理 \ref{cohomology-lemma-ML} 的证明得出结论。
\end{proof}

\begin{remark}
\label{cohomology-remark-compare-ML}
设 $(X, \mathcal{O}_X)$ 是环化空间。设
$f \in \Gamma(X, \mathcal{O}_X)$。设 $\mathcal{F}$ 是
$\mathcal{O}_X$-模。若 $\mathcal{F}$ 无 $f$-挠，则对每个
$p \geq 0$，都有逆系统的短正合列
$$
0 \to \{H^p(X, \mathcal{F})/f^nH^p(X, \mathcal{F})\}
\to \{H^p(X, \mathcal{F}/f^n\mathcal{F})\}
\to \{H^{p + 1}(X, \mathcal{F})[f^n]\} \to 0
$$
由于第一个逆系统满足 Mittag--Leffler 条件（ML），可由此得到三点：
\begin{enumerate}
\item 存在短正合列
$$
0 \to \widehat{H^p(X, \mathcal{F})} \to
\lim H^p(X, \mathcal{F}/f^n\mathcal{F}) \to
T_f(H^{p + 1}(X, \mathcal{F})) \to 0
$$
其中 $\widehat{\ }$ 表示通常的 $f$-进完备化，而 $T_f( - )$
表示
《更多代数》例
\ref{more-algebra-example-spectral-sequence-principal}
中的 $f$-进 Tate 模。
\item 有
$R^1\lim H^p(X, \mathcal{F}/f^n\mathcal{F}) =
R^1\lim H^{p + 1}(X, \mathcal{F})[f^n]$.
\item 系统 $\{H^{p + 1}(X, \mathcal{F})[f^n]\}$ 满足 ML，当且仅当
$\{H^p(X, \mathcal{F}/f^n\mathcal{F})\}$ 满足 ML。
\end{enumerate}
参见 《同调代数》引理 \ref{homology-lemma-Mittag-Leffler} 以及
《更多代数》引理 \ref{more-algebra-lemma-six-term-Rlim} 和
\ref{more-algebra-lemma-Mittag-Leffler}。
\end{remark}










\section{导出极限}
\label{cohomology-section-derived-limits}

\noindent
设 $(X, \mathcal{O}_X)$ 是环化空间。由于三角范畴
$D(\mathcal{O}_X)$ 有积
（《内射对象》引理 \ref{injectives-lemma-derived-products}），
故 $D(\mathcal{O}_X)$ 有导出极限；参见
《导出范畴》定义 \ref{derived-definition-derived-limit}。
若 $(K_n)$ 是 $D(\mathcal{O}_X)$ 中的逆系统，则记
$R\lim K_n$ 为其导出极限。

\begin{lemma}
\label{cohomology-lemma-RGamma-commutes-with-Rlim}
设 $(X, \mathcal{O}_X)$ 是环化空间。对开集 $U \subset X$，函子
$R\Gamma(U, -)$ 与 $R\lim$ 交换。此外，对 $D(\mathcal{O}_X)$
中的任意逆系统 $(K_n)$ 及任意 $m \in \mathbf{Z}$，存在短正合列
$$
0 \to
R^1\lim H^{m - 1}(U, K_n) \to H^m(U, R\lim K_n) \to
\lim H^m(U, K_n) \to 0
$$
\end{lemma}

\begin{proof}
第一个陈述由
《内射对象》引理 \ref{injectives-lemma-RF-commutes-with-Rlim}
推出。然后可将
《更多代数》注记 \ref{more-algebra-remark-compare-derived-limit}
用于 $R\lim R\Gamma(U, K_n) = R\Gamma(U, R\lim K_n)$，
从而得到上述短正合列。
\end{proof}

\begin{lemma}
\label{cohomology-lemma-Rf-commutes-with-Rlim}
设 $f : (X, \mathcal{O}_X) \to (Y, \mathcal{O}_Y)$ 是环化空间的态射。
则 $Rf_*$ 与 $R\lim$ 交换；亦即，$Rf_*$ 与导出极限交换。
\end{lemma}

\begin{proof}
设 $(K_n)$ 是 $D(\mathcal{O}_X)$ 中的逆系统。考虑定义性的特异三角
$$
R\lim K_n \to \prod K_n \to \prod K_n
$$
它位于 $D(\mathcal{O}_X)$ 中。应用正合函子 $Rf_*$，得到特异三角
$$
Rf_*(R\lim K_n) \to Rf_*\left(\prod K_n\right) \to Rf_*\left(\prod K_n\right)
$$
它位于 $D(\mathcal{O}_Y)$ 中。因此，只需证明 $Rf_*$
与导出范畴中的积交换（这些积并非仅由复形的积给出；参见
《内射对象》引理 \ref{injectives-lemma-derived-products}）。
然而，由引理 \ref{cohomology-lemma-adjoint}，$Rf_*$ 是右伴随，故结论形式地成立
（参见 《范畴》引理 \ref{categories-lemma-adjoint-exact}）。
注意：不能直接应用
《范畴》引理 \ref{categories-lemma-adjoint-exact}
，因为 $R\lim K_n$ 并不是 $D(\mathcal{O}_X)$ 中的极限。
\end{proof}

\begin{remark}
\label{cohomology-remark-discuss-derived-limit}
设 $(X, \mathcal{O}_X)$ 是环化空间。设 $(K_n)$ 是
$D(\mathcal{O}_X)$ 中的逆系统。令 $K = R\lim K_n$。对每个 $n$ 与 $m$，
令 $\mathcal{H}^m_n = H^m(K_n)$ 为 $K_n$ 的第 $m$ 个上同调层，并类似地令
$\mathcal{H}^m = H^m(K)$。记 $\underline{\mathcal{H}}^m_n$ 为预层
$$
U \longmapsto \underline{\mathcal{H}}^m_n(U) = H^m(U, K_n)
$$
类似地，令 $\underline{\mathcal{H}}^m(U) = H^m(U, K)$。
由引理 \ref{cohomology-lemma-sheafification-cohomology} 可见，
$\mathcal{H}^m_n$ 是 $\underline{\mathcal{H}}^m_n$ 的层化，而
$\mathcal{H}^m$ 是 $\underline{\mathcal{H}}^m$ 的层化。有图
$$
\xymatrix{
K \ar@{=}[d] &
\underline{\mathcal{H}}^m \ar[d] \ar[r] & 
\mathcal{H}^m \ar[d] \\
R\lim K_n &
\lim \underline{\mathcal{H}}^m_n \ar[r] & 
\lim \mathcal{H}^m_n
}
$$
一般而言，$\lim \mathcal{H}^m_n$ 未必是
$\lim \underline{\mathcal{H}}^m_n$ 的层化。若 $U \subset X$ 是开集，
则存在短正合列
\begin{equation}
\label{cohomology-equation-ses-Rlim-over-U}
0 \to
R^1\lim \underline{\mathcal{H}}^{m - 1}_n(U) \to
\underline{\mathcal{H}}^m(U) \to
\lim \underline{\mathcal{H}}^m_n(U) \to 0
\end{equation}
这由引理 \ref{cohomology-lemma-RGamma-commutes-with-Rlim} 得到。
\end{remark}

\noindent
下述引理适用于概形上转移态射满射的拟凝聚模逆系统。

\begin{lemma}
\label{cohomology-lemma-inverse-limit-is-derived-limit}
设 $(X, \mathcal{O}_X)$ 是环化空间。设 $(\mathcal{F}_n)$ 是
$\mathcal{O}_X$-模的逆系统。设 $\mathcal{B}$ 是 $X$ 的一族开集。假设
\begin{enumerate}
\item $X$ 的每个开集都有一个覆盖，其成员属于 $\mathcal{B}$；
\item 当 $p > 0$ 且 $U \in \mathcal{B}$ 时，
$H^p(U, \mathcal{F}_n) = 0$；
\item 当 $U \in \mathcal{B}$ 时，逆系统 $\mathcal{F}_n(U)$ 的
$R^1\lim$ 为零。
\end{enumerate}
则 $R\lim \mathcal{F}_n = \lim \mathcal{F}_n$，并且当 $p > 0$
且 $U \in \mathcal{B}$ 时，有
$H^p(U, \lim \mathcal{F}_n) = 0$。
\end{lemma}

\begin{proof}
令 $K_n = \mathcal{F}_n$ 且 $K = R\lim \mathcal{F}_n$。采用注记 \ref{cohomology-remark-discuss-derived-limit} 的记号并利用假设 (2)，
可见对 $U \in \mathcal{B}$，当 $m \not = 0$ 时有
$\underline{\mathcal{H}}_n^m(U) = 0$，而
$\underline{\mathcal{H}}_n^0(U) = \mathcal{F}_n(U)$。
由 Equation (\ref{cohomology-equation-ses-Rlim-over-U}) 与假设 (3)，
可见当 $m \not = 0$ 时 $\underline{\mathcal{H}}^m(U) = 0$，
而当 $m = 0$ 时它等于 $\lim \mathcal{F}_n(U)$。利用 (1) 层化，
可得当 $m \not = 0$ 时 $\mathcal{H}^m = 0$，而当 $m = 0$ 时
它等于 $\lim \mathcal{F}_n$。因此 $K = \lim \mathcal{F}_n$。
由于当 $m > 0$ 时
$H^m(U, K) = \underline{\mathcal{H}}^m(U) = 0$（见上），
第二个断言也成立。
\end{proof}

\begin{lemma}
\label{cohomology-lemma-cohomology-derived-limit-injective}
设 $(X, \mathcal{O}_X)$ 是环化空间。设 $(K_n)$ 是
$D(\mathcal{O}_X)$ 中的逆系统。设 $x \in X$ 且 $m \in \mathbf{Z}$。
假设存在整数 $n(x)$ 以及 $x$ 的开邻域基本系 $\mathfrak{U}_x$，
使得对 $U \in \mathfrak{U}_x$，
\begin{enumerate}
\item $R^1\lim H^{m - 1}(U, K_n) = 0$；并且
\item 当 $n \geq n(x)$ 时，
$H^m(U, K_n) \to H^m(U, K_{n(x)})$ 是单射。
\end{enumerate}
则茎上的态射 $H^m(R\lim K_n)_x \to H^m(K_{n(x)})_x$ 是单射。
\end{lemma}

\begin{proof}
设 $\gamma$ 是 $H^m(R\lim K_n)_x$ 中的元素，且它在
$H^m(K_{n(x)})_x$ 中映到零。由于 $H^m(R\lim K_n)$ 是预层
$U \mapsto H^m(U, R\lim K_n)$ 的层化
（由引理 \ref{cohomology-lemma-sheafification-cohomology}），
可选择 $U \in \mathfrak{U}_x$ 以及映到 $\gamma$ 的元素
$\tilde \gamma \in H^m(U, R\lim K_n)$。于是
$\tilde\gamma$ 映到
$\tilde\gamma_{n(x)} \in H^m(U, K_{n(x)})$。
再次利用引理 \ref{cohomology-lemma-sheafification-cohomology}，即
$H^m(K_{n(x)})$ 是预层 $U \mapsto H^m(U, K_{n(x)})$ 的层化，
可在缩小 $U$ 后假设 $\tilde\gamma_{n(x)} = 0$。
对这个 $U$，考虑短正合列
$$
0 \to
R^1\lim H^{m - 1}(U, K_n) \to H^m(U, R\lim K_n) \to
\lim H^m(U, K_n) \to 0
$$
它来自引理 \ref{cohomology-lemma-RGamma-commutes-with-Rlim}。
由假设 (1)，左端群为零；由假设 (2)，右端群单射到
$H^m(U, K_{n(x)})$。因此 $\tilde\gamma = 0$，从而如所需，
$\gamma = 0$。
\end{proof}

\begin{lemma}
\label{cohomology-lemma-is-limit-per-point}
设 $(X, \mathcal{O}_X)$ 是环化空间。设 $E \in D(\mathcal{O}_X)$。
假设对每个 $x \in X$，都存在函数
$p(x, -) : \mathbf{Z} \to \mathbf{Z}$ 以及 $x$ 的开邻域基本系
$\mathfrak{U}_x$，使得
$$
H^p(U, H^{m - p}(E)) = 0 \text{，其中 }
U \in \mathfrak{U}_x \text{ 且 } p > p(x, m)
$$
则
《导出范畴》注记
\ref{derived-remark-map-into-derived-limit-truncations}
中的态射 $E \to R\lim \tau_{\geq -n} E$
是 $D(\mathcal{O}_X)$ 中的同构。
\end{lemma}

\begin{proof}
令 $K_n = \tau_{\geq -n}E$ 且 $K = R\lim K_n$。典范态射
$E \to K$ 来自各典范态射 $E \to K_n = \tau_{\geq -n}E$。
必须证明 $E \to K$ 在上同调层上诱导同构
$H^m(E) \to H^m(K)$。在余下的证明中固定 $m$。若 $n \geq -m$，
则态射 $E \to \tau_{\geq -n}E = K_n$ 诱导同构
$H^m(E) \to H^m(K_n)$。为完成证明，只需说明：对每个 $x \in X$，
都存在整数 $n(x) \geq -m$，使得态射
$H^m(K)_x \to H^m(K_{n(x)})_x$ 是单射。事实上，此时复合
$$
H^m(E)_x \to H^m(K)_x \to H^m(K_{n(x)})_x
$$
是双射，而第二个箭头是单射，故第一个箭头是双射。令
$$
n(x) = 1 + \max\{-m, p(x, m - 1) - m, -1 + p(x, m) - m, -2 + p(x, m + 1) - m\}.
$$
于是总有 $n(x) \geq -m$。断言：态射
$$
H^{m - 1}(U, K_{n + 1}) \to H^{m - 1}(U, K_n)
\quad\text{且}\quad
H^m(U, K_{n + 1}) \to H^m(U, K_n)
$$
当 $n \geq n(x)$ 且 $U \in \mathfrak{U}_x$ 时都是同构。
此断言蕴含引理 \ref{cohomology-lemma-cohomology-derived-limit-injective}
的条件 (1) 与 (2) 成立，从而蕴含所需的单射性。回忆
（《导出范畴》注记
\ref{derived-remark-truncation-distinguished-triangle}）
，存在特异三角
$$
H^{-n - 1}(E)[n + 1] \to
K_{n + 1} \to K_n \to H^{-n - 1}(E)[n + 2]
$$
考察相应的上同调长正合列可知，只要
$$
H^{m + n}(U, H^{-n - 1}(E)),\quad
H^{m + n + 1}(U, H^{-n - 1}(E)),\quad
H^{m + n + 2}(U, H^{-n - 1}(E))
$$
当 $n \geq n(x)$ 且 $U \in \mathfrak{U}_x$ 时均为零，断言便成立。
而这由 $n(x)$ 的选取以及引理中的假设推出。
\end{proof}

\begin{lemma}
\label{cohomology-lemma-is-limit-spaltenstein}
\begin{reference}
\cite[Proposition 3.13]{Spaltenstein}
\end{reference}
设 $(X, \mathcal{O}_X)$ 是环化空间。设 $E \in D(\mathcal{O}_X)$。
假设对每个 $x \in X$，都存在整数 $d_x \geq 0$ 以及 $x$
的开邻域基本系 $\mathfrak{U}_x$，使得
$$
H^p(U, H^q(E)) = 0 \text{，其中 }
U \in \mathfrak{U}_x,\ p > d_x, \text{ 且 }q < 0
$$
则
《导出范畴》注记
\ref{derived-remark-map-into-derived-limit-truncations}
中的态射 $E \to R\lim \tau_{\geq -n} E$
是 $D(\mathcal{O}_X)$ 中的同构。
\end{lemma}

\begin{proof}
将引理 \ref{cohomology-lemma-is-limit-per-point} 用于
$p(x, m) = d_x + \max(0, m)$ 即得。
\end{proof}

\begin{lemma}
\label{cohomology-lemma-is-limit}
设 $(X, \mathcal{O}_X)$ 是环化空间。设 $E \in D(\mathcal{O}_X)$。
假设存在函数 $p(-) : \mathbf{Z} \to \mathbf{Z}$ 以及 $X$
的一族开集 $\mathcal{B}$，使得
\begin{enumerate}
\item $X$ 的每个开集都有一个覆盖，其成员属于 $\mathcal{B}$；并且
\item 当 $p > p(m)$ 且 $U \in \mathcal{B}$ 时，
$H^p(U, H^{m - p}(E)) = 0$。
\end{enumerate}
则
《导出范畴》注记
\ref{derived-remark-map-into-derived-limit-truncations}
中的态射 $E \to R\lim \tau_{\geq -n} E$
是 $D(\mathcal{O}_X)$ 中的同构。
\end{lemma}

\begin{proof}
将引理 \ref{cohomology-lemma-is-limit-per-point} 用于
$p(x, m) = p(m)$ 以及
$\mathfrak{U}_x = \{U \in \mathcal{B} \mid x \in U\}$。
即可。
\end{proof}

\begin{lemma}
\label{cohomology-lemma-is-limit-dimension}
设 $(X, \mathcal{O}_X)$ 是环化空间。设 $E \in D(\mathcal{O}_X)$。
假设存在整数 $d \geq 0$ 以及 $X$ 的拓扑的基 $\mathcal{B}$，使得
$$
H^p(U, H^q(E)) = 0 \text{，其中 }
U \in \mathcal{B},\ p > d, \text{ 且 }q < 0
$$
则
《导出范畴》注记
\ref{derived-remark-map-into-derived-limit-truncations}
中的态射 $E \to R\lim \tau_{\geq -n} E$
是 $D(\mathcal{O}_X)$ 中的同构。
\end{lemma}

\begin{proof}
将引理 \ref{cohomology-lemma-is-limit-spaltenstein} 用于 $d_x = d$ 以及
$\mathfrak{U}_x = \{U \in \mathcal{B} \mid x \in U\}$ 即可。
\end{proof}

\noindent
上述引理可用于在某些情形下计算上同调。

\begin{lemma}
\label{cohomology-lemma-cohomology-over-U-trivial}
设 $(X, \mathcal{O}_X)$ 是环化空间。设 $K$ 是
$D(\mathcal{O}_X)$ 的对象。设 $\mathcal{B}$ 是 $X$ 的一族开集。假设
\begin{enumerate}
\item $X$ 的每个开集都有一个覆盖，其成员属于 $\mathcal{B}$；
\item 对所有 $p > 0$、$q \in \mathbf{Z}$ 以及
$U \in \mathcal{B}$，都有 $H^p(U, H^q(K)) = 0$。
\end{enumerate}
则当 $q \in \mathbf{Z}$ 且 $U \in \mathcal{B}$ 时，
$H^q(U, K) = H^0(U, H^q(K))$。
\end{lemma}

\begin{proof}
由引理 \ref{cohomology-lemma-is-limit-dimension}（取 $d = 0$），注意有
$K = R\lim \tau_{\geq -n} K$。设 $U \in \mathcal{B}$。由
Equation (\ref{cohomology-equation-ses-Rlim-over-U})，得到短正合列
$$
0 \to R^1\lim H^{q - 1}(U, \tau_{\geq -n}K) \to
H^q(U, K) \to \lim H^q(U, \tau_{\geq -n}K) \to 0
$$
条件 (2) 蕴含
$H^q(U, \tau_{\geq -n} K) = H^0(U, H^q(\tau_{\geq -n} K))$
对所有 $q$ 都成立；这里使用例 \ref{cohomology-example-spectral-sequence} 的谱序列。
由于 $\tau_{\geq -n}K$ 下有界，该谱序列收敛。若 $n > -q$，
则 $H^q(\tau_{\geq -n}K) = H^q(K)$。因此，上述短正合列左右两端
的系统最终为常值，其值分别为 $H^0(U, H^{q - 1}(K))$ 与
$H^0(U, H^q(K))$。引理得证。
\end{proof}

\noindent
下面是另一个可以描述导出极限的情形。

\begin{lemma}
\label{cohomology-lemma-derived-limit-suitable-system}
设 $(X, \mathcal{O}_X)$ 是环化空间。设 $(K_n)$ 是
$D(\mathcal{O}_X)$ 中对象的逆系统。设 $\mathcal{B}$ 是
$X$ 的一族开集。假设
\begin{enumerate}
\item $X$ 的每个开集都有一个覆盖，其成员属于 $\mathcal{B}$；
\item 对所有 $U \in \mathcal{B}$ 与所有 $q \in \mathbf{Z}$，均有
\begin{enumerate}
\item 当 $p > 0$ 时，$H^p(U, H^q(K_n)) = 0$；
\item 逆系统 $H^0(U, H^q(K_n))$ 的 $R^1\lim$ 为零。
\end{enumerate}
\end{enumerate}
则对 $q \in \mathbf{Z}$，有
$H^q(R\lim K_n) = \lim H^q(K_n)$。
\end{lemma}

\begin{proof}
令 $K = R\lim K_n$，并取 $U \in \mathcal{B}$。由引理 \ref{cohomology-lemma-cohomology-over-U-trivial} 与 (2)(a)，有
$H^q(U, K_n) = H^0(U, H^q(K_n))$。由引理 \ref{cohomology-lemma-RGamma-commutes-with-Rlim} 与 (2)(b)，有
$H^q(U, K) = \lim H^0(U, H^q(K_n))$。因此，$H^q(U, K)$ 是各层
$H^q(K_n)$ 在 $U$ 上截面之逆极限。由于 $\lim H^q(K_n)$ 是层，
再利用假设 (1)，可见作为预层 $U \mapsto H^q(U, K)$ 的层化的
$H^q(K)$ 等于 $\lim H^q(K_n)$。引理得证。
\end{proof}








\section{构造 K-内射消解}
\label{cohomology-section-K-injective}


\noindent
设 $(X, \mathcal{O}_X)$ 是环化空间。设 $\mathcal{F}^\bullet$
是由 $\mathcal{O}_X$-模组成的复形。范畴
$\textit{Mod}(\mathcal{O}_X)$ 有足够多的内射对象，故可利用
《导出范畴》引理 \ref{derived-lemma-special-inverse-system}
构造图
$$
\xymatrix{
\ldots \ar[r] &
\tau_{\geq -2}\mathcal{F}^\bullet \ar[r] \ar[d] &
\tau_{\geq -1}\mathcal{F}^\bullet \ar[d] \\
\ldots \ar[r] & \mathcal{I}_2^\bullet \ar[r] & \mathcal{I}_1^\bullet
}
$$
它位于 $\mathcal{O}_X$-模复形的范畴中，并满足
\begin{enumerate}
\item 各竖直箭头是拟同构；
\item $\mathcal{I}_n^\bullet$ 是由内射对象组成的下有界复形；
\item 各箭头 $\mathcal{I}_{n + 1}^\bullet \to \mathcal{I}_n^\bullet$
逐项为分裂满射。
\end{enumerate}
由于 $\mathcal{O}_X$-模的范畴有极限（它们在预层层次上计算），
可形成逐项极限
$\mathcal{I}^\bullet = \lim_n \mathcal{I}_n^\bullet$。由
《导出范畴》引理
\ref{derived-lemma-bounded-below-injectives-K-injective} 和
\ref{derived-lemma-limit-K-injectives}
，这是 K-内射复形。一般而言，典范态射
\begin{equation}
\label{cohomology-equation-into-candidate-K-injective}
\mathcal{F}^\bullet \to \mathcal{I}^\bullet
\end{equation}
未必是拟同构。下述引理给出使它成为拟同构的一些条件。

\begin{lemma}
\label{cohomology-lemma-K-injective}
在上述情形中，记
$\mathcal{H}^m = H^m(\mathcal{F}^\bullet)$ 为第 $m$ 个上同调层。
设 $\mathcal{B}$ 是 $X$ 的一族开子集，且 $d \in \mathbf{N}$。假设
\begin{enumerate}
\item $X$ 的每个开集都有一个覆盖，其成员属于 $\mathcal{B}$；
\item 对每个 $U \in \mathcal{B}$，当 $p > d$ 且 $q < 0$ 时，有
$H^p(U, \mathcal{H}^q) = 0$\footnote{只需对所有 $m$ 都存在
$p(m)$，使得当 $p > p(m)$ 时，
$H^p(U. \mathcal{H}^{m - p}) = 0$；参见引理 \ref{cohomology-lemma-is-limit}。}。
\end{enumerate}
则 (\ref{cohomology-equation-into-candidate-K-injective}) 是拟同构。
\end{lemma}

\begin{proof}
由 《导出范畴》引理
\ref{derived-lemma-difficulty-K-injectives}，只需证明态射
$\mathcal{F}^\bullet \to R\lim \tau_{\geq -n} \mathcal{F}^\bullet$
是同构。这正是引理 \ref{cohomology-lemma-is-limit-dimension}。
\end{proof}

\noindent
下面给出一个技术引理，讨论层复形系统之逆极限的上同调层。
从某种意义上说，这并不是应当尝试证明的结论，因为应取导出极限，
而非实际的逆极限。

\begin{lemma}
\label{cohomology-lemma-inverse-limit-complexes}
设 $(X, \mathcal{O}_X)$ 是环化空间。设 $(\mathcal{F}_n^\bullet)$
是 $\mathcal{O}_X$-模复形的逆系统。设 $m \in \mathbf{Z}$。
假设存在 $X$ 的一族开子集 $\mathcal{B}$ 与整数 $n_0$，使得
\begin{enumerate}
\item $X$ 的每个开集都有一个覆盖，其成员属于 $\mathcal{B}$；
\item 对每个 $U \in \mathcal{B}$，
\begin{enumerate}
\item 阿贝尔群系统
$\mathcal{F}_n^{m - 2}(U)$ 与 $\mathcal{F}_n^{m - 1}(U)$
的 $R^1\lim$ 为零（例如，它们满足 Mittag--Leffler 条件）；
\item 阿贝尔群系统 $H^{m - 1}(\mathcal{F}_n^\bullet(U))$
的 $R^1\lim$ 为零（例如，它满足 Mittag--Leffler 条件）；并且
\item 对所有 $n \geq n_0$，有
$H^m(\mathcal{F}_n^\bullet(U)) = H^m(\mathcal{F}_{n_0}^\bullet(U))$
。
\end{enumerate}
\end{enumerate}
则态射
$H^m(\mathcal{F}^\bullet) \to \lim H^m(\mathcal{F}_n^\bullet) \to
H^m(\mathcal{F}_{n_0}^\bullet)$
都是层的同构，其中
$\mathcal{F}^\bullet = \lim \mathcal{F}_n^\bullet$ 是逐项逆极限。
\end{lemma}

\begin{proof}
设 $U \in \mathcal{B}$。注意，
$H^m(\mathcal{F}^\bullet(U))$ 是下列复形从左数第三处的上同调：
$$
\lim_n \mathcal{F}_n^{m - 2}(U) \to
\lim_n \mathcal{F}_n^{m - 1}(U) \to
\lim_n \mathcal{F}_n^m(U) \to
\lim_n \mathcal{F}_n^{m + 1}(U)
$$
由假设 (2)(a) 与 (2)(b)，可应用
《更多代数》引理 \ref{more-algebra-lemma-apply-Mittag-Leffler-again}
得到
$$
H^m(\mathcal{F}^\bullet(U)) = \lim H^m(\mathcal{F}_n^\bullet(U))
$$
由假设 (2)(c)，得到
$$
H^m(\mathcal{F}^\bullet(U)) = H^m(\mathcal{F}_n^\bullet(U))
$$
对所有 $n \geq n_0$ 成立。由假设 (1)，可得预层
$U \mapsto H^m(\mathcal{F}^\bullet(U))$ 的层化，等于预层
$U \mapsto H^m(\mathcal{F}_n^\bullet(U))$ 的层化，
这对所有 $n \geq n_0$ 都成立。因此，层的逆系统
$H^m(\mathcal{F}_n^\bullet)$ 从 $n \geq n_0$ 起为常值，
其值为 $H^m(\mathcal{F}^\bullet)$。引理得证。
\end{proof}












\section{逆系统与上同调，III}
\label{cohomology-section-inverse-systems-tri}

\noindent
本节利用导出极限，继续 第 \ref{cohomology-section-inverse-systems-bis} 节
中的讨论。

\begin{lemma}
\label{cohomology-lemma-formal-functions-principal}
设 $(X, \mathcal{O}_X)$ 是环化空间。设
$A \to \Gamma(X, \mathcal{O}_X)$ 是环同态，且 $f \in A$。
设 $E$ 是 $D(\mathcal{O}_X)$ 的对象。记
$$
E_n = E \otimes_{\mathcal{O}_X} (\mathcal{O}_X \xrightarrow{f^n} \mathcal{O}_X)
$$
并令 $E^\wedge = R\lim E_n$。对 $p \in \mathbf{Z}$，存在典范交换图
$$
\xymatrix{
& 0 & 0 \\
0 \ar[r] &
\widehat{H^p(X, E)} \ar[r] \ar[u] &
\lim H^p(X, E_n) \ar[r] \ar[u] &
T_f(H^{p + 1}(X, E)) \ar[r] &
0 \\
0 \ar[r] &
H^0(H^p(X, E)^\wedge) \ar[r] \ar[u] &
H^p(X, E^\wedge) \ar[r] \ar[u] &
T_f(H^{p + 1}(X, E)) \ar[r] \ar@{=}[u] &
0 \\
&
R^1\lim H^p(X, E)[f^n] \ar[u] \ar[r]^\cong &
R^1\lim H^{p - 1}(X, E_n) \ar[u] \\
& 0 \ar[u] & 0 \ar[u]
}
$$
其各行各列均正合；其中
$\widehat{H^p(X, E)} = \lim H^p(X, E)/f^n H^p(X, E)$
是通常的 $f$-进完备化，$H^p(X, E)^\wedge$ 是导出 $f$-进完备化，
而 $T_f(H^{p + 1}(X, E))$ 是 $f$-进 Tate 模；参见
《更多代数》例
\ref{more-algebra-example-spectral-sequence-principal}。
最后，有 $H^p(X, E^\wedge) = H^p(R\Gamma(X, E)^\wedge)$。
\end{lemma}

\begin{proof}
由引理 \ref{cohomology-lemma-Rf-commutes-with-Rlim}，注意有
$R\Gamma(X, E^\wedge) = R\lim R\Gamma(X, E_n)$。另一方面，
$$
R\Gamma(X, E_n) =
R\Gamma(X, E) \otimes_A^\mathbf{L} (A \xrightarrow{f^n} A)
$$
（细节从略）。可见 $R\Gamma(X, E^\wedge)$ 是导出 $f$-进完备化
$R\Gamma(X, E)^\wedge$。上述图因而由
《更多代数》引理
\ref{more-algebra-lemma-compare-completion-derived-completion}
得到。
\end{proof}

\begin{lemma}
\label{cohomology-lemma-stabilizes}
设 $\mathcal{A}$ 是阿贝尔范畴。设 $f : M \to M$ 是
$\mathcal{A}$ 中的态射。若
$M[f^n] = \Ker(f^n : M \to M)$ 稳定，则逆系统
$$
(M \xrightarrow{f^n} M)
\quad\text{与}\quad
\Coker(f^n : M \to M)
$$
在 $D(\mathcal{A})$ 中 pro-同构。
\end{lemma}

\begin{proof}
显然存在从第一个逆系统到第二个逆系统的态射。设
$M[f^c] = M[f^{c + 1}] = M[f^{c + 2}] = \ldots$.
于是，可借助下列各图，在 $D(\mathcal{A})$ 中定义反方向的逆系统态射：
$$
\xymatrix{
M/M[f^c] \ar[r]_-{f^{n + c}} \ar[d]_{f^c} &
M \ar[d]^1 \\
M \ar[r]^{f^n} & M
}
$$
由于上方水平箭头是单射，上行的复形拟同构于
$\Coker(f^{n + c} : M \to M)$。省略若干细节。
\end{proof}

\begin{example}
\label{cohomology-example-formal-functions-principal}
设 $(X, \mathcal{O}_X)$ 是环化空间。设
$A \to \Gamma(X, \mathcal{O}_X)$ 是环同态，且 $f \in A$。
设 $\mathcal{F}$ 是 $\mathcal{O}_X$-模。假设存在 $c$，使得对所有
$n \geq c$，都有
$\mathcal{F}[f^c] = \mathcal{F}[f^n]$。我们将把引理 \ref{cohomology-lemma-formal-functions-principal} 用于
$E = \mathcal{F}$。由引理 \ref{cohomology-lemma-stabilizes} 可见，
逆系统 $(E_n)$ 与逆系统 $(\mathcal{F}/f^n\mathcal{F})$ pro-同构。
因此，对 $p \in \mathbf{Z}$，得到交换图
$$
\xymatrix{
& 0 & 0 \\
0 \ar[r] &
\widehat{H^p(X, \mathcal{F})} \ar[r] \ar[u] &
\lim H^p(X, \mathcal{F}/f^n\mathcal{F}) \ar[r] \ar[u] &
T_f(H^{p + 1}(X, \mathcal{F})) \ar[r] &
0 \\
0 \ar[r] &
H^0(H^p(X, \mathcal{F})^\wedge) \ar[r] \ar[u] &
H^p(R\Gamma(X, \mathcal{F})^\wedge) \ar[r] \ar[u] &
T_f(H^{p + 1}(X, \mathcal{F})) \ar[r] \ar@{=}[u] &
0 \\
&
R^1\lim H^p(X, \mathcal{F})[f^n] \ar[u] \ar[r]^\cong &
R^1\lim H^{p - 1}(X, \mathcal{F}/f^n\mathcal{F}) \ar[u] \\
& 0 \ar[u] & 0 \ar[u]
}
$$
其各行各列均正合，其中
$\widehat{H^p(X, \mathcal{F})} =
\lim H^p(X, \mathcal{F})/f^n H^p(X, \mathcal{F})$
是通常的 $f$-进完备化，而对 $D(A)$ 中的 $M$，
$M^\wedge$ 表示导出 $f$-进完备化。
\end{example}
















\section{无界复形的 {\v C}ech 上同调}
\label{cohomology-section-cech-cohomology-of-unbounded-complexes}

\noindent
对无界复形而言，第 \ref{cohomology-section-cech-cohomology-of-complexes} 节
的构造并不“正确”。问题在于，Stacks 项目对双复形作全复形化时使用直和，
而这里必须将其替换为积。本节不这样处理，而是假设覆盖有限，并使用交替
{\v C}ech 复形。

\medskip\noindent
设 $(X, \mathcal{O}_X)$ 是环化空间。设 ${\mathcal F}^\bullet$
是 $\mathcal{O}_X$-模预层的复形。设
${\mathcal U} : X = \bigcup_{i \in I} U_i$ 是 $X$ 的一个
{\bf 有限}开覆盖。由于交替 {\v C}ech 复形
$\check{\mathcal{C}}_{alt}^\bullet(\mathcal{U}, \mathcal{F})$
（第 \ref{cohomology-section-alternating-cech} 节）
关于预层 $\mathcal{F}$ 具有函子性，可得到双复形
$\check{\mathcal{C}}^\bullet_{alt}(\mathcal{U}, \mathcal{F}^\bullet)$。
本节使用相应的全复形。构造
$\text{Tot}(\check{\mathcal{C}}^\bullet_{alt}({\mathcal U},
{\mathcal F}^\bullet))$
关于 ${\mathcal F}^\bullet$ 具有函子性。此外，还有函子变换
\begin{equation}
\label{cohomology-equation-global-sections-to-alternating-cech}
\Gamma(X, {\mathcal F}^\bullet)
\longrightarrow
\text{Tot}(\check{\mathcal{C}}^\bullet_{alt}({\mathcal U},
{\mathcal F}^\bullet))
\end{equation}
，它是复形之间的态射，并由下述规则定义：截面
$s\in \Gamma(X, {\mathcal F}^n)$ 映到元素
$\alpha = \{\alpha_{i_0\ldots i_p}\}$，其中
$\alpha_{i_0} = s|_{U_{i_0}}$，且当 $p > 0$ 时
$\alpha_{i_0\ldots i_p} = 0$。

\begin{lemma}
\label{cohomology-lemma-alternating-cech-complex-complex}
设 $(X, \mathcal{O}_X)$ 是环化空间。设
$\mathcal{U} : X = \bigcup_{i \in I} U_i$ 是有限开覆盖。
对由 $\mathcal{O}_X$-模组成的复形 $\mathcal{F}^\bullet$，
存在典范态射
$$
\text{Tot}(\check{\mathcal{C}}^\bullet_{alt}(\mathcal{U}, \mathcal{F}^\bullet))
\longrightarrow
R\Gamma(X, \mathcal{F}^\bullet)
$$
它关于 $\mathcal{F}^\bullet$ 具有函子性，并且与
(\ref{cohomology-equation-global-sections-to-alternating-cech}) 相容。
\end{lemma}

\begin{proof}
设 ${\mathcal I}^\bullet$ 是 K-内射复形，其各项均为内射
$\mathcal{O}_X$-模。对 $\mathcal{I}^\bullet$，
态射 (\ref{cohomology-equation-global-sections-to-alternating-cech}) 为
$\Gamma(X, {\mathcal I}^\bullet) \to
\text{Tot}(\check{\mathcal{C}}^\bullet_{alt}({\mathcal U},
{\mathcal I}^\bullet))$.
这是阿贝尔群复形的拟同构；此结论由
《同调代数》引理 \ref{homology-lemma-double-complex-gives-resolution}
应用于双复形
$\check{\mathcal{C}}^\bullet_{alt}({\mathcal U},
{\mathcal I}^\bullet)$，并利用引理 \ref{cohomology-lemma-injective-trivial-cech} 与 \ref{cohomology-lemma-alternating-usual}
得到。设 ${\mathcal F}^\bullet \to {\mathcal I}^\bullet$
是从 ${\mathcal F}^\bullet$ 到一个各项均内射的 K-内射复形的拟同构
（《内射对象》Theorem
\ref{injectives-theorem-K-injective-embedding-grothendieck}）。
由于 $R\Gamma(X, {\mathcal F}^\bullet)$ 由复形
$\Gamma(X, {\mathcal I}^\bullet)$ 表示，利用
$$
\text{Tot}(\check{\mathcal{C}}^\bullet_{alt}({\mathcal U},
{\mathcal F}^\bullet))
\longrightarrow
\text{Tot}(\check{\mathcal{C}}^\bullet_{alt}({\mathcal U},
{\mathcal I}^\bullet)).
$$
即可得到引理中的态射。省略函子性与相容性的验证。
\end{proof}

\begin{lemma}
\label{cohomology-lemma-alternating-cech-complex-complex-ss}
设 $(X, \mathcal{O}_X)$ 是环化空间。设
$\mathcal{U} : X = \bigcup_{i \in I} U_i$ 是有限开覆盖。
设 $\mathcal{F}^\bullet$ 是由 $\mathcal{O}_X$-模组成的复形。
设 $\mathcal{B}$ 是 $X$ 的一族开子集。假设
\begin{enumerate}
\item $X$ 的每个开集都有一个覆盖，其成员属于 $\mathcal{B}$；
\item 对所有 $i_0, \ldots, i_p \in I$，有
$U_{i_0\ldots i_p} \in \mathcal{B}$；
\item 对每个 $U \in \mathcal{B}$ 与 $p > 0$，都有
\begin{enumerate}
\item $H^p(U, \mathcal{F}^q) = 0$；
\item $H^p(U, \Coker(\mathcal{F}^{q - 1} \to \mathcal{F}^q)) = 0$；并且
\item $H^p(U, H^q(\mathcal{F})) = 0$。
\end{enumerate}
\end{enumerate}
则态射
$$
\text{Tot}(\check{\mathcal{C}}^\bullet_{alt}(\mathcal{U}, \mathcal{F}^\bullet))
\longrightarrow
R\Gamma(X, \mathcal{F}^\bullet)
$$
即引理 \ref{cohomology-lemma-alternating-cech-complex-complex} 中的态射，
是 $D(\textit{Ab})$ 中的同构。
\end{lemma}

\begin{proof}
先假设 $\mathcal{F}^\bullet$ 下有界。在这种情形下，态射
$$
\text{Tot}(\check{\mathcal{C}}^\bullet_{alt}(\mathcal{U}, \mathcal{F}^\bullet))
\longrightarrow
\text{Tot}(\check{\mathcal{C}}^\bullet(\mathcal{U}, \mathcal{F}^\bullet))
$$
由引理 \ref{cohomology-lemma-alternating-usual} 是拟同构。事实上，双复形的态射
$\check{\mathcal{C}}^\bullet_{alt}(\mathcal{U}, \mathcal{F}^\bullet) \to
\check{\mathcal{C}}^\bullet(\mathcal{U}, \mathcal{F}^\bullet)$
在这些复形所对应的第二谱序列的第一页之间诱导同构
（由 《同调代数》引理 \ref{homology-lemma-ss-double-complex}），
而这些谱序列收敛
（《同调代数》引理 \ref{homology-lemma-first-quadrant-ss}）。
因此，这种情形下的结论由引理 \ref{cohomology-lemma-cech-complex-complex-computes} 与假设 (3)(a) 得到。

\medskip\noindent
一般而言，由假设 (3)(c)，可选择如引理 \ref{cohomology-lemma-K-injective}
中的消解
$\mathcal{F}^\bullet \to \mathcal{I}^\bullet = \lim \mathcal{I}_n^\bullet$
。于是，引理中的态射变为
$$
\lim_n 
\text{Tot}(\check{\mathcal{C}}^\bullet_{alt}(\mathcal{U},
\tau_{\geq -n}\mathcal{F}^\bullet))
\longrightarrow
\Gamma(X, \mathcal{I}^\bullet) =
\lim_n \Gamma(X, \mathcal{I}_n^\bullet)
$$
这里箭头位于导出范畴中，而右侧的等式在复形层次上成立。
注意，(3)(b) 表明 $\tau_{\geq -n}\mathcal{F}^\bullet$
是满足本引理假设的下有界复形。因此，由下有界复形的情形可知，
每个态射
$$
\text{Tot}(\check{\mathcal{C}}^\bullet_{alt}(\mathcal{U},
\tau_{\geq -n}\mathcal{F}^\bullet))
\longrightarrow
\Gamma(X, \mathcal{I}_n^\bullet)
$$
都是拟同构。左侧复形在给定次数上的上同调最终为常值
（因为交替 {\v C}ech 复形有限），故右侧亦然。因此，由
《同调代数》引理 \ref{homology-lemma-apply-Mittag-Leffler-again}
（或更强的 《更多代数》引理
\ref{more-algebra-lemma-apply-Mittag-Leffler-again})
，右侧极限的上同调就是这个常值，结论得证。
\end{proof}






\section{Hom 复形}
\label{cohomology-section-hom-complexes}

\noindent
设 $(X, \mathcal{O}_X)$ 是环化空间。设
$\mathcal{L}^\bullet$ 与 $\mathcal{M}^\bullet$ 是两个
$\mathcal{O}_X$-模复形。我们构造 $\mathcal{O}_X$-模复形
$\SheafHom^\bullet(\mathcal{L}^\bullet, \mathcal{M}^\bullet)$。
具体地，对每个 $n$，令
$$
\SheafHom^n(\mathcal{L}^\bullet, \mathcal{M}^\bullet) =
\prod\nolimits_{n = p + q}
\SheafHom_{\mathcal{O}_X}(\mathcal{L}^{-q}, \mathcal{M}^p)
$$
最好将 $\SheafHom^n$ 看作如下 $\mathcal{O}_X$-模层：
从 $\mathcal{L}^\bullet$ 到 $\mathcal{M}^\bullet$
（二者看作分次 $\mathcal{O}_X$-模）的所有次数为 $n$ 的齐次
$\mathcal{O}_X$-线性态射。采用此术语，微分由规则
$$
\text{d}(f) =
\text{d}_\mathcal{M} \circ f - (-1)^n f \circ \text{d}_\mathcal{L}
$$
定义，其中
$f \in \SheafHom^n_{\mathcal{O}_X}(\mathcal{L}^\bullet,
\mathcal{M}^\bullet)$。省略 $\text{d}^2 = 0$ 的验证。
此构造是
《微分分次代数》例 \ref{dga-example-category-complexes}。
的特例。由构造立即可得
\begin{equation}
\label{cohomology-equation-cohomology-hom-complex}
H^n(\Gamma(U, \SheafHom^\bullet(\mathcal{L}^\bullet, \mathcal{M}^\bullet))) =
\Hom_{K(\mathcal{O}_U)}(\mathcal{L}^\bullet, \mathcal{M}^\bullet[n])
\end{equation}
这对所有 $n \in \mathbf{Z}$ 与每个开集 $U \subset X$ 都成立。

\begin{lemma}
\label{cohomology-lemma-compose}
设 $(X, \mathcal{O}_X)$ 是环化空间。给定 $\mathcal{O}_X$-模复形
$\mathcal{K}^\bullet,\mathcal{L}^\bullet,\mathcal{M}^\bullet$，
存在同构
$$
\SheafHom^\bullet(\mathcal{K}^\bullet,
\SheafHom^\bullet(\mathcal{L}^\bullet, \mathcal{M}^\bullet))
=
\SheafHom^\bullet(\text{Tot}(\mathcal{K}^\bullet \otimes_{\mathcal{O}_X}
\mathcal{L}^\bullet), \mathcal{M}^\bullet)
$$
它是 $\mathcal{O}_X$-模复形的同构，并关于
$\mathcal{K}^\bullet,\mathcal{L}^\bullet,\mathcal{M}^\bullet$
具有函子性。
\end{lemma}

\begin{proof}
从略。提示：证明与
《更多代数》引理 \ref{more-algebra-lemma-compose}
完全相同。
\end{proof}

\begin{lemma}
\label{cohomology-lemma-composition}
设 $(X, \mathcal{O}_X)$ 是环化空间。给定 $\mathcal{O}_X$-模复形
$\mathcal{K}^\bullet,\mathcal{L}^\bullet,\mathcal{M}^\bullet$，
存在典范态射
$$
\text{Tot}\left(
\SheafHom^\bullet(\mathcal{L}^\bullet, \mathcal{M}^\bullet)
\otimes_{\mathcal{O}_X}
\SheafHom^\bullet(\mathcal{K}^\bullet, \mathcal{L}^\bullet)
\right)
\longrightarrow
\SheafHom^\bullet(\mathcal{K}^\bullet, \mathcal{M}^\bullet)
$$
，它是 $\mathcal{O}_X$-模复形之间的态射。
\end{lemma}

\begin{proof}
从略。提示：证明与
《更多代数》引理 \ref{more-algebra-lemma-composition}
完全相同。
\end{proof}

\begin{lemma}
\label{cohomology-lemma-diagonal-better}
设 $(X, \mathcal{O}_X)$ 是环化空间。给定 $\mathcal{O}_X$-模复形
$\mathcal{K}^\bullet,\mathcal{L}^\bullet,\mathcal{M}^\bullet$，
存在典范态射
$$
\text{Tot}\left(
\mathcal{K}^\bullet \otimes_{\mathcal{O}_X}
\SheafHom^\bullet(\mathcal{M}^\bullet, \mathcal{L}^\bullet)
\right)
\longrightarrow
\SheafHom^\bullet(\mathcal{M}^\bullet,
\text{Tot}(\mathcal{K}^\bullet \otimes_{\mathcal{O}_X} \mathcal{L}^\bullet))
$$
，它是 $\mathcal{O}_X$-模复形之间的态射，并关于三个复形均具有函子性。
\end{lemma}

\begin{proof}
从略。提示：证明与
《更多代数》引理 \ref{more-algebra-lemma-diagonal-better}
完全相同。
\end{proof}

\begin{lemma}
\label{cohomology-lemma-diagonal}
设 $(X, \mathcal{O}_X)$ 是环化空间。给定 $\mathcal{O}_X$-模复形
$\mathcal{K}^\bullet,\mathcal{L}^\bullet$，存在典范态射
$$
\mathcal{K}^\bullet
\longrightarrow
\SheafHom^\bullet(\mathcal{L}^\bullet,
\text{Tot}(\mathcal{K}^\bullet \otimes_{\mathcal{O}_X} \mathcal{L}^\bullet))
$$
，它是 $\mathcal{O}_X$-模复形之间的态射，并关于两个复形均具有函子性。
\end{lemma}

\begin{proof}
从略。提示：证明与
《更多代数》引理 \ref{more-algebra-lemma-diagonal}
完全相同。
\end{proof}

\begin{lemma}
\label{cohomology-lemma-evaluate-and-more}
设 $(X, \mathcal{O}_X)$ 是环化空间。给定 $\mathcal{O}_X$-模复形
$\mathcal{K}^\bullet,\mathcal{L}^\bullet,\mathcal{M}^\bullet$，
存在典范态射
$$
\text{Tot}(\SheafHom^\bullet(\mathcal{L}^\bullet,
\mathcal{M}^\bullet) \otimes_{\mathcal{O}_X} \mathcal{K}^\bullet)
\longrightarrow
\SheafHom^\bullet(\SheafHom^\bullet(\mathcal{K}^\bullet,
\mathcal{L}^\bullet), \mathcal{M}^\bullet)
$$
，它是 $\mathcal{O}_X$-模复形之间的态射，并关于三个复形均具有函子性。
\end{lemma}

\begin{proof}
从略。提示：证明与
《更多代数》引理 \ref{more-algebra-lemma-evaluate-and-more}
完全相同。
\end{proof}

\begin{lemma}
\label{cohomology-lemma-RHom-into-K-injective}
设 $(X, \mathcal{O}_X)$ 是环化空间。设 $L$ 与 $M$ 是
$D(\mathcal{O}_X)$ 的对象。设 $\mathcal{I}^\bullet$
是表示 $M$ 的 $\mathcal{O}_X$-模 K-内射复形。设
$\mathcal{L}^\bullet$ 是表示 $L$ 的 $\mathcal{O}_X$-模复形。则
$$
H^0(\Gamma(U, \SheafHom^\bullet(\mathcal{L}^\bullet, \mathcal{I}^\bullet))) =
\Hom_{D(\mathcal{O}_U)}(L|_U, M|_U)
$$
对所有开集 $U \subset X$ 成立。
\end{lemma}

\begin{proof}
有
\begin{align*}
H^0(\Gamma(U, \SheafHom^\bullet(\mathcal{L}^\bullet, \mathcal{I}^\bullet)))
& =
\Hom_{K(\mathcal{O}_U)}(\mathcal{L}^\bullet|_U, \mathcal{I}^\bullet|_U) \\
& =
\Hom_{D(\mathcal{O}_U)}(L|_U, M|_U)
\end{align*}
第一个等号是 (\ref{cohomology-equation-cohomology-hom-complex})。
第二个等号成立，是因为由引理 \ref{cohomology-lemma-restrict-K-injective-to-open}，
$\mathcal{I}^\bullet|_U$ 是 K-内射复形。
\end{proof}

\begin{lemma}
\label{cohomology-lemma-RHom-well-defined}
设 $(X, \mathcal{O}_X)$ 是环化空间。设
$(\mathcal{I}')^\bullet \to \mathcal{I}^\bullet$
是 $\mathcal{O}_X$-模 K-内射复形之间的拟同构。设
$(\mathcal{L}')^\bullet \to \mathcal{L}^\bullet$
是 $\mathcal{O}_X$-模复形之间的拟同构。则
$$
\SheafHom^\bullet(\mathcal{L}^\bullet, (\mathcal{I}')^\bullet)
\longrightarrow
\SheafHom^\bullet((\mathcal{L}')^\bullet, \mathcal{I}^\bullet)
$$
是拟同构。
\end{lemma}

\begin{proof}
令 $M$ 为由 $\mathcal{I}^\bullet$ 与
$(\mathcal{I}')^\bullet$ 表示的 $D(\mathcal{O}_X)$ 对象。
令 $L$ 为由 $\mathcal{L}^\bullet$ 与
$(\mathcal{L}')^\bullet$ 表示的 $D(\mathcal{O}_X)$ 对象。
由引理 \ref{cohomology-lemma-RHom-into-K-injective} 可见，各层
$$
H^0(\SheafHom^\bullet(\mathcal{L}^\bullet, (\mathcal{I}')^\bullet))
\quad\text{与}\quad
H^0(\SheafHom^\bullet((\mathcal{L}')^\bullet, \mathcal{I}^\bullet))
$$
都等于预层
$$
U \longmapsto \Hom_{D(\mathcal{O}_U)}(L|_U, M|_U)
$$
所伴随的层。因此，此态射是拟同构。
\end{proof}

\begin{lemma}
\label{cohomology-lemma-RHom-from-K-flat-into-K-injective}
设 $(X, \mathcal{O}_X)$ 是环化空间。设 $\mathcal{I}^\bullet$
是 $\mathcal{O}_X$-模 K-内射复形，并设 $\mathcal{L}^\bullet$
是 $\mathcal{O}_X$-模 K-平坦复形。则
$\SheafHom^\bullet(\mathcal{L}^\bullet, \mathcal{I}^\bullet)$
是 $\mathcal{O}_X$-模 K-内射复形。
\end{lemma}

\begin{proof}
事实上，若 $\mathcal{K}^\bullet$ 是无上同调的
$\mathcal{O}_X$-模复形，则
\begin{align*}
\Hom_{K(\mathcal{O}_X)}(\mathcal{K}^\bullet,
\SheafHom^\bullet(\mathcal{L}^\bullet, \mathcal{I}^\bullet))
& =
H^0(\Gamma(X,
\SheafHom^\bullet(\mathcal{K}^\bullet,
\SheafHom^\bullet(\mathcal{L}^\bullet, \mathcal{I}^\bullet)))) \\
& =
H^0(\Gamma(X, \SheafHom^\bullet(\text{Tot}(
\mathcal{K}^\bullet \otimes_{\mathcal{O}_X} \mathcal{L}^\bullet),
\mathcal{I}^\bullet))) \\
& =
\Hom_{K(\mathcal{O}_X)}(
\text{Tot}(\mathcal{K}^\bullet \otimes_{\mathcal{O}_X} \mathcal{L}^\bullet),
\mathcal{I}^\bullet) \\
& =
0
\end{align*}
第一个等号由 (\ref{cohomology-equation-cohomology-hom-complex}) 得到。
第二个等号由引理 \ref{cohomology-lemma-compose} 得到。
第三个等号由 (\ref{cohomology-equation-cohomology-hom-complex}) 得到。
最后一个等号成立，是因为
$\text{Tot}(\mathcal{K}^\bullet \otimes_{\mathcal{O}_X} \mathcal{L}^\bullet)$
无上同调（因为 $\mathcal{L}^\bullet$ 是 K-平坦的；参见定义 \ref{cohomology-definition-K-flat}），且 $\mathcal{I}^\bullet$
是 K-内射的。
\end{proof}








\section{导出范畴中的内 Hom}
\label{cohomology-section-internal-hom}

\noindent
设 $(X, \mathcal{O}_X)$ 为环化空间，并设 $L,M$ 为
$D(\mathcal{O}_X)$ 的对象。我们希望构造 $D(\mathcal{O}_X)$ 的对象
$R\SheafHom(L, M)$，使得对 $D(\mathcal{O}_X)$ 的任意对象 $K$，
都存在典范双射
\begin{equation}
\label{cohomology-equation-internal-hom}
\Hom_{D(\mathcal{O}_X)}(K, R\SheafHom(L, M))
=
\Hom_{D(\mathcal{O}_X)}(K \otimes_{\mathcal{O}_X}^\mathbf{L} L, M)
\end{equation}
注意，由 Yoneda 引理
（《范畴》引理 \ref{categories-lemma-yoneda}），这个公式在唯一同构
意义下定义了 $R\SheafHom(L, M)$。

\medskip\noindent
为构造这样的对象，选取表示 $M$ 的 K-内射复形
$\mathcal{I}^\bullet$，并任取表示 $L$ 的
$\mathcal{O}_X$-模复形 $\mathcal{L}^\bullet$。定义
$$
R\SheafHom(L, M) = \SheafHom^\bullet(\mathcal{L}^\bullet, \mathcal{I}^\bullet)
$$
其中右端是第 \ref{cohomology-section-hom-complexes} 节构造的
$\mathcal{O}_X$-模复形。由引理 \ref{cohomology-lemma-RHom-well-defined}，
这一定义是良定的。我们得到函子
$$
D(\mathcal{O}_X)^{opp} \times D(\mathcal{O}_X) \longrightarrow D(\mathcal{O}_X),
\quad
(K, L) \longmapsto R\SheafHom(K, L)
$$
在证明式 (\ref{cohomology-equation-internal-hom}) 之前，我们先计算
$R\SheafHom(K,L)$ 的上同调群。

\begin{lemma}
\label{cohomology-lemma-section-RHom-over-U}
设 $(X,\mathcal{O}_X)$ 为环化空间，并设 $L,M$ 为
$D(\mathcal{O}_X)$ 的对象。对每个开集 $U$，有
$$
H^0(U, R\SheafHom(L, M)) =
\Hom_{D(\mathcal{O}_U)}(L|_U, M|_U)
$$
特别地，
$H^0(X,R\SheafHom(L,M))=\Hom_{D(\mathcal{O}_X)}(L,M)$。
\end{lemma}

\begin{proof}
选取表示 $M$ 的 K-内射 $\mathcal{O}_X$-模复形
$\mathcal{I}^\bullet$，以及表示 $L$ 的 K-平坦复形
$\mathcal{L}^\bullet$。于是
$\SheafHom^\bullet(\mathcal{L}^\bullet, \mathcal{I}^\bullet)$
由引理 \ref{cohomology-lemma-RHom-from-K-flat-into-K-injective} 是 K-内射的。
因此，只需在 $U$ 上取截面即可计算其在 $U$ 上的上同调，而结论由引理 \ref{cohomology-lemma-RHom-into-K-injective} 得出。
\end{proof}

\begin{lemma}
\label{cohomology-lemma-internal-hom}
设 $(X,\mathcal{O}_X)$ 为环化空间，并设 $K,L,M$ 为
$D(\mathcal{O}_X)$ 的对象。按上述构造，存在典范同构
$$
R\SheafHom(K, R\SheafHom(L, M)) =
R\SheafHom(K \otimes_{\mathcal{O}_X}^\mathbf{L} L, M)
$$
它位于 $D(\mathcal{O}_X)$ 中，并且关于 $K,L,M$ 具有函子性；
取 $H^0(X,-)$ 即恢复式 (\ref{cohomology-equation-internal-hom})。
\end{lemma}

\begin{proof}
选取表示 $M$ 的 K-内射复形 $\mathcal{I}^\bullet$，以及表示 $L$ 的
K-平坦 $\mathcal{O}_X$-模复形 $\mathcal{L}^\bullet$。任取表示 $K$ 的
$\mathcal{O}_X$-模复形 $\mathcal{K}^\bullet$。于是有
$$
\SheafHom^\bullet(\mathcal{K}^\bullet,
\SheafHom^\bullet(\mathcal{L}^\bullet, \mathcal{I}^\bullet))
=
\SheafHom^\bullet(
\text{Tot}(\mathcal{K}^\bullet \otimes_{\mathcal{O}_X} \mathcal{L}^\bullet),
\mathcal{I}^\bullet)
$$
这是引理 \ref{cohomology-lemma-compose} 的结论。注意，左端表示
$R\SheafHom(K,R\SheafHom(L,M))$（使用引理
\ref{cohomology-lemma-RHom-from-K-flat-into-K-injective}），而右端表示
$R\SheafHom(K \otimes_{\mathcal{O}_X}^\mathbf{L} L, M)$。
这就证明了引理中的显示公式。取整体截面并使用引理 \ref{cohomology-lemma-section-RHom-over-U}，即得式
(\ref{cohomology-equation-internal-hom})。
\end{proof}

\begin{lemma}
\label{cohomology-lemma-restriction-RHom-to-U}
设 $(X,\mathcal{O}_X)$ 为环化空间，并设 $K,L$ 为
$D(\mathcal{O}_X)$ 的对象。$R\SheafHom(K,L)$ 的构造与限制到开集的
操作交换；也就是说，对每个开集 $U$，有
$R\SheafHom(K|_U, L|_U) = R\SheafHom(K, L)|_U$.
\end{lemma}

\begin{proof}
这由构造和引理 \ref{cohomology-lemma-restrict-K-injective-to-open} 立即可得。
\end{proof}

\begin{lemma}
\label{cohomology-lemma-RHom-triangulated}
设 $(X,\mathcal{O}_X)$ 为环化空间。双函子 $R\SheafHom(-,-)$
在两个变元中都把特异三角变为特异三角。
\end{lemma}

\begin{proof}
只需注意，赋值
$$
(\mathcal{L}^\bullet, \mathcal{M}^\bullet) \longmapsto
\SheafHom^\bullet(\mathcal{L}^\bullet, \mathcal{M}^\bullet)
$$
在任一变元中都把逐项分裂的复形短正合列变为逐项分裂的短正合列。
细节从略。
\end{proof}

\begin{lemma}
\label{cohomology-lemma-internal-hom-composition}
设 $(X,\mathcal{O}_X)$ 为环化空间。给定 $D(\mathcal{O}_X)$ 中的
$K,L,M$，存在典范态射
$$
R\SheafHom(L, M) \otimes_{\mathcal{O}_X}^\mathbf{L} R\SheafHom(K, L)
\longrightarrow R\SheafHom(K, M)
$$
它位于 $D(\mathcal{O}_X)$ 中，并且关于 $K,L,M$ 具有函子性。
\end{lemma}

\begin{proof}
选取表示 $M$ 的 K-内射复形 $\mathcal{I}^\bullet$、表示 $L$ 的
K-内射复形 $\mathcal{J}^\bullet$，并任取表示 $K$ 的
$\mathcal{O}_X$-模复形 $\mathcal{K}^\bullet$。由引理
\ref{cohomology-lemma-composition}，存在复形态射
$$
\text{Tot}\left(
\SheafHom^\bullet(\mathcal{J}^\bullet, \mathcal{I}^\bullet)
\otimes_{\mathcal{O}_X}
\SheafHom^\bullet(\mathcal{K}^\bullet, \mathcal{J}^\bullet)
\right)
\longrightarrow
\SheafHom^\bullet(\mathcal{K}^\bullet, \mathcal{I}^\bullet)
$$
这些 $\mathcal{O}_X$-模复形
$\SheafHom^\bullet(\mathcal{J}^\bullet, \mathcal{I}^\bullet)$、
$\SheafHom^\bullet(\mathcal{K}^\bullet, \mathcal{J}^\bullet)$ 与
$\SheafHom^\bullet(\mathcal{K}^\bullet, \mathcal{I}^\bullet)$
分别表示 $R\SheafHom(L,M)$、$R\SheafHom(K,L)$ 和
$R\SheafHom(K,M)$。若选取 K-平坦复形 $\mathcal{H}^\bullet$
以及拟同构
$\mathcal{H}^\bullet \to
\SheafHom^\bullet(\mathcal{K}^\bullet, \mathcal{J}^\bullet)$,
则存在态射
$$
\text{Tot}\left(
\SheafHom^\bullet(\mathcal{J}^\bullet, \mathcal{I}^\bullet)
\otimes_{\mathcal{O}_X} \mathcal{H}^\bullet
\right)
\longrightarrow
\text{Tot}\left(
\SheafHom^\bullet(\mathcal{J}^\bullet, \mathcal{I}^\bullet)
\otimes_{\mathcal{O}_X}
\SheafHom^\bullet(\mathcal{K}^\bullet, \mathcal{J}^\bullet)
\right)
$$
其源表示
$R\SheafHom(L, M) \otimes_{\mathcal{O}_X}^\mathbf{L} R\SheafHom(K, L)$。
将上述两个显示态射复合，即得所需态射。构造具有函子性的证明从略。
\end{proof}

\begin{lemma}
\label{cohomology-lemma-internal-hom-diagonal-better}
设 $(X,\mathcal{O}_X)$ 为环化空间。给定 $D(\mathcal{O}_X)$ 中的
$K,L,M$，存在典范态射
$$
K \otimes_{\mathcal{O}_X}^\mathbf{L} R\SheafHom(M, L)
\longrightarrow
R\SheafHom(M, K \otimes_{\mathcal{O}_X}^\mathbf{L} L)
$$
它位于 $D(\mathcal{O}_X)$ 中，并且关于 $K,L,M$ 具有函子性。
\end{lemma}

\begin{proof}
选取表示 $K$ 的 K-平坦复形 $\mathcal{K}^\bullet$、表示 $L$ 的
K-内射复形 $\mathcal{I}^\bullet$，并任取表示 $M$ 的
$\mathcal{O}_X$-模复形 $\mathcal{M}^\bullet$。选取拟同构
$\text{Tot}(\mathcal{K}^\bullet \otimes_{\mathcal{O}_X} \mathcal{I}^\bullet)
\to \mathcal{J}^\bullet$
其中 $\mathcal{J}^\bullet$ 是 K-内射的。然后使用态射
$$
\text{Tot}\left(
\mathcal{K}^\bullet \otimes_{\mathcal{O}_X}
\SheafHom^\bullet(\mathcal{M}^\bullet, \mathcal{I}^\bullet)
\right)
\to
\SheafHom^\bullet(\mathcal{M}^\bullet,
\text{Tot}(\mathcal{K}^\bullet \otimes_{\mathcal{O}_X} \mathcal{I}^\bullet))
\to
\SheafHom^\bullet(\mathcal{M}^\bullet, \mathcal{J}^\bullet)
$$
其中第一个态射是引理 \ref{cohomology-lemma-diagonal-better} 中的态射。
\end{proof}

\begin{lemma}
\label{cohomology-lemma-internal-hom-diagonal}
设 $(X,\mathcal{O}_X)$ 为环化空间。给定 $D(\mathcal{O}_X)$ 中的
$K,L$，存在典范态射
$$
K \longrightarrow R\SheafHom(L, K \otimes_{\mathcal{O}_X}^\mathbf{L} L)
$$
它位于 $D(\mathcal{O}_X)$ 中，并且关于 $K$ 与 $L$ 均具有函子性。
\end{lemma}

\begin{proof}
选取表示 $K$ 的 K-平坦复形 $\mathcal{K}^\bullet$，并任取表示 $L$ 的
$\mathcal{O}_X$-模复形 $\mathcal{L}^\bullet$。选取 K-内射复形
$\mathcal{J}^\bullet$ 以及拟同构
$\text{Tot}(\mathcal{K}^\bullet \otimes_{\mathcal{O}_X} \mathcal{L}^\bullet)
\to \mathcal{J}^\bullet$。然后使用
$$
\mathcal{K}^\bullet \to
\SheafHom^\bullet(\mathcal{L}^\bullet,
\text{Tot}(\mathcal{K}^\bullet \otimes_{\mathcal{O}_X} \mathcal{L}^\bullet))
\to
\SheafHom^\bullet(\mathcal{L}^\bullet, \mathcal{J}^\bullet)
$$
其中第一个态射来自引理 \ref{cohomology-lemma-diagonal}。
\end{proof}

\begin{lemma}
\label{cohomology-lemma-dual}
设 $(X,\mathcal{O}_X)$ 为环化空间，并设 $L$ 为
$D(\mathcal{O}_X)$ 的对象。置
$L^\vee=R\SheafHom(L,\mathcal{O}_X)$。对 $D(\mathcal{O}_X)$ 中的
$M$，存在典范态射
\begin{equation}
\label{cohomology-equation-eval}
M \otimes^\mathbf{L}_{\mathcal{O}_X} L^\vee
\longrightarrow
R\SheafHom(L, M)
\end{equation}
它诱导典范态射
$$
H^0(X, M \otimes^\mathbf{L}_{\mathcal{O}_X} L^\vee)
\longrightarrow
\Hom_{D(\mathcal{O}_X)}(L, M)
$$
并且关于 $D(\mathcal{O}_X)$ 中的 $M$ 具有函子性。
\end{lemma}

\begin{proof}
利用等同 $M=R\SheafHom(\mathcal{O}_X,M)$，态射
(\ref{cohomology-equation-eval}) 是引理
\ref{cohomology-lemma-internal-hom-composition} 的特殊情形。
\end{proof}

\begin{lemma}
\label{cohomology-lemma-internal-hom-evaluate}
设 $(X,\mathcal{O}_X)$ 为环化空间，并设 $K,L,M$ 为
$D(\mathcal{O}_X)$ 的对象。存在典范态射
$$
R\SheafHom(L, M) \otimes_{\mathcal{O}_X}^\mathbf{L} K
\longrightarrow
R\SheafHom(R\SheafHom(K, L), M)
$$
它位于 $D(\mathcal{O}_X)$ 中，并且关于 $K,L,M$ 具有函子性。
\end{lemma}

\begin{proof}
选取表示 $M$ 的 K-内射复形 $\mathcal{I}^\bullet$、表示 $L$ 的
K-内射复形 $\mathcal{J}^\bullet$，以及表示 $K$ 的 K-平坦复形
$\mathcal{K}^\bullet$。所需态射由引理 \ref{cohomology-lemma-evaluate-and-more}
中的态射
$$
\text{Tot}(\SheafHom^\bullet(\mathcal{J}^\bullet,
\mathcal{I}^\bullet) \otimes_{\mathcal{O}_X} \mathcal{K}^\bullet)
\longrightarrow
\SheafHom^\bullet(\SheafHom^\bullet(\mathcal{K}^\bullet,
\mathcal{J}^\bullet), \mathcal{I}^\bullet)
$$
定义。根据我们对复形的特定选取，左端表示
$R\SheafHom(L, M) \otimes_{\mathcal{O}_X}^\mathbf{L} K$，右端表示
$R\SheafHom(R\SheafHom(K,L),M)$。它关于
$D(\mathcal{O}_X)$ 中三个对象均具有函子性的证明从略。
\end{proof}

\begin{remark}
\label{cohomology-remark-tensor-internal-hom}
设 $(X,\mathcal{O}_X)$ 为环化空间。对 $D(\mathcal{O}_X)$ 中的
$K,K',M,M'$，存在典范态射
$$
R\SheafHom(K, K') \otimes_{\mathcal{O}_X}^\mathbf{L}
R\SheafHom(M, M')
\longrightarrow
R\SheafHom(K \otimes_{\mathcal{O}_X}^\mathbf{L} M,
K' \otimes_{\mathcal{O}_X}^\mathbf{L} M')
$$
事实上，由式 (\ref{cohomology-equation-internal-hom})，给出这个态射等价于给出态射
$$
R\SheafHom(K, K') \otimes_{\mathcal{O}_X}^\mathbf{L}
R\SheafHom(M, M') \otimes_{\mathcal{O}_X}^\mathbf{L}
K \otimes_{\mathcal{O}_X}^\mathbf{L} M
\longrightarrow
K' \otimes_{\mathcal{O}_X}^\mathbf{L} M'
$$
为此，可以先交换中间两个因子（符号规则见 《更多代数》第 \ref{more-algebra-section-sign-rules} 节），再使用态射
$$
R\SheafHom(K, K') \otimes_{\mathcal{O}_X}^\mathbf{L} K \to K'
\quad\text{且}\quad
R\SheafHom(M, M') \otimes_{\mathcal{O}_X}^\mathbf{L} M \to M'
$$
这里应用引理 \ref{cohomology-lemma-internal-hom-composition}，并使用等同
$K = R\SheafHom(\mathcal{O}_X, K)$；对 $K'$、$M$ 和 $M'$
也作类似处理。
\end{remark}

\begin{remark}
\label{cohomology-remark-projection-formula-for-internal-hom}
设 $f:X\to Y$ 为环化空间的态射，并设 $K,L$ 为
$D(\mathcal{O}_X)$ 的对象。我们断言存在典范态射
$$
Rf_*R\SheafHom(L, K) \longrightarrow R\SheafHom(Rf_*L, Rf_*K)
$$
事实上，由式 (\ref{cohomology-equation-internal-hom})，给出这个态射等价于给出态射
$Rf_*R\SheafHom(L, K) \otimes_{\mathcal{O}_Y}^\mathbf{L} Rf_*L \to Rf_*K$.
为此，可以使用复合
$$
Rf_*R\SheafHom(L, K) \otimes_{\mathcal{O}_Y}^\mathbf{L} Rf_*L \to
Rf_*(R\SheafHom(L, K) \otimes_{\mathcal{O}_X}^\mathbf{L} L) \to
Rf_*K
$$
其中第一个箭头是相对杯积（注记 \ref{cohomology-remark-cup-product}），
第二个箭头则是对来自引理
\ref{cohomology-lemma-internal-hom-composition} 的典范态射
$R\SheafHom(L, K) \otimes_{\mathcal{O}_X}^\mathbf{L} L \to K$
应用 $Rf_*$ 所得（其中一个位置取 $\mathcal{O}_X$）。
\end{remark}

\begin{remark}
\label{cohomology-remark-relative-cup-and-composition}
设 $h:X\to Y$ 为环化空间的态射，并设 $K,M$ 为
$D(\mathcal{O}_Y)$ 的对象。图表
$$
\xymatrix{
Rf_*R\SheafHom_{\mathcal{O}_X}(K, M)
\otimes_{\mathcal{O}_Y}^\mathbf{L} Rf_*K
\ar[r] \ar[d] &
Rf_*\left(R\SheafHom_{\mathcal{O}_X}(K, M)
\otimes_{\mathcal{O}_X}^\mathbf{L} K\right)
\ar[d] \\
R\SheafHom_{\mathcal{O}_Y}(Rf_*K, Rf_*M) \otimes_{\mathcal{O}_Y}^\mathbf{L}
Rf_*K \ar[r] &
Rf_*M
}
$$
是交换的。这里，左侧竖直箭头来自注记
\ref{cohomology-remark-projection-formula-for-internal-hom}，顶部水平箭头是注记 \ref{cohomology-remark-cup-product} 中的箭头。另两个箭头是引理
\ref{cohomology-lemma-internal-hom-composition} 中态射的实例（把其中一个项换成
$\mathcal{O}_X$ 或 $\mathcal{O}_Y$）。
\end{remark}

\begin{remark}
\label{cohomology-remark-prepare-fancy-base-change}
设 $h:X\to Y$ 为环化空间的态射，并设 $K,L$ 为
$D(\mathcal{O}_Y)$ 的对象。我们断言存在典范态射
$$
Lh^*R\SheafHom(K, L) \longrightarrow R\SheafHom(Lh^*K, Lh^*L)
$$
它位于 $D(\mathcal{O}_X)$ 中。事实上，由式
(\ref{cohomology-equation-internal-hom})（该式见引理 \ref{cohomology-lemma-internal-hom}），
给出这样的态射等价于给出态射
$$
Lh^*R\SheafHom(K, L) \otimes^\mathbf{L} Lh^*K \longrightarrow Lh^*L
$$
由引理 \ref{cohomology-lemma-pullback-tensor-product}，这个箭头的源为
$Lh^*(\SheafHom(K,L)\otimes^\mathbf{L}K)$，故只需构造典范态射
$$
R\SheafHom(K, L) \otimes^\mathbf{L} K \longrightarrow L.
$$
为此，取经式 (\ref{cohomology-equation-internal-hom}) 与恒等态射
$$
\text{id} :
R\SheafHom(K, L)
\longrightarrow
R\SheafHom(K, L)
$$
相对应的箭头。
\end{remark}

\begin{remark}
\label{cohomology-remark-fancy-base-change}
设
$$
\xymatrix{
X' \ar[r]_h \ar[d]_{f'} &
X \ar[d]^f \\
S' \ar[r]^g &
S
}
$$
是环化空间的交换图表。设 $K,L$ 为 $D(\mathcal{O}_X)$ 的对象。
我们断言存在典范基变换态射
$$
Lg^*Rf_*R\SheafHom(K, L)
\longrightarrow
R(f')_*R\SheafHom(Lh^*K, Lh^*L)
$$
它位于 $D(\mathcal{O}_{S'})$ 中。具体地，取下列复合的伴随态射：
\begin{align*}
L(f')^*Lg^*Rf_*R\SheafHom(K, L)
& =
Lh^*Lf^*Rf_*R\SheafHom(K, L) \\
& \to
Lh^*R\SheafHom(K, L) \\
& \to
R\SheafHom(Lh^*K, Lh^*L)
\end{align*}
其中第一个箭头使用伴随态射 $Lf^*Rf_*\to\text{id}$，第二个箭头是注记 \ref{cohomology-remark-prepare-fancy-base-change} 中构造的典范态射。
\end{remark}





\section{Ext 层}
\label{cohomology-section-ext}

\noindent
设 $(X,\mathcal{O}_X)$ 为环化空间，并设
$K,L\in D(\mathcal{O}_X)$。利用导出范畴中内 Hom 的构造，取
$D(\mathcal{O}_X)$ 的对象 $R\SheafHom(K,L)$ 的第 $n$ 个上同调层，
得到良定的 $\mathcal{O}_X$-模层
$$
\SheafExt^n(K, L) = H^n(R\SheafHom(K, L))
$$
有时将这个对象写成 $\SheafExt^n_{\mathcal{O}_X}(K,L)$。
由引理 \ref{cohomology-lemma-section-RHom-over-U}，这个 $\SheafExt^n$-层是下列
对应的层化：
$$
U \longmapsto \Ext^n_{D(\mathcal{O}_U)}(K|_U, L|_U)
$$
由例 \ref{cohomology-example-spectral-sequence}，总存在谱序列
$$
E_2^{p, q} = H^p(X, \SheafExt^q(K, L))
$$
它在适当条件下收敛到 $\Ext^{p + q}_{D(\mathcal{O}_X)}(K,L)$
（例如 $L$ 下有界而 $K$ 上有界时）。





\section{整体导出 Hom}
\label{cohomology-section-global-RHom}

\noindent
设 $(X,\mathcal{O}_X)$ 为环化空间，并设
$K,L\in D(\mathcal{O}_X)$。利用导出范畴中内 Hom 的构造，得到良定对象
$$
R\Hom_X(K, L) = R\Gamma(X, R\SheafHom(K, L))
$$
它位于 $D(\Gamma(X,\mathcal{O}_X))$ 中。有时将这个对象写成
$R\Hom_{\mathcal{O}_X}(K,L)$。由引理
\ref{cohomology-lemma-section-RHom-over-U}，有
$$
H^0(R\Hom_X(K, L)) = \Hom_{D(\mathcal{O}_X)}(K, L),
\quad
H^p(R\Hom_X(K, L)) = \Ext_{D(\mathcal{O}_X)}^p(K, L)
$$
若 $f:Y\to X$ 是环化空间的态射，则存在典范态射
$$
R\Hom_X(K, L) \longrightarrow R\Hom_Y(Lf^*K, Lf^*L)
$$
它位于 $D(\Gamma(X,\mathcal{O}_X))$ 中，是对注记
\ref{cohomology-remark-prepare-fancy-base-change} 中定义的态射取整体截面所得。








\section{复形的粘合}
\label{cohomology-section-glueing-complexes}

\noindent
复形是可以粘合的！更准确地说，在某些情形下，可以把局部给定的
导出范畴对象粘成整体对象。我们先证明一些简单情形，然后在拓扑空间
与开覆盖的框架下证明十分一般的 \cite[Theorem 3.2.4]{BBD}。

\begin{lemma}
\label{cohomology-lemma-glue}
设 $(X,\mathcal{O}_X)$ 为环化空间，并设 $X=U\cup V$ 是 $X$ 的两个
开子空间之并。假设给定
\begin{enumerate}
\item $D(\mathcal{O}_U)$ 的对象 $A$，
\item $D(\mathcal{O}_V)$ 的对象 $B$，以及
\item 同构 $c:A|_{U\cap V}\to B|_{U\cap V}$。
\end{enumerate}
则存在 $D(\mathcal{O}_X)$ 的对象 $F$ 以及同构
$f:F|_U\to A$、$g:F|_V\to B$，使得
$c=g|_{U\cap V}\circ f^{-1}|_{U\cap V}$。此外，若给定
\begin{enumerate}
\item $D(\mathcal{O}_X)$ 的对象 $E$，
\item $D(\mathcal{O}_U)$ 中的态射 $a:A\to E|_U$，
\item $D(\mathcal{O}_V)$ 中的态射 $b:B\to E|_V$，
\end{enumerate}
满足
$$
a|_{U \cap V}  = b|_{U \cap V} \circ c.
$$
则存在 $D(\mathcal{O}_X)$ 中的态射 $F\to E$，其在 $U$ 上的限制为
$a\circ f$，在 $V$ 上的限制为 $b\circ g$。
\end{lemma}

\begin{proof}
以 $j_U$、$j_V$、$j_{U\cap V}$ 表示相应的开浸入。选取特异三角
$$
F \to Rj_{U, *}A \oplus Rj_{V, *}B \to Rj_{U \cap V, *}(B|_{U \cap V})
\to F[1]
$$
其中态射 $Rj_{V,*}B\to Rj_{U\cap V,*}(B|_{U\cap V})$ 是显然的，
而 $Rj_{U,*}A\to Rj_{U\cap V,*}(B|_{U\cap V})$ 是
$Rj_{U,*}A\to Rj_{U\cap V,*}(A|_{U\cap V})$ 与
$Rj_{U\cap V,*}c$ 的复合。限制到 $U$，得到
$$
F|_U \to A \oplus (Rj_{V, *}B)|_U \to (Rj_{U \cap V, *}(B|_{U \cap V}))|_U
\to F|_U[1]
$$
记 $j:U\cap V\to U$。限制到开集与上同调的相容性表明，
$(Rj_{V,*}B)|_U$ 与 $(Rj_{U\cap V,*}(B|_{U\cap V}))|_U$
都典范同构于 $Rj_*(B|_{U\cap V})$。因此，上一个显示图表的第二个
箭头有截面，从而态射 $F|_U\to A$ 是同构。类似地，
$F|_V\to B$ 也是同构。态射 $F\to E$ 的存在性来自 $\Hom$ 的
Mayer--Vietoris 序列；见引理 \ref{cohomology-lemma-mayer-vietoris-hom}。
\end{proof}

\begin{lemma}
\label{cohomology-lemma-vanishing-and-glueing}
设 $f:(X,\mathcal{O}_X)\to(Y,\mathcal{O}_Y)$ 为环化空间的态射，
并设 $\mathcal{B}$ 为 $Y$ 的拓扑基。
\begin{enumerate}
\item 假设 $K$ 属于 $D(\mathcal{O}_X)$，并且对 $V\in\mathcal{B}$，
当 $i<0$ 时有 $H^i(f^{-1}(V),K)=0$。则 $Rf_*K$ 的负次数上同调层
为零；对所有开集 $V\subset Y$，当 $i<0$ 时有
$H^i(f^{-1}(V),K)=0$；并且对应
$V\mapsto H^0(f^{-1}V,K)$ 是 $Y$ 上的层。
\item 假设 $K,L$ 属于 $D(\mathcal{O}_X)$，并且对
$V\in\mathcal{B}$，当 $i<0$ 时有
$\Ext^i(K|_{f^{-1}V},L|_{f^{-1}V})=0$。则对所有开集
$V\subset Y$，当 $i<0$ 时有
$\Ext^i(K|_{f^{-1}V},L|_{f^{-1}V})=0$；并且对应
$V\mapsto\Hom(K|_{f^{-1}V},L|_{f^{-1}V})$ 是 $Y$ 上的层。
\end{enumerate}
\end{lemma}

\begin{proof}
引理 \ref{cohomology-lemma-unbounded-describe-higher-direct-images} 告诉我们，
$H^i(Rf_*K)$ 是预层
$V\mapsto H^i(f^{-1}(V),K)=H^i(V,Rf_*K)$ 所伴随的层。
(1) 的假设蕴含 $Rf_*K$ 的次数 $<0$ 的上同调层为零。因此，对任意
开集 $V\subset Y$，上同调群 $H^i(V,Rf_*K)$ 在 $i<0$ 时为零，
在 $i=0$ 时等于 $H^0(V,H^0(Rf_*K))$。这证明了 (1)。

\medskip\noindent
为证明 (2)，对复形 $R\SheafHom(K,L)$ 应用 (1)，并使用引理
\ref{cohomology-lemma-section-RHom-over-U} 作相应转换。
\end{proof}

\begin{situation}
\label{cohomology-situation-locally-given}
设 $(X,\mathcal{O}_X)$ 为环化空间。给定
\begin{enumerate}
\item $X$ 的一族开集 $\mathcal{B}$；
\item 对每个 $U\in\mathcal{B}$，$D(\mathcal{O}_U)$ 中的对象 $K_U$；
\item 对满足 $V\subset U$ 且 $V,U\in\mathcal{B}$ 的每一对开集，
$D(\mathcal{O}_V)$ 中的同构 $\rho^U_V:K_U|_V\to K_V$；
\end{enumerate}
并要求每当 $W \subset V \subset U$ 且 $U, V, W$ 属于
$\mathcal{B}$ 时，
都有 $\rho^U_W=\rho^V_W\circ\rho^U_V|_W$。
\end{situation}

\noindent
若不再加假设，我们无法由此证明任何结论。一个有趣的情形是
$\mathcal{B}$ 为 $X$ 的拓扑基。另一个情形是给定拓扑空间的态射
$f:X\to Y$，而 $\mathcal{B}$ 的元素是 $Y$ 的某个拓扑基中各元素的
逆像。

\medskip\noindent
在情形 \ref{cohomology-situation-locally-given} 中，一个{\it 解}是二元组
$(K,\rho_U)$，其中 $K$ 是 $D(\mathcal{O}_X)$ 的对象，而对
$U\in\mathcal{B}$，$\rho_U:K|_U\to K_U$ 是同构，并且对所有
$V\subset U$、$U,V\in\mathcal{B}$，都有
$\rho^U_V\circ\rho_U|_V=\rho_V$。在某些情形下，解是唯一的。

\begin{lemma}
\label{cohomology-lemma-uniqueness}
在情形 \ref{cohomology-situation-locally-given} 中，假设
\begin{enumerate}
\item $X = \bigcup_{U \in \mathcal{B}} U$，并且
对 $U,V\in\mathcal{B}$ 有
$U \cap V = \bigcup_{W \in \mathcal{B}, W \subset U \cap V} W$,
\item 对任意 $U\in\mathcal{B}$，当 $i<0$ 时有
$\Ext^i(K_U,K_U)=0$。
\end{enumerate}
若解 $(K,\rho_U)$ 存在，则它在唯一同构意义下唯一；此外，当
$i<0$ 时有 $\Ext^i(K,K)=0$。
\end{lemma}

\begin{proof}
设 $(K,\rho_U)$ 与 $(K',\rho'_U)$ 是两个解。令 $f:X\to Y$ 为
《拓扑》引理 \ref{topology-lemma-create-map-from-subcollection}
中构造的连续映射，并置 $\mathcal{O}_Y=f_*\mathcal{O}_X$。
于是 $K,K'$ 和 $\mathcal{B}$ 满足引理
\ref{cohomology-lemma-vanishing-and-glueing} (2) 的情形。因此，$K$ 的负次数
Ext 群为零，并且对应
$$
V \longmapsto \Hom(K|_{f^{-1}V}, K'|_{f^{-1}V})
$$
是 $Y$ 上的层。由于 $(K,\rho_U)$ 与 $(K',\rho'_U)$ 都是解，故
$$
(\rho'_U)^{-1} \circ \rho_U : K|_U \longrightarrow K'|_U
$$
这些定义在 $U=f^{-1}(f(U))$ 上的态射在交叠处一致。因此得到上述层的
唯一整体截面，它定义所需同构 $K\to K'$，并与所有已有结构相容。
\end{proof}

\begin{remark}
\label{cohomology-remark-uniqueness}
采用引理 \ref{cohomology-lemma-uniqueness} 的记号与假设。设
$U,V\in\mathcal{B}$，并令 $\mathcal{B}'$ 为 $\mathcal{B}$ 中包含于
$U\cap V$ 的元素之集。则
$$
(\{K_{U'}\}_{U' \in \mathcal{B}'},
\{\rho_{V'}^{U'}\}_{V' \subset U'\text{ 且 }U', V' \in \mathcal{B}'})
$$
是环化空间 $U\cap V$ 上满足引理 \ref{cohomology-lemma-uniqueness} 假设的系统。
而且，$(K_U|_{U\cap V},\rho^U_{U'})$ 与
$(K_V|_{U\cap V},\rho^V_{U'})$ 都是该系统的解。由此引理，得到唯一同构
$$
\rho_{U, V} : K_U|_{U \cap V} \longrightarrow K_V|_{U \cap V}
$$
使得对每个 $U'\subset U\cap V$、$U'\in\mathcal{B}$，图表
$$
\xymatrix{
K_U|_{U'} \ar[rr]_{\rho_{U, V}|_{U'}} \ar[rd]_{\rho^U_{U'}} & &
K_V|_{U'} \ar[ld]^{\rho^V_{U'}} \\
& K_{U'}
}
$$
交换。再取第三个元素 $W\in\mathcal{B}$。得到同构
$\rho_{U, W} : K_U|_{U \cap W} \to K_W|_{U \cap W}$ 和
$\rho_{V, W} : K_U|_{V \cap W} \to K_W|_{V \cap W}$，
它们满足与 $\rho_{U,V}$ 类似的性质。最后，有
$$
\rho_{U, W}|_{U \cap V \cap W} =
\rho_{V, W}|_{U \cap V \cap W} \circ \rho_{U, V}|_{U \cap V \cap W}
$$
这由引理中的唯一性成立，因为等式两端都是与态射
$\rho^U_{U''}$ 和 $\rho^W_{U''}$ 相容的唯一同构，其中
$U''\subset U\cap V\cap W$、$U''\in\mathcal{B}$。略去一些次要细节。
集合 $(K_U,\rho_{U,V})$ 是 $X$ 的开覆盖
$\mathcal{U}:X=\bigcup_{U\in\mathcal{B}}U$ 在导出范畴中的下降数据。
用这种语言来说，寻找解的存在性，就是寻找“下降数据的有效性”。
\end{remark}

\begin{lemma}
\label{cohomology-lemma-solution-in-finite-case}
在情形 \ref{cohomology-situation-locally-given} 中，假设
\begin{enumerate}
\item $X = U_1 \cup \ldots \cup U_n$，其中 $U_i \in \mathcal{B}$；
\item 对 $U,V\in\mathcal{B}$，有
$U \cap V = \bigcup_{W \in \mathcal{B}, W \subset U \cap V} W$,
\item 对任意 $U\in\mathcal{B}$，当 $i<0$ 时有
$\Ext^i(K_U,K_U)=0$。
\end{enumerate}
则解存在，并且在唯一同构意义下唯一。
\end{lemma}

\begin{proof}
唯一性已由引理 \ref{cohomology-lemma-uniqueness} 得到。我们对 $n$ 作归纳来证明
本引理。$n=1$ 的情形是立即的。

\medskip\noindent
设 $n=2$。考虑注记 \ref{cohomology-remark-uniqueness} 中构造的同构
$\rho_{U_1, U_2} : K_{U_1}|_{U_1 \cap U_2} \to K_{U_2}|_{U_1 \cap U_2}$
。由引理 \ref{cohomology-lemma-glue}，得到 $D(\mathcal{O}_X)$ 中的对象 $K$，
以及与 $\rho_{U_1,U_2}$ 相容的同构
$\rho_{U_1}:K|_{U_1}\to K_{U_1}$ 与
$\rho_{U_2}:K|_{U_2}\to K_{U_2}$。取 $U\in\mathcal{B}$。
我们将构造同构 $\rho_U:K|_U\to K_U$，并把验证
$(K,\rho_U)$ 是解留给读者。令 $\mathcal{B}'$ 为 $\mathcal{B}$ 中
包含于 $U\cap U_1$ 或包含于 $U\cap U_2$ 的元素之集。于是
$(K_U,\rho^U_{U'})$ 是环化空间 $U$ 上系统
$(\{K_{U'}\}_{U' \in \mathcal{B}'},
\{\rho_{V'}^{U'}\}_{V' \subset U'\text{ 且 }U', V' \in \mathcal{B}'})$
的一个解。我们断言 $(K|_U,\tau_{U'})$ 是另一个解，其中对
$U'\in\mathcal{B}'$，如下选取 $\tau_{U'}$：若 $U'\subset U_1$，
则取复合
$$
K|_{U'} \xrightarrow{\rho_{U_1}|_{U'}}
K_{U_1}|_{U'} \xrightarrow{\rho^{U_1}_{U'}}
K_{U'}
$$
若 $U'\subset U_2$，则取复合
$$
K|_{U'} \xrightarrow{\rho_{U_2}|_{U'}}
K_{U_2}|_{U'} \xrightarrow{\rho^{U_2}_{U'}}
K_{U'}.
$$
为验证这是一个解，使用注记 \ref{cohomology-remark-uniqueness} 中所述
$\rho_{U_1,U_2}$ 的性质，以及 $\rho_{U_1}$、$\rho_{U_2}$ 与
$\rho_{U_1,U_2}$ 的相容性。由引理 \ref{cohomology-lemma-uniqueness}，得到唯一
同构 $K|_{U'}\to K_{U'}$；对 $U'\in\mathcal{B}'$，它与态射
$\tau_{U'}$ 和 $\rho^U_{U'}$ 相容。

\medskip\noindent
设 $n>2$。考虑开子空间 $X'=U_1\cup\ldots\cup U_{n-1}$，并令
$\mathcal{B}'$ 为 $\mathcal{B}$ 中包含于 $X'$ 的元素之集。于是得到
环化空间 $X'$ 上的系统
$(\{K_U\}_{U \in \mathcal{B}'}, \{\rho_V^U\}_{U, V \in \mathcal{B}'})$
，可以对它应用归纳假设。我们得到解 $(K_{X'},\rho^{X'}_U)$。
然后考虑 $X$ 的开集族
$\mathcal{B}^*=\mathcal{B}\cup\{X'\}$，并得到系统
$(\{K_U\}_{U \in \mathcal{B}^*},
\{\rho_V^U\}_{V \subset U\text{ 且 }U, V \in \mathcal{B}^*})$。
把引理 \ref{cohomology-lemma-uniqueness} 应用于解 $K_{X'}$，可知这个新系统
也满足条件 (3)。对该系统有 $X=X'\cup U_n$，因而归结到上面已经
处理的 $n=2$ 情形。
\end{proof}

\begin{lemma}
\label{cohomology-lemma-glueing-increasing-union}
设 $X$ 为环化空间，$E$ 为良序集，并设
$$
X = \bigcup\nolimits_{\alpha \in E} W_\alpha
$$
是一个开覆盖，满足 $W_\alpha\subset W_{\alpha+1}$，并且当 $\alpha$
不是后继元时，
$W_\alpha=\bigcup_{\beta<\alpha}W_\beta$。设 $K_\alpha$ 为
$D(\mathcal{O}_{W_\alpha})$ 的对象，且在 $i<0$ 时有
$\Ext^i(K_\alpha,K_\alpha)=0$。假设对所有 $\beta<\alpha$ 给定
同构
$\rho_\beta^\alpha :  K_\alpha|_{W_\beta} \to K_\beta$，
它位于 $D(\mathcal{O}_{W_\beta})$ 中；并且当
$\gamma<\beta<\alpha$ 时有
$\rho_\gamma^\alpha = \rho_\gamma^\beta \circ \rho^\alpha_\beta|_{W_\gamma}$
。则存在 $D(\mathcal{O}_X)$ 的对象 $K$，以及对
$\alpha\in E$ 的同构 $K|_{W_\alpha}\to K_\alpha$，它们与同构
$\rho_\beta^\alpha$ 相容。
\end{lemma}

\begin{proof}
在本证明中，$\alpha,\beta,\gamma,\ldots$ 表示 $E$ 的元素。
在 $W_\alpha$ 上选取表示 $K_\alpha$ 的 K-内射复形
$I_\alpha^\bullet$。对 $\beta<\alpha$，以
$j_{\beta,\alpha}:W_\beta\to W_\alpha$ 表示包含态射。我们将用
超限递归，对所有 $\beta<\alpha$ 构造复形态射
$$
\tau_{\beta, \alpha} :
(j_{\beta, \alpha})_!I_\beta^\bullet
\longrightarrow
I_\alpha^\bullet
$$
它表示同构
$\rho^\alpha_\beta:K_\alpha|_{W_\beta}\to K_\beta$ 之逆的伴随态射。
此外，我们将使得当 $\gamma<\beta<\alpha$ 时有
$$
\tau_{\gamma, \alpha} = \tau_{\beta, \alpha} \circ
(j_{\beta, \alpha})_!\tau_{\gamma, \beta}
$$
作为复形态射的等式。具体地，假设已对所有
$\gamma<\beta<\alpha$ 给定能正确复合的 $\tau_{\gamma,\beta}$。
若 $\alpha=\alpha'+1$ 是后继元，则任取复形态射
$$
(j_{\alpha', \alpha})_!I_{\alpha'}^\bullet \to I_\alpha^\bullet
$$
它是同构
$\rho^\alpha_{\alpha'}:K_\alpha|_{W_{\alpha'}}\to K_{\alpha'}$
之逆的伴随态射（由于 $I_\alpha^\bullet$ 是 K-内射的，这样选取是
可能的），并且对任意 $\beta<\alpha'$，置
$$
\tau_{\beta, \alpha} = \tau_{\alpha', \alpha} \circ
(j_{\alpha', \alpha})_!\tau_{\beta, \alpha'}
$$
若 $\alpha$ 不是后继元，则考虑 $W_\alpha$ 上的复形
$$
C^\bullet = \colim_{\beta < \alpha} (j_{\beta, \alpha})_!I_\beta^\bullet
$$
（逐项余极限），其中该系统的转移态射由
$\beta'<\beta<\alpha$ 时的态射 $\tau_{\beta',\beta}$ 给出。
我们断言 $C^\bullet$ 表示 $K_\alpha$。事实上，对
$\beta<\alpha$，余投影
$(j_{\beta,\alpha})_!I_\beta^\bullet\to C^\bullet$ 的限制给出态射
$$
\sigma_\beta : I_\beta^\bullet \longrightarrow C^\bullet|_{W_\beta}
$$
它是拟同构：若 $x\in W_\beta$，则考察茎得到
$$
(C^\bullet)_x =
\colim_{\beta' < \alpha}
\left((j_{\beta', \alpha})_!I_{\beta'}^\bullet\right)_x =
\colim_{\beta \leq \beta' < \alpha} (I_{\beta'}^\bullet)_x
\longleftarrow
(I_\beta^\bullet)_x
$$
这是拟同构。这里使用了取茎与余极限交换、滤过余极限正合，以及当
$\beta\leq\beta'<\alpha$ 时态射
$(I_\beta^\bullet)_x\to(I_{\beta'}^\bullet)_x$ 是拟同构。
因此，$(C^\bullet,\sigma_\beta^{-1})$ 是系统
$(\{K_\beta\}_{\beta<\alpha},
\{\rho^\beta_{\beta'}\}_{\beta'<\beta<\alpha})$ 的解。
由于 $(K_\alpha,\rho^\alpha_\beta)$ 是另一个解，我们得到
$D(\mathcal{O}_{W_\alpha})$ 中与所有态射相容的唯一同构
$\sigma:K_\alpha\to C^\bullet$；见引理
\ref{cohomology-lemma-solution-in-finite-case}（正是在这里用到了负次数 Ext 群
为零）。选取表示 $\sigma$ 的复形态射
$\tau:C^\bullet\to I_\alpha^\bullet$，并置
$$
\tau_{\beta, \alpha} = \tau|_{W_\beta} \circ \sigma_\beta
$$
由此得到所需态射。最后，取 $K$ 为下列复形所表示的导出范畴对象：
$$
K^\bullet = \colim_{\alpha \in E} (W_\alpha \to X)_!I_\alpha^\bullet
$$
其中转移态射由我们对 $\beta<\alpha$ 精心构造的
$\tau_{\beta,\alpha}$ 给出。与上面完全相同的论证表明，对所有
$\alpha$，余投影的限制确定同构
$$
K|_{W_\alpha} \longrightarrow K_\alpha
$$
它与给定态射 $\rho^\alpha_\beta$ 相容。
\end{proof}

\noindent
利用超限归纳，我们可以证明一般情形下的结论。

\begin{theorem}[BBD 粘合引理]
\label{cohomology-theorem-glueing-bbd-general}
\begin{reference}
\cite[Theorem 3.2.4]{BBD} 在不加有界性假设时的特殊情形。
\end{reference}
在情形 \ref{cohomology-situation-locally-given} 中，假设
\begin{enumerate}
\item $X = \bigcup_{U \in \mathcal{B}} U$,
\item 对 $U,V\in\mathcal{B}$，有
$U \cap V = \bigcup_{W \in \mathcal{B}, W \subset U \cap V} W$,
\item 对任意 $U\in\mathcal{B}$，当 $i<0$ 时有
$\Ext^i(K_U,K_U)=0$。
\end{enumerate}
则存在 $D(\mathcal{O}_X)$ 的对象 $K$，以及对
$U\in\mathcal{B}$ 的 $D(\mathcal{O}_U)$ 中的同构
$\rho_U:K|_U\to K_U$，使得对所有满足
$V\subset U$、$U,V\in\mathcal{B}$ 的开集，都有
$\rho^U_V\circ\rho_U|_V=\rho_V$。二元组 $(K,\rho_U)$ 在唯一同构
意义下唯一。
\end{theorem}

\begin{proof}
上文把二元组 $(K,\rho_U)$ 称为解。唯一性来自引理
\ref{cohomology-lemma-uniqueness}。若 $X$ 有由 $\mathcal{B}$ 中元素组成的
有限覆盖（例如 $X$ 拟紧），则本定理是引理
\ref{cohomology-lemma-solution-in-finite-case} 的推论。在一般情形下，用超限归纳
和引理 \ref{cohomology-lemma-glueing-increasing-union} 作完全相同的论证。

\medskip\noindent
首先用超限递归，对任意序数 $\alpha$ 选取开集
$W_\alpha\subset X$。具体地，置 $W_0=\emptyset$。若
$\alpha=\beta+1$ 是后继序数，则当 $W_\beta=X$ 时置
$W_\alpha=X$；当 $W_\beta \not = X$ 时，取不包含于 $W_\beta$ 的
$U_\alpha\in\mathcal{B}$，并置
$W_\alpha=W_\beta\cup U_\alpha$。若 $\alpha$ 是极限序数，则置
$W_\alpha=\bigcup_{\beta<\alpha}W_\beta$。于是当 $\alpha$ 足够大时，
$W_\alpha=X$。注意，对每个 $\alpha$，开集 $W_\alpha$ 都是
$\mathcal{B}$ 中元素之并。因此，若
$\mathcal{B}_\alpha=\{U\in\mathcal{B},U\subset W_\alpha\}$，则
$$
S_\alpha = (\{K_U\}_{U \in \mathcal{B}_\alpha},
\{\rho_V^U\}_{V \subset U\text{ 且 }U, V \in \mathcal{B}_\alpha})
$$
是环化空间 $W_\alpha$ 上引理 \ref{cohomology-lemma-uniqueness} 所述的系统。

\medskip\noindent
我们用超限归纳证明每个 $\alpha$ 的系统 $S_\alpha$ 都有解。
由于 $\alpha$ 足够大时该系统就是定理中给定的系统，这将证明定理。

\medskip\noindent
设 $\alpha=\beta+1$ 为后继序数。（上面引理的证明已经处理过这个
情形，但为清楚起见，我们重复该论证。）此时
$W_\alpha=W_\beta\cup U_\alpha$，其中
$U_\alpha\in\mathcal{B}$。由归纳假设，系统 $S_\beta$ 有解
$(K_{W_\beta}, \{\rho^{W_\beta}_U\}_{U \in \mathcal{B}_\beta})$
。于是可以考虑 $W_\alpha$ 的开集族
$\mathcal{B}_\alpha^* = \mathcal{B}_\alpha \cup \{W_\beta\}$
，并得到系统
$(\{K_U\}_{U \in \mathcal{B}_\alpha^*},
\{\rho_V^U\}_{V \subset U\text{ 且 }U, V \in \mathcal{B}_\alpha^*})$。
把引理 \ref{cohomology-lemma-uniqueness} 应用于解 $K_{W_\beta}$，可知这个
新系统也满足条件 (3)。对该系统有
$W_\alpha=W_\beta\cup U_\alpha$，因而归结到引理
\ref{cohomology-lemma-solution-in-finite-case} 已处理的情形。

\medskip\noindent
设 $\alpha$ 为极限序数。回忆此时
$W_\alpha=\bigcup_{\beta<\alpha}W_\beta$。对 $\beta<\alpha$，令
$(K_{W_\beta}, \{\rho^{W_\beta}_U\}_{U \in \mathcal{B}_\beta})$
为 $S_\beta$ 的解。当 $\gamma<\beta<\alpha$ 时，限制
$K_{W_\beta}|_{W_\gamma}$ 连同对
$U\in\mathcal{B}_\gamma$ 的态射 $\rho^{W_\beta}_U$，构成
$S_\gamma$ 的解。由唯一性，得到唯一同构
$\rho_{W_\gamma}^{W_\beta} : K_{W_\beta}|_{W_\gamma} \to K_{W_\gamma}$
它与对 $U\in\mathcal{B}_\gamma$ 的态射 $\rho^{W_\beta}_U$ 和
$\rho^{W_\gamma}_U$ 相容。这些态射以正确方式复合；也就是说，当
$\delta<\gamma<\beta<\alpha$ 时，
$\rho_{W_\delta}^{W_\gamma}\circ
\rho_{W_\gamma}^{W_\beta}|_{W_\delta}
=\rho^{W_\delta}_{W_\beta}$。因此，可以应用引理
\ref{cohomology-lemma-glueing-increasing-union}（注意，把引理
\ref{cohomology-lemma-uniqueness} 应用于解 $K_{W_\beta}$，可知
$K_{W_\beta}$ 的负次数 Ext 群为零），得到 $K_{W_\alpha}$ 和同构
$$
\rho_{W_\beta}^{W_\alpha} :
K_{W_\alpha}|_{W_\beta}
\longrightarrow
K_{W_\beta}
$$
它们与 $\gamma<\beta<\alpha$ 时的态射
$\rho_{W_\gamma}^{W_\beta}$ 相容。

\medskip\noindent
为证明 $K_{W_\alpha}$ 是解，还需对 $U\in\mathcal{B}_\alpha$
构造满足某些相容性的同构
$\rho_U^{W_\alpha}:K_{W_\alpha}|_U\to K_U$。我们把
$\rho_U^{W_\alpha}$ 选为唯一满足如下条件的态射：对任意
$\beta<\alpha$ 以及任意满足 $V\subset U$ 的
$V\in\mathcal{B}_\beta$，图表
$$
\xymatrix{
K_{W_\alpha}|_V \ar[r]_{\rho_U^{W_\alpha}|_V}
\ar[d]_{\rho_{W_\beta}^{W_\alpha}|_V}
& K_U|_V \ar[d]^{\rho_U^V} \\
K_{W_\beta} \ar[r]^{\rho_V^{W_\beta}}
& K_V
}
$$
交换。这一选取有意义，因为
$$
(\{K_V\}_{V \subset U, V \in \mathcal{B}_\beta\text{ 对某个 }\beta < \alpha},
\{\rho_V^{V'}\}_{V \subset V'\text{ 且 }V, V' \subset U
\text{，且 }V, V' \in \mathcal{B}_\beta\text{ 对某个 }\beta < \alpha})
$$
是环化空间 $U$ 上引理 \ref{cohomology-lemma-uniqueness} 所述的系统，并且
$(K_U,\rho^U_V)$ 与
$(K_{W_\alpha}|_U,  \rho_V^{W_\beta}\circ \rho_{W_\beta}^{W_\alpha}|_V)$
都是该系统的解。这给出存在性和唯一性。略去这些态射满足所需相容性的
证明（这只是记号整理）。
\end{proof}



\section{严格完美复形}
\label{cohomology-section-strictly-perfect}

\noindent
模的严格完美复形将在后面用于定义伪凝聚复形和完美复形。其定义如下。

\begin{definition}
\label{cohomology-definition-strictly-perfect}
设 $(X,\mathcal{O}_X)$ 为环化空间，并设 $\mathcal{E}^\bullet$
为 $\mathcal{O}_X$-模复形。若仅有有限多个 $i$ 使
$\mathcal{E}^i$ 非零，并且对所有 $i$，$\mathcal{E}^i$ 都是某个
有限自由 $\mathcal{O}_X$-模的直和项，则称
$\mathcal{E}^\bullet$ 是{\it 严格完美的}。
\end{definition}

\noindent
警告：由于我们没有假设 $X$ 是局部环化空间，有限自由
$\mathcal{O}_X$-模的直和项未必是有限局部自由的。

\begin{lemma}
\label{cohomology-lemma-cone}
严格完美复形之间态射的锥是严格完美的。
\end{lemma}

\begin{proof}
这由定义立即可得。
\end{proof}

\begin{lemma}
\label{cohomology-lemma-tensor}
两个严格完美复形的张量积所对应的全复形是严格完美的。
\end{lemma}

\begin{proof}
从略。
\end{proof}

\begin{lemma}
\label{cohomology-lemma-strictly-perfect-pullback}
设 $f:(X,\mathcal{O}_X)\to(Y,\mathcal{O}_Y)$ 为环化空间的态射。
若 $\mathcal{F}^\bullet$ 是严格完美的 $\mathcal{O}_Y$-模复形，
则 $f^*\mathcal{F}^\bullet$ 是严格完美的
$\mathcal{O}_X$-模复形。
\end{lemma}

\begin{proof}
有限自由模的拉回仍是有限自由的。函子 $f^*$ 是加性函子，因而保持
直和项。引理得证。
\end{proof}

\begin{lemma}
\label{cohomology-lemma-local-lift-map}
设 $(X,\mathcal{O}_X)$ 为环化空间。给定 $\mathcal{O}_X$-模的实线图表
$$
\xymatrix{
\mathcal{E} \ar@{..>}[dr] \ar[r] & \mathcal{F} \\
& \mathcal{G} \ar[u]_p
}
$$
其中 $\mathcal{E}$ 是某个有限自由 $\mathcal{O}_X$-模的直和项，
并且 $p$ 满射，则在 $X$ 上局部存在使图表交换的虚线箭头。
\end{lemma}

\begin{proof}
可以假设对某个 $n$ 有 $\mathcal{E}=\mathcal{O}_X^{\oplus n}$。
此时，寻找虚线箭头等价于提升基元素在
$\Gamma(X,\mathcal{F})$ 中的像。由层满射的刻画（《层》第 \ref{sheaves-section-exactness-points} 节），这在局部是可行的。
\end{proof}

\begin{lemma}
\label{cohomology-lemma-local-homotopy}
设 $(X,\mathcal{O}_X)$ 为环化空间。
\begin{enumerate}
\item 设 $\alpha:\mathcal{E}^\bullet\to\mathcal{F}^\bullet$
为 $\mathcal{O}_X$-模复形的态射，其中 $\mathcal{E}^\bullet$
严格完美，而 $\mathcal{F}^\bullet$ 无上同调。则 $\alpha$
在 $X$ 上局部零同伦。
\item 设 $\alpha:\mathcal{E}^\bullet\to\mathcal{F}^\bullet$
为 $\mathcal{O}_X$-模复形的态射，其中 $\mathcal{E}^\bullet$
严格完美，且当 $i<a$ 时有 $\mathcal{E}^i=0$，当 $i\geq a$ 时有
$H^i(\mathcal{F}^\bullet)=0$。则 $\alpha$ 在 $X$ 上局部零同伦。
\end{enumerate}
\end{lemma}

\begin{proof}
第一条由第二条推出，故只证明 (2)。我们对复形
$\mathcal{E}^\bullet$ 的长度作归纳。若对某个有限自由
$\mathcal{O}_X$-模的直和项 $\mathcal{E}$ 及整数 $n\geq a$，有
$\mathcal{E}^\bullet\cong\mathcal{E}[-n]$，则结论来自引理
\ref{cohomology-lemma-local-lift-map} 以及如下事实：
$\mathcal{F}^{n - 1} \to \Ker(\mathcal{F}^n \to \mathcal{F}^{n + 1})$
由所假设的 $H^n(\mathcal{F}^\bullet)$ 为零而满射。若
$\mathcal{E}^i$ 仅在 $i\in[a,b]$ 时可能非零，则有复形的分裂正合列
$$
0 \to \mathcal{E}^b[-b] \to \mathcal{E}^\bullet \to
\sigma_{\leq b - 1}\mathcal{E}^\bullet \to 0
$$
它在 $K(\mathcal{O}_X)$ 中确定一个特异三角。因此，由三角范畴的
公理得到正合列
$$
\Hom_{K(\mathcal{O}_X)}(
\sigma_{\leq b - 1}\mathcal{E}^\bullet, \mathcal{F}^\bullet)
\to
\Hom_{K(\mathcal{O}_X)}(\mathcal{E}^\bullet, \mathcal{F}^\bullet)
\to
\Hom_{K(\mathcal{O}_X)}(\mathcal{E}^b[-b], \mathcal{F}^\bullet)
$$
。复合 $\mathcal{E}^b[-b]\to\mathcal{F}^\bullet$ 在局部零同伦，
故可以假设我们的态射来自上述显示正合列左端的一个元素。由归纳假设，
这个元素在局部为零。
\end{proof}

\begin{lemma}
\label{cohomology-lemma-lift-through-quasi-isomorphism}
设 $(X,\mathcal{O}_X)$ 为环化空间。给定
$\mathcal{O}_X$-模复形的实线图表
$$
\xymatrix{
\mathcal{E}^\bullet \ar@{..>}[dr] \ar[r]_\alpha & \mathcal{F}^\bullet \\
& \mathcal{G}^\bullet \ar[u]_f
}
$$
其中 $\mathcal{E}^\bullet$ 严格完美，当 $j<a$ 时
$\mathcal{E}^j=0$，而 $H^j(f)$ 在 $j>a$ 时为同构、在 $j=a$ 时
满射，则在 $X$ 上局部存在使图表在同伦意义下交换的虚线箭头。
\end{lemma}

\begin{proof}
对 $f$ 的假设蕴含锥 $C(f)^\bullet$ 的次数 $\geq a$ 的上同调层
为零。因此，由引理 \ref{cohomology-lemma-local-homotopy}，存在开覆盖
$X=\bigcup U_i$，使得复合
$\mathcal{E}^\bullet \to \mathcal{F}^\bullet \to C(f)^\bullet$
在 $U_i$ 上零同伦。由于
$$
\mathcal{G}^\bullet \to \mathcal{F}^\bullet \to C(f)^\bullet \to
\mathcal{G}^\bullet[1]
$$
限制后成为 $K(\mathcal{O}_{U_i})$ 中的特异三角，故可把
$\alpha|_{U_i}$ 在同伦意义下提升为所需态射
$\alpha_i:\mathcal{E}^\bullet|_{U_i}\to\mathcal{G}^\bullet|_{U_i}$。
\end{proof}

\begin{lemma}
\label{cohomology-lemma-local-actual}
设 $(X,\mathcal{O}_X)$ 为环化空间。设
$\mathcal{E}^\bullet$、$\mathcal{F}^\bullet$ 为
$\mathcal{O}_X$-模复形，并且 $\mathcal{E}^\bullet$ 严格完美。
\begin{enumerate}
\item 对任意元素
$\alpha \in \Hom_{D(\mathcal{O}_X)}(\mathcal{E}^\bullet, \mathcal{F}^\bullet)$
，存在开覆盖 $X=\bigcup U_i$，使得 $\alpha|_{U_i}$ 由复形态射
$\alpha_i:\mathcal{E}^\bullet|_{U_i}\to
\mathcal{F}^\bullet|_{U_i}$ 给出。
\item 给定复形态射
$\alpha:\mathcal{E}^\bullet\to\mathcal{F}^\bullet$，若它在群
$\Hom_{D(\mathcal{O}_X)}(\mathcal{E}^\bullet, \mathcal{F}^\bullet)$
中的像为零，则存在开覆盖 $X=\bigcup U_i$，使得
$\alpha|_{U_i}$ 零同伦。
\end{enumerate}
\end{lemma}

\begin{proof}
证明 (1)。由导出范畴的构造，可以找到拟同构
$f:\mathcal{F}^\bullet\to\mathcal{G}^\bullet$ 以及复形态射
$\beta:\mathcal{E}^\bullet\to\mathcal{G}^\bullet$，使得
$\alpha=f^{-1}\beta$。因此，结论来自引理
\ref{cohomology-lemma-lift-through-quasi-isomorphism}。(2) 的证明从略。
\end{proof}

\begin{lemma}
\label{cohomology-lemma-Rhom-strictly-perfect}
设 $(X,\mathcal{O}_X)$ 为环化空间。设
$\mathcal{E}^\bullet$、$\mathcal{F}^\bullet$ 为
$\mathcal{O}_X$-模复形，并且 $\mathcal{E}^\bullet$ 严格完美。
则内 Hom $R\SheafHom(\mathcal{E}^\bullet,\mathcal{F}^\bullet)$
由复形 $\mathcal{H}^\bullet$ 表示，其各项为
$$
\mathcal{H}^n =
\bigoplus\nolimits_{n = p + q}
\SheafHom_{\mathcal{O}_X}(\mathcal{E}^{-q}, \mathcal{F}^p)
$$
，微分如第 \ref{cohomology-section-hom-complexes} 节所述。
\end{lemma}

\begin{proof}
选取到 K-内射复形的拟同构
$\mathcal{F}^\bullet\to\mathcal{I}^\bullet$。令
$(\mathcal{H}')^\bullet$ 为各项如下的复形：
$$
(\mathcal{H}')^n =
\prod\nolimits_{n = p + q}
\SheafHom_{\mathcal{O}_X}(\mathcal{E}^{-q}, \mathcal{I}^p)
$$
由第 \ref{cohomology-section-internal-hom} 节的构造，它表示
$R\SheafHom(\mathcal{E}^\bullet,\mathcal{F}^\bullet)$。只需证明态射
$$
\mathcal{H}^\bullet \longrightarrow (\mathcal{H}')^\bullet
$$
是拟同构。给定开集 $U\subset X$，直接考察可得
$$
H^0(\mathcal{H}^\bullet(U)) =
\Hom_{K(\mathcal{O}_U)}(\mathcal{E}^\bullet|_U, \mathcal{I}^\bullet|_U)
\to
H^0((\mathcal{H}')^\bullet(U)) =
\Hom_{D(\mathcal{O}_U)}(\mathcal{E}^\bullet|_U, \mathcal{I}^\bullet|_U)
$$
由引理 \ref{cohomology-lemma-local-actual}，对应
$U\mapsto H^0(\mathcal{H}^\bullet(U))$ 的层化等于对应
$U\mapsto H^0((\mathcal{H}')^\bullet(U))$ 的层化。其他上同调层也可作
类似论证。因此 $\mathcal{H}^\bullet$ 与
$(\mathcal{H}')^\bullet$ 拟同构，引理得证。
\end{proof}

\begin{lemma}
\label{cohomology-lemma-Rhom-strictly-perfect-K-flat}
在引理 \ref{cohomology-lemma-Rhom-strictly-perfect} 的情形下，若
$\mathcal{F}^\bullet$ 是 K-平坦的，则
$\mathcal{H}^\bullet$ 也是 K-平坦的。
\end{lemma}

\begin{proof}
注意，$\mathcal{H}^\bullet$ 就是 Hom 复形
$\SheafHom^\bullet(\mathcal{E}^\bullet, \mathcal{F}^\bullet)$
，因为严格完美复形 $\mathcal{E}^\bullet$ 的有界性保证 Hom 复形
定义中的乘积变成直和。设 $\mathcal{K}^\bullet$ 为无上同调的
$\mathcal{O}_X$-模复形。考虑引理 \ref{cohomology-lemma-diagonal-better} 中的态射
$$
\gamma :
\text{Tot}(\mathcal{K}^\bullet \otimes
\SheafHom^\bullet(\mathcal{E}^\bullet, \mathcal{F}^\bullet))
\longrightarrow
\SheafHom^\bullet(\mathcal{E}^\bullet,
\text{Tot}(\mathcal{K}^\bullet \otimes \mathcal{F}^\bullet))
$$
。由于 $\mathcal{F}^\bullet$ 是 K-平坦的，复形
$\text{Tot}(\mathcal{K}^\bullet\otimes\mathcal{F}^\bullet)$ 无上同调；
因而由引理 \ref{cohomology-lemma-local-actual}（或者也可用引理
\ref{cohomology-lemma-Rhom-strictly-perfect}），$\gamma$ 的目标也无上同调。
所以，为证明本引理，只需证明 $\gamma$ 是复形同构。为此，可以对
复形 $\mathcal{E}^\bullet$ 的长度作归纳。若长度 $\leq1$，则对某个
$n\geq0$ 和 $k\in\mathbf{Z}$，$\mathcal{E}^\bullet$ 是
$\mathcal{O}_X^{\oplus n}[k]$ 的直和项，此时直接考察即可得到结论。
若长度 $>1$，则考虑逐项分裂的复形短正合列
$$
0 \to \sigma_{\geq a + 1} \mathcal{E}^\bullet \to
\mathcal{E}^\bullet \to
\sigma_{\leq a} \mathcal{E}^\bullet \to 0
$$
，其中 $a\in\mathbf{Z}$ 适当；见 《同调代数》第 \ref{homology-section-truncations} 节。于是 $\gamma$ 嵌入逐项
分裂的复形短正合列之间的态射。由归纳假设，$\gamma$ 对
$\sigma_{\geq a+1}\mathcal{E}^\bullet$ 和
$\sigma_{\leq a}\mathcal{E}^\bullet$ 都是同构，因而对
$\mathcal{E}^\bullet$ 的结论也成立。略去一些细节。
\end{proof}

\begin{lemma}
\label{cohomology-lemma-Rhom-complex-of-direct-summands-finite-free}
设 $(X,\mathcal{O}_X)$ 为环化空间。设
$\mathcal{E}^\bullet$、$\mathcal{F}^\bullet$ 为
$\mathcal{O}_X$-模复形，满足
\begin{enumerate}
\item 当 $n\ll0$ 时，$\mathcal{F}^n=0$；
\item 当 $n\gg0$ 时，$\mathcal{E}^n=0$；并且
\item $\mathcal{E}^n$ 同构于某个有限自由
$\mathcal{O}_X$-模的直和项。
\end{enumerate}
则内 Hom $R\SheafHom(\mathcal{E}^\bullet,\mathcal{F}^\bullet)$
由复形 $\mathcal{H}^\bullet$ 表示，其各项为
$$
\mathcal{H}^n =
\bigoplus\nolimits_{n = p + q}
\SheafHom_{\mathcal{O}_X}(\mathcal{E}^{-q}, \mathcal{F}^p)
$$
，微分如第 \ref{cohomology-section-internal-hom} 节所述。
\end{lemma}

\begin{proof}
选取拟同构 $\mathcal{F}^\bullet\to\mathcal{I}^\bullet$，其中
$\mathcal{I}^\bullet$ 是下有界的内射复形。注意
$\mathcal{I}^\bullet$ 是 K-内射的（《导出范畴》引理
\ref{derived-lemma-bounded-below-injectives-K-injective}）。因此，
第 \ref{cohomology-section-internal-hom} 节的构造表明
$R\SheafHom(\mathcal{E}^\bullet,\mathcal{F}^\bullet)$ 由复形
$(\mathcal{H}')^\bullet$ 表示，其各项为
$$
(\mathcal{H}')^n =
\prod\nolimits_{n = p + q}
\SheafHom_{\mathcal{O}_X}(\mathcal{E}^{-q}, \mathcal{I}^p) =
\bigoplus\nolimits_{n = p + q}
\SheafHom_{\mathcal{O}_X}(\mathcal{E}^{-q}, \mathcal{I}^p)
$$
（之所以相等，是因为仅有有限多个非零项）。注意
$\mathcal{H}^\bullet$ 是各项为
$\SheafHom_{\mathcal{O}_X}(\mathcal{E}^{-q}, \mathcal{F}^p)$
的双复形所对应的全复形；$(\mathcal{H}')^\bullet$ 也类似。
自然态射 $(\mathcal{H}')^\bullet\to\mathcal{H}^\bullet$
来自双复形的态射。因此，为证明这个态射是拟同构，可以使用双复形的
谱序列（《同调代数》引理 \ref{homology-lemma-first-quadrant-ss}）
$$
{}'E_1^{p, q} =
H^p(\SheafHom_{\mathcal{O}_X}(\mathcal{E}^{-q}, \mathcal{F}^\bullet))
$$
，它收敛到 $H^{p+q}(\mathcal{H}^\bullet)$；
$(\mathcal{H}')^\bullet$ 也类似。为完成本引理的证明，只需证明
$\mathcal{F}^\bullet\to\mathcal{I}^\bullet$ 诱导上同调层的同构
$$
H^p(\SheafHom_{\mathcal{O}_X}(\mathcal{E}, \mathcal{F}^\bullet))
\longrightarrow
H^p(\SheafHom_{\mathcal{O}_X}(\mathcal{E}, \mathcal{I}^\bullet))
$$
，其中 $\mathcal{E}$ 是某个有限自由 $\mathcal{O}_X$-模的直和项。
当 $\mathcal{E}$ 有限自由时这是显然的，故结论成立。
\end{proof}





\section{伪凝聚模}
\label{cohomology-section-pseudo-coherent}

\noindent
本节讨论伪凝聚复形。

\begin{definition}
\label{cohomology-definition-pseudo-coherent}
设 $(X,\mathcal{O}_X)$ 为环化空间，$\mathcal{E}^\bullet$ 为
$\mathcal{O}_X$-模复形，并设 $m\in\mathbf{Z}$。
\begin{enumerate}
\item 若存在开覆盖 $X=\bigcup U_i$，并且对每个 $i$ 存在复形态射
$\alpha_i : \mathcal{E}_i^\bullet \to \mathcal{E}^\bullet|_{U_i}$
，其中 $\mathcal{E}_i^\bullet$ 在 $U_i$ 上严格完美，并且
$H^j(\alpha_i)$ 在 $j>m$ 时为同构，而 $H^m(\alpha_i)$ 满射，
则称 $\mathcal{E}^\bullet$ 是{\it $m$-伪凝聚的}。
\item 若 $\mathcal{E}^\bullet$ 对所有 $m$ 都是 $m$-伪凝聚的，
则称它是{\it 伪凝聚的}。
\item 当且仅当 $D(\mathcal{O}_X)$ 的对象 $E$ 可由一个
$m$-伪凝聚（相应地，伪凝聚）的 $\mathcal{O}_X$-模复形表示时，
称该对象是{\it $m$-伪凝聚的}（相应地，{\it 伪凝聚的}）。
\end{enumerate}
\end{definition}

\noindent
若 $X$ 拟紧，则 $D(\mathcal{O}_X)$ 的 $m$-伪凝聚对象属于
$D^-(\mathcal{O}_X)$。但若 $X$ 不拟紧，则未必如此。

\begin{lemma}
\label{cohomology-lemma-pseudo-coherent-independent-representative}
设 $(X,\mathcal{O}_X)$ 为环化空间，并设 $E$ 为
$D(\mathcal{O}_X)$ 的对象。
\begin{enumerate}
\item 若存在开覆盖 $X=\bigcup U_i$、$U_i$ 上的严格完美复形
$\mathcal{E}_i^\bullet$，以及 $D(\mathcal{O}_{U_i})$ 中的态射
$\alpha_i:\mathcal{E}_i^\bullet\to E|_{U_i}$，使得
$H^j(\alpha_i)$ 在 $j>m$ 时为同构，而 $H^m(\alpha_i)$ 满射，
则 $E$ 是 $m$-伪凝聚的。
\item 若 $E$ 是 $m$-伪凝聚的，则表示 $E$ 的任意复形都是
$m$-伪凝聚的。
\end{enumerate}
\end{lemma}

\begin{proof}
任取表示 $E$ 的复形 $\mathcal{F}^\bullet$，并设
$X=\bigcup U_i$ 及
$\alpha_i:\mathcal{E}_i^\bullet\to E|_{U_i}$ 如 (1) 所述。
我们将证明复形 $\mathcal{F}^\bullet$ 是 $m$-伪凝聚的，这将同时证明
(1) 和 (2)。由引理 \ref{cohomology-lemma-local-actual}，细化开覆盖
$X=\bigcup U_i$ 后，可以用复形态射
$\alpha_i:\mathcal{E}_i^\bullet\to\mathcal{F}^\bullet|_{U_i}$
表示各态射 $\alpha_i$。由假设，$H^j(\alpha_i)$ 在 $j>m$ 时为同构，
而 $H^m(\alpha_i)$ 满射，故 $\mathcal{F}^\bullet$ 是
$m$-伪凝聚的。
\end{proof}

\begin{lemma}
\label{cohomology-lemma-pseudo-coherent-pullback}
设 $f:(X,\mathcal{O}_X)\to(Y,\mathcal{O}_Y)$ 为环化空间的态射，
并设 $E$ 为 $D(\mathcal{O}_Y)$ 的对象。若 $E$ 是
$m$-伪凝聚的，则 $Lf^*E$ 也是 $m$-伪凝聚的。
\end{lemma}

\begin{proof}
用 $\mathcal{O}_Y$-模复形 $\mathcal{E}^\bullet$ 表示 $E$，并按定义
\ref{cohomology-definition-pseudo-coherent} 选取开覆盖 $Y=\bigcup V_i$ 和态射
$\alpha_i : \mathcal{E}_i^\bullet \to \mathcal{E}^\bullet|_{V_i}$
。置 $U_i=f^{-1}(V_i)$。由引理
\ref{cohomology-lemma-pseudo-coherent-independent-representative}，只需证明
$Lf^*\mathcal{E}^\bullet|_{U_i}$ 是 $m$-伪凝聚的。选取特异三角
$$
\mathcal{E}_i^\bullet \to
\mathcal{E}^\bullet|_{V_i} \to
C \to
\mathcal{E}_i^\bullet[1]
$$
对 $\alpha_i$ 的假设恰好意味着对所有 $j\geq m$，上同调层
$H^j(C)$ 为零。以 $f_i:U_i\to V_i$ 表示 $f$ 的限制。注意
$Lf^*\mathcal{E}^\bullet|_{U_i}=
Lf_i^*(\mathcal{E}|_{V_i})$。应用 $Lf_i^*$，得到特异三角
$$
Lf_i^*\mathcal{E}_i^\bullet \to
Lf_i^*\mathcal{E}|_{V_i} \to
Lf_i^*C \to
Lf_i^*\mathcal{E}_i^\bullet[1]
$$
由 $Lf_i^*$ 作为左导出函子的构造可知，对 $j\geq m$，
$H^j(Lf_i^*C)=0$（使用 《导出范畴》引理
\ref{derived-lemma-negative-vanishing} 的对偶）。因此
$H^j(Lf_i^*\alpha_i)$ 在 $j>m$ 时为同构，而
$H^m(Lf^*\alpha_i)$ 满射。另一方面，
$Lf_i^*\mathcal{E}_i^\bullet = f_i^*\mathcal{E}_i^\bullet$.
由引理 \ref{cohomology-lemma-strictly-perfect-pullback}，它严格完美。结论得证。
\end{proof}

\begin{lemma}
\label{cohomology-lemma-cone-pseudo-coherent}
设 $(X,\mathcal{O}_X)$ 为环化空间，$m\in\mathbf{Z}$，并设
$(K,L,M,f,g,h)$ 为 $D(\mathcal{O}_X)$ 中的特异三角。
\begin{enumerate}
\item 若 $K$ 是 $(m+1)$-伪凝聚的，而 $L$ 是 $m$-伪凝聚的，
则 $M$ 是 $m$-伪凝聚的。
\item 若 $K$ 与 $M$ 都是 $m$-伪凝聚的，则 $L$ 是 $m$-伪凝聚的。
\item 若 $L$ 是 $(m+1)$-伪凝聚的，而 $M$ 是 $m$-伪凝聚的，
则 $K$ 是 $(m+1)$-伪凝聚的。
\end{enumerate}
\end{lemma}

\begin{proof}
证明 (1)。选取开覆盖 $X=\bigcup U_i$，以及
$D(\mathcal{O}_{U_i})$ 中的态射
$\alpha_i:\mathcal{K}_i^\bullet\to K|_{U_i}$，使
$\mathcal{K}_i^\bullet$ 严格完美，且 $H^j(\alpha_i)$ 在
$j>m+1$ 时为同构、在 $j=m+1$ 时满射。可以把
$\mathcal{K}_i^\bullet$ 换成
$\sigma_{\geq m+1}\mathcal{K}_i^\bullet$，故可以假设当 $j<m+1$ 时
$\mathcal{K}_i^j=0$。细化开覆盖后，可以选取
$D(\mathcal{O}_{U_i})$ 中的态射
$\beta_i:\mathcal{L}_i^\bullet\to L|_{U_i}$，使
$\mathcal{L}_i^\bullet$ 严格完美，且 $H^j(\beta)$ 在 $j>m$ 时
为同构、在 $j=m$ 时满射。由引理
\ref{cohomology-lemma-lift-through-quasi-isomorphism}，再次细化覆盖后，可以找到
复形态射 $\gamma_i:\mathcal{K}^\bullet\to\mathcal{L}^\bullet$，
使图表
$$
\xymatrix{
K|_{U_i} \ar[r] & L|_{U_i} \\
\mathcal{K}_i^\bullet \ar[u]^{\alpha_i} \ar[r]^{\gamma_i} &
\mathcal{L}_i^\bullet \ar[u]_{\beta_i}
}
$$
在 $D(\mathcal{O}_{U_i})$ 中交换（这要求用实际复形态射表示
$\alpha_i$、$\beta_i$ 以及 $K|_{U_i}\to L|_{U_i}$；略去一些细节）。
锥 $C(\gamma_i)^\bullet$ 严格完美（引理 \ref{cohomology-lemma-cone}）。
图表的交换性蕴含存在特异三角之间的态射
$$
(\mathcal{K}_i^\bullet, \mathcal{L}_i^\bullet, C(\gamma_i)^\bullet)
\longrightarrow
(K|_{U_i}, L|_{U_i}, M|_{U_i}).
$$
由所诱导的上同调长正合列之间的态射，以及 《同调代数》引理
\ref{homology-lemma-four-lemma} 和 \ref{homology-lemma-five-lemma}，
可知 $C(\gamma_i)^\bullet\to M|_{U_i}$ 在次数 $>m$ 的上同调上诱导
同构，在次数 $m$ 上诱导满射。因此，由引理
\ref{cohomology-lemma-pseudo-coherent-independent-representative}，$M$ 是
$m$-伪凝聚的。

\medskip\noindent
把特异三角旋转后，由 (1) 即得断言 (2) 和 (3)。
\end{proof}

\begin{lemma}
\label{cohomology-lemma-tensor-pseudo-coherent}
设 $(X,\mathcal{O}_X)$ 为环化空间，并设 $K,L$ 为
$D(\mathcal{O}_X)$ 的对象。
\begin{enumerate}
\item 若 $K$ 是 $n$-伪凝聚的，且当 $i>a$ 时有 $H^i(K)=0$；
而 $L$ 是 $m$-伪凝聚的，且当 $j>b$ 时有 $H^j(L)=0$，则
$K\otimes_{\mathcal{O}_X}^\mathbf{L}L$ 是 $t$-伪凝聚的，其中
$t=\max(m+a,n+b)$。
\item 若 $K$ 和 $L$ 都是伪凝聚的，则
$K\otimes_{\mathcal{O}_X}^\mathbf{L}L$ 是伪凝聚的。
\end{enumerate}
\end{lemma}

\begin{proof}
证明 (1)。用开覆盖的各成员替换 $X$ 后，可以假设存在严格完美复形
$\mathcal{K}^\bullet$ 和 $\mathcal{L}^\bullet$，以及态射
$\alpha:\mathcal{K}^\bullet\to K$ 和
$\beta:\mathcal{L}^\bullet\to L$，使得 $H^i(\alpha)$ 在 $i>n$
时为同构、在 $i=n$ 时满射，而 $H^i(\beta)$ 在 $i>m$ 时为同构、
在 $i=m$ 时满射。于是态射
$$
\alpha \otimes^\mathbf{L} \beta :
\text{Tot}(\mathcal{K}^\bullet \otimes_{\mathcal{O}_X} \mathcal{L}^\bullet)
\to K \otimes_{\mathcal{O}_X}^\mathbf{L} L
$$
在 $i>t$ 时的第 $i$ 个上同调层上诱导同构，在 $i=t$ 时诱导满射。
这由 Tor 的谱序列得出（细节从略）。

\medskip\noindent
证明 (2)。先用某个开覆盖的各成员替换 $X$，便可归结到 $K$ 和 $L$
上有界的情形。此时断言立即由 (1) 得出。
\end{proof}

\begin{lemma}
\label{cohomology-lemma-summands-pseudo-coherent}
设 $(X,\mathcal{O}_X)$ 为环化空间，$m\in\mathbf{Z}$。若
$K\oplus L$ 在 $D(\mathcal{O}_X)$ 中是 $m$-伪凝聚的
（相应地，伪凝聚的），则 $K$ 和 $L$ 也如此。
\end{lemma}

\begin{proof}
假设 $K\oplus L$ 是 $m$-伪凝聚的。用开覆盖的各成员替换 $X$ 后，
可以假设 $K\oplus L\in D^-(\mathcal{O}_X)$，从而
$L\in D^-(\mathcal{O}_X)$。注意，存在特异三角
$$
(K \oplus L, K \oplus L, L \oplus L[1]) =
(K, K, 0) \oplus (L, L, L \oplus L[1])
$$
；见 《导出范畴》引理
\ref{derived-lemma-direct-sum-triangles}。由引理
\ref{cohomology-lemma-cone-pseudo-coherent}，$L\oplus L[1]$ 是
$m$-伪凝聚的，故 $L[1]\oplus L[2]$ 也为 $m$-伪凝聚。归纳可知
$L[n]\oplus L[n+1]$ 是 $m$-伪凝聚的。由于 $L$ 上有界，当 $n$
足够大时，$L[n]$ 是 $m$-伪凝聚的。因此，利用特异三角
$$
(L[n], L[n] \oplus L[n - 1], L[n - 1])
$$
反向递推，便得 $L[n-1],L[n-2],\ldots,L$ 都是 $m$-伪凝聚的，
如所需。
\end{proof}

\begin{lemma}
\label{cohomology-lemma-complex-pseudo-coherent-modules}
设 $(X,\mathcal{O}_X)$ 为环化空间，$m\in\mathbf{Z}$。设
$\mathcal{F}^\bullet$ 为（局部）上有界的 $\mathcal{O}_X$-模复形，
并且对所有 $i$，$\mathcal{F}^i$ 都是 $(m-i)$-伪凝聚的。
则 $\mathcal{F}^\bullet$ 是 $m$-伪凝聚的。
\end{lemma}

\begin{proof}
从略。提示：使用引理 \ref{cohomology-lemma-cone-pseudo-coherent} 和截断，
如 《更多代数》引理
\ref{more-algebra-lemma-complex-pseudo-coherent-modules} 的证明所示。
\end{proof}

\begin{lemma}
\label{cohomology-lemma-cohomology-pseudo-coherent}
设 $(X,\mathcal{O}_X)$ 为环化空间，$m\in\mathbf{Z}$，并设
$E$ 为 $D(\mathcal{O}_X)$ 的对象。若 $E$（局部）上有界，并且对所有
$i$，$H^i(E)$ 都是 $(m-i)$-伪凝聚的，则 $E$ 是 $m$-伪凝聚的。
\end{lemma}

\begin{proof}
从略。提示：使用引理 \ref{cohomology-lemma-cone-pseudo-coherent} 和截断，
如 《更多代数》引理
\ref{more-algebra-lemma-cohomology-pseudo-coherent} 的证明所示。
\end{proof}

\begin{lemma}
\label{cohomology-lemma-finite-cohomology}
设 $(X,\mathcal{O}_X)$ 为环化空间，$K$ 为
$D(\mathcal{O}_X)$ 的对象，并设 $m\in\mathbf{Z}$。
\begin{enumerate}
\item 若 $K$ 是 $m$-伪凝聚的，且当 $i>m$ 时有 $H^i(K)=0$，
则 $H^m(K)$ 是有限型 $\mathcal{O}_X$-模。
\item 若 $K$ 是 $m$-伪凝聚的，且当 $i>m+1$ 时有 $H^i(K)=0$，
则 $H^{m+1}(K)$ 是有限表示的 $\mathcal{O}_X$-模。
\end{enumerate}
\end{lemma}

\begin{proof}
证明 (1)。可以在 $X$ 上局部工作。因此，可以假设存在严格完美复形
$\mathcal{E}^\bullet$ 以及态射
$\alpha:\mathcal{E}^\bullet\to K$，它在次数 $>m$ 的上同调上诱导
同构，在次数 $m$ 上诱导满射。只需对 $\mathcal{E}^\bullet$
证明结论。令 $n$ 为使 $\mathcal{E}^n\not=0$ 的最大整数。若
$n=m$，则 $H^m(\mathcal{E}^\bullet)$ 是 $\mathcal{E}^n$ 的商，
结论显然。若 $n>m$，则由于 $H^n(E^\bullet)=0$，
$\mathcal{E}^{n-1}\to\mathcal{E}^n$ 满射。由引理
\ref{cohomology-lemma-local-lift-map}，可以在局部找到该满射的截面，并写成
$\mathcal{E}^{n-1}=\mathcal{E}'\oplus\mathcal{E}^n$。因此，只需对
复形 $(\mathcal{E}')^\bullet$ 证明结论；它与
$\mathcal{E}^\bullet$ 相同，只是次数 $n-1$ 的项换成
$\mathcal{E}'$，次数 $n$ 的项换成 $0$。对 $n$ 作归纳即得结论。

\medskip\noindent
证明 (2)。可以在 $X$ 上局部工作。因此，可以假设存在严格完美复形
$\mathcal{E}^\bullet$ 以及态射
$\alpha:\mathcal{E}^\bullet\to K$，它在次数 $>m$ 的上同调上诱导
同构，在次数 $m$ 上诱导满射。如 (1) 的证明那样，可以归结到当
$i>m+1$ 时 $\mathcal{E}^i=0$ 的情形。于是
$H^{m + 1}(K) \cong H^{m + 1}(\mathcal{E}^\bullet) =
\Coker(\mathcal{E}^m \to \mathcal{E}^{m + 1})$
，这是有限表示的。
\end{proof}

\begin{lemma}
\label{cohomology-lemma-n-pseudo-module}
设 $(X,\mathcal{O}_X)$ 为环化空间，并设 $\mathcal{F}$ 为
$\mathcal{O}_X$-模层。
\begin{enumerate}
\item 把 $\mathcal{F}$ 视为 $D(\mathcal{O}_X)$ 的对象时，它是
$0$-伪凝聚的，当且仅当 $\mathcal{F}$ 是有限型
$\mathcal{O}_X$-模；并且
\item 把 $\mathcal{F}$ 视为 $D(\mathcal{O}_X)$ 的对象时，它是
$(-1)$-伪凝聚的，当且仅当 $\mathcal{F}$ 是有限表示的
$\mathcal{O}_X$-模。
\end{enumerate}
\end{lemma}

\begin{proof}
一个方向的蕴含使用引理 \ref{cohomology-lemma-finite-cohomology}，另一个方向
使用引理 \ref{cohomology-lemma-cohomology-pseudo-coherent}。
\end{proof}





\section{Tor 维数}
\label{cohomology-section-tor}

\noindent
本节将更仔细地考察由平坦模给出的消解。

\begin{definition}
\label{cohomology-definition-tor-amplitude}
令 $(X, \mathcal{O}_X)$ 为环化空间。
令 $E$ 为 $D(\mathcal{O}_X)$ 的对象。
令 $a, b \in \mathbf{Z}$ 且 $a \leq b$。
\begin{enumerate}
\item 称 $E$ 的 {\it Tor 振幅位于 $[a, b]$}，
如果对每个 $\mathcal{O}_X$-模 $\mathcal{F}$ 以及每个 $i \not \in [a, b]$，都有
$H^i(E \otimes_{\mathcal{O}_X}^\mathbf{L} \mathcal{F}) = 0$。
\item 如果 $E$ 的 Tor 振幅位于 $[a, b]$，其中 $a, b$ 可适当选取，
则称其具有{\it 有限 Tor 维数}。
\item 如果存在开覆盖 $X = \bigcup U_i$，使得
对每个 $i$，$E|_{U_i}$ 都具有有限 Tor 维数，
则称 $E$ {\it 局部具有有限 Tor 维数}。
\end{enumerate}
如果把 $\mathcal{O}_X$-模 $\mathcal{F}$ 的 $\mathcal{F}[0]$
视为 $D(\mathcal{O}_X)$ 的对象时，其 Tor 振幅位于 $[-d, 0]$，
则称该模的 {\it Tor 维数 $\leq d$}。
\end{definition}

\noindent
注意，如果定义中的 $E$ 具有有限 Tor 维数，
那么 $E$ 是 $D^b(\mathcal{O}_X)$ 的对象；
在上述定义中取 $\mathcal{F} = \mathcal{O}_X$ 即可看出这一点。

\begin{lemma}
\label{cohomology-lemma-last-one-flat}
令 $(X, \mathcal{O}_X)$ 为环化空间。
令 $\mathcal{E}^\bullet$ 为由平坦
$\mathcal{O}_X$-模组成的上有界复形，其 Tor 振幅位于 $[a, b]$。
则 $\Coker(d_{\mathcal{E}^\bullet}^{a - 1})$ 是平坦
$\mathcal{O}_X$-模。
\end{lemma}

\begin{proof}
由于 $\mathcal{E}^\bullet$ 是由平坦模组成的上有界复形，对任意
$\mathcal{O}_X$-模 $\mathcal{F}$，都有
$\mathcal{E}^\bullet \otimes_{\mathcal{O}_X} \mathcal{F} =
\mathcal{E}^\bullet \otimes_{\mathcal{O}_X}^{\mathbf{L}} \mathcal{F}$。
因此，对每个 $\mathcal{O}_X$-模 $\mathcal{F}$，序列
$$
\mathcal{E}^{a - 2} \otimes_{\mathcal{O}_X} \mathcal{F} \to
\mathcal{E}^{a - 1} \otimes_{\mathcal{O}_X} \mathcal{F} \to
\mathcal{E}^a \otimes_{\mathcal{O}_X} \mathcal{F}
$$
在中间正合。由于
$\mathcal{E}^{a - 2} \to \mathcal{E}^{a - 1} \to \mathcal{E}^a \to
\Coker(d^{a - 1}) \to 0$
是平坦消解，故
$\text{Tor}_1^{\mathcal{O}_X}(\Coker(d^{a - 1}), \mathcal{F}) = 0$
对所有 $\mathcal{O}_X$-模 $\mathcal{F}$ 成立。这意味着
$\Coker(d^{a - 1})$ 平坦，见引理 \ref{cohomology-lemma-flat-tor-zero}。
\end{proof}

\begin{lemma}
\label{cohomology-lemma-tor-amplitude}
令 $(X, \mathcal{O}_X)$ 为环化空间。
令 $E$ 为 $D(\mathcal{O}_X)$ 的对象。
令 $a, b \in \mathbf{Z}$ 且 $a \leq b$。下列条件等价：
\begin{enumerate}
\item $E$ 的 Tor 振幅位于 $[a, b]$。
\item $E$ 可由平坦 $\mathcal{O}_X$-模组成的复形
$\mathcal{E}^\bullet$ 表示，并且
当 $i \not \in [a, b]$ 时 $\mathcal{E}^i = 0$。
\end{enumerate}
\end{lemma}

\begin{proof}
如果 (2) 成立，则可计算
$E \otimes_{\mathcal{O}_X}^\mathbf{L} \mathcal{F} =
\mathcal{E}^\bullet \otimes_{\mathcal{O}_X} \mathcal{F}$，
显然 (1) 成立。

\medskip\noindent
假设 (1) 成立。根据第 \ref{cohomology-section-flat} 节，
可用由平坦 $\mathcal{O}_X$-模组成的上有界复形
$\mathcal{K}^\bullet$ 表示 $E$。
令 $n$ 为满足 $\mathcal{K}^n \not = 0$ 的最大整数。
如果 $n > b$，则由 $H^n(\mathcal{K}^\bullet) = 0$ 可知
$\mathcal{K}^{n - 1} \to \mathcal{K}^n$ 是满射。由于 $\mathcal{K}^n$ 平坦，
$\Ker(\mathcal{K}^{n - 1} \to \mathcal{K}^n)$ 也是平坦的
（模，引理 \ref{modules-lemma-flat-ses}）。
因此可将 $\mathcal{K}^\bullet$ 替换为
$\tau_{\leq n - 1}\mathcal{K}^\bullet$。于是对 $n$ 归纳后，
可约化到如下情形：$K^\bullet$ 是由平坦
$\mathcal{O}_X$-模组成的复形，并且当 $i > b$ 时 $\mathcal{K}^i = 0$。

\medskip\noindent
令 $\mathcal{E}^\bullet = \tau_{\geq a}\mathcal{K}^\bullet$。
除 $\mathcal{E}^a$ 的平坦性以外，其余都很清楚；
这一平坦性立即由引理 \ref{cohomology-lemma-last-one-flat}
和定义得到。
\end{proof}

\begin{lemma}
\label{cohomology-lemma-tor-amplitude-pullback}
令 $f : (X, \mathcal{O}_X) \to (Y, \mathcal{O}_Y)$ 为环化空间的态射。
令 $E$ 为 $D(\mathcal{O}_Y)$ 的对象。
如果 $E$ 的 Tor 振幅位于 $[a, b]$，那么 $Lf^*E$ 的 Tor 振幅也位于
$[a, b]$。
\end{lemma}

\begin{proof}
假设 $E$ 的 Tor 振幅位于 $[a, b]$。由引理 \ref{cohomology-lemma-tor-amplitude}，
可用由平坦 $\mathcal{O}$-模组成的复形
$\mathcal{E}^\bullet$ 表示 $E$，且
当 $i \not \in [a, b]$ 时 $\mathcal{E}^i = 0$。于是
$f^*\mathcal{E}^\bullet$ 表示 $Lf^*E$。
由模，引理 \ref{modules-lemma-base-change-flat}，
各模 $f^*\mathcal{E}^i$ 都是平坦的。
因此由引理 \ref{cohomology-lemma-tor-amplitude}，
$Lf^*E$ 的 Tor 振幅位于 $[a, b]$。
\end{proof}

\begin{lemma}
\label{cohomology-lemma-tor-amplitude-stalk}
令 $(X, \mathcal{O}_X)$ 为环化空间。
令 $E$ 为 $D(\mathcal{O}_X)$ 的对象。
令 $a, b \in \mathbf{Z}$ 且 $a \leq b$。下列条件等价：
\begin{enumerate}
\item $E$ 的 Tor 振幅位于 $[a, b]$。
\item 对每个 $x \in X$，$D(\mathcal{O}_{X, x})$ 的对象 $E_x$
的 Tor 振幅位于 $[a, b]$。
\end{enumerate}
\end{lemma}

\begin{proof}
在 $x$ 处取茎，等同于沿环化空间态射
$(x, \mathcal{O}_{X, x}) \to (X, \mathcal{O}_X)$ 拉回。
因此 (1) $\Rightarrow$ (2) 由引理 \ref{cohomology-lemma-tor-amplitude-pullback} 得出。
反之，注意取茎与张量积可交换
（模，引理 \ref{modules-lemma-stalk-tensor-product}）。因此
$$
(E \otimes_{\mathcal{O}_X}^\mathbf{L} \mathcal{F})_x =
E_x \otimes_{\mathcal{O}_{X, x}}^\mathbf{L} \mathcal{F}_x
$$
另一方面，取茎是正合的，所以
$$
H^i(E \otimes_{\mathcal{O}_X}^\mathbf{L} \mathcal{F})_x =
H^i((E \otimes_{\mathcal{O}_X}^\mathbf{L} \mathcal{F})_x) =
H^i(E_x \otimes_{\mathcal{O}_{X, x}}^\mathbf{L} \mathcal{F}_x)
$$
而要检验
$H^i(E \otimes_{\mathcal{O}_X}^\mathbf{L} \mathcal{F})$ 是否为零，
只需检验其所有茎是否为零
（模，引理 \ref{modules-lemma-abelian}）。故 (2) 推出 (1)。
\end{proof}

\begin{lemma}
\label{cohomology-lemma-cone-tor-amplitude}
令 $(X, \mathcal{O}_X)$ 为环化空间。
令 $(K, L, M, f, g, h)$ 为 $D(\mathcal{O}_X)$ 中的特异
三角。令 $a, b \in \mathbf{Z}$。
\begin{enumerate}
\item 如果 $K$ 的 Tor 振幅位于 $[a + 1, b + 1]$，并且
$L$ 的 Tor 振幅位于 $[a, b]$，那么 $M$ 的
Tor 振幅位于 $[a, b]$。
\item 如果 $K$ 和 $M$ 的 Tor 振幅位于 $[a, b]$，那么
$L$ 的 Tor 振幅位于 $[a, b]$。
\item 如果 $L$ 的 Tor 振幅位于 $[a + 1, b + 1]$，
并且 $M$ 的 Tor 振幅位于 $[a, b]$，那么
$K$ 的 Tor 振幅位于 $[a + 1, b + 1]$。
\end{enumerate}
\end{lemma}

\begin{proof}
略去证明。提示：这直接来自特异三角所伴随的上同调长正合列，
以及
$- \otimes_{\mathcal{O}_X}^{\mathbf{L}} \mathcal{F}$
保持特异三角这一事实。
最容易证明的是 (2)，其余结论由平移得到。
\end{proof}

\begin{lemma}
\label{cohomology-lemma-tensor-tor-amplitude}
令 $(X, \mathcal{O}_X)$ 为环化空间。令 $K, L$ 为
$D(\mathcal{O}_X)$ 的对象。如果 $K$ 的 Tor 振幅位于 $[a, b]$，且
$L$ 的 Tor 振幅位于 $[c, d]$，那么 $K \otimes_{\mathcal{O}_X}^\mathbf{L} L$
的 Tor 振幅位于 $[a + c, b + d]$。
\end{lemma}

\begin{proof}
略去证明。提示：使用 Tor 的谱序列。
\end{proof}

\begin{lemma}
\label{cohomology-lemma-summands-tor-amplitude}
令 $(X, \mathcal{O}_X)$ 为环化空间。令 $a, b \in \mathbf{Z}$。
对于 $D(\mathcal{O}_X)$ 的对象 $K$、$L$，如果 $K \oplus L$ 的 Tor
振幅位于 $[a, b]$，那么 $K$ 和 $L$ 也都如此。
\end{lemma}

\begin{proof}
由 Tor 函子的可加性立即可得。
\end{proof}







\section{完美复形}
\label{cohomology-section-perfect}

\noindent
本节讨论环化空间上完美复形的
性质。

\begin{definition}
\label{cohomology-definition-perfect}
令 $(X, \mathcal{O}_X)$ 为环化空间。
令 $\mathcal{E}^\bullet$ 为 $\mathcal{O}_X$-模复形。
如果存在开覆盖 $X = \bigcup U_i$，使得对每个 $i$
都存在复形态射
$\mathcal{E}_i^\bullet \to \mathcal{E}^\bullet|_{U_i}$，
它是拟同构，且 $\mathcal{E}_i^\bullet$
是 $\mathcal{O}_{U_i}$-模的严格完美复形，
则称 $\mathcal{E}^\bullet$ 是{\it 完美的}。
如果 $D(\mathcal{O}_X)$ 的对象 $E$ 可由
完美的 $\mathcal{O}_X$-模复形表示，则称其是{\it 完美的}。
\end{definition}

\noindent
如果 $X$ 拟紧，那么 $D(\mathcal{O}_X)$ 的完美对象
属于 $D^b(\mathcal{O}_X)$。但当
$X$ 非拟紧时，未必如此。

\begin{lemma}
\label{cohomology-lemma-perfect-independent-representative}
令 $(X, \mathcal{O}_X)$ 为环化空间。
令 $E$ 为 $D(\mathcal{O}_X)$ 的对象。
\begin{enumerate}
\item 如果存在开覆盖 $X = \bigcup U_i$，以及 $U_i$ 上的
严格完美复形 $\mathcal{E}_i^\bullet$，使得
$\mathcal{E}_i^\bullet$ 在 $D(\mathcal{O}_{U_i})$ 中表示 $E|_{U_i}$，
那么 $E$ 是完美的。
\item 如果 $E$ 是完美的，那么表示 $E$ 的任意复形都是完美的。
\end{enumerate}
\end{lemma}

\begin{proof}
与引理 \ref{cohomology-lemma-pseudo-coherent-independent-representative}
的证明相同。
\end{proof}

\begin{lemma}
\label{cohomology-lemma-perfect-on-locally-ringed}
令 $(X, \mathcal{O}_X)$ 为环化空间。令 $E$ 为
$D(\mathcal{O}_X)$ 的对象。假设所有茎 $\mathcal{O}_{X, x}$
都是局部环。则下列条件等价：
\begin{enumerate}
\item $E$ 是完美的；
\item 存在开覆盖 $X = \bigcup U_i$，使得
$E|_{U_i}$ 可由有限局部自由
$\mathcal{O}_{U_i}$-模组成的有限复形表示；以及
\item 存在开覆盖 $X = \bigcup U_i$，使得
$E|_{U_i}$ 可由有限自由
$\mathcal{O}_{U_i}$-模组成的有限复形表示。
\end{enumerate}
\end{lemma}

\begin{proof}
这由引理 \ref{cohomology-lemma-perfect-independent-representative}
以及如下事实得出：在 $X$ 上，有限自由模的每个直和项
都是有限局部自由的。见模，引理
\ref{modules-lemma-direct-summand-of-locally-free-is-locally-free}。
\end{proof}

\begin{lemma}
\label{cohomology-lemma-perfect-precise}
令 $(X, \mathcal{O}_X)$ 为环化空间。
令 $E$ 为 $D(\mathcal{O}_X)$ 的对象。
令 $a \leq b$ 为整数。如果 $E$ 的 Tor 振幅位于 $[a, b]$
并且是 $(a - 1)$-伪凝聚的，那么 $E$ 是完美的。
\end{lemma}

\begin{proof}
用开覆盖的成员替换 $X$ 后，可假设存在
严格完美复形 $\mathcal{E}^\bullet$ 和态射
$\alpha : \mathcal{E}^\bullet \to E$，使得 $H^i(\alpha)$
对 $i \geq a$ 是同构。我们可以并且确实将 $\mathcal{E}^\bullet$ 替换为
$\sigma_{\geq a - 1}\mathcal{E}^\bullet$。选取特异三角
$$
\mathcal{E}^\bullet \to E \to C \to \mathcal{E}^\bullet[1]
$$
由 $E$ 和 $\mathcal{E}^\bullet$ 的上同调层消失性
以及关于 $\alpha$ 的假设，得到 $C \cong \mathcal{K}[2 - a]$，其中
$\mathcal{K} = \Ker(\mathcal{E}^{a - 1} \to \mathcal{E}^a)$。
令 $\mathcal{F}$ 为 $\mathcal{O}_X$-模。
应用 $- \otimes_{\mathcal{O}_X}^\mathbf{L} \mathcal{F}$ 后，
$E$ 的 Tor 振幅位于 $[a, b]$ 这一假设
推出 $\mathcal{K} \otimes_{\mathcal{O}_X} \mathcal{F} \to
\mathcal{E}^{a - 1} \otimes_{\mathcal{O}_X} \mathcal{F}$ 的像等于
$\Ker(\mathcal{E}^{a - 1} \otimes_{\mathcal{O}_X} \mathcal{F}
\to \mathcal{E}^a \otimes_{\mathcal{O}_X} \mathcal{F})$。
由此 $\text{Tor}_1^{\mathcal{O}_X}(\mathcal{E}', \mathcal{F}) = 0$，
其中 $\mathcal{E}' = \Coker(\mathcal{E}^{a - 1} \to \mathcal{E}^a)$。
因此 $\mathcal{E}'$ 平坦（引理 \ref{cohomology-lemma-flat-tor-zero}）。
于是由模，引理 \ref{modules-lemma-flat-locally-finite-presentation}，
$\mathcal{E}'$ 局部为某个有限自由模的直和项。
故复形
$$
\mathcal{E}' \to \mathcal{E}^{a + 1} \to \ldots \to \mathcal{E}^b
$$
局部拟同构于 $E$，从而 $E$ 是完美的。
\end{proof}

\begin{lemma}
\label{cohomology-lemma-perfect}
令 $(X, \mathcal{O}_X)$ 为环化空间。
令 $E$ 为 $D(\mathcal{O}_X)$ 的对象。
下列条件等价：
\begin{enumerate}
\item $E$ 是完美的；以及
\item $E$ 是伪凝聚的，并且局部具有有限 Tor 维数。
\end{enumerate}
\end{lemma}

\begin{proof}
假设 (1)。按定义，这意味着存在开覆盖
$X = \bigcup U_i$，使得 $E|_{U_i}$ 可由
严格完美复形表示。因此，由引理 \ref{cohomology-lemma-pseudo-coherent-independent-representative}，$E$ 是伪凝聚的（即
对所有 $m$ 都是 $m$-伪凝聚的）。
此外，有限自由模的直和项是平坦的，所以
由引理 \ref{cohomology-lemma-tor-amplitude}，
$E|_{U_i}$ 具有有限 Tor 维数。因此 (2) 成立。

\medskip\noindent
假设 (2)。用开覆盖的成员替换 $X$ 后，
可假设存在整数 $a \leq b$，使得 $E$
的 Tor 振幅位于 $[a, b]$。由于 $E$ 对所有 $m$
都是 $m$-伪凝聚的，应用引理 \ref{cohomology-lemma-perfect-precise} 即得结论。
\end{proof}

\begin{lemma}
\label{cohomology-lemma-perfect-pullback}
令 $f : (X, \mathcal{O}_X) \to (Y, \mathcal{O}_Y)$ 为环化
空间的态射。令 $E$ 为 $D(\mathcal{O}_Y)$ 的对象。如果 $E$ 在
$D(\mathcal{O}_Y)$ 中完美，那么 $Lf^*E$ 在 $D(\mathcal{O}_X)$ 中完美。
\end{lemma}

\begin{proof}
这由引理 \ref{cohomology-lemma-perfect}、
\ref{cohomology-lemma-tor-amplitude-pullback} 和
\ref{cohomology-lemma-pseudo-coherent-pullback} 得出。
（另一种证明是照搬引理
\ref{cohomology-lemma-pseudo-coherent-pullback} 的证明。）
\end{proof}

\begin{lemma}
\label{cohomology-lemma-two-out-of-three-perfect}
令 $(X, \mathcal{O}_X)$ 为环化空间。令 $(K, L, M, f, g, h)$
为 $D(\mathcal{O}_X)$ 中的特异三角。如果 $K, L, M$
三者中有两个完美，那么第三个也完美。
\end{lemma}

\begin{proof}
第一种证明：合并使用引理 \ref{cohomology-lemma-perfect}、\ref{cohomology-lemma-cone-pseudo-coherent} 和
\ref{cohomology-lemma-cone-tor-amplitude}。
第二种证明（略述）：设 $K$ 和 $L$ 完美。用开覆盖的成员
替换 $X$ 后，可假设 $K$ 和 $L$
分别由严格完美复形 $\mathcal{K}^\bullet$
和 $\mathcal{L}^\bullet$ 表示。再次用开覆盖的成员
替换 $X$ 后，可假设态射 $K \to L$ 由
复形态射 $\alpha : \mathcal{K}^\bullet \to \mathcal{L}^\bullet$ 给出，
见引理 \ref{cohomology-lemma-local-actual}。
于是 $M$ 同构于 $\alpha$ 的锥，而由引理 \ref{cohomology-lemma-cone}，
该锥是严格完美的。
\end{proof}

\begin{lemma}
\label{cohomology-lemma-tensor-perfect}
令 $(X, \mathcal{O}_X)$ 为环化空间。
如果 $K, L$ 是 $D(\mathcal{O}_X)$ 的完美对象，那么
$K \otimes_{\mathcal{O}_X}^\mathbf{L} L$ 也是完美对象。
\end{lemma}

\begin{proof}
这由引理 \ref{cohomology-lemma-perfect}、\ref{cohomology-lemma-tensor-pseudo-coherent} 和
\ref{cohomology-lemma-tensor-tor-amplitude} 得出。
\end{proof}

\begin{lemma}
\label{cohomology-lemma-summands-perfect}
令 $(X, \mathcal{O}_X)$ 为环化空间。
如果 $K \oplus L$ 是 $D(\mathcal{O}_X)$ 的完美对象，那么
$K$ 和 $L$ 也都是完美对象。
\end{lemma}

\begin{proof}
这由引理 \ref{cohomology-lemma-perfect}、\ref{cohomology-lemma-summands-pseudo-coherent} 和
\ref{cohomology-lemma-summands-tor-amplitude} 得出。
\end{proof}

\begin{lemma}
\label{cohomology-lemma-pushforward-perfect}
令 $(X, \mathcal{O}_X)$ 为环化空间。令 $j : U \to X$ 为
开子空间。令 $E$ 为 $D(\mathcal{O}_U)$ 的完美对象，其上同调
层支撑在闭子集 $T \subset U$ 上，并且 $j(T)$
在 $X$ 中闭。那么 $Rj_*E$ 是 $D(\mathcal{O}_X)$ 的完美对象。
\end{lemma}

\begin{proof}
完美复形这一性质在 $X$ 上是局部的。因此只需检验
$Rj_*E$ 限制到 $U$ 和 $V = X \setminus j(T)$ 后都是完美的。
我们有 $Rj_*E|_U = E$，它是完美的。又有
$Rj_*E|_V = 0$，因为 $E|_{U \setminus T} = 0$。
\end{proof}

\begin{lemma}
\label{cohomology-lemma-perfect-max-open-coh-loc-free}
令 $(X, \mathcal{O}_X)$ 为环化空间。令 $D(\mathcal{O}_X)$ 中的 $E$
为完美对象。假设所有茎 $\mathcal{O}_{X, x}$ 都是局部环。
那么集合
$$
U =
\{x \in X \mid
H^i(E)_x\text{ 是有限自由 }
\mathcal{O}_{X, x}\text{-模，对所有 }i\in \mathbf{Z}\}
$$
在 $X$ 中开，并且是满足如下性质的极大开集 $U \subset X$：
对所有 $i \in \mathbf{Z}$，$H^i(E)|_U$ 都是有限局部自由的。
\end{lemma}

\begin{proof}
注意，如果 $V \subset X$ 是某个开集，使得对所有 $i \in \mathbf{Z}$，
$H^i(E)|_V$ 都是有限局部自由的，那么 $V \subset U$。
令 $x \in U$。我们将证明 $x$ 的某个开邻域包含于 $U$，
并且对所有 $i$，$H^i(E)$ 在该邻域上都是有限局部自由的。
这将完成证明。
在证明过程中，我们可以（有限多次）用 $x$ 的开邻域
替换 $X$。
因此可假设 $E$ 由严格完美复形
$\mathcal{E}^\bullet$ 表示。设当 $i \not \in [a, b]$ 时
$\mathcal{E}^i = 0$。我们将对 $b - a$ 归纳证明结论。模
$H^b(E) = \Coker(d^{b - 1} : \mathcal{E}^{b - 1} \to \mathcal{E}^b)$
是有限表示的。由于 $H^b(E)_x$ 有限自由，
由模，引理 \ref{modules-lemma-finite-presentation-stalk-free}，
$H^b(E)$ 在 $x$ 的某个开邻域上有限自由。
因此，用一个（可能更小的）开邻域替换 $X$ 后，
可假设有直和分解
$\mathcal{E}^b = \Im(d^{b - 1}) \oplus H^b(E)$，并且
$H^b(E)$ 有限自由，见引理 \ref{cohomology-lemma-local-lift-map}。
再次使用同一论证，可假设
$\mathcal{E}^{b - 1} = \Ker(d^{b - 1}) \oplus \Im(d^{b - 1})$。
复形
$\mathcal{E}^a \to \ldots \to \mathcal{E}^{b - 2} \to \Ker(d^{b - 1})$
是严格完美复形，表示完美
对象 $E'$，且当 $i \not = b$ 时 $H^i(E) = H^i(E')$。
因此由归纳假设得到结论。
\end{proof}







\section{对偶}
\label{cohomology-section-duals}

\noindent
本节刻画复形范畴和导出范畴中的
可对偶对象。
特别地，我们将看到，$D(\mathcal{O}_X)$ 的对象
有对偶，当且仅当它是完美的（这由例 \ref{cohomology-example-dual-derived} 和引理 \ref{cohomology-lemma-left-dual-derived} 得出）。

\begin{lemma}
\label{cohomology-lemma-symmetric-monoidal-cat-complexes}
令 $(X, \mathcal{O}_X)$ 为环化空间。以
$\mathcal{F}^\bullet \otimes \mathcal{G}^\bullet =
\text{Tot}(\mathcal{F}^\bullet \otimes_{\mathcal{O}_X} \mathcal{G}^\bullet)$
定义张量积后，$\mathcal{O}_X$-模复形的范畴
是对称幺半范畴（符号规则见
更多代数，第 \ref{more-algebra-section-sign-rules} 节）。
\end{lemma}

\begin{proof}
略去证明。提示：取复形 $\mathbf{1}$ 为单位对象，
它在次数 $0$ 处为 $\mathcal{O}_X$，在其余次数处为零，并带有显然的同构
$\text{Tot}(\mathbf{1} \otimes_{\mathcal{O}_X} \mathcal{G}^\bullet) =
\mathcal{G}^\bullet$ 和
$\text{Tot}(\mathcal{F}^\bullet \otimes_{\mathcal{O}_X} \mathbf{1}) =
\mathcal{F}^\bullet$。
为证明此引理，需要检验若干图的交换性，见
范畴，定义
\ref{categories-definition-monoidal-category} 和
\ref{categories-definition-symmetric-monoidal-category}。
每种情形下的验证都很直接。
\end{proof}

\begin{example}
\label{cohomology-example-dual}
令 $(X, \mathcal{O}_X)$ 为环化空间。令 $\mathcal{F}^\bullet$
为 $\mathcal{O}_X$-模的局部有界复形，使得每个
$\mathcal{F}^n$ 都局部为某个有限自由
$\mathcal{O}_X$-模的直和项。换言之，存在开覆盖
$X = \bigcup U_i$，使得 $\mathcal{F}^\bullet|_{U_i}$ 是严格
完美复形。考虑第 \ref{cohomology-section-hom-complexes} 节中的复形
$$
\mathcal{G}^\bullet = \SheafHom^\bullet(\mathcal{F}^\bullet, \mathcal{O}_X)
$$
令
$$
\eta :
\mathcal{O}_X
\to
\text{Tot}(\mathcal{F}^\bullet \otimes_{\mathcal{O}_X} \mathcal{G}^\bullet)
\quad\text{且}\quad
\epsilon :
\text{Tot}(\mathcal{G}^\bullet \otimes_{\mathcal{O}_X} \mathcal{F}^\bullet)
\to
\mathcal{O}_X
$$
分别为 $\eta = \sum \eta_n$ 和 $\epsilon = \sum \epsilon_n$，
其中 $\eta_n : \mathcal{O}_X \to
\mathcal{F}^n \otimes_{\mathcal{O}_X} \mathcal{G}^{-n}$，
并且
$\epsilon_n : \mathcal{G}^{-n} \otimes_{\mathcal{O}_X} \mathcal{F}^n
\to \mathcal{O}_X$ 如模，例 \ref{modules-example-dual} 中所定义。
那么按范畴，定义 \ref{categories-definition-dual}，
$\mathcal{G}^\bullet, \eta, \epsilon$
是 $\mathcal{F}^\bullet$ 的左对偶。
我们略去下列等式的验证：
$(1 \otimes \epsilon) \circ (\eta \otimes 1) = \text{id}_{\mathcal{F}^\bullet}$
以及
$(\epsilon \otimes 1) \circ (1 \otimes \eta) =
\text{id}_{\mathcal{G}^\bullet}$。请比较
更多代数，引理 \ref{more-algebra-lemma-left-dual-complex}。
\end{example}

\begin{lemma}
\label{cohomology-lemma-left-dual-complex}
令 $(X, \mathcal{O}_X)$ 为环化空间。令 $\mathcal{F}^\bullet$
为 $\mathcal{O}_X$-模复形。如果 $\mathcal{F}^\bullet$
在 $\mathcal{O}_X$-模复形的幺半范畴中
有左对偶
（范畴，定义 \ref{categories-definition-dual}），
那么 $\mathcal{F}^\bullet$ 是局部有界复形，其各项都
局部为有限自由 $\mathcal{O}_X$-模的直和项，
并且该左对偶就是例 \ref{cohomology-example-dual} 中的构造。
\end{lemma}

\begin{proof}
由左对偶的唯一性
（范畴，注记 \ref{categories-remark-left-dual-adjoint}），
只要证明 $\mathcal{F}^\bullet$ 具有所述性质，便得到最后一个断言。
令 $\mathcal{G}^\bullet, \eta, \epsilon$ 为左对偶。
写成 $\eta = \sum \eta_n$ 和 $\epsilon = \sum \epsilon_n$，
其中 $\eta_n : \mathcal{O}_X \to
\mathcal{F}^n \otimes_{\mathcal{O}_X} \mathcal{G}^{-n}$，
并且
$\epsilon_n : \mathcal{G}^{-n} \otimes_{\mathcal{O}_X} \mathcal{F}^n
\to \mathcal{O}_X$。由于
$(1 \otimes \epsilon) \circ (\eta \otimes 1) = \text{id}_{\mathcal{F}^\bullet}$
以及
$(\epsilon \otimes 1) \circ (1 \otimes \eta) = \text{id}_{\mathcal{G}^\bullet}$，
由范畴，定义 \ref{categories-definition-dual}，立即可见
$(1 \otimes \epsilon_n) \circ (\eta_n \otimes 1) = \text{id}_{\mathcal{F}^n}$
以及
$(\epsilon_n \otimes 1) \circ (1 \otimes \eta_n) =
\text{id}_{\mathcal{G}^{-n}}$。
因此由模，引理 \ref{modules-lemma-left-dual-module}，
$\mathcal{F}^n$ 局部为某个有限自由
$\mathcal{O}_X$-模的直和项。
由于和 $\eta = \sum \eta_n$ 局部有限，可知
$\mathcal{F}^\bullet$ 局部有界。
\end{proof}

\begin{lemma}
\label{cohomology-lemma-internal-hom-evaluate-isom}
令 $(X, \mathcal{O}_X)$ 为环化空间。令 $K, L, M \in D(\mathcal{O}_X)$。
如果 $K$ 完美，那么引理 \ref{cohomology-lemma-internal-hom-evaluate} 中的态射
$$
R\SheafHom(L, M) \otimes_{\mathcal{O}_X}^\mathbf{L} K
\longrightarrow
R\SheafHom(R\SheafHom(K, L), M)
$$
是同构。
\end{lemma}

\begin{proof}
由于该态射整体定义，并且等式的左、右两端的构造都与局部化
交换（见引理 \ref{cohomology-lemma-restriction-RHom-to-U}），为证明断言可在
$X$ 上局部工作。因此可假设 $K$ 由严格
完美复形 $\mathcal{E}^\bullet$ 表示。

\medskip\noindent
如果 $K_1 \to K_2 \to K_3$ 是 $D(\mathcal{O}_X)$ 中的特异三角，
那么得到特异三角
$$
R\SheafHom(L, M) \otimes_{\mathcal{O}_X}^\mathbf{L} K_1 \to
R\SheafHom(L, M) \otimes_{\mathcal{O}_X}^\mathbf{L} K_2 \to
R\SheafHom(L, M) \otimes_{\mathcal{O}_X}^\mathbf{L} K_3
$$
以及
$$
R\SheafHom(R\SheafHom(K_1, L), M) \to
R\SheafHom(R\SheafHom(K_2, L), M)
R\SheafHom(R\SheafHom(K_3, L), M)
$$
见第 \ref{cohomology-section-flat} 节和引理 \ref{cohomology-lemma-RHom-triangulated}。引理 \ref{cohomology-lemma-internal-hom-evaluate} 的箭头对 $K$ 具有函子性，
因而得到这两个特异三角之间的态射。
于是，如果结论对 $K_1$ 和 $K_3$ 成立，则由导出范畴，引理
\ref{derived-lemma-third-isomorphism-triangle}，
结论对 $K_2$ 也成立。

\medskip\noindent
结合上述观察与朴素截断的特异三角
$$
\sigma_{\geq n}\mathcal{E}^\bullet \to \mathcal{E}^\bullet \to
\sigma_{\leq n - 1}\mathcal{E}^\bullet
$$
可约化到 $K$ 是放在某个次数上的
有限自由 $\mathcal{O}_X$-模的直和项这一情形。
再利用该问题与直和的显然相容性
（类似于上文所述）以及移位，可进一步约化到
$K = \mathcal{O}_X^{\oplus n}$ 的情形，其中 $n$ 是某个整数。
此时结论显然。
\end{proof}

\begin{lemma}
\label{cohomology-lemma-dual-perfect-complex}
令 $(X, \mathcal{O}_X)$ 为环化空间。令 $K$ 为
$D(\mathcal{O}_X)$ 的完美对象。那么
$K^\vee = R\SheafHom(K, \mathcal{O}_X)$ 也是完美对象，
且 $(K^\vee)^\vee \cong K$。存在函子性同构
$$
M \otimes^\mathbf{L}_{\mathcal{O}_X} K^\vee = R\SheafHom(K, M)
$$
以及
$$
H^0(X, M \otimes^\mathbf{L}_{\mathcal{O}_X} K^\vee) =
\Hom_{D(\mathcal{O}_X)}(K, M)
$$
其中 $M$ 属于 $D(\mathcal{O}_X)$。
\end{lemma}

\begin{proof}
由引理 \ref{cohomology-lemma-internal-hom-evaluate}，存在典范态射
$$
K = R\SheafHom(\mathcal{O}_X, \mathcal{O}_X)
\otimes_{\mathcal{O}_X}^\mathbf{L} K \longrightarrow
R\SheafHom(R\SheafHom(K, \mathcal{O}_X), \mathcal{O}_X) =
(K^\vee)^\vee
$$
由引理 \ref{cohomology-lemma-internal-hom-evaluate-isom}，它是同构。
为检验其余断言，我们将不再特别说明地使用如下事实：
内 Hom 的构造与到开集的限制可交换
（引理 \ref{cohomology-lemma-restriction-RHom-to-U}）。
可在 $X$ 上局部检验 $K^\vee$ 是完美的。
由引理 \ref{cohomology-lemma-dual}，
为看出最后一个断言，只需检验态射
(\ref{cohomology-equation-eval})
$$
M \otimes^\mathbf{L}_{\mathcal{O}_X} K^\vee
\longrightarrow
R\SheafHom(K, M)
$$
是同构。这也可在 $X$ 上局部检验。
因此，只需在 $K$ 由严格完美复形表示时，
证明这两个断言。

\medskip\noindent
假设 $K$ 由严格完美复形
$\mathcal{E}^\bullet$ 表示。由引理 \ref{cohomology-lemma-Rhom-strictly-perfect} 可知，
$K^\vee$ 由如下复形表示：其次数 $n$ 的项为
$(\mathcal{E}^{-n})^\vee =
\SheafHom_{\mathcal{O}_X}(\mathcal{E}^{-n}, \mathcal{O}_X)$。
由于 $\mathcal{E}^{-n}$ 是某个有限自由
$\mathcal{O}_X$-模的直和项，$(\mathcal{E}^{-n})^\vee$ 亦然。
因此 $K^\vee$ 也由严格完美复形表示，
从而可见 $K^\vee$ 完美。
为看出 (\ref{cohomology-equation-eval}) 是同构，用复形
$\mathcal{F}^\bullet$ 表示 $M$。
由引理 \ref{cohomology-lemma-Rhom-strictly-perfect}，复形
$R\SheafHom(K, M)$ 由以下各项组成的复形表示：
$$
\bigoplus\nolimits_{n = p + q}
\SheafHom_{\mathcal{O}_X}(\mathcal{E}^{-q}, \mathcal{F}^p)
$$
另一方面，对象 $M \otimes^\mathbf{L}_{\mathcal{O}_X} K^\vee$
由以下各项组成的复形表示：
$$
\bigoplus\nolimits_{n = p + q}
\mathcal{F}^p \otimes_{\mathcal{O}_X} (\mathcal{E}^{-q})^\vee
$$
因此，(\ref{cohomology-equation-eval}) 是同构这一断言
约化为如下典范态射
$$
\mathcal{F}
\otimes_{\mathcal{O}_X}
\SheafHom_{\mathcal{O}_X}(\mathcal{E}, \mathcal{O}_X)
\longrightarrow
\SheafHom_{\mathcal{O}_X}(\mathcal{E}, \mathcal{F})
$$
在 $\mathcal{E}$ 是某个有限自由 $\mathcal{O}_X$-模的直和项、
且 $\mathcal{F}$ 是任意 $\mathcal{O}_X$-模时为同构这一断言。
这立即由 $\mathcal{E}$ 有限自由时的相应断言得出。
\end{proof}


\begin{lemma}
\label{cohomology-lemma-symmetric-monoidal-derived}
令 $(X, \mathcal{O}_X)$ 为环化空间。导出范畴
$D(\mathcal{O}_X)$ 以导出张量积为张量积，并带有通常的结合
约束和交换约束，从而是对称幺半范畴（符号规则见
更多代数，第 \ref{more-algebra-section-sign-rules} 节）。
\end{lemma}

\begin{proof}
略去证明。可与引理 \ref{cohomology-lemma-symmetric-monoidal-cat-complexes} 比较。
\end{proof}

\begin{example}
\label{cohomology-example-dual-derived}
令 $(X, \mathcal{O}_X)$ 为环化空间。令 $K$ 为
$D(\mathcal{O}_X)$ 的完美对象。按引理
\ref{cohomology-lemma-dual-perfect-complex}，令
$K^\vee = R\SheafHom(K, \mathcal{O}_X)$。
则态射
$$
K \otimes_{\mathcal{O}_X}^\mathbf{L} K^\vee \longrightarrow R\SheafHom(K, K)
$$
是同构（由该引理）。记
$$
\eta :
\mathcal{O}_X
\longrightarrow
K \otimes_{\mathcal{O}_X}^\mathbf{L} K^\vee
$$
为如下态射：它将 $1$ 送到在上述同构下
对应于 $\text{id}_K$ 的截面。
记
$$
\epsilon : 
K^\vee
\otimes_{\mathcal{O}_X}^\mathbf{L} K
\longrightarrow
\mathcal{O}_X
$$
为求值态射（例如可用引理 \ref{cohomology-lemma-internal-hom-composition} 构造）。那么
按范畴，定义 \ref{categories-definition-dual}，
$K^\vee, \eta, \epsilon$ 是 $K$ 的左对偶。
我们略去下列等式的验证：
$(1 \otimes \epsilon) \circ (\eta \otimes 1) = \text{id}_K$
以及
$(\epsilon \otimes 1) \circ (1 \otimes \eta) =
\text{id}_{K^\vee}$。
\end{example}

\begin{lemma}
\label{cohomology-lemma-left-dual-derived}
令 $(X, \mathcal{O}_X)$ 为环化空间。令 $M$ 为
$D(\mathcal{O}_X)$ 的对象。如果 $M$ 在幺半范畴
$D(\mathcal{O}_X)$ 中有左对偶（范畴，定义
\ref{categories-definition-dual}），那么 $M$ 完美，并且该左对偶
就是例 \ref{cohomology-example-dual-derived} 中的构造。
\end{lemma}

\begin{proof}
令 $x \in X$。只需找到 $x$ 的开邻域 $U$，
使得 $M$ 限制到 $U$ 后为完美复形。因此在证明过程中，
可以（有限多次）用 $x$ 的开邻域替换 $X$。
令 $N, \eta, \epsilon$ 为左对偶。

\medskip\noindent
我们将多次使用如下论证。任取
表示 $M$ 的 $\mathcal{O}_X$-模复形
$\mathcal{M}^\bullet$。选取表示 $N$ 的 K-平坦复形
$\mathcal{N}^\bullet$，其各项都是平坦
$\mathcal{O}_X$-模，见引理 \ref{cohomology-lemma-K-flat-resolution}。
考虑态射
$$
\eta : \mathcal{O}_X \to
\text{Tot}(\mathcal{M}^\bullet \otimes_{\mathcal{O}_X} \mathcal{N}^\bullet)
$$
缩小 $X$ 后，可找到整数 $N$，以及对
$i = 1, \ldots, N$ 的整数 $n_i \in \mathbf{Z}$ 和
$\mathcal{M}^{n_i}$、$\mathcal{N}^{-n_i}$ 的截面
$f_i$、$g_i$，使得
$$
\eta(1) = \sum\nolimits_i f_i \otimes g_i
$$
令 $\mathcal{K}^\bullet \subset \mathcal{M}^\bullet$ 为任意
包含各截面 $f_i$（$i = 1, \ldots, N$）的
$\mathcal{O}_X$-模子复形。
由于各模 $\mathcal{N}^n$ 平坦，有
$\text{Tot}(\mathcal{K}^\bullet \otimes_{\mathcal{O}_X} \mathcal{N}^\bullet)
\subset
\text{Tot}(\mathcal{M}^\bullet \otimes_{\mathcal{O}_X} \mathcal{N}^\bullet)$，
故 $\eta$ 经由
$$
\tilde \eta :
\mathcal{O}_X \to
\text{Tot}(\mathcal{K}^\bullet \otimes_{\mathcal{O}_X} \mathcal{N}^\bullet)
$$
分解。记由 $\mathcal{K}^\bullet$ 表示的
$D(\mathcal{O}_X)$ 对象为 $K$，得到交换图
$$
\xymatrix{
M \ar[rr]_-{\eta \otimes 1} \ar[rrd]_{\tilde \eta \otimes 1} & &
M \otimes^\mathbf{L} N \otimes^\mathbf{L} M
\ar[r]_-{1 \otimes \epsilon} &
M \\
& &
K \otimes^\mathbf{L} N \otimes^\mathbf{L} M
\ar[u] \ar[r]^-{1 \otimes \epsilon} &
K \ar[u]
}
$$
由于上行的复合是 $M$ 上的恒等态射，
可知 $M$ 是 $D(\mathcal{O}_X)$ 中 $K$ 的直和项。

\medskip\noindent
第一次应用上述论证时，可选择子复形
$\mathcal{K}^\bullet = \sigma_{\geq a} \tau_{\leq b}\mathcal{M}^\bullet$，
其中对 $i = 1, \ldots, N$ 有 $a < n_i < b$。因此
$M$ 是 $D(\mathcal{O}_X)$ 中某个有界复形的直和项，
从而可假设 $M$ 属于 $D^b(\mathcal{O}_X)$。（回忆上述过程
涉及缩小 $X$。）

\medskip\noindent
由于 $M$ 属于 $D^b(\mathcal{O}_X)$，可选择
$\mathcal{M}^\bullet$ 为由平坦模组成的上有界复形
（由模，引理 \ref{modules-lemma-module-quotient-flat} 和
导出范畴，引理 \ref{derived-lemma-subcategory-left-resolution}）。
于是，在上述论证中可选择
$\mathcal{K}^\bullet = \sigma_{\geq a}\mathcal{M}^\bullet$，
其中对 $i = 1, \ldots, N$ 有 $a < n_i$。
由此可假设 $M$ 是
$D(\mathcal{O}_X)$ 中某个由平坦模组成的有界复形的直和项。
特别地，$M$ 的 Tor 振幅有限。

\medskip\noindent
设 $M$ 的 Tor 振幅位于 $[a, b]$。假设 $M$ 是 $m$-伪凝聚的，
我们将证明（缩小 $X$ 后）可假设 $M$
是 $(m - 1)$-伪凝聚的。结合引理 \ref{cohomology-lemma-perfect-precise} 以及 $M$
无论如何都是 $(b + 1)$-伪凝聚的这一事实，这将完成证明。
缩小 $X$ 后，可假设存在严格完美
复形 $\mathcal{E}^\bullet$ 以及 $D(\mathcal{O}_X)$ 中的态射
$\alpha : \mathcal{E}^\bullet \to M$，使得 $H^i(\alpha)$ 在
$i > m$ 时为同构，在 $i = m$ 时为满射。可以并且确实假设
当 $i < m$ 时 $\mathcal{E}^i = 0$。选取特异三角
$$
\mathcal{E}^\bullet \to M \to L \to \mathcal{E}^\bullet[1]
$$
注意，当 $i \geq m$ 时 $H^i(L) = 0$。因此可用满足
当 $i \geq m$ 时 $\mathcal{L}^i = 0$ 的复形
$\mathcal{L}^\bullet$ 表示 $L$。态射
$L \to \mathcal{E}^\bullet[1]$ 由复形态射
$\mathcal{L}^\bullet \to \mathcal{E}^\bullet[1]$ 给出；
该态射在除次数 $m - 1$ 之外的所有次数上为零，而在该次数
处得到态射 $\mathcal{L}^{m - 1} \to \mathcal{E}^m$，见
导出范畴，引理 \ref{derived-lemma-negative-exts}。
于是 $M$ 由复形
$$
\mathcal{M}^\bullet :
\ldots \to
\mathcal{L}^{m - 2} \to
\mathcal{L}^{m - 1} \to
\mathcal{E}^m \to
\mathcal{E}^{m + 1} \to \ldots
$$
表示。把第二段中的讨论应用于此复形，得到
$\mathcal{M}^{n_i}$ 的截面 $f_i$，其中 $i = 1, \ldots, N$。
对 $n < m$，令 $\mathcal{K}^n \subset \mathcal{L}^n$
为由满足 $n_i = n$ 的截面
$f_i$ 以及满足 $n_i = n - 1$ 的 $d(f_i)$ 生成的
$\mathcal{O}_X$-子模。
对 $n \geq m$，令 $\mathcal{K}^n = \mathcal{E}^n$。
显然有特异三角的态射
$$
\xymatrix{
\mathcal{E}^\bullet \ar[r] &
\mathcal{M}^\bullet \ar[r] &
\mathcal{L}^\bullet \ar[r] &
\mathcal{E}^\bullet[1] \\
\mathcal{E}^\bullet \ar[r] \ar[u] &
\mathcal{K}^\bullet \ar[r] \ar[u] &
\sigma_{\leq m - 1}\mathcal{K}^\bullet \ar[r] \ar[u] &
\mathcal{E}^\bullet[1] \ar[u]
}
$$
其中所有态射都如上所述。
记与复形 $\mathcal{K}^\bullet$ 对应的
$D(\mathcal{O}_X)$ 对象为 $K$。
由证明第二段中的论证，得到 $D(\mathcal{O}_X)$ 中的
态射 $s : M \to K$，使得复合
$M \to K \to M$ 是 $M$ 上的恒等态射。我们不知道图
$$
\xymatrix{
\mathcal{E}^\bullet \ar[r] &
\mathcal{K}^\bullet \ar@{=}[r] &
K \\
\mathcal{E}^\bullet \ar[u]^{\text{id}} \ar[r]^i &
\mathcal{M}^\bullet \ar@{=}[r] &
M \ar[u]_s
}
$$
是否交换，但知道它与态射 $K \to M$ 复合后交换。
由引理 \ref{cohomology-lemma-local-actual}，缩小 $X$ 后，可假设
$s \circ i$ 由复形态射
$\sigma : \mathcal{E}^\bullet \to \mathcal{K}^\bullet$ 给出。
再次使用同一引理，可假设 $\sigma$ 与含入
$\mathcal{K}^\bullet \subset \mathcal{M}^\bullet$ 的复合
由某个同伦
$\{h^i : \mathcal{E}^i \to \mathcal{M}^{i - 1}\}$
同伦于零。因此，把 $\mathcal{K}^{m - 1}$ 替换为
$\mathcal{K}^{m - 1} + \Im(h^m)$ 后（注意替换后
$\mathcal{K}^{m - 1}$ 仍由有限多个整体截面生成），
可见 $\sigma$ 本身同伦于零！
这意味着有如下实线部分交换的图：
$$
\xymatrix{
\mathcal{E}^\bullet \ar[r] &
M \ar[r] &
\mathcal{L}^\bullet \ar[r] &
\mathcal{E}^\bullet[1] \\
\mathcal{E}^\bullet \ar[r] \ar[u] &
K \ar[r] \ar[u] &
\sigma_{\leq m - 1}\mathcal{K}^\bullet \ar[r] \ar[u] &
\mathcal{E}^\bullet[1] \ar[u] \\
\mathcal{E}^\bullet \ar[r] \ar[u] &
M \ar[r] \ar[u]^s &
\mathcal{L}^\bullet \ar[r] \ar@{..>}[u] &
\mathcal{E}^\bullet[1] \ar[u]
}
$$
由三角范畴的公理，得到可嵌入该图的虚线
箭头。
考察次数 $m - 1$ 的上同调层，得到
$$
\xymatrix{
H^{m - 1}(M) \ar[r] &
H^{m - 1}(\mathcal{L}^\bullet) \ar[r] &
H^m(\mathcal{E}^\bullet) \\
H^{m - 1}(K) \ar[r] \ar[u] &
H^{m - 1}(\sigma_{\leq m - 1}\mathcal{K}^\bullet) \ar[r] \ar[u] &
H^m(\mathcal{E}^\bullet) \ar[u] \\
H^{m - 1}(M) \ar[r] \ar[u] &
H^{m - 1}(\mathcal{L}^\bullet) \ar[r] \ar[u] &
H^m(\mathcal{E}^\bullet) \ar[u]
}
$$
由于左列和右列中的竖直复合都是恒等态射，
可知中间的竖直复合
$H^{m - 1}(\mathcal{L}^\bullet) \to
H^{m - 1}(\sigma_{\leq m - 1}\mathcal{K}^\bullet) \to
H^{m - 1}(\mathcal{L}^\bullet)$ 是满射！
特别地，$H^{m - 1}(\sigma_{\leq m - 1}\mathcal{K}^\bullet) \to
H^{m - 1}(\mathcal{L}^\bullet)$ 是满射。
利用上述三角态射所诱导的上同调层长正合列态射，
追图可见，这推出 $H^i(K) \to H^i(M)$ 在
$i \geq m$ 时为同构，在 $i = m - 1$ 时为满射。
按构造，可选择 $r \geq 0$ 和满射
$\mathcal{O}_X^{\oplus r} \to \mathcal{K}^{m - 1}$。于是复合
$$
(\mathcal{O}_X^{\oplus r} \to \mathcal{E}^m \to
\mathcal{E}^{m + 1} \to \ldots ) \longrightarrow
K \longrightarrow M
$$
在次数 $\geq m$ 的上同调层上诱导同构，
并在次数 $m - 1$ 上诱导满射，证明完成。
\end{proof}






\section{杂项}
\label{cohomology-section-misc}

\noindent
一些无法归入别处的结果。

\begin{lemma}
\label{cohomology-lemma-colim-and-lim-of-duals}
令 $(X, \mathcal{O}_X)$ 为环化空间。令
$(K_n)_{n \in \mathbf{N}}$ 为 $D(\mathcal{O}_X)$ 中完美对象组成的系统。
令 $K = \text{hocolim} K_n$ 为导出余极限
（导出范畴，定义 \ref{derived-definition-derived-colimit}）。
那么对 $D(\mathcal{O}_X)$ 的任意对象 $E$，都有
$$
R\SheafHom(K, E) = R\lim E \otimes^\mathbf{L}_{\mathcal{O}_X} K_n^\vee
$$
其中 $(K_n^\vee)$ 是对偶完美复形组成的逆系统。
\end{lemma}

\begin{proof}
由引理 \ref{cohomology-lemma-dual-perfect-complex}，有
$R\lim E \otimes^\mathbf{L}_{\mathcal{O}_X} K_n^\vee =
R\lim R\SheafHom(K_n, E)$，
它嵌入特异三角
$$
R\lim R\SheafHom(K_n, E) \to
\prod R\SheafHom(K_n, E) \to
\prod R\SheafHom(K_n, E)
$$
由于 $K$ 同样嵌入特异三角
$\bigoplus K_n \to \bigoplus K_n \to K$，只需证明
$\prod R\SheafHom(K_n, E) = R\SheafHom(\bigoplus K_n, E)$。
这是 (\ref{cohomology-equation-internal-hom}) 以及导出张量积
与直和可交换的形式推论。
\end{proof}

\begin{lemma}
\label{cohomology-lemma-ext-composition-is-cup}
令 $(X, \mathcal{O}_X)$ 为环化空间。令 $K$ 和 $E$ 为
$D(\mathcal{O}_X)$ 的对象，其中 $E$ 完美。图
$$
\xymatrix{
H^0(X, K \otimes_{\mathcal{O}_X}^\mathbf{L} E^\vee) \times H^0(X, E)
\ar[r] \ar[d] &
H^0(X, K \otimes_{\mathcal{O}_X}^\mathbf{L} E^\vee
\otimes_{\mathcal{O}_X}^\mathbf{L} E) \ar[d] \\
\Hom_X(E, K) \times H^0(X, E) \ar[r] &
H^0(X, K)
}
$$
交换，其中上方水平箭头是杯积；右侧竖直箭头使用
$\epsilon : E^\vee \otimes_{\mathcal{O}_X}^\mathbf{L} E \to \mathcal{O}_X$
（例 \ref{cohomology-example-dual-derived}）；左侧竖直箭头使用引理 \ref{cohomology-lemma-dual-perfect-complex}；下方水平
箭头是显然的箭头。
\end{lemma}

\begin{proof}
我们将简写 $\otimes = \otimes_{\mathcal{O}_X}^\mathbf{L}$
和 $\mathcal{O} = \mathcal{O}_X$。我们把 $E$、$K$
分别与 $R\SheafHom(\mathcal{O}, E)$、$R\SheafHom(\mathcal{O}, K)$
等同，并把 $E^\vee$ 与 $R\SheafHom(E, \mathcal{O})$ 等同。

\medskip\noindent
令 $\xi \in H^0(X, K \otimes E^\vee)$ 且 $\eta \in H^0(X, E)$。
记 $D(\mathcal{O})$ 中相应的态射为
$\tilde \xi : \mathcal{O} \to K \otimes E^\vee$ 和
$\tilde \eta : \mathcal{O} \to E$。由引理
\ref{cohomology-lemma-second-cup-equals-first}，杯积
$\xi \cup \eta$ 对应于
$\tilde \xi \otimes \tilde \eta : \mathcal{O} \to
K \otimes E^\vee \otimes E$。

\medskip\noindent
我们断言，由引理 \ref{cohomology-lemma-dual-perfect-complex} 与 $\xi$
对应的态射 $\xi' : E \to K$ 是复合
$$
E = \mathcal{O} \otimes E
\xrightarrow{\tilde \xi \otimes 1_E}
K \otimes E^\vee \otimes E
\xrightarrow{1_K \otimes \epsilon}
K
$$
引理 \ref{cohomology-lemma-dual-perfect-complex} 中的构造
使用求值态射 (\ref{cohomology-equation-eval})；后者又使用
$E$ 与 $R\SheafHom(\mathcal{O}, E)$ 的等同，以及引理 \ref{cohomology-lemma-internal-hom-composition} 中构造的复合
$\underline{\circ}$。
因此 $\xi'$ 是复合
\begin{align*}
E = \mathcal{O} \otimes
R\SheafHom(\mathcal{O}, E)
& \xrightarrow{\tilde \xi \otimes 1}
R\SheafHom(\mathcal{O}, K) \otimes
R\SheafHom(E, \mathcal{O}) \otimes
R\SheafHom(\mathcal{O}, E) \\
& \xrightarrow{\underline{\circ} \otimes 1}
R\SheafHom(E, K) \otimes R\SheafHom(\mathcal{O}, E) \\
& \xrightarrow{\underline{\circ}}
R\SheafHom(\mathcal{O}, K) = K
\end{align*}
由此以及引理 \ref{cohomology-lemma-internal-hom-composition} 中构造的复合
$\underline{\circ}$ 具有结合性
（在此插入未来引用），并且例 \ref{cohomology-example-dual-derived} 中
$\epsilon$ 定义为复合
$\underline{\circ} : E^\vee \otimes E \to \mathcal{O}$，
立即得到该断言。

\medskip\noindent
使用前两段的结果，可知引理的断言等价于
$(1_K \otimes \epsilon) \circ (\tilde \xi \otimes \tilde \eta)$
等于
$(1_K \otimes \epsilon) \circ (\tilde \xi \otimes 1_E)
\circ (1_\mathcal{O} \otimes \tilde \eta)$，
这是立即的。
\end{proof}

\begin{lemma}
\label{cohomology-lemma-pullback-internal-hom}
令 $h : X \to Y$ 为环化空间的态射。
令 $K, M$ 为 $D(\mathcal{O}_Y)$ 的对象。注记
\ref{cohomology-remark-prepare-fancy-base-change} 中的典范态射
$$
Lh^*R\SheafHom(K, M) \longrightarrow R\SheafHom(Lh^*K, Lh^*M)
$$
在下列情形中是同构：
\begin{enumerate}
\item $K$ 完美；
\item $h$ 平坦，$K$ 伪凝聚，且 $M$（局部）下有界；
\item $\mathcal{O}_X$ 作为 $h^{-1}\mathcal{O}_Y$-模具有有限 Tor 维数，
$K$ 伪凝聚，且 $M$（局部）下有界。
\end{enumerate}
\end{lemma}

\begin{proof}
证明 (1)。问题在 $Y$ 上是局部的，故可假设
$K$ 由严格完美复形
$\mathcal{E}^\bullet$ 表示，见第 \ref{cohomology-section-perfect} 节。
选取表示 $M$ 的 K-平坦复形 $\mathcal{F}^\bullet$。
应用引理 \ref{cohomology-lemma-Rhom-strictly-perfect} 可知，
$R\SheafHom(K, L)$ 由复形
$\mathcal{H}^\bullet =
\SheafHom^\bullet(\mathcal{E}^\bullet, \mathcal{F}^\bullet)$ 表示，
其项为 $\mathcal{H}^n = \bigoplus\nolimits_{n = p + q}
\SheafHom_{\mathcal{O}_X}(\mathcal{E}^{-q}, \mathcal{F}^p)$。
由第 \ref{cohomology-section-derived-pullback} 节中 $Lh^*$ 的构造，
可知 $Lh^*K$ 由严格完美复形
$h^*\mathcal{E}^\bullet$ 表示（引理
\ref{cohomology-lemma-strictly-perfect-pullback}）。
类似地，对象 $Lh^*M$ 由
复形 $h^*\mathcal{F}^\bullet$ 表示。
最后，由引理 \ref{cohomology-lemma-Rhom-strictly-perfect-K-flat}，
$\mathcal{H}^\bullet$ 是 K-平坦的，故对象
$Lh^*R\SheafHom(K, M)$ 由 $h^*\mathcal{H}^\bullet$ 表示。
因此为完成证明，只需说明
$h^*\mathcal{H}^\bullet =
\SheafHom^\bullet(h^*\mathcal{E}^\bullet, h^*\mathcal{F}^\bullet)$。
为此只需注意，只要 $\mathcal{E}$ 是有限自由
$\mathcal{O}_X$-模的直和项，就有
$h^*\SheafHom(\mathcal{E}, \mathcal{F}) =
\SheafHom(h^*\mathcal{E}, \mathcal{F})$。

\medskip\noindent
证明 (2)。由于 $h$ 平坦，可以在任意 $\mathcal{O}_Y$-模复形上
直接使用 $h^*$ 来计算 $Lh^*$。
特别地，对所有 $i \in \mathbf{Z}$，有 $H^i(Lh^*K) = h^*H^i(K)$。
设当 $i < a$ 时 $H^i(M) = 0$。令 $K' \to K$ 为
$D(\mathcal{O}_Y)$ 中的态射，它对所有 $i \geq b$ 诱导同构
$H^i(K') \to H^i(K)$。那么相应的态射
$$
R\SheafHom(K, M) \to R\SheafHom(K', M)
$$
以及
$$
R\SheafHom(Lh^*K, Lh^*M) \to R\SheafHom(Lh^*K', Lh^*M)
$$
在次数 $< a - b$ 的上同调层上都是同构（略去细节）。
因此，为证明引理陈述中的态射在次数 $< a - b$ 的
上同调层上诱导同构，只需在这些次数上对 $K'$ 证明结论。
此外，与 (1) 的证明相同，问题在 $Y$ 上是局部的。
因此可假设 $K$ 由严格完美复形表示，见
第 \ref{cohomology-section-pseudo-coherent} 节。这就约化到情形 (1)。

\medskip\noindent
证明 (3)。证明与 (2) 相同，只是使用
$Lh^*$ 的上同调维数有界来得到所需的消失性。
略去细节。
\end{proof}

\begin{lemma}
\label{cohomology-lemma-stalk-internal-hom}
令 $X$ 为环化空间。令 $K, M$ 为 $D(\mathcal{O}_X)$ 的对象。
令 $x \in X$。典范态射
$$
R\SheafHom(K, M)_x \longrightarrow
R\Hom_{\mathcal{O}_{X, x}}(K_x, M_x)
$$
在下列情形中是同构：
\begin{enumerate}
\item $K$ 完美；
\item $K$ 伪凝聚且 $M$（局部）下有界。
\end{enumerate}
\end{lemma}

\begin{proof}
令 $Y = \{x\}$ 为单点环化空间，其结构层
由 $\mathcal{O}_{X, x}$ 给出。然后把引理
\ref{cohomology-lemma-pullback-internal-hom}
应用于平坦含入态射 $Y \to X$。
\end{proof}







\section{导出范畴中的可逆对象}
\label{cohomology-section-invertible-D-or-R}

\noindent
我们刻画环化空间的导出范畴中的可逆对象
（包括结构层的茎是局部环和不是局部环的情形）。

\begin{lemma}
\label{cohomology-lemma-category-summands-finite-free}
令 $(X, \mathcal{O}_X)$ 为环化空间。
置 $R = \Gamma(X, \mathcal{O}_X)$。作为有限自由
$\mathcal{O}_X$-模直和项的 $\mathcal{O}_X$-模所成的范畴，
等价于有限投射 $R$-模所成的范畴。
\end{lemma}

\begin{proof}
注意，有限投射 $R$-模等同于
有限自由 $R$-模的直和项。
该等价由函子 $\mathcal{E} \mapsto
\Gamma(X, \mathcal{E})$ 给出。逆函子由
模，引理 \ref{modules-lemma-construct-quasi-coherent-sheaves}
中的构造给出。
\end{proof}

\begin{lemma}
\label{cohomology-lemma-invertible-derived}
令 $(X, \mathcal{O}_X)$ 为环化空间。令 $M$ 为
$D(\mathcal{O}_X)$ 的对象。下列条件等价：
\begin{enumerate}
\item $M$ 在 $D(\mathcal{O}_X)$ 中可逆，见
范畴，定义 \ref{categories-definition-invertible}；以及
\item 存在局部有限的乘积分解
$$
\mathcal{O}_X = \prod\nolimits_{n \in \mathbf{Z}} \mathcal{O}_n
$$
并且对每个 $n$，存在可逆 $\mathcal{O}_n$-模
$\mathcal{H}^n$（模，定义 \ref{modules-definition-invertible}），
使得在 $D(\mathcal{O}_X)$ 中 $M = \bigoplus \mathcal{H}^n[-n]$。
\end{enumerate}
如果 (1) 和 (2) 成立，那么 $M$ 是 $D(\mathcal{O}_X)$ 的完美对象。如果
对所有 $x \in X$，$\mathcal{O}_{X, x}$ 都是局部环，那么这些条件
还等价于
\begin{enumerate}
\item[(3)] 存在开覆盖 $X = \bigcup U_i$，
并且对每个 $i$ 存在整数 $n_i$，使得 $M|_{U_i}$
由放在次数 $n_i$ 上的可逆 $\mathcal{O}_{U_i}$-模表示。
\end{enumerate}
\end{lemma}

\begin{proof}
假设 (2)。考虑对象 $R\SheafHom(M, \mathcal{O}_X)$
和复合态射
$$
R\SheafHom(M, \mathcal{O}_X) \otimes_{\mathcal{O}_X}^\mathbf{L} M \to
\mathcal{O}_X
$$
为证明这是同构，可在局部工作。因此可假设
$\mathcal{O}_X = \prod_{a \leq n \leq b} \mathcal{O}_n$
且 $M = \bigoplus_{a \leq n \leq b} \mathcal{H}^n[-n]$。
于是只需证明
$$
R\SheafHom(\mathcal{H}^m, \mathcal{O}_X)
\otimes_{\mathcal{O}_X}^\mathbf{L} \mathcal{H}^n
$$
在 $n \not = m$ 时为零，在 $n = m$ 时等于 $\mathcal{O}_n$。
$n \not = m$ 的情形由如下事实得出：$\mathcal{O}_n$ 和
$\mathcal{O}_m$ 都是平坦 $\mathcal{O}_X$-代数，并且
$\mathcal{O}_n \otimes_{\mathcal{O}_X} \mathcal{O}_m = 0$。
利用可逆 $\mathcal{O}_X$-模的局部结构
（模，引理 \ref{modules-lemma-invertible}）并在局部工作，
$n = m$ 情形中的同构可直接得到；
略去细节。由于 $D(\mathcal{O}_X)$ 是对称幺半范畴，
可知 $M$ 可逆。

\medskip\noindent
假设 (1)。(2) 中的描述给出了
$\mathcal{O}_n$ 的候选，即
$\SheafHom_{\mathcal{O}_X}(H^n(M), H^n(M))$。
如果这是一个局部有限的环层族，
并且 $\mathcal{O}_X = \prod \mathcal{O}_n$，那么利用
$\mathcal{O}_X$ 中来自该乘积分解的幂等元，立即得到
直和分解 $M = \bigoplus H^n(M)[-n]$。
这说明，为证明 (2)，可在 $X$ 上局部工作。

\medskip\noindent
选取 $D(\mathcal{O}_X)$ 的对象 $N$
以及同构
$M \otimes_{\mathcal{O}_X}^\mathbf{L} N \cong \mathcal{O}_X$。
令 $x \in X$。
那么 $N$ 是幺半范畴 $D(\mathcal{O}_X)$ 中 $M$ 的左对偶，
由引理 \ref{cohomology-lemma-left-dual-derived} 可知 $M$ 完美。
由对称性可见 $N$ 完美。用 $x$ 的开邻域替换 $X$ 后，
可假设 $M$ 和 $N$ 分别由严格完美
复形 $\mathcal{E}^\bullet$ 和 $\mathcal{F}^\bullet$ 表示。
于是 $M \otimes_{\mathcal{O}_X}^\mathbf{L} N$ 由
$\text{Tot}(\mathcal{E}^\bullet \otimes_{\mathcal{O}_X} \mathcal{F}^\bullet)$ 表示。
再次缩小 $X$ 后，可假设互逆同构
$\mathcal{O}_X \to M \otimes_{\mathcal{O}_X}^\mathbf{L} N$ 和
$M \otimes_{\mathcal{O}_X}^\mathbf{L} N \to \mathcal{O}_X$
由复形态射
$$
\alpha : \mathcal{O}_X \to
\text{Tot}(\mathcal{E}^\bullet \otimes_{\mathcal{O}_X} \mathcal{F}^\bullet)
\quad\text{且}\quad
\beta :
\text{Tot}(\mathcal{E}^\bullet \otimes_{\mathcal{O}_X} \mathcal{F}^\bullet)
\to \mathcal{O}_X
$$
给出。见引理 \ref{cohomology-lemma-local-actual}。于是作为复形态射，
$\beta \circ \alpha = 1$；而作为
$D(\mathcal{O}_X)$ 中的态射，$\alpha \circ \beta = 1$。
缩小 $X$ 后，由同一引理，可假设复合
$\alpha \circ \beta$ 通过某个同伦 $\theta$ 同伦于 $1$，
其分量为
$$
\theta^n :
\text{Tot}^n(\mathcal{E}^\bullet \otimes_{\mathcal{O}_X} \mathcal{F}^\bullet)
\to
\text{Tot}^{n - 1}(
\mathcal{E}^\bullet \otimes_{\mathcal{O}_X} \mathcal{F}^\bullet)
$$
置 $R = \Gamma(X, \mathcal{O}_X)$。由引理 \ref{cohomology-lemma-category-summands-finite-free}，
得到
\begin{enumerate}
\item $M^\bullet = \Gamma(X, \mathcal{E}^\bullet)$ 是由
有限投射 $R$-模组成的有界复形；
\item $N^\bullet = \Gamma(X, \mathcal{F}^\bullet)$ 是由
有限投射 $R$-模组成的有界复形；
\item $\alpha$ 和 $\beta$ 对应于复形态射
$a : R \to \text{Tot}(M^\bullet \otimes_R N^\bullet)$ 和
$b : \text{Tot}(M^\bullet \otimes_R N^\bullet) \to R$；
\item $\theta^n$ 对应于态射
$h^n : \text{Tot}^n(M^\bullet \otimes_R N^\bullet) \to
\text{Tot}^{n - 1}(M^\bullet \otimes_R N^\bullet)$；以及
\item $b \circ a = 1$ 且 $b \circ a - 1 = dh + hd$。
\end{enumerate}
由此 $M^\bullet$ 和 $N^\bullet$ 定义
$D(R)$ 中互逆的对象。由
更多代数，引理 \ref{more-algebra-lemma-invertible-derived}，
得到乘积分解 $R = \prod_{a \leq n \leq b} R_n$
以及可逆 $R_n$-模 $H^n$，使得
$M^\bullet \cong \bigoplus_{a \leq n \leq b} H^n[-n]$。
由于每个 $H^n$ 都是投射 $R$-模，$D(R)$ 中的该同构
可提升为复形态射
$$
\bigoplus H^n[-n] \longrightarrow M^\bullet
$$
相应地，再次使用引理 \ref{cohomology-lemma-category-summands-finite-free}，
得到态射
$$
\bigoplus H^n \otimes_R \mathcal{O}_X[-n] \to \mathcal{E}^\bullet
$$
它在 $D(\mathcal{O}_X)$ 中为同构。置
$\mathcal{O}_n = R_n \otimes_R \mathcal{O}_X$，可知 (2) 成立。

\medskip\noindent
如果 $\mathcal{O}_X$ 的所有茎都是局部环，那么
(2) 与 (3) 的等价性很容易证明。略去细节。
\end{proof}







\section{紧对象}
\label{cohomology-section-compact}

\noindent
本节研究环化空间上的模的导出范畴中的紧对象。
回忆，紧对象定义于
导出范畴，定义 \ref{derived-definition-compact-object}。
在适当的环化空间上，完美对象是紧的。

\begin{lemma}
\label{cohomology-lemma-when-jshriek-compact}
令 $X$ 为环化空间。令 $j : U \to X$ 为
开集的含入。如果存在整数 $d$，使得
\begin{enumerate}
\item 对所有 $p > d$，$H^p(U, \mathcal{F}) = 0$；并且
\item 函子 $\mathcal{F} \mapsto H^p(U, \mathcal{F})$
与直和可交换，
\end{enumerate}
那么 $\mathcal{O}_X$-模 $j_!\mathcal{O}_U$ 是
$D(\mathcal{O}_X)$ 的紧对象。
\end{lemma}

\begin{proof}
假设 (1) 和 (2)。由层，引理
\ref{sheaves-lemma-j-shriek-modules}，
$\Hom(j_!\mathcal{O}_U, \mathcal{F}) = \mathcal{F}(U)$，
故对 $D(\mathcal{O}_X)$ 中的 $K$，有
$\Hom(j_!\mathcal{O}_U, K) = R\Gamma(U, K)$。
因此必须证明 $R\Gamma(U, -)$ 与直和可交换。
第一个假设意味着函子
$F = H^0(U, -)$ 具有有限上同调维数。此外，第二个
假设意味着内射模的任意直和都对 $F$ 无环。
令 $K_i$ 为 $D(\mathcal{O}_X)$ 的一族对象。
选取表示 $K_i$ 且各项内射的 K-内射代表
$I_i^\bullet$，见内射对象，定理
\ref{injectives-theorem-K-injective-embedding-grothendieck}。
由于可通过把 $F$ 应用于任意无环对象复形来计算 $RF$
（导出范畴，引理 \ref{derived-lemma-unbounded-right-derived}），
并且 $\bigoplus K_i$ 由 $\bigoplus I_i^\bullet$ 表示
（内射对象，引理 \ref{injectives-lemma-derived-products}），
可知 $R\Gamma(U, \bigoplus K_i)$ 由
$\bigoplus H^0(U, I_i^\bullet)$ 表示。因此
$R\Gamma(U, -)$ 如所需地与直和可交换。
\end{proof}

\begin{lemma}
\label{cohomology-lemma-perfect-is-compact}
令 $X$ 为环化空间。假设 $X$ 的底拓扑空间
具有下列性质：
\begin{enumerate}
\item $X$ 拟紧；
\item 存在由拟紧开子集组成的基；并且
\item 任意两个拟紧开集的交仍拟紧。
\end{enumerate}
令 $K$ 为 $D(\mathcal{O}_X)$ 的完美对象。那么
\begin{enumerate}
\item[(a)] $K$ 在如下意义下是 $D^+(\mathcal{O}_X)$ 的紧对象：
如果 $M = \bigoplus_{i \in I} M_i$ 下有界，那么
$\Hom(K, M) = \bigoplus_{i \in I} \Hom(K, M_i)$。
\item[(b)] 如果 $X$ 具有有限上同调维数，即存在
$d$，使得当 $i > d$ 时 $H^i(X, \mathcal{F}) = 0$，那么
$K$ 是 $D(\mathcal{O}_X)$ 的紧对象。
\end{enumerate}
\end{lemma}

\begin{proof}
令 $K^\vee$ 为 $K$ 的对偶，见引理 \ref{cohomology-lemma-dual-perfect-complex}。那么对
$D(\mathcal{O}_X)$ 中的 $M$，函子性地有
$$
\Hom_{D(\mathcal{O}_X)}(K, M) =
H^0(X, K^\vee \otimes_{\mathcal{O}_X}^\mathbf{L} M)
$$
由于 $K^\vee \otimes_{\mathcal{O}_X}^\mathbf{L} -$ 与
直和可交换，只需证明 $R\Gamma(X, -)$
与相应的直和可交换。

\medskip\noindent
证明 (b)。由于 $R\Gamma(X, K) = R\Hom(\mathcal{O}_X, K)$，
并且由引理 \ref{cohomology-lemma-quasi-separated-cohomology-colimit}，
$H^p(X, -)$ 与直和可交换，这是引理 \ref{cohomology-lemma-when-jshriek-compact} 的特例。

\medskip\noindent
证明 (a)。令 $\mathcal{I}_i$（$i \in I$）为一族内射
$\mathcal{O}_X$-模。由引理 \ref{cohomology-lemma-quasi-separated-cohomology-colimit}，
可见
$$
H^p(X, \bigoplus\nolimits_{i \in I} \mathcal{I}_i) =
\bigoplus\nolimits_{i \in I} H^p(X, \mathcal{I}_i) = 0
$$
对所有 $p$ 成立。现在若 $M = \bigoplus M_i$ 如 (a) 中所述，
则存在 $a \in \mathbf{Z}$，使得当 $n < a$ 时
$H^n(M_i) = 0$。因此可选取表示 $M_i$ 的内射
$\mathcal{O}_X$-模复形 $\mathcal{I}_i^\bullet$，
使得当 $n < a$ 时 $\mathcal{I}_i^n = 0$，见
导出范畴，引理 \ref{derived-lemma-injective-resolutions-exist}。
由内射对象，引理 \ref{injectives-lemma-derived-products}，
直和复形 $\bigoplus \mathcal{I}_i^\bullet$
表示 $M$。由勒雷无环性
（导出范畴，引理 \ref{derived-lemma-leray-acyclicity}），
可见
$$
R\Gamma(X, M) = \Gamma(X, \bigoplus \mathcal{I}_i^\bullet) =
\bigoplus \Gamma(X, \bigoplus \mathcal{I}_i^\bullet) =
\bigoplus R\Gamma(X, M_i)
$$
如所需。
\end{proof}



















\section{投影公式}
\label{cohomology-section-projection-formula}

\noindent
本节汇集投影公式的若干变体。
最基本的版本是引理 \ref{cohomology-lemma-projection-formula}。
在陈述并证明它之后，我们将讨论一个涉及
完美复形的更一般版本。

\begin{lemma}
\label{cohomology-lemma-injective-tensor-finite-locally-free}
令 $X$ 为环化空间。
令 $\mathcal{I}$ 为内射 $\mathcal{O}_X$-模。
令 $\mathcal{E}$ 为 $\mathcal{O}_X$-模。
假设 $\mathcal{E}$ 在 $X$ 上有限局部自由，见
模，定义 \ref{modules-definition-locally-free}。
那么 $\mathcal{E} \otimes_{\mathcal{O}_X} \mathcal{I}$
是内射 $\mathcal{O}_X$-模。
\end{lemma}

\begin{proof}
这是因为在引理的假设下，有
$$
\Hom_{\mathcal{O}_X}(\mathcal{F},
\mathcal{E} \otimes_{\mathcal{O}_X} \mathcal{I})
=
\Hom_{\mathcal{O}_X}(
\mathcal{F} \otimes_{\mathcal{O}_X} \mathcal{E}^\vee, \mathcal{I})
$$
其中
$\mathcal{E}^\vee = \SheafHom_{\mathcal{O}_X}(\mathcal{E}, \mathcal{O}_X)$
是 $\mathcal{E}$ 的对偶，它同样有限局部自由。由于与有限局部自由层
作张量积是正合函子，由同调，引理
\ref{homology-lemma-characterize-injectives} 即得结论。
\end{proof}

\begin{lemma}
\label{cohomology-lemma-projection-formula}
令 $f : X \to Y$ 为环化空间的态射。
令 $\mathcal{F}$ 为 $\mathcal{O}_X$-模。
令 $\mathcal{E}$ 为 $\mathcal{O}_Y$-模。
假设 $\mathcal{E}$ 在 $Y$ 上有限局部自由，见
模，定义 \ref{modules-definition-locally-free}。
那么对所有 $q \geq 0$，存在同构
$$
\mathcal{E} \otimes_{\mathcal{O}_Y} R^qf_*\mathcal{F}
\longrightarrow
R^qf_*(f^*\mathcal{E} \otimes_{\mathcal{O}_X} \mathcal{F})
$$
事实上，在 $D^{+}(Y)$ 中存在关于 $\mathcal{F}$
具有函子性的同构
$$
\mathcal{E} \otimes_{\mathcal{O}_Y} Rf_*\mathcal{F}
\longrightarrow
Rf_*(f^*\mathcal{E} \otimes_{\mathcal{O}_X} \mathcal{F})
$$
\end{lemma}

\begin{proof}
在 $X$ 上选取内射消解 $\mathcal{F} \to \mathcal{I}^\bullet$。
注意，$f^*\mathcal{E}$ 也有限局部自由，故得到消解
$$
f^*\mathcal{E} \otimes_{\mathcal{O}_X} \mathcal{F}
\longrightarrow
f^*\mathcal{E} \otimes_{\mathcal{O}_X} \mathcal{I}^\bullet
$$
由引理 \ref{cohomology-lemma-injective-tensor-finite-locally-free}，
这是内射消解。
应用 $f_*$，得到
$$
Rf_*(f^*\mathcal{E} \otimes_{\mathcal{O}_X} \mathcal{F})
=
f_*(f^*\mathcal{E} \otimes_{\mathcal{O}_X} \mathcal{I}^\bullet).
$$
因此，只要能证明
$f_*(f^*\mathcal{E} \otimes_{\mathcal{O}_X} \mathcal{F}) =
\mathcal{E} \otimes_{\mathcal{O}_Y} f_*(\mathcal{F})$
关于 $\mathcal{O}_X$-模 $\mathcal{F}$ 具有函子性，引理即得证。
当 $\mathcal{E} = \mathcal{O}_Y^{\oplus n}$ 时这是显然的；
一般情形通过在 $Y$ 上局部工作得到。略去细节。
\end{proof}

\noindent
令 $f : X \to Y$ 为环化空间的态射。
令 $E \in D(\mathcal{O}_X)$ 且 $K \in D(\mathcal{O}_Y)$。
不作任何进一步假设，存在态射
\begin{equation}
\label{cohomology-equation-projection-formula-map}
Rf_*E \otimes^\mathbf{L}_{\mathcal{O}_Y} K
\longrightarrow
Rf_*(E \otimes^\mathbf{L}_{\mathcal{O}_X} Lf^*K)
\end{equation}
也就是说，它伴随于典范态射
$$
Lf^*(Rf_*E \otimes^\mathbf{L}_{\mathcal{O}_Y} K) =
Lf^*Rf_*E \otimes^\mathbf{L}_{\mathcal{O}_X} Lf^*K
\longrightarrow
E \otimes^\mathbf{L}_{\mathcal{O}_X} Lf^*K
$$
该态射来自 $Lf^*Rf_*E \to E$ 以及引理 \ref{cohomology-lemma-pullback-tensor-product} 和 \ref{cohomology-lemma-adjoint}。
投影公式一个相当一般的版本如下。

\begin{lemma}
\label{cohomology-lemma-projection-formula-perfect}
令 $f : X \to Y$ 为环化空间的态射。
令 $E \in D(\mathcal{O}_X)$ 且 $K \in D(\mathcal{O}_Y)$。
如果 $K$ 完美，那么在 $D(\mathcal{O}_Y)$ 中
$$
Rf_*E \otimes^\mathbf{L}_{\mathcal{O}_Y} K =
Rf_*(E \otimes^\mathbf{L}_{\mathcal{O}_X} Lf^*K)
$$
\end{lemma}

\begin{proof}
为检验 (\ref{cohomology-equation-projection-formula-map}) 是同构，
可在 $Y$ 上局部工作，也就是说，要找到开覆盖
$\{V_j \to Y\}$，使得该态射限制到 $V_j$ 后为同构。
按完美对象的定义，这意味着可假设 $K$ 由
$\mathcal{O}_Y$-模的严格完美复形表示。
注意，完全一般地，断言对
$K = K_1 \oplus K_2$ 成立，当且仅当它对
$K_1$ 和 $K_2$ 都成立。因此可假设 $K$ 是由
有限自由 $\mathcal{O}_Y$-模组成的有限复形。
此时，涉及朴素截断的简单论证把断言约化到
$K$ 由有限自由 $\mathcal{O}_Y$-模表示的情形。
由于断言关于变量 $K$ 在有限直和项下不变，可知
只需对 $K = \mathcal{O}_Y[n]$ 证明它，
而在该情形中结论平凡。
\end{proof}

\noindent
以下是一种投影公式完全一般地成立的情形。

\begin{lemma}
\label{cohomology-lemma-projection-formula-closed-immersion}
令 $f : X \to Y$ 为环化空间的态射，并且 $f$
是到某个闭子集上的同胚。那么
(\ref{cohomology-equation-projection-formula-map}) 总是同构。
\end{lemma}

\begin{proof}
由于 $f$ 是到某个闭子集上的同胚，函子 $f_*$
是正合的（模，引理 \ref{modules-lemma-i-star-exact}）。
因此，可通过把 $f_*$ 应用于任意代表复形来计算 $Rf_*$。
选取表示 $K$ 的 K-平坦 $\mathcal{O}_Y$-模复形
$\mathcal{K}^\bullet$，并任取表示 $E$ 的
$\mathcal{O}_X$-模复形 $\mathcal{E}^\bullet$。那么
$f^*\mathcal{K}^\bullet$ 表示 $Lf^*K$，并且它是
K-平坦 $\mathcal{O}_X$-模复形
（引理 \ref{cohomology-lemma-pullback-K-flat}）。因此
(\ref{cohomology-equation-projection-formula-map}) 的右端由
$$
f_*\text{Tot}(\mathcal{E}^\bullet
\otimes_{\mathcal{O}_X} f^*\mathcal{K}^\bullet)
$$
表示。同理，左端由
$$
\text{Tot}(f_*\mathcal{E}^\bullet \otimes_{\mathcal{O}_Y} \mathcal{K}^\bullet)
$$
表示。由于 $f_*$ 与直和可交换
（模，引理 \ref{modules-lemma-i-star-right-adjoint}），
只需证明
$$
f_*(\mathcal{E} \otimes_{\mathcal{O}_X} f^*\mathcal{K}) =
f_*\mathcal{E} \otimes_{\mathcal{O}_Y} \mathcal{K}
$$
对任意 $\mathcal{O}_X$-模 $\mathcal{E}$ 和 $\mathcal{O}_Y$-模
$\mathcal{K}$ 成立。我们将在茎上检验这一点。
令 $y \in Y$。如果 $y \not \in f(X)$，那么两端的茎
都为零。如果 $y = f(x)$，那么要证明
$$
\mathcal{E}_x \otimes_{\mathcal{O}_{X, x}}
(\mathcal{O}_{X, x} \otimes_{\mathcal{O}_{Y, y}} \mathcal{F}_y) =
\mathcal{E}_x \otimes_{\mathcal{O}_{Y, y}} \mathcal{F}_y
$$
（使用层，引理 \ref{sheaves-lemma-stalks-closed-pushforward}
和引理 \ref{sheaves-lemma-stalk-pullback-modules}）。
该等式成立，故引理得证。
\end{proof}

\begin{remark}
\label{cohomology-remark-compatible-with-diagram}
态射 (\ref{cohomology-equation-projection-formula-map}) 与注记
\ref{cohomology-remark-base-change} 中的基变换态射具有如下相容性。
具体而言，假设
$$
\xymatrix{
X' \ar[r]_{g'} \ar[d]_{f'} &
X \ar[d]^f \\
Y' \ar[r]^g &
Y
}
$$
是环化空间的交换图。
令 $E \in D(\mathcal{O}_X)$ 且 $K \in D(\mathcal{O}_Y)$。
那么图
$$
\xymatrix{
Lg^*(Rf_*E \otimes^\mathbf{L}_{\mathcal{O}_Y} K) \ar[r]_p \ar[d]_t &
Lg^*Rf_*(E \otimes^\mathbf{L}_{\mathcal{O}_X} Lf^*K) \ar[d]_b \\
Lg^*Rf_*E \otimes^\mathbf{L}_{\mathcal{O}_{Y'}} Lg^*K \ar[d]_b &
Rf'_*L(g')^*(E \otimes^\mathbf{L}_{\mathcal{O}_X} Lf^*K) \ar[d]_t \\
Rf'_*L(g')^*E \otimes^\mathbf{L}_{\mathcal{O}_{Y'}} Lg^*K \ar[rd]_p &
Rf'_*(L(g')^*E \otimes^\mathbf{L}_{\mathcal{O}_{Y'}} L(g')^*Lf^*K) \ar[d]_c \\
& Rf'_*(L(g')^*E \otimes^\mathbf{L}_{\mathcal{O}_{Y'}} L(f')^*Lg^*K)
}
$$
交换。这里，标记为 $t$ 的箭头由应用引理 \ref{cohomology-lemma-pullback-tensor-product} 得到，标记为 $b$ 的箭头
由应用注记 \ref{cohomology-remark-base-change} 得到，标记为 $p$ 的箭头
由应用 (\ref{cohomology-equation-projection-formula-map}) 得到，而
$c$ 来自 $L(g')^* \circ Lf^* = L(f')^* \circ Lg^*$。
略去验证。
\end{remark}

















\section{Berthelot 与 Ogus 引入的算子}
\label{cohomology-section-eta}

\noindent
本节继续更多代数，第
\ref{more-algebra-section-eta} 节中开始的讨论。
我们强烈建议读者先阅读该节。

\begin{lemma}
\label{cohomology-lemma-invertible-ideal-sheaf}
令 $(X, \mathcal{O}_X)$ 为环化空间。令
$\mathcal{I} \subset \mathcal{O}_X$ 为理想层。
考虑下列两个条件：
\begin{enumerate}
\item 对每个 $x \in X$，存在 $x$ 的开邻域
$U \subset X$ 和 $f \in \mathcal{I}(U)$，使得
$\mathcal{I}|_U = \mathcal{O}_U \cdot f$，并且
$f : \mathcal{O}_U \to \mathcal{O}_U$ 是单射；以及
\item $\mathcal{I}$ 作为 $\mathcal{O}_X$-模可逆。
\end{enumerate}
那么 (1) 推出 (2)；如果结构层的所有茎
$\mathcal{O}_{X, x}$ 都是局部环，则逆命题也成立。
\end{lemma}

\begin{proof}
略去证明。提示：使用模，引理
\ref{modules-lemma-invertible-is-locally-free-rank-1}。
\end{proof}

\begin{situation}
\label{cohomology-situation-eta}
令 $(X, \mathcal{O}_X)$ 为环化空间。令
$\mathcal{I} \subset \mathcal{O}_X$ 为满足引理 \ref{cohomology-lemma-invertible-ideal-sheaf} 中条件 (1)
的理想层\footnote{本节的讨论可推广到只要求
$\mathcal{I}$ 是模，第 \ref{modules-section-invertible} 节所定义的
可逆 $\mathcal{O}_X$-模这一情形。}。
\end{situation}

\begin{lemma}
\label{cohomology-lemma-I-torsion-free}
在情形 \ref{cohomology-situation-eta} 中，
令 $\mathcal{F}$ 为 $\mathcal{O}_X$-模。
下列条件等价：
\begin{enumerate}
\item 由被 $\mathcal{I}$ 零化的截面组成的子层
$\mathcal{F}[\mathcal{I}] \subset \mathcal{F}$ 为零；
\item 对所有 $n \geq 1$，子层 $\mathcal{F}[\mathcal{I}^n]$ 为零；
\item 乘法态射
$\mathcal{I} \otimes_{\mathcal{O}_X} \mathcal{F} \to \mathcal{F}$
是单射；
\item 对每个满足
$\mathcal{I}|_U = \mathcal{O}_U \cdot f$
（其中 $f \in \mathcal{I}(U)$）的开集 $U \subset X$，
态射 $f : \mathcal{F}|_U \to \mathcal{F}|_U$ 都是单射；
\item 对每个 $x \in X$ 以及理想
$\mathcal{I}_x \subset \mathcal{O}_{X, x}$ 的生成元 $f$，
元素 $f$ 在茎 $\mathcal{F}_x$ 上是非零因子。
\end{enumerate}
\end{lemma}

\begin{proof}
略去证明。
\end{proof}

\noindent
在情形 \ref{cohomology-situation-eta} 中，
令 $\mathcal{F}$ 为 $\mathcal{O}_X$-模。
如果引理 \ref{cohomology-lemma-I-torsion-free} 的等价条件成立，
则称 $\mathcal{F}$ {\it $\mathcal{I}$-无挠}。
在这种情况下，对任意 $i \in \mathbf{Z}$，记
$$
\mathcal{I}^i\mathcal{F} =
\mathcal{I}^{\otimes i} \otimes_{\mathcal{O}_X} \mathcal{F}
$$
于是有含入
$$
\ldots \subset
\mathcal{I}^{i + 1}\mathcal{F} \subset
\mathcal{I}^i\mathcal{F} \subset
\mathcal{I}^{i - 1}\mathcal{F} \subset \ldots
$$
各模 $\mathcal{I}^i\mathcal{F}$ 作为 $\mathcal{O}_X$-模
局部同构于 $\mathcal{F}$，但整体上未必同构。

\medskip\noindent
令 $\mathcal{F}^\bullet$ 为由 $\mathcal{I}$-无挠
$\mathcal{O}_X$-模组成的复形，其微分为
$d^i : \mathcal{F}^i \to \mathcal{F}^{i + 1}$。
此时定义 $\eta_\mathcal{I}\mathcal{F}^\bullet$
为具有下列各项的复形：
\begin{align*}
(\eta_\mathcal{I}\mathcal{F})^i
& =
\Ker\left(
d^i, -1 :
\mathcal{I}^i\mathcal{F}^i \oplus \mathcal{I}^{i + 1}\mathcal{F}^{i + 1}
\to
\mathcal{I}^i\mathcal{F}^{i + 1}
\right) \\
& =
\Ker\left(d^i :
\mathcal{I}^i\mathcal{F}^i
\to
\mathcal{I}^i\mathcal{F}^{i + 1}/
\mathcal{I}^{i + 1}\mathcal{F}^{i + 1}
\right)
\end{align*}
其微分由 $d^i$ 诱导。换言之，
$(\eta_\mathcal{I}\mathcal{F})^i$ 的局部截面 $s$
等同于 $\mathcal{I}^i\mathcal{F}^i$ 的局部截面 $s$，
且它在 $\mathcal{I}^i\mathcal{F}^{i + 1}$ 中的像 $d^i(s)$
属于子层 $\mathcal{I}^{i + 1}\mathcal{F}^{i + 1}$。
注意，$\eta_\mathcal{I}\mathcal{F}^\bullet$
仍是由 $\mathcal{I}$-无挠模组成的复形。

\medskip\noindent
令 $a^\bullet : \mathcal{F}^\bullet \to \mathcal{G}^\bullet$ 为
由 $\mathcal{I}$-无挠 $\mathcal{O}_X$-模组成的复形之间的态射。
于是由态射
$\mathcal{I}^i\mathcal{F}^i \to \mathcal{I}^i\mathcal{G}^i$
诱导复形态射
$$
\eta_\mathcal{I} a^\bullet :
\eta_\mathcal{I}\mathcal{F}^\bullet
\longrightarrow
\eta_\mathcal{I}\mathcal{G}^\bullet
$$
读者可检验，这给出由
$\mathcal{I}$-无挠 $\mathcal{O}_X$-模组成的复形范畴上的
自函子。

\medskip\noindent
如果 $a^\bullet, b^\bullet : \mathcal{F}^\bullet \to \mathcal{G}^\bullet$
是由 $\mathcal{I}$-无挠 $\mathcal{O}_X$-模组成的
复形之间的两个态射，并且
$h = \{h^i : \mathcal{F}^i \to \mathcal{G}^{i - 1}\}$
是 $a^\bullet$ 与 $b^\bullet$ 之间的同伦，那么定义
$\eta_\mathcal{I}h$ 为态射族
$(\eta_\mathcal{I}h)^i : (\eta_\mathcal{I}\mathcal{F})^i \to
(\eta_\mathcal{I}\mathcal{G})^{i - 1}$，
它把 $x$ 送到 $h^i(x)$；这是有意义的，因为
$x$ 是 $\mathcal{I}^i\mathcal{F}^i$ 的局部截面
意味着 $h^i(x)$ 是 $\mathcal{I}^i\mathcal{G}^{i - 1}$ 的局部截面，
而后者当然包含于 $(\eta_\mathcal{I}\mathcal{G})^{i - 1}$。
读者可检验，$\eta_\mathcal{I}h$ 是
$\eta_\mathcal{I}a^\bullet$ 与 $\eta_\mathcal{I}b^\bullet$ 之间的同伦。
综上，在 $\mathcal{I}$-无挠 $\mathcal{O}_X$-模的
加性范畴的同伦范畴上，得到函子
$$
\eta_f :
K(\mathcal{I}\text{-无挠 }\mathcal{O}_X\text{-模})
\longrightarrow
K(\mathcal{I}\text{-无挠 }\mathcal{O}_X\text{-模})
$$
（导出范畴，第 \ref{derived-section-homotopy} 节）。
$\eta_\mathcal{I}$ 在任何意义下都不是三角范畴的正合函子；
可与更多代数，例
\ref{more-algebra-example-eta-not-distinguished} 比较。

\begin{lemma}
\label{cohomology-lemma-eta-stalks}
在情形 \ref{cohomology-situation-eta} 中，
令 $\mathcal{F}^\bullet$ 为由 $\mathcal{I}$-无挠
$\mathcal{O}_X$-模组成的复形。
对 $x \in X$，选取生成元 $f \in \mathcal{I}_x$。那么茎
$(\eta_\mathcal{I}\mathcal{F}^\bullet)_x$
典范同构于更多代数，第 \ref{more-algebra-section-eta} 节中
构造的复形 $\eta_f\mathcal{F}^\bullet_x$。
\end{lemma}

\begin{proof}
略去证明。
\end{proof}

\begin{lemma}
\label{cohomology-lemma-eta-first-property}
在情形 \ref{cohomology-situation-eta} 中，
令 $\mathcal{F}^\bullet$ 为由 $\mathcal{I}$-无挠
$\mathcal{O}_X$-模组成的复形。存在上同调层的典范同构
$$
\mathcal{I}^{\otimes i} \otimes_{\mathcal{O}_X}
\left(
H^i(\mathcal{F}^\bullet)/H^i(\mathcal{F}^\bullet)[\mathcal{I}]
\right)
\longrightarrow H^i(\eta_\mathcal{I}\mathcal{F}^\bullet)
$$
\end{lemma}

\begin{proof}
如下定义态射：
$$
\mathcal{I}^{\otimes i} \otimes_{\mathcal{O}_X} H^i(\mathcal{F}^\bullet)
\longrightarrow
H^i(\eta_\mathcal{I}\mathcal{F}^\bullet)
$$
令 $g$ 为 $\mathcal{I}^{\otimes i}$ 的局部截面，
令 $\overline{s}$ 为 $H^i(\mathcal{F}^\bullet)$ 的局部截面。
那么 $\overline{s}$（局部地）是
$\Ker(d^i : \mathcal{F}^i \to \mathcal{F}^{i + 1})$
的局部截面 $s$ 的类。
我们把 $g \otimes \overline{s}$ 送到
$(\eta_\mathcal{I}\mathcal{F})^i \subset \mathcal{I}^i\mathcal{F}$
的局部截面 $gs$。
当然，$gs$ 属于 $\eta_\mathcal{I}\mathcal{F}^\bullet$ 上
$d^i$ 的核，因而定义
$H^i(\eta_\mathcal{I}\mathcal{F}^\bullet)$ 的局部截面。
检验其适定性没有困难。
我们断言，该态射经由引理中给出的同构分解。
这可在茎上检验，因此通过引理 \ref{cohomology-lemma-eta-stalks}，
它转化为更多代数，引理
\ref{more-algebra-lemma-eta-first-property} 的结果。
\end{proof}

\begin{lemma}
\label{cohomology-lemma-eta-qis}
在情形 \ref{cohomology-situation-eta} 中，
令 $\mathcal{F}^\bullet \to \mathcal{G}^\bullet$ 为
由 $\mathcal{I}$-无挠 $\mathcal{O}_X$-模组成的复形之间的态射。
那么诱导态射
$\eta_\mathcal{I}\mathcal{F}^\bullet \to \eta_\mathcal{I}\mathcal{G}^\bullet$
也是拟同构。
\end{lemma}

\begin{proof}
这是因为引理 \ref{cohomology-lemma-eta-first-property} 的同构
与复形态射相容。
\end{proof}


\begin{lemma}
\label{cohomology-lemma-Leta}
在情形 \ref{cohomology-situation-eta} 中，存在加性
函子\footnote{注意，此函子并不正合，即不把
特异三角变为特异三角。}
$L\eta_\mathcal{I} : D(\mathcal{O}_X) \to D(\mathcal{O}_X)$，
使得如果 $D(\mathcal{O}_X)$ 中的 $M$ 由
$\mathcal{I}$-无挠 $\mathcal{O}_X$-模复形
$\mathcal{F}^\bullet$ 表示，那么
$L\eta_\mathcal{I}M = \eta_\mathcal{I}\mathcal{F}^\bullet$。
态射也同样如此。
\end{lemma}

\begin{proof}
记 $\mathcal{T} \subset \textit{Mod}(\mathcal{O}_X)$ 为
$\mathcal{I}$-无挠 $\mathcal{O}_X$-模组成的全子范畴。
相应地，有含入
$$
K(\mathcal{T})
\quad\subset\quad
K(\textit{Mod}(\mathcal{O}_X)) = K(\mathcal{O}_X)
$$
它把 $K(\mathcal{T})$ 视为 $K(\mathcal{O}_X)$ 的全三角子范畴。
令 $S \subset \text{Arrows}(K(\mathcal{T}))$ 为拟同构组成的集合。
我们将应用
导出范畴，引理 \ref{derived-lemma-localization-subcategory}
证明态射
$$
S^{-1}K(\mathcal{T}) \longrightarrow D(\mathcal{O}_X)
$$
是三角范畴的等价。该引理说明，只需证明：
给定 $\mathcal{O}_X$-模复形 $\mathcal{G}^\bullet$，
存在拟同构 $\mathcal{F}^\bullet \to \mathcal{G}^\bullet$，
其中 $\mathcal{F}^\bullet$ 是由 $\mathcal{I}$-无挠
$\mathcal{O}_X$-模组成的复形。
由引理 \ref{cohomology-lemma-K-flat-resolution}，可找到拟同构
$\mathcal{F}^\bullet \to \mathcal{G}^\bullet$，使得复形
$\mathcal{F}^\bullet$ 是 K-平坦的（我们不会用到这一点），并且
由平坦 $\mathcal{O}_X$-模 $\mathcal{F}^i$ 组成。
由引理 \ref{cohomology-lemma-I-torsion-free} 的第三个刻画，
平坦 $\mathcal{O}_X$-模是
$\mathcal{I}$-无挠 $\mathcal{O}_X$-模，证明完成。

\medskip\noindent
有了这些准备，就可以定义 $L\eta_f$。
具体地，由本节引理 \ref{cohomology-lemma-I-torsion-free} 后面的讨论，
我们已经有良定义函子
$$
K(\mathcal{T}) \xrightarrow{\eta_f} K(\mathcal{T}) \to
K(\mathcal{O}_X) \to D(\mathcal{O}_X)
$$
按照引理 \ref{cohomology-lemma-eta-qis}，它把拟同构
送到拟同构。因此，由范畴，引理
\ref{categories-lemma-properties-left-localization}，
该函子经由
$S^{-1}K(\mathcal{T}) = D(\mathcal{O}_X)$ 分解。
\end{proof}

\noindent
在情形 \ref{cohomology-situation-eta} 中，我们构造 Bockstein 算子。
首先注意有交换图
$$
\xymatrix{
0 \ar[r] &
\mathcal{I}^{i + 1} \ar[r] \ar[d] &
\mathcal{I}^i \ar[r] \ar[d] &
\mathcal{I}^i/\mathcal{I}^{i + 1}
\ar[r] \ar@{=}[d] &
0 \\
0 \ar[r] &
\mathcal{I}^{i + 1}/\mathcal{I}^{i + 2} \ar[r] &
\mathcal{I}^i/\mathcal{I}^{i + 2} \ar[r] &
\mathcal{I}^i/\mathcal{I}^{i + 1} \ar[r] &
0
}
$$
其两行都是 $\mathcal{O}_X$-模的短正合列。
令 $M$ 为 $D(\mathcal{O}_X)$ 的对象。
把上述图与 $M$ 作张量积，得到特异三角的态射
$$
\xymatrix{
M \otimes^\mathbf{L} \mathcal{I}^{i + 1} \ar[r] \ar[d] &
M \otimes^\mathbf{L} \mathcal{I}^i \ar[r] \ar[d] &
M \otimes^\mathbf{L} \mathcal{I}^i/\mathcal{I}^{i + 1} \ar[d]^{\text{id}} \\
M \otimes^\mathbf{L} \mathcal{I}^{i + 1}/\mathcal{I}^{i + 2} \ar[r] &
M \otimes^\mathbf{L} \mathcal{I}^i/\mathcal{I}^{i + 2} \ar[r] &
M \otimes^\mathbf{L} \mathcal{I}^i/\mathcal{I}^{i + 1}
}
$$
尤其是，与下方三角相伴的上同调层长正合列
确定{\it Bockstein 算子}
$$
\beta = \beta^i :
H^i(M \otimes^\mathbf{L} \mathcal{I}^i/\mathcal{I}^{i + 1})
\longrightarrow
H^{i + 1}(M \otimes^\mathbf{L} \mathcal{I}^{i + 1}/\mathcal{I}^{i + 2})
$$
其中 $i \in \mathbf{Z}$。为供后用，我们在此记录：
由上述交换图，Bockstein 算子有分解
\begin{equation}
\label{cohomology-equation-factorization-bockstein}
\vcenter{
\xymatrix{
H^i(M \otimes^\mathbf{L} \mathcal{I}^i/\mathcal{I}^{i + 1})
\ar[r]_\delta \ar[rd]_\beta &
H^{i + 1}(M \otimes^\mathbf{L} \mathcal{I}^{i + 1}) \ar[d] \\
&
H^{i + 1}(M \otimes^\mathbf{L} \mathcal{I}^{i + 1}/\mathcal{I}^{i + 2})
}
}
\end{equation}
其中 $\delta$ 是来自上述交换图中上方特异三角的边界算子。
得到复形
\begin{equation}
\label{cohomology-equation-complex-bocksteins}
H^\bullet(M/\mathcal{I}) =
\left[
\begin{matrix}
\ldots \\
\downarrow \\
H^{i - 1}(M \otimes^\mathbf{L} \mathcal{I}^{i - 1}/\mathcal{I}^i) \\
\downarrow \beta \\
H^i(M \otimes^\mathbf{L} \mathcal{I}^i/\mathcal{I}^{i + 1}) \\
\downarrow \beta \\
H^{i + 1}(M \otimes^\mathbf{L} \mathcal{I}^{i + 1}/\mathcal{I}^{i + 2}) \\
\downarrow \\
\ldots
\end{matrix}
\right]
\end{equation}
也就是说，$\beta \circ \beta = 0$。这可以在茎上检验，
此时可由更多代数，第 \ref{more-algebra-section-eta} 节中
证明的相应代数结果推出。
另一种证明：短正合列
$0 \to \mathcal{I}^{i + 1}/\mathcal{I}^{i + 2}
\to \mathcal{I}^i/\mathcal{I}^{i + 2}
\to \mathcal{I}^i/\mathcal{I}^{i + 1} \to 0$
在 $D(\mathcal{O}_X)$ 中定义态射
$b^i : \mathcal{I}^i/\mathcal{I}^{i + 1} \to
(\mathcal{I}^{i + 1}/\mathcal{I}^{i + 2})[1]$；
这些态射与 $M$ 作张量积并取上同调层后，
诱导上述态射 $\beta$。
然后利用导出范畴，引理
\ref{derived-lemma-cup-ext-1-zero} 中的判据，以及
模 $\mathcal{I}^i/\mathcal{I}^{i + 3}$
是 $\mathcal{I}^{i + 1}/\mathcal{I}^{i + 3}$ 被
$\mathcal{I}^i/\mathcal{I}^{i + 1}$ 的扩张这一事实，可证明复合
$b^{i + 1}[1] \circ b^i : \mathcal{I}^i/\mathcal{I}^{i + 1} \to
(\mathcal{I}^{i + 1}/\mathcal{I}^{i + 2})[1] \to
(\mathcal{I}^{i + 2}/\mathcal{I}^{i + 3})[2]$
在 $D(\mathcal{O}_X)$ 中为零。


\begin{lemma}
\label{cohomology-lemma-eta-second-property}
在情形 \ref{cohomology-situation-eta} 中，令 $M$ 为
$D(\mathcal{O}_X)$ 的对象。存在 $D(\mathcal{O}_X)$ 中的典范同构
$$
L\eta_\mathcal{I}M \otimes^\mathbf{L} \mathcal{O}_X/\mathcal{I}
\longrightarrow
H^\bullet(M/\mathcal{I})
$$
其中右端是复形
(\ref{cohomology-equation-complex-bocksteins})。
\end{lemma}

\begin{proof}
根据引理 \ref{cohomology-lemma-eta-qis} 中 $L\eta_\mathcal{I}$ 的构造，
可假设 $M$ 由 $\mathcal{I}$-无挠
$\mathcal{O}_X$-模复形 $\mathcal{F}^\bullet$ 表示。那么
$L\eta_\mathcal{I}M$ 由复形
$\eta_\mathcal{I}\mathcal{F}^\bullet$ 表示，而后者同样
由 $\mathcal{I}$-无挠 $\mathcal{O}_X$-模组成。
因此 $L\eta_\mathcal{I}M \otimes^\mathbf{L} \mathcal{O}_X/\mathcal{I}$
由复形
$\eta_\mathcal{I}\mathcal{F}^\bullet \otimes \mathcal{O}_X/\mathcal{I}$
表示。类似地，复形 $H^\bullet(M/\mathcal{I})$ 的各项为
$H^i(\mathcal{F}^\bullet \otimes \mathcal{I}^i/\mathcal{I}^{i + 1})$。

\medskip\noindent
令 $f$ 为 $\mathcal{I}$ 的局部生成元。
令 $s$ 为 $(\eta_\mathcal{I}\mathcal{F})^i$ 的局部截面。
那么可写成 $s = f^is'$，其中 $s'$ 是
$\mathcal{F}^i$ 的局部截面；类似地，
$d^i(s) = f^{i + 1}t$，其中 $t$ 是
$\mathcal{F}^{i + 1}$ 的局部截面。因此 $d^i$ 把 $f^is'$
映到 $\mathcal{F}^{i + 1} \otimes \mathcal{I}^i/\mathcal{I}^{i + 1}$
中的零。
故可把 $s$ 映到
$H^i(\mathcal{F}^\bullet \otimes \mathcal{I}^i/\mathcal{I}^{i + 1})$
中 $f^is'$ 的类。
该规则定义 $\mathcal{O}_X$-模态射
$$
(\eta_\mathcal{I}\mathcal{F})^i \otimes \mathcal{O}_X/\mathcal{I}
\longrightarrow
H^i(\mathcal{F}^\bullet \otimes \mathcal{I}^i/\mathcal{I}^{i + 1})
$$
计算表明，这些态射与微分相容（本质上因为
$\beta$ 把 $f^is'$ 的类送到 $f^{i + 1}t$ 的类），
从而得到一个表示引理陈述中箭头的复形态射。

\medskip\noindent
为完成证明，注意上一段给出的构造在茎上
与更多代数，引理
\ref{more-algebra-lemma-eta-second-property}
中构造的态射一致，结论随即得出。
\end{proof}

\begin{lemma}
\label{cohomology-lemma-eta-tensor-invertible}
在情形 \ref{cohomology-situation-eta} 中，
令 $\mathcal{F}^\bullet$ 为由
$\mathcal{I}$-无挠 $\mathcal{O}_X$-模组成的复形。
令 $\mathcal{L}$ 为可逆 $\mathcal{O}_X$-模。
那么 $\eta_\mathcal{I}(\mathcal{F}^\bullet \otimes \mathcal{L}) =
(\eta_\mathcal{I}\mathcal{F}^\bullet) \otimes \mathcal{L}$。
\end{lemma}

\begin{proof}
由构造立即可得。
\end{proof}

\begin{lemma}
\label{cohomology-lemma-eta-cohomology-locally-free}
在情形 \ref{cohomology-situation-eta} 中，令 $M$ 为 $D(\mathcal{O}_X)$ 的对象。
令 $x \in X$ 且 $\mathcal{O}_{X, x}$ 非零。如果 $H^i(M)_x$
是有限自由 $\mathcal{O}_{X, x}$-模，那么
$H^i(L\eta_\mathcal{I}M)_x$ 是同秩的
有限自由 $\mathcal{O}_{X, x}$-模。
\end{lemma}

\begin{proof}
具体地，设 $f \in \mathcal{O}_{X, x}$ 生成茎 $\mathcal{I}_x$。
那么 $f$ 是 $\mathcal{O}_{X, x}$ 中的非零因子，因而
$H^i(M)_x[f] = 0$。因此由引理 \ref{cohomology-lemma-eta-first-property}，
$H^i(L\eta_\mathcal{I}M)_x$ 同构于
$\mathcal{I}^i_x \otimes_{\mathcal{O}_{X, x}} H^i(M)_x$，
后者如所需地是同秩自由模。
\end{proof}
